% !TeX program = xelatex
% Compile twice: xelatex -interaction=nonstopmode main.tex
\documentclass[UTF8,10.5pt,a4paper,oneside,fontset=fandol,scheme=plain]{ctexbook}
% ==================== Art-Book Editorial Preamble for AI Series Volume 5 ====================
\usepackage[top=24mm,bottom=24mm,left=23mm,right=23mm,headheight=22pt,headsep=7mm]{geometry}
\usepackage{amsmath,amssymb,mathtools,amsthm}
\usepackage{booktabs,tabularx,array,multirow,longtable}
\usepackage{graphicx}
\xeCJKDeclareCharClass{CJK}{"2460 -> "2462}
\usepackage{xcolor}
\usepackage{tikz}
\usepackage{pgfplots}
\pgfplotsset{compat=1.18}
\usetikzlibrary{arrows.meta,positioning,shapes.geometric,calc,decorations.pathreplacing,backgrounds}
\usepackage[most]{tcolorbox}
\usepackage{fancyhdr}
\usepackage{listings}
\usepackage{enumitem}
\usepackage{hyperref}

% ==================== Curated Gallery Color Palette ====================
\definecolor{ink}{HTML}{1B2430}        % Deep charcoal-slate ink
\definecolor{teal}{HTML}{0B5555}       % Refined verdigris teal
\definecolor{tealDark}{HTML}{083B3B}   % Deep gallery spine teal
\definecolor{tealLight}{HTML}{F2F8F8}  % Delicate ice-teal parchment
\definecolor{tealLine}{HTML}{A3C6C5}   % Soft architectural rule
\definecolor{tealMist}{HTML}{DCEBEB}   % Watermark chapter numeral tint
\definecolor{navy}{HTML}{1D3B5E}       % Muted Prussian indigo
\definecolor{navyLight}{HTML}{F2F5FA}  % Soft indigo tint
\definecolor{ochre}{HTML}{9E4A22}      % Burnished terracotta-ochre
\definecolor{gold}{HTML}{B97A24}       % Fine burnished amber-gold accent
\definecolor{ochreLight}{HTML}{FCF5F0} % Warm parchment tint
\definecolor{plum}{HTML}{5B3A6B}       % Muted editorial plum
\definecolor{plumLight}{HTML}{F8F4FB}  % Soft plum tint
\definecolor{slate}{HTML}{526071}      % Muted architectural gray
\definecolor{codebg}{HTML}{F7F9FB}     % Clean light code canvas
\definecolor{codeborder}{HTML}{D4DCE5}

\hypersetup{
  unicode=true,
  colorlinks=true,
  urlcolor=teal,
  linkcolor=teal,
  citecolor=teal,
  pdftitle={天书中的智能 · 第5卷：深度学习、自动微分与现代架构},
  pdfauthor={天书中的智能：人工智能的七卷手记}
}

% ==================== Editorial Chapter Openers & Typography ====================
\newcommand{\chapWatermarkNum}{\ifnum\value{chapter}<10 0\fi\thechapter}
\ctexset{
  contentsname={目\quad 录},
  chapter={
    name={第,章},
    number=\arabic{chapter},
    format=\LARGE\sffamily\bfseries\color{tealDark},
    nameformat=\small\sffamily\bfseries\color{gold},
    aftername=\hfill{\fontsize{38}{38}\selectfont\sffamily\bfseries\color{tealMist}\chapWatermarkNum}\par\vspace{-4mm},
    aftertitle=\par\vspace{2.8mm}{\noindent\color{gold}\rule{26mm}{1.5pt}}{\color{tealLine}\rule{\dimexpr\linewidth-26mm\relax}{0.55pt}}\par,
    beforeskip=6pt,
    afterskip=10pt
  },
  section={
    format=\Large\sffamily\bfseries\color{teal},
    beforeskip=14pt plus 3pt minus 2pt,
    afterskip=6pt plus 1pt
  },
  subsection={
    format=\large\sffamily\bfseries\color{ink},
    beforeskip=10pt plus 2pt minus 1pt,
    afterskip=4pt plus 1pt
  },
  subsubsection={
    format=\normalsize\sffamily\bfseries\color{navy},
    beforeskip=7pt,
    afterskip=2pt
  }
}
\renewcommand{\lstlistingname}{代码}
\renewcommand{\figurename}{图}
\renewcommand{\tablename}{表}

% Artisanal Chapter Epigraph macro (2 args: quote + right-aligned attribution)
\newcommand{\chapEpigraph}[2]{%
  \vspace{-2mm}%
  \begin{tcolorbox}[
    enhanced,
    colback=ochreLight!45!white,
    colframe=ochreLight!45!white,
    borderline west={1.8pt}{0pt}{gold},
    boxrule=0pt,
    arc=0.8mm,
    left=4mm,right=3.5mm,top=2mm,bottom=2mm,
    before skip=0pt,after skip=8pt
  ]
    {\kaishu\normalsize\color{ink!85}“#1”}\par\vspace{1mm}%
    \hfill{\footnotesize\sffamily\color{slate}——\,#2}%
  \end{tcolorbox}%
}

% Exercise tier header macro
\newcommand{\exTier}[2]{%
  \par\vspace{3.8mm}\noindent
  \colorbox{tealLight}{\strut\;\sffamily\bfseries\footnotesize\color{tealDark}#1\;}\;%
  {\small\sffamily\bfseries\color{slate}#2}\par\nopagebreak\vspace{1.2mm}%
}
\newcommand{\exHint}[1]{\par\smallskip\noindent{\footnotesize\sffamily\color{navy}\textbf{【提示】}\,#1}}

\setlength{\parindent}{0pt}
\setlength{\parskip}{5pt plus 1.5pt minus 0.5pt}
\linespread{1.26}
\setlist[itemize]{leftmargin=1.8em,topsep=3pt,itemsep=2.5pt,parsep=1pt}
\setlist[enumerate]{leftmargin=2.0em,topsep=3pt,itemsep=2.5pt,parsep=1pt}
\setlength{\tabcolsep}{5.5pt}
\renewcommand{\arraystretch}{1.22}

% ==================== Running Headers & Footers ====================
\pagestyle{fancy}
\fancyhf{}
\fancyhead[L]{\small\sffamily\color{slate}天书中的智能 \;\textcolor{gold}{$\cdot$}\; 第 5 卷}
\fancyhead[R]{\small\sffamily\color{slate}\nouppercase{\leftmark}}
\fancyfoot[C]{\small\sffamily\color{slate}\textcolor{gold}{$\diamond$}\quad\textcolor{tealDark}{\thepage}\quad\textcolor{gold}{$\diamond$}}
\renewcommand{\headrulewidth}{0.4pt}
\renewcommand{\headrule}{\hbox to\headwidth{\color{tealLine}\leaders\hrule height \headrulewidth\hfill}}
\renewcommand{\chaptermark}[1]{\markboth{第 \thechapter\ 章\quad #1}{}}

\fancypagestyle{plain}{
  \fancyhf{}
  \fancyfoot[C]{\small\sffamily\color{slate}\textcolor{gold}{$\diamond$}\quad\textcolor{tealDark}{\thepage}\quad\textcolor{gold}{$\diamond$}}
  \renewcommand{\headrulewidth}{0pt}
}

% ==================== Theorem & Pedagogical Box Environments ====================
\theoremstyle{definition}
\newtheorem{theorem}{定理}[chapter]
\newtheorem{lemma}[theorem]{引理}
\newtheorem{proposition}[theorem]{命题}
\newtheorem{corollary}[theorem]{推论}
\newtheorem{definition}[theorem]{定义}
\newtheorem{example}[theorem]{例}
\newtheorem{remark}[theorem]{注记}
\renewcommand{\proofname}{\textbf{\sffamily\color{teal}证明}}

\newtcolorbox{thmBox}[1][]{
  enhanced,breakable,
  colback=tealLight,
  colframe=tealLine,
  borderline west={2.4pt}{0pt}{teal},
  boxrule=0.45pt,
  arc=1.5mm,
  left=4mm,right=3.8mm,top=2.8mm,bottom=2.8mm,
  before skip=6.5pt,after skip=6.5pt,
  #1
}

\newtcolorbox{defBox}[1][]{
  enhanced,breakable,
  colback=navyLight,
  colframe=navy!30,
  borderline west={2.4pt}{0pt}{navy},
  boxrule=0.45pt,
  arc=1.5mm,
  left=4mm,right=3.8mm,top=2.8mm,bottom=2.8mm,
  before skip=6.5pt,after skip=6.5pt,
  #1
}

\newtcolorbox{insightBox}[1]{
  enhanced,breakable,
  colback=white,
  colframe=navy!35,
  borderline west={2.4pt}{0pt}{gold},
  coltitle=navy,
  fonttitle=\sffamily\bfseries\small,
  attach title to upper={\quad},
  title={【直觉与下游坐标】#1},
  boxrule=0.45pt,
  arc=1.2mm,
  left=4mm,right=3.8mm,top=2.6mm,bottom=2.6mm,
  before skip=6.5pt,after skip=6.5pt
}

\newtcolorbox{recallBox}[1]{
  enhanced,breakable,
  colback=codebg,
  colframe=codeborder,
  borderline west={2.2pt}{0pt}{slate},
  coltitle=slate,
  fonttitle=\sffamily\bfseries\small,
  attach title to upper={\quad},
  title={【前置速忆】#1},
  boxrule=0.4pt,
  arc=1.2mm,
  left=3.8mm,right=3.5mm,top=2.4mm,bottom=2.4mm,
  before skip=5.5pt,after skip=5.5pt
}

\newtcolorbox{failBox}[1]{
  enhanced,breakable,
  colback=ochreLight,
  colframe=ochre!35,
  borderline west={2.4pt}{0pt}{ochre},
  coltitle=ochre,
  fonttitle=\sffamily\bfseries\small,
  attach title to upper={\quad},
  title={【反例、数值陷阱与排错】#1},
  boxrule=0.45pt,
  arc=1.2mm,
  left=4mm,right=3.8mm,top=2.6mm,bottom=2.6mm,
  before skip=6.5pt,after skip=6.5pt
}

\newtcolorbox{extBox}[1]{
  enhanced,breakable,
  colback=plumLight,
  colframe=plum!30,
  borderline west={2.4pt}{0pt}{plum},
  coltitle=plum,
  fonttitle=\sffamily\bfseries\small,
  attach title to upper={\quad},
  title={【选读拓展 · 初读可跳过】#1},
  boxrule=0.45pt,
  arc=1.2mm,
  left=4mm,right=3.8mm,top=2.6mm,bottom=2.6mm,
  before skip=6.5pt,after skip=6.5pt
}

\newtcolorbox{chapIntro}[1]{
  enhanced,
  colback=tealLight,
  colframe=tealLine,
  borderline west={2.8pt}{0pt}{tealDark},
  boxrule=0.45pt,
  arc=1.5mm,
  left=4.2mm,right=4mm,top=3mm,bottom=3mm,
  before skip=2pt,after skip=10pt,
  title={\sffamily\bfseries\small\color{tealDark}#1},
  colbacktitle=tealLight,
  attach boxed title to top left={yshift=-1mm,xshift=2mm},
  boxed title style={empty,size=minimal}
}

% ==================== Code Listing Style ====================
\lstdefinestyle{pycode}{
  language=Python,
  basicstyle=\footnotesize\ttfamily\color{ink},
  keywordstyle=\bfseries\color{teal},
  commentstyle=\itshape\color{slate},
  stringstyle=\color{ochre},
  backgroundcolor=\color{codebg},
  frame=single,
  rulecolor=\color{codeborder},
  framesep=5pt,
  xleftmargin=4pt,
  xrightmargin=4pt,
  numbers=left,
  numberstyle=\tiny\ttfamily\color{slate!80},
  numbersep=6pt,
  columns=fullflexible,
  keepspaces=true,
  showstringspaces=false,
  breaklines=true,
  aboveskip=6.5pt,
  belowskip=6.5pt
}
\lstset{style=pycode}

% ==================== Math Shortcuts ====================
\newcommand{\R}{\mathbb{R}}
\newcommand{\C}{\mathbb{C}}
\newcommand{\N}{\mathbb{N}}
\newcommand{\E}{\mathbb{E}}
\newcommand{\Prob}{\mathbb{P}}
\newcommand{\Var}{\operatorname{Var}}
\newcommand{\Cov}{\operatorname{Cov}}
\newcommand{\DKL}{D_{\mathrm{KL}}}
\newcommand{\softmax}{\operatorname{softmax}}
\newcommand{\ReLU}{\operatorname{ReLU}}
\newcommand{\LN}{\operatorname{LN}}
\newcommand{\BN}{\operatorname{BN}}
\newcommand{\RMSN}{\operatorname{RMSNorm}}
\newcommand{\Attn}{\operatorname{Attn}}
\newcommand{\diag}{\operatorname{diag}}
\newcommand{\tr}{\operatorname{tr}}
\newcommand{\argmin}{\operatorname*{arg\,min}}
\newcommand{\argmax}{\operatorname*{arg\,max}}
\newcommand{\norm}[1]{\left\|#1\right\|}
\newcommand{\abs}[1]{\left|#1\right|}
\newcommand{\inner}[2]{\left\langle #1,\, #2 \right\rangle}
\newcommand{\had}{\odot}

\makeatletter
\setlength{\@fpsep}{16pt}
\makeatother


\begin{document}
\color{ink}

% ==================== Cover Page (Gallery Art-Book Edition) ====================
\frontmatter
\thispagestyle{empty}
\begin{tikzpicture}[remember picture, overlay]
  % Architectural deep teal spine + burnished gold hairline + vertical monograph label
  \fill[tealDark] (current page.north west) rectangle ([xshift=11mm]current page.south west);
  \fill[gold] ([xshift=11mm]current page.north west) rectangle ([xshift=12.2mm]current page.south west);
  \node[rotate=90, anchor=center, font=\sffamily\bfseries\scriptsize, color=tealMist, opacity=0.85]
    at ([xshift=5.5mm]current page.west) {INTELLIGENCE FROM THE BOOK \quad $\cdot$ \quad VOLUME V};

  % Bespoke fine-art geometric Transformer Attention Web & Diffusion Wave motif (top-right)
  \begin{scope}[shift={([xshift=-95mm, yshift=-20mm]current page.north east)}, scale=0.56, transform shape]
    % Faint architectural attention arcs connecting token nodes
    \foreach \i/\x in {1/-2.2, 2/-1.1, 3/0.0, 4/1.1, 5/2.2} {
      \node[circle, fill=teal, opacity=0.45, inner sep=2.2pt] (b\i) at (\x, -1.2) {};
      \node[circle, fill=navy, opacity=0.50, inner sep=2.2pt] (t\i) at (\x, 1.2) {};
    }
    % Weighted self-attention lines
    \draw[teal, opacity=0.14, line width=0.6pt] (b1) -- (t2) (b1) -- (t3) (b2) -- (t4) (b3) -- (t1) (b4) -- (t2) (b5) -- (t3);
    \draw[gold, opacity=0.58, line width=1.25pt] (b1) to[out=70,in=-110] (t4);
    \draw[gold, opacity=0.45, line width=1.0pt] (b2) to[out=75,in=-105] (t3);
    \draw[ochre, opacity=0.48, line width=1.1pt] (b4) to[out=105,in=-75] (t2);
    \draw[teal, opacity=0.42, line width=0.95pt] (b3) to[out=85,in=-95] (t5);
    \draw[teal, opacity=0.35, line width=0.85pt] (b5) to[out=100,in=-80] (t4);
    % Residual highway arc on the right
    \draw[navy, opacity=0.30, line width=0.9pt, dashed] (2.45, -1.2) to[out=35, in=-35] (2.45, 1.2);
    % Subtle computational graph / reverse-mode halo rings
    \draw[teal, opacity=0.15, line width=0.7pt] (0, 0) ellipse (2.65cm and 0.65cm);
    \draw[gold, opacity=0.22, line width=0.6pt] (0, 0) circle (1.45cm);
  \end{scope}
\end{tikzpicture}
\vspace*{1mm}
\noindent{\small\sffamily\bfseries\color{teal}天书中的智能 \;\textcolor{gold}{$\cdot$}\; 人工智能的七卷手记 \hfill 第 5 卷 \quad 2026 年 11 月版}
\vspace{1.5mm}
\noindent{\color{gold}\rule{30mm}{1.8pt}}{\color{tealLine}\rule{\dimexpr\linewidth-30mm\relax}{0.7pt}}
\vspace{6.5mm}

{\fontsize{25}{33}\selectfont\sffamily\bfseries\color{ink}深度学习、自动微分与现代架构}\par
\vspace{2.5mm}
{\large\sffamily\color{teal}从反向模式自动微分与训练动力学，走向 Transformer、预训练与扩散生成}\par
\vspace{4.5mm}
{\small\color{slate}对标 Stanford CS231n / CS224n / CS336 $\cdot$ UC Berkeley CS182 $\cdot$ CMU 11-785 $\cdot$ 南大《深度学习》}\par

\vspace{5.5mm}

\begin{tcolorbox}[
  enhanced,
  colback=tealLight,
  colframe=tealLine,
  borderline west={2.6pt}{0pt}{tealDark},
  arc=2mm,
  boxrule=0.5pt,
  left=4.5mm,right=4.5mm,top=3mm,bottom=3mm
]
\textbf{\sffamily\color{tealDark}本卷定位：彻底拆开现代深度学习的黑箱，从零构造核心引擎}\par\vspace{1.2mm}
\small
告别“调库炼丹”的模糊直觉，本卷用第一性原理与极简可验证代码，将现代神经网络从最底层的微积分计算图一路搭到当代大模型与扩散生成架构：
\begin{itemize}
  \item \textbf{自动微分与训练动力学}：亲手编写反向模式自动微分引擎（VJP）并与 PyTorch 精度互验；严格推导 Kaiming 初始化方差守恒、BatchNorm/LayerNorm/RMSNorm 差异与 AdamW 解耦权重衰减。
  \item \textbf{视觉与序列归纳偏置}：剖析 2D 卷积的平移等变与 \texttt{im2col} 矩阵本质、ResNet 残差梯度高速公路的数学证明，以及 RNN/LSTM 门控对长程依赖的化解。
  \item \textbf{注意力机制与迷你 GPT}：证明 $\sqrt{d_k}$ 缩放与 RoPE 旋转位置编码的相对位置性质，完整手写因果自注意力、KV Cache 复杂度分析与可运行迷你 GPT 实测。
  \item \textbf{自监督预训练与深度生成}：打通 BPE、InfoNCE 互信息下界与 Chinchilla 缩放定律；从第 4 卷 EM 的 ELBO 无缝推导出 VAE 重参数化、DDPM 扩散去噪与流匹配（Flow Matching）。
\end{itemize}
\end{tcolorbox}

\vspace{4.5mm}

\noindent\textbf{\sffamily\color{tealDark}第 5 卷章节导览}
\vspace{1.5mm}

{\footnotesize\renewcommand{\arraystretch}{1.15}
\begin{tabularx}{\linewidth}{@{}p{15mm}p{47mm}X@{}}
\toprule
\textbf{章节} & \textbf{标题} & \textbf{核心理论、完整证明与实测实验} \\
\midrule
序\quad 章 & 本卷坐标与八校考点覆盖对照 & 从经典机器学习迈入深层表征、八校对标表与本卷知识依赖 DAG \\
第 1 章 & 多层感知机与反向模式自动微分 & ReLU 分段线性折叠、计算图与 VJP、60 行手写引擎、PyTorch 互验与非线性分类 \\
第 2 章 & 深层训练动力学：初始化与优化 & 梯度消失/爆炸、Kaiming 方差证明、BN/LN/RMSNorm、AdamW 与 20 层梯度实测 \\
第 3 章 & 卷积、残差网络与循环序列建模 & 2D 卷积等变性、ResNet 梯度高速公路证明、RNN/LSTM 门控与 24 层深度实测 \\
第 4 章 & 注意力机制与 Transformer 架构 & $\sqrt{d_k}$ 缩放证明、多头因果注意力、RoPE、KV Cache 与迷你 GPT 序列实测 \\
第 5 章 & 自监督预训练、对比学习与缩放律 & BPE 分词、BERT/GPT、InfoNCE 互信息下界证明、Chinchilla 缩放律与对比实测 \\
第 6 章 & 深度生成：VAE、扩散与流匹配 & 重参数化 VAE、DDPM 前向/后验严格推导、分数匹配、流匹配与 2D 扩散实测 \\
附\quad 录 & 全卷 48 道分层习题完整精解 & 第 1--6 章 A/B/C/D 四层习题分步推导、反例验算与代码自检基准 \\
\bottomrule
\end{tabularx}}

\vfill
\noindent{\color{tealLine}\rule{\linewidth}{0.5pt}}
{\footnotesize\color{slate}配套实验脚本 \texttt{code/vol5\_experiments.py} 与实测数据随本卷提供。上一卷：第 4 卷；下一卷：第 6 卷《强化学习、大模型后训练与 AI 系统》。}

% ==================== Table of Contents ====================
\clearpage
\tableofcontents

% ==================== Main Matter ====================
\mainmatter
\chapter*{序章：深层表征的崛起与八校考点对照}
\addcontentsline{toc}{chapter}{序章：深层表征的崛起与八校考点对照}
\markboth{序章\quad 本卷坐标与八校考点对照}{}

\chapEpigraph{你不必相信上帝，但你应当相信“天书”的存在。}{Paul Erdős（保罗·埃尔德什，谈论记录最美证明的 \emph{THE BOOK}）}

\begin{chapIntro}{本章定位与从第 4 卷的跨越}
在前四卷中，特征向量 $x \in \R^d$ 往往由人工设计，我们只在其上学习浅层决策边界或核映射。从本卷（第 5 卷）开始，我们进入现代深度学习的核心革命：\textbf{让模型通过端到端反向传播，自动从原始数据中学习多层次的内部表征（representation）}。

本卷拒绝“调包黑箱”，坚持将每一个标志性架构——从 60 行手写自动微分、Kaiming 方差守恒、ResNet 残差高速公路、因果自注意力迷你 GPT，到对比学习 InfoNCE 与 DDPM 扩散去噪——都拆解到第一性原理的数学推导与可在笔记本电脑上秒级复现的实测代码。
\end{chapIntro}

\section*{0.1\quad 从经典机器学习迈入深度学习的四个认知跃迁}
\addcontentsline{toc}{section}{0.1\quad 从经典机器学习迈入深度学习的四个认知跃迁}

\begin{enumerate}
  \item \textbf{从手工求导到计算图自动微分（第 1 章）}：面对千万乃至百亿参数，解析推导每个偏导既不现实也易出错；反向模式自动微分（VJP）以一次前向加一次反向的 $O(1)$ 倍常数开销，精确给出标量损失对全部参数的梯度。
  \item \textbf{从凸函数收敛到非凸训练动力学（第 2--3 章）}：深度网络损失面高度非凸，训练能否成功首先取决于信号与梯度能否稳定穿透数十层——这催生了 Kaiming 初始化、LayerNorm/RMSNorm、AdamW 以及改变深度学习历史的 ResNet 残差跳跃连接。
  \item \textbf{从局部归纳偏置到全局内容寻址（第 4--5 章）}：卷积利用平移等变、循环网络利用时间递推，而 Transformer 用自注意力让任意两个 token 以 $O(1)$ 路径长度直接交互，并与自监督预训练和缩放定律共同铸就了当代基础模型范式。
  \item \textbf{从判别分类到深度概率生成（第 6 章）}：将第 4 卷 EM 算法的 ELBO 下界与神经网络结合，便诞生了变分自编码器（VAE）；再将“一步隐变量映射”拆解为渐进加噪与反向去噪，便通向当今图像与视频生成的核心引擎——扩散模型（DDPM）与流匹配（Flow Matching）。
\end{enumerate}

\section*{0.2\quad 第 5 卷八校核心考点覆盖对照表}
\addcontentsline{toc}{section}{0.2\quad 第 5 卷八校核心考点覆盖对照表}

\begin{center}
\footnotesize
\begin{tabularx}{\linewidth}{@{}p{34mm}p{35mm}p{15mm}X@{}}
\toprule
\textbf{八校核心考点} & \textbf{主要对标课程} & \textbf{本卷位置} & \textbf{下游直接用途} \\
\midrule
MLP、计算图、VJP 与手写反向模式自动微分 & CMU 10-414 / Stanford CS231n & 第 1 章 & 彻底理解 PyTorch \texttt{autograd} 与自定义算子梯度 \\
梯度消失/爆炸、Kaiming 初始化、BN/LN 与 AdamW & CS231n / Berkeley CS182 & 第 2 章 & 深层网络稳定训练与大模型优化器配置诊断 \\
2D 卷积、\texttt{im2col}、ResNet 与 RNN/LSTM & CS231n / 南大深度学习 & 第 3 章 & 视觉骨干理解、残差流直觉与序列长程依赖建模 \\
缩放点积注意力、多头、RoPE、Transformer 与 KV Cache & Stanford CS224n / CS336 & 第 4 章 & 现代 LLM 架构核心、第 6 卷大模型后训练与推理优化 \\
BPE、BERT/GPT、InfoNCE 对比学习与缩放定律 & CS224n / CS336 / 智班 & 第 5 章 & 多模态对齐（CLIP）、算力配比与预训练范式理解 \\
VAE（ELBO+重参数化）、DDPM 扩散与流匹配 & CS231n / MIT 6.S184 & 第 6 章 & 现代视觉/音频/分子生成模型的统一数学框架 \\
\bottomrule
\end{tabularx}
\end{center}

\section*{0.3\quad 第 5 卷知识阶梯（零先修倒挂）}
\addcontentsline{toc}{section}{0.3\quad 第 5 卷知识阶梯（零先修倒挂）}

\begin{center}
\resizebox{0.98\linewidth}{!}{%
\begin{tikzpicture}[
  node distance=6.5mm and 7.5mm,
  ch/.style={
    rectangle, rounded corners=1.8mm, draw=teal, thick, fill=tealLight,
    text width=34mm, minimum height=14.5mm, align=center, font=\sffamily\footnotesize, text=ink, inner sep=4pt
  },
  goal/.style={
    rectangle, rounded corners=1.8mm, draw=ochre, very thick, fill=ochreLight,
    text width=37mm, minimum height=14.5mm, align=center, font=\sffamily\footnotesize, text=ink, inner sep=4pt
  },
  arr/.style={-{Stealth[length=2.2mm]}, thick, teal, preaction={draw=white, line width=2.5pt, arrows=-}}
]
  \node[ch] (c1) {\textbf{第 1 章}\\MLP 与反向模式\\自动微分（VJP）};
  \node[ch, right=of c1] (c2) {\textbf{第 2 章}\\训练动力学：初始化\\归一化与 AdamW};
  \node[ch, right=of c2] (c3) {\textbf{第 3 章}\\卷积、ResNet 残差\\与循环网络};

  \node[ch, below=11mm of c1] (c4) {\textbf{第 4 章}\\注意力、Transformer\\与手写迷你 GPT};
  \node[ch, below=11mm of c2] (c5) {\textbf{第 5 章}\\自监督预训练、\\InfoNCE 与缩放律};
  \node[goal, below=11mm of c3] (c6) {\textbf{\color{ochre}第 6 章（承接第 4 卷 EM）}\\深度生成：VAE、\\DDPM 与流匹配};

  \draw[arr] (c1) -- (c2);
  \draw[arr] (c2) -- (c3);
  \draw[arr] ([xshift=-9mm]c2.south) -- (c4.north);
  \draw[arr] ([xshift=-9mm]c3.south) -- ([xshift=9mm]c4.north);
  \draw[arr] (c4) -- (c5);
  \draw[arr] ([xshift=9mm]c2.south) -- ([xshift=-9mm]c6.north);
  \draw[arr] (c4.south) -- ++(0,-4.8mm) -| (c6.south);
\end{tikzpicture}}
\end{center}

\chapter[多层感知机与反向模式自动微分]{多层感知机与反向模式自动微分}
\label{ch:mlp-autodiff}

\chapEpigraph{凡是我不能亲手创造的，我就不算真正理解。}{Richard P. Feynman（理查德·费曼，黑板遗言，1988）}

\begin{chapIntro}{本章导读：打开深度学习的第一块基石}
在前面四卷中，我们的模型大多拥有显式而紧凑的数学结构：线性回归、逻辑回归、支持向量机、决策树与高斯混合模型。只要特征空间固定，我们就能用纸笔写出损失函数对每个参数的闭合梯度。然而，当问题变为“从像素识别猫狗”或“从上下文预测下一个词”时，人工设计线性可分特征几乎成为不可能的任务。

本章完成从“浅层线性模型”到“深度可微计算”的决定性一跃：
\begin{enumerate}
  \item \textbf{表达力之源}：分析线性模型为何在异或（XOR）与双月（Two-Moons）数据上必然折戟，揭示 ReLU 多层感知机（MLP）通过\textbf{分段线性空间折叠}实现万能逼近的构造性原理。
  \item \textbf{自动微分的核心引擎}：严格区分符号微分、数值差分与自动微分；从雅可比矩阵出发对比前向模式（JVP）与反向模式（VJP），证明为何反向模式仅用 $O(1)$ 倍前向计算代价即可一次求出标量损失对百万乃至百亿参数的全部梯度。
  \item \textbf{亲手构造与张量向量化}：逐节点手算一个非线性计算图的前向与反向过程，用约 60 行纯 Python 写出反向模式自动微分引擎并与 PyTorch \texttt{autograd} 互验至机器精度（误差 \texttt{0.00e+00}）；随后推导出矩阵仿射层、ReLU 与数值稳定融合交叉熵的张量 VJP，并在双月数据集上用纯 NumPy 实测验证。
\end{enumerate}
\end{chapIntro}

\section{线性模型的天花板与 ReLU 折叠几何}
\label{sec:mlp-geometry}

\subsection{为什么线性模型注定无法解开异或与双月}

一个二分类线性模型（如第 4 卷中的逻辑回归或线性感知机）对输入 $x \in \R^d$ 计算仿射得分 $z = w^\top x + b$，其决策边界是单张超平面
\[
  \mathcal{H} = \{x \in \R^d : w^\top x + b = 0\},
\]
它将整个输入空间 $\R^d$ 切成两个凸半空间 $\mathcal{H}^+ = \{w^\top x + b \ge 0\}$ 与 $\mathcal{H}^- = \{w^\top x + b < 0\}$。这意味着：\textbf{线性分类器所能表示的任何一类决策区域都必须是凸集}。

\begin{thmBox}
\begin{proposition}[异或问题（XOR）的线性不可分性]
\label{prop:xor-inseparable}
设二维数据集包含四个点：$x^{(1)} = (0,0)^\top, x^{(4)} = (1,1)^\top$ 属于负类（$y=0$），$x^{(2)} = (0,1)^\top, x^{(3)} = (1,0)^\top$ 属于正类（$y=1$）。不存在任何 $w \in \R^2$ 与 $b \in \R$，使得对正类样本严格有 $w^\top x + b > 0$ 而对负类样本严格有 $w^\top x + b < 0$。
\end{proposition}
\begin{proof}
反证法。假设存在 $(w_1, w_2, b) \in \R^3$ 满足上述严格分离条件，则四个点分别给出四个不等式：
\begin{align*}
  (0,0):\quad & b < 0, & (1,1):\quad & w_1 + w_2 + b < 0, \\
  (0,1):\quad & w_2 + b > 0, & (1,0):\quad & w_1 + b > 0.
\end{align*}
将前两式相加得 $w_1 + w_2 + 2b < 0$；将后两式相加得 $w_1 + w_2 + 2b > 0$，矛盾。从凸几何视角看，正类线段与负类线段在中心点 $(\frac{1}{2}, \frac{1}{2})$ 处相交，两类凸包不分离，故不可能被任何单一超平面分开。
\end{proof}
\end{thmBox}

Minsky 与 Papert（1969）在《感知机》一书中将这一简洁的凸几何事实推至台前，曾让单层神经网络陷入第一次低谷。在更贴近现实数据的\textbf{双月（Two-Moons）}分布中，两个半环形流形相互咬合，任何单一直线都无法绕过弯钩将两者完整切开（后文实验~\ref{sec:exp-moons} 中线性逻辑回归的准确率将被死死卡在 $89.17\%$）。

打破这一天花板的核心思想异常朴素：\textbf{在最终的线性分类器之前，先学习一个非线性坐标变换 $\phi_\theta: \R^d \to \R^m$，把原始空间中互相纠缠的流形“拉直、折叠”到新的表示空间，使其线性可分}。如果将这样的“仿射变换 + 逐元素非线性”复合 $L$ 层，就得到了\textbf{多层感知机（Multilayer Perceptron, MLP）}。

\begin{defBox}
\begin{definition}[多层感知机与激活函数]
设输入为 $h^{(0)} = x \in \R^{d_0}$。一个 $L$ 层全连接前馈网络（MLP）由权重矩阵 $W^{(l)} \in \R^{d_l \times d_{l-1}}$、偏置向量 $b^{(l)} \in \R^{d_l}$ 与逐元素非线性激活函数 $\sigma: \R \to \R$ 定义，逐层递推：
\begin{equation}
  z^{(l)} = W^{(l)} h^{(l-1)} + b^{(l)}, \qquad h^{(l)} = \sigma\!\left(z^{(l)}\right), \qquad l = 1, \dots, L-1,
  \label{eq:mlp-def}
\end{equation}
最终输出层通常为线性映射 $\hat{y} = z^{(L)} = W^{(L)} h^{(L-1)} + b^{(L)} \in \R^{d_L}$。若省略非线性（即令 $\sigma(z)=z$），则无论堆叠多少层，复合映射 $\hat{y} = W^{(L)} \cdots W^{(1)} x + \tilde{b}$ 仍退化为单层仿射变换，毫无表达力增益。
\end{definition}
\end{defBox}

现代深度学习中最基础、也最具代表性的激活函数是\textbf{修正线性单元（ReLU）}：
\[
  \ReLU(z) = \max(0, z).
\]
与早期常用的 $\mathrm{sigmoid}(z) = (1+e^{-z})^{-1}$ 或 $\tanh(z)$ 相比，ReLU 在正半轴导数恒为 $1$，不会因输入幅度增大而饱和（第~\ref{ch:train-dynamics} 章将定量展示其对梯度传播的巨大优势），且计算只需一次比较与置零。

\subsection{ReLU 的几何本质：空间折叠与分段线性区域}

理解 ReLU 网络最深刻的方式不是把它看成神秘的“生物神经元”，而是把它看成\textbf{折纸几何}。

考虑隐藏层中的第 $j$ 个神经元 $h_j = \ReLU(w_j^\top x + b_j)$。它在输入空间 $\R^d$ 中画出一张超平面 $w_j^\top x + b_j = 0$：超平面一侧的点保持仿射值不变，另一侧的点全部被压扁到 $0$。这就好比沿一条折痕将纸面的一侧折平。若第一隐藏层有 $m$ 个 ReLU 神经元，它们在 $\R^d$ 中画出 $m$ 张超平面，将空间划分为若干凸多面体区域；在每个区域内部，每个 ReLU 的开关状态（激活或关闭）是恒定的，因此整个网络在该区域内是一个\textbf{纯粹的线性函数}。

用手工构造就能瞬间解决异或问题：设隐藏层含两个 ReLU 神经元
\[
  h_1 = \ReLU(x_1 + x_2 - 0.5), \qquad h_2 = \ReLU(x_1 + x_2 - 1.5),
\]
再令输出层为 $z = h_1 - 2 h_2 - 0.25$。逐点代入四个角点即可验算：
\begin{itemize}
  \item $(0,0) \mapsto (h_1, h_2) = (0, 0) \mapsto z = -0.25 < 0$；
  \item $(0,1)$ 与 $(1,0) \mapsto (h_1, h_2) = (0.5, 0) \mapsto z = +0.25 > 0$；
  \item $(1,1) \mapsto (h_1, h_2) = (1.5, 0.5) \mapsto z = 1.5 - 1.0 - 0.25 = +0.25$？注意修正系数：若取 $z = h_1 - 2 h_2 - 0.25$，在 $(1,1)$ 处 $h_1 - 2h_2 = 0.5$ 仍为正！再检查一次：在 $(1,1)$ 处 $h_1 = 1.5, h_2 = 0.5$，我们需要 $z(1,1) < 0$ 且 $z(0,1) > 0$。令 $h_1 = \ReLU(x_1 + x_2 - 0.5)$、$h_2 = \ReLU(x_1 + x_2 - 1.0)$，则四个点映射为：
  \[
    (0,0)\mapsto (0,0),\quad (0,1),(1,0)\mapsto (0.5, 0),\quad (1,1)\mapsto (1.5, 1.0).
  \]
  取输出 $z = 2h_1 - 4h_2 - 0.5$：在 $(0,0)$ 处 $z = -0.5 < 0$；在 $(0,1),(1,0)$ 处 $z = 1.0 - 0 - 0.5 = +0.5 > 0$；在 $(1,1)$ 处 $z = 3.0 - 4.0 - 0.5 = -1.5 < 0$！两个负类点在隐藏空间中被折叠到了同一侧，单张直线 $2h_1 - 4h_2 = 0.5$ 便将 XOR 完美分开。
\end{itemize}

更令人惊叹的是\textbf{深度带来的指数级折叠效应}。根据超平面排列的 Zaslavsky 计数定理，$m$ 张超平面在 $\R^d$ 中最多分出 $\sum_{k=0}^{d} \binom{m}{k} = O(m^d)$ 个线性区域；因此宽度为 $m$ 的单隐藏层 ReLU 网络在线性区域数上随 $m$ 呈多项式增长。然而，当我们将 $L$ 层宽度为 $m$（$m \ge d$）的 ReLU 层串联时，每一层都能把上一层所有线性区域\emph{同步再折叠一次}，Montúfar 等人（2014）证明其最大线性区域数可达 $\Omega\!\left((m/d)^{(L-1)d} m^d\right)$——随深度 $L$ 呈\textbf{指数增长}！这从几何上解释了为何“深而窄”的现代架构往往比“浅而宽”的网络更高效地表示复杂函数。

\subsection{万能逼近定理的一维构造性直觉}

单隐藏层神经网络究竟能逼近多大范围的函数？Cybenko（1989）与 Hornik（1991）证明了著名的\textbf{万能逼近定理（Universal Approximation Theorem）}。与其抽象地援引 Hahn-Banach 定理，不如直接用 ReLU 搭出“三角形帐篷（帽函数）”，在十行之内亲眼看见它的构造性证明。

\begin{thmBox}
\begin{theorem}[一维 ReLU 双层网络的构造性万能逼近]
\label{thm:uat-1d}
设 $f: [0, 1] \to \R$ 是闭区间上的任意连续函数。对任意精度 $\epsilon > 0$，存在一个仅有单个隐藏层的 ReLU 神经网络
\[
  \hat{f}(x) = c_0 + \sum_{k=1}^{M} a_k\,\ReLU(x - t_k),
\]
使得 $\sup_{x \in [0, 1]} \abs{f(x) - \hat{f}(x)} < \epsilon$。
\end{theorem}
\begin{proof}
由于 $f$ 在紧集 $[0, 1]$ 上连续，故 $f$ 一致连续：存在步长 $h = 1/N$（$N \in \N^+$），使得只要 $\abs{x - x'} \le h$，就有 $\abs{f(x) - f(x')} < \epsilon/2$。记等距网格点 $x_k = k h$（$k = 0, 1, \dots, N$）。

\textbf{第一步（构造分段线性插值）}：在每个子区间 $[x_k, x_{k+1}]$ 上，用连接 $(x_k, f(x_k))$ 与 $(x_{k+1}, f(x_{k+1}))$ 的直线段构造连续分段线性函数 $g(x)$。对任意 $x \in [x_k, x_{k+1}]$，存在 $\lambda \in [0, 1]$ 使 $g(x) = (1-\lambda)f(x_k) + \lambda f(x_{k+1})$，因此
\[
  \abs{f(x) - g(x)} \le (1-\lambda)\abs{f(x) - f(x_k)} + \lambda \abs{f(x) - f(x_{k+1})} < \frac{\epsilon}{2} < \epsilon.
\]

\textbf{第二步（用 ReLU 组合精确表示 $g(x)$）}：利用标准三角帽基函数（每个帽函数由 3 个 ReLU 线性组合而成）：
\[
  \phi_k(x) = \frac{1}{h}\Big(\ReLU(x - x_{k-1}) - 2\,\ReLU(x - x_k) + \ReLU(x - x_{k+1})\Big),
\]
不难验证 $\phi_k(x_k) = 1$，在 $[x_{k-1}, x_k]$ 上线性上升、在 $[x_k, x_{k+1}]$ 上线性下降，且在区间 $[x_{k-1}, x_{k+1}]$ 外恒为 $0$。或者更紧凑地，利用斜率差分直接写出：设 $s_k = \frac{f(x_{k+1}) - f(x_k)}{h}$ 为第 $k$ 段斜率，则对任意 $x \in [0, 1]$，
\[
  g(x) = f(0) + s_0\,\ReLU(x - x_0) + \sum_{k=1}^{N-1} (s_k - s_{k-1})\,\ReLU(x - x_k).
\]
每当 $x$ 跨过折点 $x_k$ 时，新增的 ReLU 恰好给函数斜率补上跳变量 $s_k - s_{k-1}$，使总斜率变为 $s_k$；而在端点 $x=0$ 处 $g(0) = f(0)$。由此 $g(x)$ 就是一个含 $N$ 个隐藏 ReLU 神经元的单隐藏层网络，满足 $\sup_{x \in [0,1]}\abs{f(x) - g(x)} < \epsilon$。
\end{proof}
\end{thmBox}

\begin{insightBox}{万能逼近定理\emph{没有}承诺什么？}
定理~\ref{thm:uat-1d} 证明的只是\textbf{表示存在性（Expressivity）}，而不是\textbf{优化可达性}与\textbf{样本泛化性}。在高维 $\R^d$ 中，若强行用单隐藏层拼出 $d$ 维网格帐篷，所需神经元数 $N^d$ 会遭遇维度灾难；更重要的是，即使网络架构足以表达目标函数，我们还需要一种高效算法来求得梯度、并在非凸损失面上把参数优化到位——这就是本章自动微分与第~\ref{ch:train-dynamics} 章训练动力学的使命。
\end{insightBox}

\section{三种求导范式与计算图抽象}
\label{sec:three-diff}

训练一个拥有 $P$ 个参数的神经网络 $\theta \in \R^P$，核心瓶颈是在每一次迭代中计算标量损失 $\mathcal{L}(\theta) \in \R$ 对全体参数的梯度向量
\[
  \nabla_\theta \mathcal{L}(\theta) = \left(\frac{\partial \mathcal{L}}{\partial \theta_1}, \frac{\partial \mathcal{L}}{\partial \theta_2}, \dots, \frac{\partial \mathcal{L}}{\partial \theta_P}\right)^\top \in \R^P.
\]
在现代大模型中，$P$ 动辄达到 $10^7 \sim 10^{11}$。计算机科学历史上先后出现过三种计算导数的思路，它们的复杂度与精度天差地别：

\begin{enumerate}
  \item \textbf{数值微分（Numerical Differentiation）}：利用中心有限差分公式逐坐标近似
  \[
    \frac{\partial \mathcal{L}}{\partial \theta_i} \approx \frac{\mathcal{L}(\theta + h e_i) - \mathcal{L}(\theta - h e_i)}{2h}.
  \]
  它有两个致命缺陷：第一，每个参数都至少需要额外执行 $1\sim 2$ 次完整前向传播，计算全部 $P$ 个偏导需要 $O(P)$ 次前向计算，对百万参数网络完全不可行；第二，步长 $h$ 面临\textbf{截断误差} $O(h^2)$ 与浮点\textbf{舍入误差} $O(\epsilon_{\text{mach}}/h)$ 的两难权衡（在双精度下 $h \approx 10^{-5}\sim 10^{-6}$ 时最优相对误差仅约 $10^{-11}$）。因此数值微分在深度学习中绝不用于训练，但它是检验手写反向传播代码的黄金工具——\textbf{梯度检验（Gradient Checking）}。

  \item \textbf{符号微分（Symbolic Differentiation）}：像 SymPy 或 Mathematica 一样，对数学表达式树反复套用乘积法则与链式法则，输出封闭的符号公式。问题在于\textbf{表达式膨胀（Expression Swell）}：考虑嵌套复合函数 $f(x) = g_L(g_{L-1}(\cdots g_1(x)\cdots))$ 或带有共享子表达式的计算，如果不显式复用中间变量，展开后的符号导数长度可能随网络深度呈指数级爆炸，且难以直接处理程序中的循环与条件分支。

  \item \textbf{自动微分（Automatic Differentiation, AD）}：它既不是数值近似，也不是整式符号展开，而是\textbf{沿着程序实际执行的初等算子序列，在具体数值点上按链式法则精确传递导数}。每一个初等算子（加、乘、$\exp$、$\tanh$、矩阵乘）的局部导数早已写成固定规则，中间变量被缓存复用，既彻底消除表达式膨胀，又达到机器浮点精度（无截断误差）。
\end{enumerate}

\begin{defBox}
\begin{definition}[计算图（Computational DAG）]
任何由初等可微算子复合而成的程序执行轨迹，都可以表示为一个\textbf{有向无环图（DAG）} $\mathcal{G} = (\mathcal{V}, \mathcal{E})$，其顶点按拓扑序编号为 $v_1, v_2, \dots, v_K$：
\begin{itemize}
  \item 源节点 $v_1, \dots, v_P$ 为输入变量（参数或数据）；
  \item 每个中间节点 $v_k$ 由其直接前驱集合 $\mathrm{pa}(k) \subset \{1, \dots, k-1\}$ 经由一个初等可微算子 $\phi_k$ 算出：$v_k = \phi_k(v_{\mathrm{pa}(k)})$；
  \item 汇节点 $v_K = \mathcal{L}$ 为最终输出（在训练中为标量损失）。
\end{itemize}
对任意有向边 $(i \to k) \in \mathcal{E}$，定义其\textbf{局部导数（局部雅可比）}为 $J_{k,i} = \frac{\partial \phi_k}{\partial v_i}$。
\end{definition}
\end{defBox}

按拓扑序从 $v_1$ 计算到 $v_K$ 的过程称为\textbf{前向传播（Forward Pass）}。自动微分的全部奥秘在于：利用同一个计算图 $\mathcal{G}$，我们既可以顺着拓扑序传递导数（前向模式），也可以逆着拓扑序回传导数（反向模式）。

\section{前向模式（JVP）与反向模式（VJP）的数学本质}
\label{sec:jvp-vjp}

设复合函数 $f = f_L \circ f_{L-1} \circ \cdots \circ f_1 : \R^{n_0} \to \R^{n_L}$，其中第 $l$ 步输出 $v_l = f_l(v_{l-1}) \in \R^{n_l}$，对应的局部雅可比矩阵为 $J_l = \frac{\partial v_l}{\partial v_{l-1}} \in \R^{n_l \times n_{l-1}}$。根据多元微积分链式法则，整体雅可比矩阵是各层局部雅可比的连乘：
\begin{equation}
  J_f = \frac{\partial v_L}{\partial v_0} = J_L J_{L-1} \cdots J_2 J_1 \in \R^{n_L \times n_0}.
  \label{eq:chain-matrix}
\end{equation}
矩阵乘法满足结合律，但\textbf{不同的加括号顺序会带来天壤之别的计算复杂度}！

\subsection{前向模式自动微分：雅可比-向量积（JVP）}

若从右向左结合式~\eqref{eq:chain-matrix}，并给定一个输入扰动方向（切向量）$\dot{v}_0 \in \R^{n_0}$，定义每个节点的\textbf{切线值（tangent）}为 $\dot{v}_k = \frac{\partial v_k}{\partial v_0} \dot{v}_0$。在拓扑序遍历中，前向模式同步推进原值与切线值：
\begin{equation}
  \dot{v}_k = \sum_{i \in \mathrm{pa}(k)} \frac{\partial \phi_k}{\partial v_i}\,\dot{v}_i = \operatorname{JVP}_{\phi_k}(v_{\mathrm{pa}(k)}, \dot{v}_{\mathrm{pa}(k)}).
  \label{eq:jvp}
\end{equation}
如果令 $\dot{v}_0 = e_j$（第 $j$ 个标准基向量），一次前向模式传播便算出整体雅可比矩阵 $J_f$ 的\textbf{第 $j$ 列} $\frac{\partial v_L}{\partial (v_0)_j} \in \R^{n_L}$。因此，要算齐 $J_f \in \R^{n_L \times n_0}$ 的全部列，前向模式需要运行 $n_0$ 次！在数学上，前向模式等价于用\textbf{对偶数（Dual Numbers）} $a + b\varepsilon$（满足 $\varepsilon^2 = 0$）重载算术运算，因为泰勒展开立刻给出 $f(a + b\varepsilon) = f(a) + f'(a)b\,\varepsilon$。

\subsection{反向模式自动微分：向量-雅可比积（VJP）}

现在改从左向右结合式~\eqref{eq:chain-matrix}。给定一个输出端余切向量（对标量损失 $\mathcal{L} \in \R$ 而言，它就是标量 $1$）$\bar{v}_L \in \R^{n_L}$，为计算图中的每个节点 $v_i$ 定义其\textbf{伴随量（adjoint）}：
\[
  \bar{v}_i := \left(\frac{\partial v_L}{\partial v_i}\right)^\top \bar{v}_L \in \R^{n_i}.
\]
当 $v_L = \mathcal{L} \in \R$ 且置 $\bar{v}_L = 1$ 时，$\bar{v}_i = \nabla_{v_i} \mathcal{L}$ 恰好就是损失函数对节点 $v_i$ 的完整梯度！

如何逆着拓扑序算出所有节点的伴随量？由多元链式法则，节点 $v_i$ 对最终输出 $\mathcal{L}$ 的全部影响，必须流经它的所有直接后继节点 $\mathrm{ch}(i) = \{k : (i \to k) \in \mathcal{E}\}$：
\begin{equation}
  \bar{v}_i = \sum_{k \in \mathrm{ch}(i)} \left(\frac{\partial \phi_k}{\partial v_i}\right)^\top \bar{v}_k = \sum_{k \in \mathrm{ch}(i)} \operatorname{VJP}_{\phi_k}(v_{\mathrm{pa}(k)}, \bar{v}_k)_i.
  \label{eq:vjp}
\end{equation}
式~\eqref{eq:vjp} 有一个至关重要的工程含义：\textbf{如果一个变量在计算图中被多处复用（出度 $\abs{\mathrm{ch}(i)} > 1$），来自各个分支的反向梯度必须累加（\texttt{+=}），绝不能覆盖（\texttt{=}）}！

\begin{thmBox}
\begin{theorem}[反向模式自动微分的复杂度优势（廉价梯度原理）]
\label{thm:cheap-gradient}
设标量损失函数 $\mathcal{L}: \R^P \to \R$ 由计算图 $\mathcal{G} = (\mathcal{V}, \mathcal{E})$ 计算，前向求值耗时记为 $\operatorname{Cost}(\mathcal{L})$。假设每个初等算子 $\phi_k$ 最多有常数个输入，且单次局部 VJP 的算术代价不超过其前向代价的常数倍 $c$（对加、减、乘、除及常见初等函数，$c \le 3$）。则：
\begin{enumerate}
  \item \textbf{反向模式}仅需\textbf{一次前向传播 + 一次反向遍历}，即可同时精确求出全部 $P$ 个偏导数 $\nabla_\theta \mathcal{L} \in \R^P$，总算术代价满足
  \[
    \operatorname{Cost}\!\left(\mathcal{L}, \nabla_\theta \mathcal{L}\right) \le (1 + c)\,\operatorname{Cost}(\mathcal{L}) \le 4\,\operatorname{Cost}(\mathcal{L}),
  \]
  \textbf{与参数量 $P$ 完全无关}（这正是 Baur-Strassen 定理的现代自动微分表述）。
  \item \textbf{前向模式}若要计算完整梯度 $\nabla_\theta \mathcal{L}$，需要分别以 $e_1, \dots, e_P$ 为切向量运行 $P$ 趟，总代价为 $\Theta(P \cdot \operatorname{Cost}(\mathcal{L}))$。
\end{enumerate}
一般地，对向量映射 $f: \R^{n_0} \to \R^{n_L}$，前向模式计算完整雅可比需 $n_0$ 趟，反向模式需 $n_L$ 趟。因此：当 $n_0 \gg n_L$（深度学习训练的常态：$P \sim 10^9, n_L = 1$）时，反向模式占压倒性优势；当 $n_0 \ll n_L$（如偏微分方程灵敏度分析或计算雅可比方向导数）时，前向模式反而更优。
\end{theorem}
\begin{proof}
在反向模式中，前向传播按拓扑序访问每个节点 $v_k$ 一次，耗时 $\operatorname{Cost}(\mathcal{L})$，并将中间数值缓存于内存（称为\textbf{磁带 / tape}）。反向传播严格按逆拓扑序访问每个节点 $v_k$ 一次：访问 $v_k$ 时，其所有后继节点均已处理完毕，故 $\bar{v}_k$ 已经累加完整；随后只需为每个 $i \in \mathrm{pa}(k)$ 计算局部 VJP 项 $(\partial \phi_k / \partial v_i)^\top \bar{v}_k$ 并累加到 $\bar{v}_i$。每条边 $(i \to k)$ 在全过程中恰好被正向和反向各遍历一次。对基本算术算子，乘法 $u \cdot v$ 前向需 1 次乘法，反向需 2 次乘法与 2 次累加；加法前向 1 次加法，反向仅 2 次累加。综上，反向遍历总算术运算数不超过 $c \cdot \operatorname{Cost}(\mathcal{L})$（对标准算术图 $c \le 3$，总计 $\le 4\times$），且一次遍历结束时所有源节点 $v_1, \dots, v_P$ 的伴随量 $\bar{v}_1, \dots, \bar{v}_P$ 已全部就位。
\end{proof}
\end{thmBox}

\begin{center}
\small
\begin{tabularx}{\linewidth}{@{}p{26mm}p{26mm}p{24mm}p{26mm}X@{}}
\toprule
\textbf{求导方式} & \textbf{算力代价（$\R^P \to \R$）} & \textbf{数值精度} & \textbf{显存需求} & \textbf{典型用途} \\
\midrule
数值有限差分 & $O(P) \times$ 前向 & 截断+舍入误差 & $O(1)$ & 单元测试梯度检验 \\
符号微分 & 易指数膨胀 & 解析精确 & 表达式易爆炸 & 小型公式化简 \\
前向模式（JVP） & $O(P) \times$ 前向 & 机器精度 & $O(1)$（无需存激活） & $n_0 \ll n_L$、Hessian-向量积 \\
反向模式（VJP） & $\le 4 \times$ 前向 & 机器精度 & $O(\text{前向激活数})$ & \textbf{深度学习核心训练引擎} \\
\bottomrule
\end{tabularx}
\end{center}

\section{手算计算图与 60 行自动微分引擎}
\label{sec:micrograd}

\subsection{一个非线性计算图的逐节点手工推演}

为了彻底消除对“反向传播”的神秘感，我们严格按照配套脚本 \texttt{code/vol5\_experiments.py} 中的精确数值，手算一个单神经元均方损失计算图的前向与反向过程：
\[
  y = \tanh(x_1 w_1 + x_2 w_2 + b), \qquad L = (y - 0.7)^2,
\]
取输入与参数为 $x_1 = 2.0,\; x_2 = -1.5,\; w_1 = -0.8,\; w_2 = 1.2,\; b = 0.5$。

\textbf{第一步：按拓扑序执行前向传播}（保留 6 位小数）：
\begin{align*}
  u_1 &= x_1 w_1 = 2.0 \times (-0.8) = -1.600000, \\
  u_2 &= x_2 w_2 = (-1.5) \times 1.2 = -1.800000, \\
  s   &= u_1 + u_2 = -1.600000 + (-1.800000) = -3.400000, \\
  z   &= s + b = -3.400000 + 0.500000 = -2.900000, \\
  y   &= \tanh(-2.900000) = -0.993963, \\
  r   &= y - 0.7 = -1.693963, \\
  L   &= r \cdot r = (-1.693963)^2 = 2.869511.
\end{align*}

\textbf{第二步：按逆拓扑序回传伴随量 $\bar{v} = \partial L / \partial v$}：
\begin{enumerate}
  \item 终点初始化：$\bar{L} = 1.000000$。
  \item 平方节点 $L = r \cdot r$（注意变量 $r$ 在乘法中出现两次，两条边梯度累加！）：
  \[
    \bar{r} = r \cdot \bar{L} + r \cdot \bar{L} = 2r \bar{L} = 2 \times (-1.693963) = -3.387926.
  \]
  \item 减常数节点 $r = y + (-0.7)$：$\bar{y} = 1 \cdot \bar{r} = -3.387926$。
  \item 双曲正切节点 $y = \tanh(z)$，局部导数为 $\frac{dy}{dz} = 1 - y^2 = 1 - (-0.993963)^2 = 0.012037$：
  \[
    \bar{z} = (1 - y^2)\,\bar{y} = 0.012037 \times (-3.387926) = -0.040781.
  \]
  （顺便留意：由于 $z = -2.9$ 已落入 $\tanh$ 饱和区，$1-y^2 \approx 0.012$ 把上游梯度瞬间衰减了将近两个数量级——这正是第~\ref{ch:train-dynamics} 章讨论梯度消失的微型缩影！）
  \item 加法节点 $z = s + b$ 与 $s = u_1 + u_2$（加法节点像“梯度分流器”，将上游伴随量原样复制给每个输入）：
  \[
    \bar{s} = \bar{b} = \bar{z} = -0.040781, \qquad \bar{u}_1 = \bar{u}_2 = \bar{s} = -0.040781.
  \]
  \item 乘法节点 $u_1 = x_1 w_1$ 与 $u_2 = x_2 w_2$（乘法节点像“交叉交换器”，将上游伴随量乘以\emph{另一个}操作数的前向值）：
  \begin{align*}
    \bar{w}_1 &= x_1 \bar{u}_1 = 2.0 \times (-0.040781) = -0.081562, &
    \bar{x}_1 &= w_1 \bar{u}_1 = (-0.8) \times (-0.040781) = +0.032625, \\
    \bar{w}_2 &= x_2 \bar{u}_2 = (-1.5) \times (-0.040781) = +0.061172, &
    \bar{x}_2 &= w_2 \bar{u}_2 = 1.2 \times (-0.040781) = -0.048937.
  \end{align*}
\end{enumerate}

图~\ref{fig:comp-graph} 完整绘制了该计算图的前向数据流（青色实线）与反向伴随流（赭色虚线）。读者只需沿着虚线箭头把每条边上的局部导数乘起来，便能直观看到链式法则如何在图上自动展开。

\begin{figure}[htbp]
\centering
\resizebox{\linewidth}{!}{%
\begin{tikzpicture}[
  node distance=10mm and 12mm,
  var/.style={rectangle, rounded corners=1.5mm, draw=navy, thick, fill=navyLight,
              minimum width=23mm, minimum height=11mm, align=center, font=\scriptsize, text=ink},
  op/.style={rectangle, rounded corners=2mm, draw=teal, thick, fill=tealLight,
             minimum width=24mm, minimum height=12mm, align=center, font=\scriptsize, text=ink},
  loss/.style={rectangle, rounded corners=2mm, draw=ochre, very thick, fill=ochreLight,
               minimum width=23mm, minimum height=12mm, align=center, font=\scriptsize, text=ink},
  fwd/.style={-{Stealth[length=2.2mm]}, thick, teal},
  bwd/.style={-{Stealth[length=2.2mm]}, thick, ochre, dashed}
]
  % Input leaf nodes
  \node[var] (x1) {$x_1 = 2.0$\\[-1pt]\textcolor{ochre}{$\bar{x}_1 = +0.0326$}};
  \node[var, below=4mm of x1] (w1) {$w_1 = -0.8$\\[-1pt]\textcolor{ochre}{$\bar{w}_1 = -0.0816$}};
  \node[var, below=6mm of w1] (x2) {$x_2 = -1.5$\\[-1pt]\textcolor{ochre}{$\bar{x}_2 = -0.0489$}};
  \node[var, below=4mm of x2] (w2) {$w_2 = 1.2$\\[-1pt]\textcolor{ochre}{$\bar{w}_2 = +0.0612$}};
  \node[var, below=6mm of w2] (b)  {$b = 0.5$\\[-1pt]\textcolor{ochre}{$\bar{b} = -0.0408$}};

  % Mul nodes
  \node[op, right=21mm of $(x1)!0.5!(w1)$] (u1) {$u_1 = x_1 w_1$\\$=-1.6000$\\[-1pt]\textcolor{ochre}{$\bar{u}_1 = -0.0408$}};
  \node[op, right=21mm of $(x2)!0.5!(w2)$] (u2) {$u_2 = x_2 w_2$\\$=-1.8000$\\[-1pt]\textcolor{ochre}{$\bar{u}_2 = -0.0408$}};

  % Add s
  \node[op, right=24mm of $(u1)!0.5!(u2)$] (s) {$s = u_1 + u_2$\\$=-3.4000$\\[-1pt]\textcolor{ochre}{$\bar{s} = -0.0408$}};

  % Add z
  \node[op, right=11mm of s] (z) {$z = s + b$\\$=-2.9000$\\[-1pt]\textcolor{ochre}{$\bar{z} = -0.0408$}};

  % Tanh y
  \node[op, right=11mm of z] (y) {$y = \tanh(z)$\\$=-0.9940$\\[-1pt]\textcolor{ochre}{$\bar{y} = -3.3879$}};

  % Loss L
  \node[loss, right=12mm of y] (L) {$L=(y-0.7)^2$\\$=2.8695$\\[-1pt]\textcolor{ochre}{$\bar{L} = 1.0000$}};

  % Forward arrows
  \draw[fwd] (x1.east) -- (u1.west);
  \draw[fwd] (w1.east) -- (u1.west);
  \draw[fwd] (x2.east) -- (u2.west);
  \draw[fwd] (w2.east) -- (u2.west);
  \draw[fwd] (u1.east) -- (s.west);
  \draw[fwd] (u2.east) -- (s.west);
  \draw[fwd] (s.east) -- (z.west);
  \draw[fwd] (b.east) -| (z.south);
  \draw[fwd] (z.east) -- (y.west);
  \draw[fwd] (y.east) -- (L.west);

  % Backward annotations along main spine (cleanly routed outside all boxes)
  \draw[bwd] (L.north) to[bend right=26] node[above, font=\tiny, text=ochre] {$\times\,2(y-0.7)$} (y.north);
  \draw[bwd] (y.north) to[bend right=26] node[above, font=\tiny, text=ochre] {$\times\,(1-y^2)=0.0120$} (z.north);
  \draw[bwd] (z.north) to[bend right=26] node[above, font=\tiny, text=ochre] {$\times\,1$} (s.north);
  \draw[bwd] (s.north) to[out=130, in=25] node[above right, font=\tiny, text=ochre, inner sep=1pt] {$\times\,1$} (u1.east);
  \draw[bwd] (u1.south) to[out=-115, in=-25] node[below=2pt, font=\tiny, text=ochre, fill=white, inner sep=0.8pt] {$\times\,x_1=2.0$} (w1.south east);

  % Legend
  \node[anchor=south west, font=\scriptsize, text=slate] at ($(b.south east)+(16mm,-2mm)$)
    {\textcolor{teal}{\textbf{—— 前向传播（按拓扑序求值并缓存）}} \qquad \textcolor{ochre}{\textbf{- - - 反向模式 VJP（逆拓扑序累加伴随量）}}};
\end{tikzpicture}}
\caption{标量计算图的前向求值与反向模式伴随量传播。每个节点上方为前向数值，下方赭色为精确梯度 $\bar{v} = \partial L/\partial v$。}
\label{fig:comp-graph}
\end{figure}

\subsection{从零编写 60 行反向模式自动微分引擎}

下面这段不足 60 行的纯 Python 代码（取自配套实验脚本 \texttt{vol5\_experiments.py} 中的标量自动微分实验 \texttt{exp1\_micrograd\_and\_mlp}）就是 PyTorch \texttt{autograd} 与 JAX \texttt{grad} 的极简原型。它通过运算符重载动态构建计算图（Define-by-Run），并在 \texttt{backward()} 中执行深度优先拓扑排序，再逆序调用每个节点注册的局部 VJP 闭包 \texttt{\_backward}：

\begin{lstlisting}[caption={极简反向模式自动微分引擎 \texttt{Value} 与 PyTorch 精度互验}]
import math, torch

class Value:
    def __init__(self, data, _children=()):
        self.data = float(data)
        self.grad = 0.0
        self._backward = lambda: None
        self._prev = set(_children)

    def __add__(self, other):
        other = other if isinstance(other, Value) else Value(other)
        out = Value(self.data + other.data, (self, other))
        def _b():
            self.grad += out.grad       # 多分支复用必须用 += 累加
            other.grad += out.grad
        out._backward = _b
        return out

    def __mul__(self, other):
        other = other if isinstance(other, Value) else Value(other)
        out = Value(self.data * other.data, (self, other))
        def _b():
            self.grad += other.data * out.grad
            other.grad += self.data * out.grad
        out._backward = _b
        return out

    def tanh(self):
        t = math.tanh(self.data)
        out = Value(t, (self,))
        def _b():
            self.grad += (1.0 - t * t) * out.grad
        out._backward = _b
        return out

    def backward(self):
        topo, visited = [], set()
        def build(v):
            if v not in visited:
                visited.add(v)
                for child in v._prev:
                    build(child)
                topo.append(v)          # 后序压栈 => topo 为合法拓扑序
        build(self)
        self.grad = 1.0
        for node in reversed(topo):     # 逆拓扑序调用局部 VJP
            node._backward()
\end{lstlisting}

运行 \texttt{code/vol5\_experiments.py} 中的实测校验，将上述 \texttt{Value} 引擎与双精度 PyTorch \texttt{autograd} 对比，输出为：
\begin{itemize}
  \item $\partial L / \partial w_1 = -0.081562$，$\partial L / \partial w_2 = +0.061172$；
  \item 所有叶节点梯度的最大绝对误差 \texttt{max\_grad\_diff = 0.00e+00}——与工业级框架 PyTorch 达到\textbf{比特级一致}！
\end{itemize}
留意代码第 43 行的精妙之处：在表达式 \texttt{(y + Value(-0.7)) * (y + Value(-0.7))} 中，同一个临时节点 \texttt{r} 作为左右两个孩子传入 \texttt{\_\_mul\_\_}，其 \texttt{self.grad +=} 与 \texttt{other.grad +=} 会对同一对象执行两次累加，自动得到 $2r\bar{L}$。这正是为什么式~\eqref{eq:vjp} 要求累加而非赋值。

\section{矩阵与张量算子的 VJP 推导}
\label{sec:tensor-vjp}

标量级 \texttt{Value} 清晰揭示了原理，但现实中把百万次乘法拆成百万个 Python 对象会带来巨大的解释器开销。工业深度学习框架（PyTorch/JAX/Tinygrad）的计算图节点是\textbf{多维张量（Tensor）}，局部 VJP 直接调用底层 BLAS/CUDA 矩阵乘法。本节推导神经网络最核心的三大张量算子 VJP。

\begin{recallBox}{矩阵微分的“同形伴随”原则}
为避免四阶张量下标混乱，我们始终保持\textbf{梯度与原变量同形}：对任意矩阵 $A \in \R^{m \times n}$ 和标量损失 $\mathcal{L}$，定义 $\bar{A} := \frac{\partial \mathcal{L}}{\partial A} \in \R^{m \times n}$，其第 $(i,j)$ 元素为 $\frac{\partial \mathcal{L}}{\partial A_{ij}}$。利用 Frobenius 内积 $\inner{U}{V}_F = \tr(U^\top V) = \sum_{i,j} U_{ij} V_{ij}$，一阶全微分恒有
\[
  d\mathcal{L} = \inner{\bar{A}}{dA}_F = \tr(\bar{A}^\top dA).
\]
只要把 $d\mathcal{L} = \inner{\bar{Y}}{dY}_F$ 化为 $\inner{\text{某矩阵}}{dX}_F$，该矩阵就必为 $\bar{X}$——极其干净且绝不会转置出错！
\end{recallBox}

\subsection{批量仿射层（Linear Layer）：$Y = XW + \mathbf{1}_N b^\top$}

设小批量输入 $X \in \R^{N \times d_{\text{in}}}$，权重 $W \in \R^{d_{\text{in}} \times d_{\text{out}}}$，偏置 $b \in \R^{d_{\text{out}}}$，输出 $Y = XW + \mathbf{1}_N b^\top \in \R^{N \times d_{\text{out}}}$（其中 $\mathbf{1}_N \in \R^N$ 为全 1 列向量，对应广播加法）。已知上游传入的梯度 $\bar{Y} \in \R^{N \times d_{\text{out}}}$。

\begin{thmBox}
\begin{proposition}[批量仿射层的张量 VJP]
\label{prop:linear-vjp}
对于 $Y = XW + \mathbf{1}_N b^\top$，关于 $X, W, b$ 的反向模式 VJP 为：
\begin{equation}
  \bar{X} = \bar{Y} W^\top \in \R^{N \times d_{\text{in}}}, \qquad
  \bar{W} = X^\top \bar{Y} \in \R^{d_{\text{in}} \times d_{\text{out}}}, \qquad
  \bar{b} = \bar{Y}^\top \mathbf{1}_N = \sum_{n=1}^{N} \bar{Y}_{n,:}^\top \in \R^{d_{\text{out}}}.
  \label{eq:linear-vjp}
\end{equation}
\end{proposition}
\begin{proof}
对 $Y = XW + \mathbf{1}_N b^\top$ 取全微分得 $dY = (dX) W + X (dW) + \mathbf{1}_N (db)^\top$。代入 Frobenius 迹恒等式，并利用迹的循环不变性 $\tr(ABC) = \tr(CAB)$：
\begin{align*}
  d\mathcal{L} = \tr(\bar{Y}^\top dY)
  &= \tr\!\left(\bar{Y}^\top (dX) W\right) + \tr\!\left(\bar{Y}^\top X (dW)\right) + \tr\!\left(\bar{Y}^\top \mathbf{1}_N (db)^\top\right) \\
  &= \tr\!\left((W \bar{Y}^\top) dX\right) + \tr\!\left((\bar{Y}^\top X) dW\right) + \tr\!\left((\mathbf{1}_N^\top \bar{Y}) db\right) \\
  &= \tr\!\left((\bar{Y} W^\top)^\top dX\right) + \tr\!\left((X^\top \bar{Y})^\top dW\right) + \inner{\bar{Y}^\top \mathbf{1}_N}{db}.
\end{align*}
逐项比对系数即得式~\eqref{eq:linear-vjp}。
\end{proof}
\end{thmBox}

一个极实用的心算验证法是\textbf{维度匹配}：$\bar{X}$ 必须与 $X$ 同形（$N \times d_{\text{in}}$），由 $\bar{Y}$ ($N \times d_{\text{out}}$) 和 $W$ ($d_{\text{in}} \times d_{\text{out}}$) 拼出 $N \times d_{\text{in}}$ 的唯一合法乘积就是 $\bar{Y} W^\top$；同理拼出 $\bar{W}$ ($d_{\text{in}} \times d_{\text{out}}$) 的唯一乘积就是 $X^\top \bar{Y}$；而前向传播把 $b$ 沿批量维\emph{广播复制}了 $N$ 份，反向传播则必须沿批量维\textbf{求和归约}。

\subsection{逐元素激活层：ReLU 的次梯度 VJP}

对于 $H = \ReLU(Z) = \max(0, Z)$，由于算子逐元素独立作用，给定上游梯度 $\bar{H}$，有
\begin{equation}
  \bar{Z} = \bar{H} \had \mathbb{I}[Z > 0],
  \label{eq:relu-vjp}
\end{equation}
其中 $\had$ 表示 Hadamard 逐元素乘积，$\mathbb{I}[Z > 0]$ 是与 $Z$ 同形的 $0/1$ 掩码矩阵。尽管 $\ReLU$ 在原点 $z = 0$ 处不可导，其 Clarke 次微分集合为 $[0, 1]$，实践中框架按约定取 $0$（在连续浮点分布下精确命中 $0$ 的概率几乎为零，且次梯度下降理论依然适用）。注意：为了在反向传播时计算式~\eqref{eq:relu-vjp}，前向传播必须缓存 $Z > 0$ 的布尔掩码（或缓存 $H$），这正是反向模式显存占用的直接来源。

\subsection{数值稳定的融合 LogSumExp 与交叉熵 VJP}

在多分类任务中，网络最后一层输出未归一化的逻辑得分（logits）$z \in \R^C$，经由 Softmax 变为概率分布 $p_k = \frac{e^{z_k}}{\sum_{j=1}^C e^{z_j}}$，再与独热标签 $y \in \{0, 1\}^C$（$\sum_k y_k = 1$）计算交叉熵损失 $\ell(z, y) = -\sum_{k=1}^C y_k \log p_k$。

若天真地先算 \texttt{exp(z)} 再取 \texttt{log}，当 $z_k = 100$ 时，单精度浮点 \texttt{exp(100)} 直接上溢成 \texttt{inf}，导致 \texttt{inf / inf = NaN}。解决方法是利用平移不变性：令 $m = \max_j z_j$，则
\begin{equation}
  \log \sum_{j=1}^C e^{z_j} = m + \log \sum_{j=1}^C e^{z_j - m}.
  \label{eq:lse-trick}
\end{equation}
由于 $z_j - m \le 0$，所有指数项严格落在 $(0, 1]$ 内，且至少有一项等于 $1$，彻底杜绝上溢与 $\log(0)$ 下溢！将 Softmax 与交叉熵\textbf{融合（Fuse）}后，单样本损失写为
\[
  \ell(z, y) = -\sum_{k=1}^C y_k \left(z_k - \log \sum_{j=1}^C e^{z_j}\right) = \log \sum_{j=1}^C e^{z_j} - \sum_{k=1}^C y_k z_k.
\]

\begin{thmBox}
\begin{proposition}[融合 Softmax + 交叉熵的简洁 VJP]
\label{prop:ce-vjp}
对批量大小 $N$、样本 $n = 1, \dots, N$ 的平均交叉熵损失 $\mathcal{L} = \frac{1}{N}\sum_{n=1}^N \ell(z^{(n)}, y^{(n)})$，其关于 logits 矩阵 $Z \in \R^{N \times C}$ 的梯度为：
\begin{equation}
  \frac{\partial \mathcal{L}}{\partial Z} = \frac{1}{N}(P - Y), \qquad P_{n,k} = \softmax(z^{(n)})_k.
  \label{eq:ce-vjp}
\end{equation}
\end{proposition}
\begin{proof}
对单个样本 $z \in \R^C$ 求第 $k$ 个分量的偏导：
\[
  \frac{\partial \ell}{\partial z_k} = \frac{\partial}{\partial z_k}\left(\log \sum_{j=1}^C e^{z_j}\right) - y_k = \frac{e^{z_k}}{\sum_{j=1}^C e^{z_j}} - y_k = p_k - y_k.
\]
对批量中 $N$ 个样本取算术平均，即得 $\bar{Z} = \frac{1}{N}(P - Y)$。二分类 Logistic 损失 $\ell(z, y) = \log(1 + e^z) - y z$ 是 $C=2$ 时的特例，其梯度同样是极其优雅的 $\frac{1}{N}(\sigma(z) - y)$。
\end{proof}
\end{thmBox}

式~\eqref{eq:ce-vjp} 是深度学习中最赏心悦目的抵消之一：Softmax 的雅可比矩阵原本是复杂的 $C \times C$ 稠密矩阵 $\diag(p) - pp^\top$，交叉熵的导数带有危险的 $-y_k / p_k$ 除法（当 $p_k \to 0$ 时爆炸），两者分开求导既慢又不稳定；而将它们融合后，所有危险的除法在代数上精确抵消，只剩“预测概率减去真实标签”这一完全有界（落在 $[-1, 1]$）的误差信号！这就是 PyTorch 强烈推荐使用 \texttt{F.cross\_entropy(logits, target)} 而非手动串联 \texttt{F.softmax} 与 \texttt{torch.log} 的根本原因。

\section{实测验证：纯 NumPy 手写 MLP 攻克双月流形}
\label{sec:exp-moons}

现在我们把式~\eqref{eq:linear-vjp}、\eqref{eq:relu-vjp} 与~\eqref{eq:ce-vjp} 拼装成一个\textbf{零框架依赖}的纯 NumPy 两层 ReLU MLP（结构 $2 \to 32 \to 1$），并在经典的二维双月数据集（$N = 240$ 个样本，噪声标准差 $\sigma = 0.10$，随机种子 $2026$）上与线性逻辑回归同台竞技。

\begin{lstlisting}[caption={纯 NumPy 2 层 ReLU MLP 的前向与手动 VJP 反向传播（摘自 \texttt{vol5\_experiments.py}）}]
# 前向传播：X (240, 2) -> W1 (2, 32) -> ReLU -> W2 (32, 1) -> sigmoid
h_pre = X @ W1 + b1
h = np.maximum(0.0, h_pre)
zm = (h @ W2 + b2).ravel()
pm = 1.0 / (1.0 + np.exp(-np.clip(zm, -30, 30)))
mlp_loss = float(np.mean(np.logaddexp(0, zm) - y * zm))  # 数值稳定二元交叉熵

# 反向传播：严格按式 (1.4)-(1.7) 逐层执行张量 VJP
dzm = ((pm - y) / len(y))[:, None]      # (N, 1)，融合 Sigmoid + BCE 梯度
dW2 = h.T @ dzm                         # (32, 1)
db2 = np.sum(dzm, axis=0)               # (1,)
dh = dzm @ W2.T                         # (N, 32)
dh_pre = dh * (h_pre > 0)               # ReLU 门控掩码 VJP
dW1 = X.T @ dh_pre                      # (2, 32)
db1 = np.sum(dh_pre, axis=0)            # (32,)
\end{lstlisting}

图~\ref{fig:mlp-moons} 直接读取实测记录文件 \texttt{data/mlp\_moons.csv}，对比了 600 轮梯度下降中两者的损失与准确率演化轨迹：
\begin{itemize}
  \item \textbf{线性逻辑回归}（学习率 $\eta = 0.40$）：前 150 轮迅速找到最优分割直线后便撞上几何天花板；第 480 至 600 轮准确率彻底停滞于 \textbf{89.17\%}，交叉熵损失停在 \textbf{0.2305}（\texttt{0.23054}）。再增加迭代次数也无济于事——这正是欠拟合与表示能力不足的典型指纹。
  \item \textbf{两层 ReLU MLP}（$2 \to 32 \to 1$，学习率 $\eta = 0.45$）：32 个隐藏神经元在平面上切出多段折线，逐渐将两个半月弯钩完整包裹分离；第 150 轮准确率已达 $95.00\%$，第 510 轮起稳定达到 \textbf{99.17\%}，最终第 600 轮损失降至 \textbf{0.0460}（\texttt{0.04603}），仅为线性基线的五分之一。
\end{itemize}

\begin{figure}[htbp]
\centering
\begin{tikzpicture}
  % Left panel: Loss
  \begin{axis}[
    name=lossplot,
    width=0.47\linewidth, height=5.8cm,
    xlabel={训练轮数（Epoch）},
    ylabel={二元交叉熵损失},
    xmin=-12, xmax=612, ymin=0, ymax=0.80,
    grid=both, grid style={tealLine!45, dashed},
    axis line style={slate},
    tick label style={font=\scriptsize},
    label style={font=\footnotesize\sffamily, color=ink},
    title={\small\sffamily\bfseries (a) 训练损失对比},
    legend style={at={(0.97,0.96)}, anchor=north east, font=\scriptsize, fill=white, draw=tealLine}
  ]
    \addplot[ochre, very thick, mark=square*, mark size=1.4pt, mark repeat=3]
      table[x=epoch, y=lin_loss, col sep=comma]{data/mlp_moons.csv};
    \addlegendentry{线性模型（终值 0.2305）}
    \addplot[teal, very thick, mark=*, mark size=1.4pt, mark repeat=3]
      table[x=epoch, y=mlp_loss, col sep=comma]{data/mlp_moons.csv};
    \addlegendentry{ReLU MLP（终值 0.0460）}
  \end{axis}

  % Right panel: Accuracy
  \begin{axis}[
    at={(lossplot.east)}, anchor=west, xshift=1.4cm,
    width=0.47\linewidth, height=5.8cm,
    xlabel={训练轮数（Epoch）},
    ylabel={分类准确率（\%）},
    xmin=-12, xmax=612, ymin=48, ymax=103,
    grid=both, grid style={tealLine!45, dashed},
    axis line style={slate},
    tick label style={font=\scriptsize},
    label style={font=\footnotesize\sffamily, color=ink},
    title={\small\sffamily\bfseries (b) 分类准确率对比},
    legend style={at={(0.97,0.05)}, anchor=south east, font=\scriptsize, fill=white, draw=tealLine}
  ]
    \addplot[ochre, very thick, mark=square*, mark size=1.4pt, mark repeat=3]
      table[x=epoch, y=lin_acc, col sep=comma]{data/mlp_moons.csv};
    \addlegendentry{线性模型（停滞 89.17\%）}
    \addplot[teal, very thick, mark=*, mark size=1.4pt, mark repeat=3]
      table[x=epoch, y=mlp_acc, col sep=comma]{data/mlp_moons.csv};
    \addlegendentry{ReLU MLP（达到 99.17\%）}
    \draw[dashed, slate, thin] (axis cs:0, 89.17) -- (axis cs:600, 89.17);
  \end{axis}
\end{tikzpicture}
\caption{双月数据集（$N=240$）上的实测训练曲线（读取自 \texttt{data/mlp\_moons.csv}）。线性分类器受限于单一超平面在 $89.17\%$ 处停滞，而 32 单元纯 NumPy ReLU MLP 凭借分段线性折叠达到 $99.17\%$。}
\label{fig:mlp-moons}
\end{figure}

\begin{failBox}{自动微分与 MLP 实现中的三大经典踩坑}
\begin{enumerate}
  \item \textbf{死亡 ReLU（Dead ReLU）}：若某隐藏神经元因过大的学习率或负偏置，在全体训练样本上均满足 $w_j^\top x + b_j < 0$，则由式~\eqref{eq:relu-vjp} 知其梯度恒为 $0$，此后参数永远无法更新、该神经元永久“死亡”。缓解手段包括合理控制学习率、采用 Kaiming 初始化（第~\ref{ch:train-dynamics} 章），或在需要时改用 LeakyReLU（$\max(\alpha z, z)$）、GELU 或 SiLU/SwiGLU（第 4 章 Transformer 标配）。
  \item \textbf{忘记调用 \texttt{optimizer.zero\_grad()}}：正如我们在 \texttt{Value} 引擎中看到的，为了正确处理多分支复用，PyTorch 默认将每次 \texttt{loss.backward()} 算出的梯度\textbf{累加}到 \texttt{param.grad} 中！如果在训练循环中忘记清空上一轮的梯度，参数更新方向就会变成历史所有小批量梯度的总和，导致训练发散或极度诡异的震荡。（反过来，刻意在 $K$ 个微批量后再清零，就是大模型显存受限时常用的\textbf{梯度累积（Gradient Accumulation）}技巧。）
  \item \textbf{隐式广播形状陷阱与原地修改（In-place）破坏磁带}：在 NumPy/PyTorch 中，若把形状为 \texttt{(N, 1)} 的预测与形状为 \texttt{(N,)} 的标签相减，广播机制会静默生成一个 \texttt{(N, N)} 的巨大矩阵，损失看似在下降却早已算错全部样本配对！务必对中间张量加 \texttt{assert} 断言形状。此外，在 PyTorch 中对需要参与反向求导的激活值执行原地修改（如 \texttt{h += 1} 或 \texttt{relu\_} 覆盖了计算局部导数所需的旧值），会直接触发运行时版本计数器报错。
\end{enumerate}
\end{failBox}

\begin{extBox}{梯度检查点（Activation Checkpointing）：用算力换显存}
由定理~\ref{thm:cheap-gradient} 可知，反向模式自动微分虽然在\textbf{时间}上仅耗 $\le 4\times$ 前向计算，却付出了沉重的\textbf{空间}代价：它必须把全部 $L$ 层的中间激活值留在 GPU 显存里，等待反向遍历时计算局部 VJP。对于一个 $L$ 层深度网络，激活显存随深度线性增长为 $O(L)$；在训练百亿参数大模型和超长上下文时，激活显存往往远超模型参数本身。

Griewank 与 Walther（2000）以及 Chen 等人（2016）提出的\textbf{梯度检查点（Gradient Checkpointing / Activation Rematerialization）}优雅地打破了这一瓶颈：
\begin{itemize}
  \item \textbf{分段策略}：将 $L$ 层划分为 $\sqrt{L}$ 段，每段长度为 $\sqrt{L}$。前向传播时，仅在每段的边界保存检查点激活（共需 $O(\sqrt{L})$ 显存），丢弃段内所有中间激活。
  \item \textbf{按需重算}：反向传播进入某一段时，从该段起点的检查点重新做一次局部前向传播，恢复该段内的 $\sqrt{L}$ 个中间激活，再完成该段反向 VJP 后立即释放。
\end{itemize}
这一策略将峰值激活显存从 $O(L)$ 压缩至 $O(\sqrt{L})$（对 $L=100$ 层足足节省 $10\times$ 显存！），代价仅仅是额外多做\textbf{一次}前向传播——总耗时仅增加约 $33\%$。这是现代大模型分布式训练标配技术之一，我们将在第 6 卷第 5 章（AI 系统与显存优化）中结合 FlashAttention 深入探讨。
\end{extBox}

\section*{本章小结与核心脉络}
\addcontentsline{toc}{section}{本章小结与核心脉络}

本章确立了深度学习得以运转的两根支柱：在\textbf{表示侧}，ReLU MLP 通过多层分段线性折叠突破线性超平面的凸性牢笼，一维帽函数构造给出了万能逼近的直观证明；在\textbf{计算侧}，反向模式自动微分（VJP）沿着计算图逆拓扑序回传伴随量，以与参数量 $P$ 无关的 $O(1)$ 倍前向开销，精确求出全部梯度，并通过张量级 VJP（仿射、ReLU、融合交叉熵）无缝衔接硬件矩阵加速。

然而，手握精确的自动微分引擎，是否意味着只要把 MLP 堆到 20 层、50 层就能无脑提升性能？恰恰相反——在下一章中我们会惊骇地发现，未经精心设计的 20 层 ReLU 网络在第一次前向传播结束时，信号就已经要么坍缩到 $10^{-10}$，要么爆炸到 $10^3$。打开深度训练大门的第二把钥匙，正是\textbf{深层网络的训练动力学}。

\section*{本章分层习题}
\addcontentsline{toc}{section}{本章分层习题}

\exTier{A 层：概念与手算热身}{巩固计算图前向/反向规则与分段线性几何直觉}

\begin{enumerate}
  \item[\textbf{5.1.1}] \textbf{三输入计算图手算}：考虑标量函数 $f(a, b, c) = \sigma(a b + c) \cdot a$，其中 $\sigma(z) = (1 + e^{-z})^{-1}$ 为 Sigmoid 函数。在点 $(a, b, c) = (2.0, -0.5, 1.0)$ 处：
  \begin{enumerate}
    \item 画出包含所有中间节点的计算 DAG，标出每条边的前向值；
    \item 注意变量 $a$ 的出度为 $2$，请按逆拓扑序手算全部伴随量 $\bar{a}, \bar{b}, \bar{c}$。
  \end{enumerate}
  \exHint{利用 $\sigma'(z) = \sigma(z)(1-\sigma(z))$，并在计算 $\bar{a}$ 时将乘法分支与内部 $ab$ 分支的梯度相加。}

  \item[\textbf{5.1.2}] \textbf{ReLU 折纸几何与 XOR 变体}：构造一个具有最少数目隐藏 ReLU 神经元的双层网络，在二维平面上精确表示“正方形指示器”的平滑近似：在正方形 $[-1, 1]^2$ 内部输出严格为正，在正方形 $[-2, 2]^2$ 外部输出严格为负。分析每个隐藏神经元对应的超平面位置。
  \exHint{先用两个一维 ReLU 组合出区间 $[-1, 1]$ 上的梯形平台函数，再对 $x_1$ 与 $x_2$ 两个方向求和并减去阈值。}
\end{enumerate}

\exTier{B 层：定理推导与严谨证明}{掌握张量 VJP、复合算子求导与复杂度界}

\begin{enumerate}
  \item[\textbf{5.1.3}] \textbf{双线性形式与二次型的张量 VJP}：设 $x \in \R^d, A \in \R^{d \times d}$，标量函数为 $f(x, A) = x^\top A x$，另有批量矩阵乘法 $Z = X W_1 W_2$。使用 Frobenius 迹恒等式推导 $\bar{x}, \bar{A}$ 以及给定 $\bar{Z}$ 时的 $\bar{X}, \bar{W}_1, \bar{W}_2$。当 $A$ 并非对称矩阵时，$\nabla_x (x^\top A x)$ 是否仍等于 $2Ax$？
  \exHint{展开 $d(x^\top A x) = (dx)^\top A x + x^\top (dA) x + x^\top A (dx)$，利用标量的转置等于自身将第一项改写为 $x^\top A^\top dx$。}

  \item[\textbf{5.1.4}] \textbf{带标签平滑（Label Smoothing）的融合交叉熵 VJP}：现代视觉与语言模型常将独热标签 $y$ 替换为平滑目标 $\tilde{y} = (1-\varepsilon)y + \frac{\varepsilon}{C}\mathbf{1}_C$（$\varepsilon \in (0,1)$）。证明融合损失 $\ell_{\text{ls}}(z, y) = -\sum_{k=1}^C \tilde{y}_k \log \softmax(z)_k$ 关于 logits $z$ 的梯度仍为 $p - \tilde{y}$，并说明为何它等价于在原交叉熵上加了一个惩罚项 $\varepsilon \DKL(u \,\|\, p)$（$u$ 为均匀分布）。
  \exHint{利用 $\sum_k \tilde{y}_k = 1$ 直接套用命题~\ref{prop:ce-vjp} 的证明步骤，再拆开 $-\sum_k u_k \log p_k$ 与 KL 散度比较。}
\end{enumerate}

\exTier{C 层：代码实现与数值实验}{扩展手写自动微分引擎并用数值差分检验梯度}

\begin{enumerate}
  \item[\textbf{5.1.5}] \textbf{为 \texttt{Value} 引擎补齐运算符并实现自动梯度检验}：扩展第~\ref{sec:micrograd} 节的 \texttt{Value} 类，支持幂运算 \texttt{\_\_pow\_\_}、指数 \texttt{exp()}、对数 \texttt{log()}、除法 \texttt{\_\_truediv\_\_} 与 \texttt{relu()}。编写通用函数 \texttt{grad\_check(f, xs, h=1e-6)}，用中心有限差分检验自定义 2 层标量 MLP 的所有梯度，验证相对误差小于 $10^{-8}$。
  \exHint{除法可通过 $a / b = a \cdot b^{-1}$ 复用乘法与幂运算；注意处理相对误差分母接近 $0$ 时的平滑项 $\max(\abs{g_{\text{ad}}} + \abs{g_{\text{num}}}, 10^{-8})$。}

  \item[\textbf{5.1.6}] \textbf{双月数据集上的深度 vs 宽度消融}：修改 \texttt{code/vol5\_experiments.py} 中的 \texttt{exp1\_micrograd\_and\_mlp}，在固定总参数量约 $100$ 个的条件下，比较单隐藏层（$2 \to 25 \to 1$）与双隐藏层（$2 \to 9 \to 8 \to 1$）在双月数据与同心圆数据上的收敛速度与决策边界形状。
  \exHint{在二维网格 $[-1.5, 2.5] \times [-1.0, 1.5]$ 上评估模型输出概率并画等高线，可直观看到分段线性折痕。}
\end{enumerate}

\exTier{D 层：开放探究与前沿迁移}{连接高阶导数与计算图系统设计}

\begin{enumerate}
  \item[\textbf{5.1.7}] \textbf{Pearlmutter 技巧：免构造 Hessian 的 Hessian-向量积（HVP）}：对于拥有 $P = 10^7$ 个参数的神经网络，显式存储 $P \times P$ 的 Hessian 矩阵 $H = \nabla^2_\theta \mathcal{L}$ 需要数百 TB 内存。证明借助\textbf{一次反向模式 + 一次前向模式}（Forward-over-Reverse），可以在 $O(\operatorname{Cost}(\mathcal{L}))$ 时间与内存内精确计算任意方向 $v \in \R^P$ 上的 Hessian-向量积 $Hv = \nabla_\theta (\nabla_\theta \mathcal{L}(\theta)^\top v)$。为什么在反向-反向（Reverse-over-Reverse）模式下也能算出 $Hv$，二者常数开销有何差异？
  \exHint{令 $g(\theta) = \nabla_\theta \mathcal{L}(\theta)$，则 $Hv$ 就是映射 $g: \R^P \to \R^P$ 沿方向 $v$ 的 JVP！也可对标量 $\phi(\theta) = g(\theta)^\top v$ 再执行一次 VJP。}

  \item[\textbf{5.1.8}] \textbf{隐函数微分（Implicit Differentiation）与深度平衡模型}：假设某层的输出由不动点方程 $z^* = f_\theta(z^*, x)$ 隐式定义（如迭代求解到收敛），若直接对 100 步不动点迭代展开反向传播，需要缓存 100 份激活。利用隐函数定理证明：
  \[
    \frac{\partial z^*}{\partial \theta} = \left(I - \frac{\partial f_\theta(z^*, x)}{\partial z^*}\right)^{-1} \frac{\partial f_\theta(z^*, x)}{\partial \theta},
  \]
  并说明如何通过求解一个线性方程组，使该层反向传播的显存占用降为 $O(1)$（与前向迭代步数无关）。
  \exHint{对等式 $z^*(\theta) = f_\theta(z^*(\theta), x)$ 两端关于 $\theta$ 求全导数并整理，再将左乘上游伴随量 $\bar{z}^*$ 转化为求解线性系统 $(I - J_z^\top) u = \bar{z}^*$。}
\end{enumerate}

\chapter[深层训练动力学：初始化与优化]{深层网络的训练动力学：初始化、归一化与自适应优化}
\label{ch:train-dynamics}

\chapEpigraph{给我四个参数，我就能画出一头大象；给我五个，我就能让它甩起鼻子。}{John von Neumann（冯·诺依曼，费米转述）}

\begin{chapIntro}{本章导读：驯服 20 层乃至 100 层网络的信号与梯度}
第~\ref{ch:mlp-autodiff} 章赋予了我们两件利器：深层 ReLU 网络的指数级折叠表达力，以及反向模式自动微分的精确梯度。然而在深度学习历史上，直到 2010 年代中期之前，训练一个 20 层以上的纯前馈网络都如同走钢丝——稍不留神，前向激活就会在十几层内坍缩成全零，或者爆炸成 \texttt{NaN}；反向梯度也随之呈指数级消失或发散。

本章从第一性原理剖析深层网络的\textbf{训练动力学（Training Dynamics）}，逐一拆解让现代深度网络稳定收敛的四大核心支柱：
\begin{enumerate}
  \item \textbf{信号传播与初始化定理}：定量刻画 $L$ 层连乘如何导致方差呈 $\sigma^{2L}$ 指数失控；严格证明 \textbf{Xavier（Glorot）} 与 \textbf{Kaiming（He）} 方差守恒初始化定理（尤其揭示 ReLU 将二阶矩精确砍半的 $\frac{2}{n_{\text{in}}}$ 因子）。
  \item \textbf{归一化家族的张量几何}：统一比较 \textbf{BatchNorm}、\textbf{LayerNorm} 与 \textbf{RMSNorm} 的数学定义、归一化轴切面与推理陷阱，并用 20 层 ReLU 实测（\texttt{data/init\_prop.csv}）直观展示四条激活传播曲线的十三个数量级悬殊。
  \item \textbf{自适应优化与解耦权重衰减}：沿着 $\text{SGD+Momentum} \to \text{AdaGrad} \to \text{RMSProp} \to \text{Adam} \to \text{AdamW}$ 的脉络完整推导；严格证明 Adam 的偏差校正公式，并数学证明为何 \textbf{Adam + $L_2$ 正则化不等价于 AdamW 的解耦权重衰减}。
  \item \textbf{正则化与学习率调度}：证明反向 Dropout 的期望守恒性，剖析线性预热（Warmup）与余弦退火（Cosine Annealing）背后的动力学直觉。
\end{enumerate}
\end{chapIntro}

\section{深层信号与梯度的指数病理}
\label{sec:vanishing-exploding}

考虑一个 $L$ 层无偏置前馈网络，每层宽度均为 $n$，输入向量 $x^{(0)} \in \R^n$ 满足各分量独立同分布、零均值且单位方差：$\E[x^{(0)}_j] = 0,\; \Var(x^{(0)}_j) = 1$。为先剥离非线性干扰，考察线性链
\[
  x^{(l)} = W^{(l)} x^{(l-1)}, \qquad l = 1, 2, \dots, L,
\]
其中每个权重矩阵 $W^{(l)} \in \R^{n \times n}$ 的元素独立抽自零均值、方差为 $\sigma_w^2$ 的分布。

\subsection{前向方差的指数坍缩与爆炸}

考察第 $l$ 层任意单个神经元 $x^{(l)}_i = \sum_{j=1}^{n} W^{(l)}_{ij} x^{(l-1)}_j$。利用权重与上一层激活的相互独立性以及零均值条件 $\E[W^{(l)}_{ij}] = 0$，我们立刻得到：
\begin{align*}
  \E\!\left[x^{(l)}_i\right] &= \sum_{j=1}^{n} \E\!\left[W^{(l)}_{ij}\right] \E\!\left[x^{(l-1)}_j\right] = 0, \\
  \Var\!\left(x^{(l)}_i\right) &= \E\!\left[\left(\sum_{j=1}^{n} W^{(l)}_{ij} x^{(l-1)}_j\right)^2\right]
  = \sum_{j=1}^{n} \E\!\left[\left(W^{(l)}_{ij}\right)^2\right] \E\!\left[\left(x^{(l-1)}_j\right)^2\right]
  = n \sigma_w^2 \Var\!\left(x^{(l-1)}_j\right).
\end{align*}
其中交叉项 $\E[W^{(l)}_{ij} W^{(l)}_{ik} x^{(l-1)}_j x^{(l-1)}_k]$ 对所有 $j \neq k$ 均因权重独立且零均值而严格消失。沿着 $L$ 层递推，输出方差满足指数律：
\begin{equation}
  \Var\!\left(x^{(L)}_i\right) = \left(n \sigma_w^2\right)^L \Var\!\left(x^{(0)}_i\right).
  \label{eq:linear-var-exp}
\end{equation}
式~\eqref{eq:linear-var-exp} 的后果触目惊心：记增益因子 $\alpha = n \sigma_w^2$。只要 $\alpha$ 略微偏离临界点 $1$，深度 $L$ 就会将其指数放大：
\begin{itemize}
  \item 若 $\sigma_w = 0.5 / \sqrt{n}$（即 $\alpha = 0.25$），在 $L = 20$ 层后方差缩小为 $0.25^{20} \approx 9.1 \times 10^{-13}$，标准差跌至 $10^{-6}$ 量级；若再叠加上 ReLU 每次砍半的能量衰减，激活值直接跌入 $10^{-10}$ 的浮点噪声泥潭！
  \item 若 $\sigma_w = 2.2 / \sqrt{n}$（即 $\alpha = 4.84$），在 $L = 20$ 层后标准差剧增至 $2.2^{20} \approx 7.0 \times 10^6$（即使加上 ReLU 砍半也仍达数千倍），激活与梯度迅速上溢为 \texttt{inf}/\texttt{NaN}。
\end{itemize}

\subsection{反向梯度的对称灾难与饱和激活函数}

反向传播与前向传播在数学上完全对偶。由第~\ref{ch:mlp-autodiff} 章的仿射层 VJP，设反向误差信号（伴随量）为 $\delta^{(l)} = \frac{\partial \mathcal{L}}{\partial z^{(l)}} \in \R^{n_l}$，则
\[
  \delta^{(l-1)} = \left(W^{(l)}\right)^\top \delta^{(l)} \had \sigma'\!\left(z^{(l-1)}\right).
\]
从第 $L$ 层一路回传到第 $1$ 层，梯度要连乘 $L-1$ 个矩阵和 $L-1$ 个激活导数对角阵：
\begin{equation}
  \frac{\partial z^{(L)}}{\partial z^{(1)}} = \prod_{l=2}^{L} W^{(l)} \diag\!\left(\sigma'\!\left(z^{(l-1)}\right)\right).
  \label{eq:grad-chain}
\end{equation}
这就解释了 1990 年代 Sigmoid 网络为何几乎不可能训练超过 5 层：因为 Sigmoid 导数的全局最大值仅为
\[
  \max_{z \in \R} \sigma'(z) = \sigma(0)(1-\sigma(0)) = \frac{1}{4}!
\]
哪怕权重矩阵都是完美的正交矩阵（谱范数恒为 $1$），仅 Sigmoid 导数一项就会把梯度每层至少衰减至 $\frac{1}{4}$，20 层之后梯度就被压缩到 $4^{-20} \approx 9.1 \times 10^{-13}$——底层权重完全收不到任何有效更新信号，这就是历史上著名的\textbf{梯度消失问题（Vanishing Gradient Problem）}。

\section{Xavier 与 Kaiming 方差守恒初始化定理}
\label{sec:init-theory}

既然式~\eqref{eq:linear-var-exp} 表明信号能否平稳穿透深层网络取决于每层的方差缩放系数是否精确等于 $1$，我们就可以把这一条件写成严格的数学定理，直接解出不同激活函数下最优的权重初始化方差！

\begin{thmBox}
\begin{theorem}[Xavier/Glorot 与 Kaiming/He 方差守恒初始化]
\label{thm:xavier-kaiming}
考虑第 $l$ 层仿射变换 $z^{(l)} = W^{(l)} x^{(l-1)} + b^{(l)}$，其中 $W^{(l)} \in \R^{n_{\text{out}} \times n_{\text{in}}}$ 的元素独立同分布于零均值、对称分布（方差为 $\sigma_w^2 = \Var(w)$），偏置 $b^{(l)} = 0$，输入 $x^{(l-1)} = \sigma(z^{(l-1)})$ 与 $W^{(l)}$ 相互独立，且上一层预激活 $z^{(l-1)}$ 各分量独立同分布、关于原点对称（均值 $0$，方差 $\Var(z^{(l-1)})$）。
\begin{enumerate}
  \item \textbf{线性或原点附近 $\tanh$ 激活（Xavier / Glorot, 2010）}：当 $\sigma(z) \approx z$ 时，
  \[
    \Var\!\left(z^{(l)}\right) = n_{\text{in}} \sigma_w^2\,\Var\!\left(z^{(l-1)}\right), \qquad
    \Var\!\left(\delta^{(l-1)}\right) = n_{\text{out}} \sigma_w^2\,\Var\!\left(\delta^{(l)}\right).
  \]
  前向方差守恒要求 $\sigma_w^2 = \frac{1}{n_{\text{in}}}$，反向方差守恒要求 $\sigma_w^2 = \frac{1}{n_{\text{out}}}$。当 $n_{\text{in}} \neq n_{\text{out}}$ 时取调和平均折中，即得 \textbf{Xavier 初始化}：
  \begin{equation}
    \Var(w) = \frac{2}{n_{\text{in}} + n_{\text{out}}}.
    \label{eq:xavier-init}
  \end{equation}
  \item \textbf{ReLU 激活（Kaiming / He, 2015）}：当 $x^{(l-1)} = \ReLU(z^{(l-1)})$ 时，有精确的二阶矩关系
  \begin{equation}
    \E\!\left[\ReLU\!\left(z^{(l-1)}\right)^2\right] = \frac{1}{2}\,\Var\!\left(z^{(l-1)}\right),
    \label{eq:relu-half-moment}
  \end{equation}
  由此前向预激活方差递推满足 $\Var(z^{(l)}) = \frac{1}{2} n_{\text{in}} \sigma_w^2\,\Var(z^{(l-1)})$，反向梯度方差递推满足 $\Var(\delta^{(l-1)}) = \frac{1}{2} n_{\text{out}} \sigma_w^2\,\Var(\delta^{(l)})$。因此保持前向（fan-in）或反向（fan-out）信号方差守恒的 \textbf{Kaiming 初始化}为：
  \begin{equation}
    \Var(w) = \frac{2}{n_{\text{in}}} \quad (\text{fan-in 模式}) \qquad \text{或} \qquad
    \Var(w) = \frac{2}{n_{\text{out}}} \quad (\text{fan-out 模式}).
    \label{eq:kaiming-init}
  \end{equation}
\end{enumerate}
\end{theorem}
\begin{proof}
首先证明 ReLU 下的前向二阶矩公式~\eqref{eq:relu-half-moment}。由于 $W^{(l)}$ 与 $x^{(l-1)}$ 独立且 $\E[W^{(l)}_{ij}] = 0$，第 $l$ 层预激活 $z^{(l)}_i = \sum_{j=1}^{n_{\text{in}}} W^{(l)}_{ij} x^{(l-1)}_j$ 仍满足零均值 $\E[z^{(l)}_i] = 0$，并且由于 $W^{(l)}_{ij}$ 分布关于原点对称，$z^{(l)}_i$ 的分布也关于原点对称！

下面计算 $z^{(l)}_i$ 的方差。注意此时 $x^{(l-1)}_j = \ReLU(z^{(l-1)}_j)$ \textbf{不再是零均值}（它恒非负！），因此不能粗暴地套用 $\Var(WX) = \Var(W)\Var(X)$，而必须使用二阶原点矩：
\begin{align*}
  \Var\!\left(z^{(l)}_i\right)
  &= \E\!\left[\left(\sum_{j=1}^{n_{\text{in}}} W^{(l)}_{ij} x^{(l-1)}_j\right)^2\right]
  = \sum_{j=1}^{n_{\text{in}}} \E\!\left[\left(W^{(l)}_{ij}\right)^2\right] \E\!\left[\left(x^{(l-1)}_j\right)^2\right] \\
  &= n_{\text{in}} \Var(w)\, \E\!\left[\ReLU\!\left(z^{(l-1)}_j\right)^2\right].
\end{align*}
因为交叉项中含有 $\E[W^{(l)}_{ij} W^{(l)}_{ik}] = \E[W^{(l)}_{ij}]\E[W^{(l)}_{ik}] = 0$（$j \neq k$），所以依然严格为零。

接下来利用 $z := z^{(l-1)}_j$ 分布关于原点对称（概率密度满足 $p(-z) = p(z)$）且 $\E[z] = 0$，直接积分计算 $\E[\ReLU(z)^2]$：
\[
  \E\!\left[\ReLU(z)^2\right]
  = \int_{0}^{+\infty} z^2 p(z)\,dz
  = \frac{1}{2} \int_{-\infty}^{+\infty} z^2 p(z)\,dz
  = \frac{1}{2}\E[z^2] = \frac{1}{2}\Var(z).
\]
代入上式立刻得到
\[
  \Var\!\left(z^{(l)}\right) = \frac{1}{2} n_{\text{in}} \Var(w)\,\Var\!\left(z^{(l-1)}\right).
\]
令每层缩放系数 $\frac{1}{2} n_{\text{in}} \Var(w) = 1$，便得到 Kaiming fan-in 条件 $\Var(w) = \frac{2}{n_{\text{in}}}$。

再看反向传播：记 $g^{(l)}_j = \frac{\partial \mathcal{L}}{\partial x^{(l)}_j}$ 与 $\delta^{(l)}_j = \frac{\partial \mathcal{L}}{\partial z^{(l)}_j} = g^{(l)}_j \cdot \mathbb{I}[z^{(l)}_j > 0]$。由对称性，$\Prob(z^{(l)}_j > 0) = \frac{1}{2}$；近似认为掩码与上游梯度 $g^{(l)}_j$ 独立且 $\E[g^{(l)}_j] = 0$，则 $\Var(\delta^{(l)}_j) = \frac{1}{2}\Var(g^{(l)}_j)$。又因为 $g^{(l-1)}_k = \sum_{i=1}^{n_{\text{out}}} W^{(l)}_{ik} \delta^{(l)}_i$，重复完全相同的求和计算即得
\[
  \Var\!\left(g^{(l-1)}\right) = n_{\text{out}} \Var(w)\,\Var\!\left(\delta^{(l)}\right) = \frac{1}{2} n_{\text{out}} \Var(w)\,\Var\!\left(g^{(l)}\right),
\]
令缩放系数为 $1$ 即得 fan-out 条件 $\Var(w) = \frac{2}{n_{\text{out}}}$。当激活函数退化为线性 $\sigma(z)=z$ 时，式~\eqref{eq:relu-half-moment} 中无 $\frac{1}{2}$ 因子，立刻退化为 Xavier 条件。
\end{proof}
\end{thmBox}

\begin{insightBox}{一个极容易被忽略的细节：Kaiming 守恒的是 $z$ 的方差，而非 $\ReLU(z)$ 的方差！}
设 $z \sim \mathcal{N}(0, \sigma^2)$，经 ReLU 后 $x = \ReLU(z)$ 的均值和方差是多少？直接积分可得：
\[
  \E[x] = \int_0^\infty z \frac{1}{\sqrt{2\pi}\sigma} e^{-z^2/(2\sigma^2)} dz = \frac{\sigma}{\sqrt{2\pi}} \approx 0.3989\sigma,
\]
\[
  \Var(x) = \E[x^2] - (\E[x])^2 = \left(\frac{1}{2} - \frac{1}{2\pi}\right)\sigma^2 \approx 0.3408\,\sigma^2 \implies \operatorname{std}(x) \approx 0.5838\,\sigma.
\]
但是请注意！由于下一层输入 $z' = Wx$ 的每个权重 $W_{ij}$ 零均值，它会自动把 $x$ 的非零均值抵消掉，因此进入下一层预激活方差的不是 $\Var(x)$，而是完整的二阶矩 $\E[x^2] = \frac{1}{2}\sigma^2$！这也解释了为什么后文实测中（图~\ref{fig:init-prop}），在 Kaiming 初始化下 ReLU 之后的激活标准差稳定在约 $0.84 \sim 0.96$（略低于初始高斯输入 $1.0$，具体数值取决于输入首层分布与有限宽度波动），而不会像 $\sigma_w = 0.5/\sqrt{D}$ 那样指数衰减。
\end{insightBox}

如果采用更一般的带负半轴斜率 $a$ 的 LeakyReLU（$\sigma(z) = \max(az, z)$），完全仿照定理~\ref{thm:xavier-kaiming} 的积分即可推得 $\E[\sigma(z)^2] = \frac{1 + a^2}{2}\Var(z)$，因此通用的 Kaiming 公式是 $\Var(w) = \frac{2}{(1+a^2)n_{\text{in}}}$——这正是 PyTorch \texttt{nn.init.kaiming\_normal\_} 源码中的精确公式！

\section{归一化家族：BatchNorm、LayerNorm 与 RMSNorm}
\label{sec:normalization}

即便使用了完美的 Kaiming 初始化，它也只保证\textbf{第 0 步（初始化瞬间）}的方差守恒。一旦梯度下降开始更新权重，各层分布就会随着参数演化而漂移，且残差支路、注意力权重等结构会再次破坏方差平衡。现代深度学习对此给出的终极工程解药，是在网络内部显式插入\textbf{归一化层（Normalization Layers）}，把激活张量实时拉回“零均值、单位方差”的舒适区。

归一化家族的核心区别，仅仅在于\textbf{沿着张量的哪几个维度计算均值与方差}。

\subsection{三种主流归一化的精确数学定义}

\begin{defBox}
\begin{definition}[BatchNorm、LayerNorm 与 RMSNorm]
设某层激活张量在视觉任务中形状为 $(N, C, H, W)$（批量、通道、高、宽），或在序列任务中形状为 $(B, T, D)$（批量、序列长度、特征维）。取数值稳定常数 $\epsilon \approx 10^{-5}$，引入可学习的仿射缩放参数 $\gamma$ 与平移参数 $\beta$（初始化为 $\gamma = \mathbf{1}, \beta = \mathbf{0}$，保证归一化层在需要时可恢复恒等映射）。
\begin{enumerate}
  \item \textbf{批归一化（Batch Normalization, BN；Ioffe \& Szegedy, 2015）}：对每个通道 $c \in \{1,\dots,C\}$，\textbf{跨整个小批量与空间像素} $(N, H, W)$ 统计标量均值与方差：
  \begin{equation}
  \begin{aligned}
    \mu_c &= \frac{1}{NHW}\sum_{n,h,w} x_{n,c,h,w}, \qquad
    \sigma_c^2 = \frac{1}{NHW}\sum_{n,h,w} (x_{n,c,h,w} - \mu_c)^2, \\
    \BN(x)_{n,c,h,w} &= \gamma_c \frac{x_{n,c,h,w} - \mu_c}{\sqrt{\sigma_c^2 + \epsilon}} + \beta_c.
  \end{aligned}
  \label{eq:bn}
  \end{equation}
  \item \textbf{层归一化（Layer Normalization, LN；Ba et al., 2016）}：对\textbf{每个样本、每个时间步} $(b, t)$ 独立地在特征维度 $D$ 上统计均值与方差（完全不跨样本！）：
  \begin{equation}
  \begin{aligned}
    \mu_{b,t} &= \frac{1}{D}\sum_{d=1}^{D} x_{b,t,d}, \qquad
    \sigma_{b,t}^2 = \frac{1}{D}\sum_{d=1}^{D} (x_{b,t,d} - \mu_{b,t})^2, \\
    \LN(x)_{b,t,d} &= \gamma_d \frac{x_{b,t,d} - \mu_{b,t}}{\sqrt{\sigma_{b,t}^2 + \epsilon}} + \beta_d.
  \end{aligned}
  \label{eq:ln}
  \end{equation}
  \item \textbf{均方根归一化（RMSNorm；Zhang \& Sennrich, 2019）}：现代大语言模型（LLaMA、Qwen、DeepSeek、Gemma 等）的默认标配。它发现 LayerNorm 的核心收益来自\textbf{重缩放不变性（Re-scaling Invariance）}而非中心平移，因此直接去掉均值减法与偏置 $\beta$，仅除以均方根：
  \begin{equation}
    \operatorname{RMS}(x_{b,t}) = \sqrt{\frac{1}{D}\sum_{d=1}^{D} x_{b,t,d}^2 + \epsilon}, \qquad
    \RMSN(x)_{b,t,d} = \gamma_d \frac{x_{b,t,d}}{\operatorname{RMS}(x_{b,t})}.
    \label{eq:rmsnorm}
  \end{equation}
\end{enumerate}
\end{definition}
\end{defBox}

图~\ref{fig:norm-slices} 用直观的张量切面对比了 BatchNorm、LayerNorm 与 RMSNorm 在 $(B, T, D)$ / $(N, C, L)$ 空间中的统计范围：BN 沿着批量轴 $N$ 纵向切割，因而把同一批次内不同样本耦合在了一起；而 LN 与 RMSNorm 沿着特征轴 $D$ 横向切片，每个 token 完全自给自足。

\begin{figure}[htbp]
\centering
\resizebox{0.96\linewidth}{!}{%
\begin{tikzpicture}[
  font=\sffamily\footnotesize,
  axis/.style={-{Stealth[length=2mm]}, thick, slate}
]
  % Helper macro to draw a 3D tensor grid: (x0, y0, highlight mode)
  % Mode 1: BatchNorm (slice across N and L for one C)
  \begin{scope}[shift={(0,0)}]
    \node[font=\sffamily\bfseries\small, color=tealDark] at (1.5, 3.3) {(a) BatchNorm};
    \node[font=\scriptsize, color=slate] at (1.5, 2.85) {沿批量 $N$ 与空间 $L$ 归一化};
    % Draw 4x3 front face (C=4 cols, N=3 rows) and depth L=2
    \foreach \n in {0,1,2} {
      \foreach \c in {0,1,2,3} {
        \ifnum\c=1
          \fill[teal!55] (\c*0.65, \n*0.65) rectangle ++(0.65, 0.65);
        \else
          \fill[tealLight] (\c*0.65, \n*0.65) rectangle ++(0.65, 0.65);
        \fi
        \draw[slate!60] (\c*0.65, \n*0.65) rectangle ++(0.65, 0.65);
      }
    }
    % Top face (depth L)
    \foreach \c in {0,1,2,3} {
      \ifnum\c=1
        \fill[teal!35] (\c*0.65, 1.95) -- ++(0.45, 0.35) -- ++(0.65, 0) -- ++(-0.45, -0.35) -- cycle;
      \else
        \fill[navyLight] (\c*0.65, 1.95) -- ++(0.45, 0.35) -- ++(0.65, 0) -- ++(-0.45, -0.35) -- cycle;
      \fi
      \draw[slate!60] (\c*0.65, 1.95) -- ++(0.45, 0.35) -- ++(0.65, 0) -- ++(-0.45, -0.35) -- cycle;
    }
    % Right face
    \foreach \n in {0,1,2} {
      \fill[navyLight!70] (2.60, \n*0.65) -- ++(0.45, 0.35) -- ++(0, 0.65) -- ++(-0.45, -0.35) -- cycle;
      \draw[slate!60] (2.60, \n*0.65) -- ++(0.45, 0.35) -- ++(0, 0.65) -- ++(-0.45, -0.35) -- cycle;
    }
    % Axes labels
    \draw[axis] (-0.25, 0) -- (-0.25, 2.05) node[above, font=\scriptsize] {$N$ (批量)};
    \draw[axis] (0, -0.25) -- (2.7, -0.25) node[right, font=\scriptsize] {$C$ (通道)};
    \draw[axis] (2.75, -0.1) -- (3.25, 0.3) node[right, font=\scriptsize] {$H\!\times\!W$};
    \node[font=\scriptsize, color=tealDark] at (1.5, -0.85) {依赖批量大小；需维护滑动统计量};
  \end{scope}

  % Mode 2: LayerNorm (slice across D for each (b, t))
  \begin{scope}[shift={(5.4,0)}]
    \node[font=\sffamily\bfseries\small, color=navy] at (1.5, 3.3) {(b) LayerNorm};
    \node[font=\scriptsize, color=slate] at (1.5, 2.85) {单样本单位置沿特征 $D$ 归一化};
    \foreach \n in {0,1,2} {
      \foreach \c in {0,1,2,3} {
        \ifnum\n=1
          \fill[navy!50] (\c*0.65, \n*0.65) rectangle ++(0.65, 0.65);
        \else
          \fill[navyLight] (\c*0.65, \n*0.65) rectangle ++(0.65, 0.65);
        \fi
        \draw[slate!60] (\c*0.65, \n*0.65) rectangle ++(0.65, 0.65);
      }
    }
    \foreach \c in {0,1,2,3} {
      \fill[navyLight!60] (\c*0.65, 1.95) -- ++(0.45, 0.35) -- ++(0.65, 0) -- ++(-0.45, -0.35) -- cycle;
      \draw[slate!60] (\c*0.65, 1.95) -- ++(0.45, 0.35) -- ++(0.65, 0) -- ++(-0.45, -0.35) -- cycle;
    }
    \foreach \n in {0,1,2} {
      \fill[navyLight!70] (2.60, \n*0.65) -- ++(0.45, 0.35) -- ++(0, 0.65) -- ++(-0.45, -0.35) -- cycle;
      \draw[slate!60] (2.60, \n*0.65) -- ++(0.45, 0.35) -- ++(0, 0.65) -- ++(-0.45, -0.35) -- cycle;
    }
    \draw[axis] (-0.25, 0) -- (-0.25, 2.05) node[above, font=\scriptsize] {$B$ (批量)};
    \draw[axis] (0, -0.25) -- (2.7, -0.25) node[right, font=\scriptsize] {$D$ (特征)};
    \draw[axis] (2.75, -0.1) -- (3.25, 0.3) node[right, font=\scriptsize] {$T$ (序列)};
    \node[font=\scriptsize, color=navy] at (1.5, -0.85) {减均值且除标准差；Transformer 标准};
  \end{scope}

  % Mode 3: RMSNorm (same slice as LN, no mean subtraction, L2 sphere projection)
  \begin{scope}[shift={(10.8,0)}]
    \node[font=\sffamily\bfseries\small, color=ochre] at (1.5, 3.3) {(c) RMSNorm};
    \node[font=\scriptsize, color=slate] at (1.5, 2.85) {仅沿特征 $D$ 按均方根重缩放};
    \foreach \n in {0,1,2} {
      \foreach \c in {0,1,2,3} {
        \ifnum\n=1
          \fill[ochre!55] (\c*0.65, \n*0.65) rectangle ++(0.65, 0.65);
        \else
          \fill[ochreLight] (\c*0.65, \n*0.65) rectangle ++(0.65, 0.65);
        \fi
        \draw[slate!60] (\c*0.65, \n*0.65) rectangle ++(0.65, 0.65);
      }
    }
    \foreach \c in {0,1,2,3} {
      \fill[ochreLight!60] (\c*0.65, 1.95) -- ++(0.45, 0.35) -- ++(0.65, 0) -- ++(-0.45, -0.35) -- cycle;
      \draw[slate!60] (\c*0.65, 1.95) -- ++(0.45, 0.35) -- ++(0.65, 0) -- ++(-0.45, -0.35) -- cycle;
    }
    \foreach \n in {0,1,2} {
      \fill[ochreLight!70] (2.60, \n*0.65) -- ++(0.45, 0.35) -- ++(0, 0.65) -- ++(-0.45, -0.35) -- cycle;
      \draw[slate!60] (2.60, \n*0.65) -- ++(0.45, 0.35) -- ++(0, 0.65) -- ++(-0.45, -0.35) -- cycle;
    }
    \draw[axis] (-0.25, 0) -- (-0.25, 2.05) node[above, font=\scriptsize] {$B$ (批量)};
    \draw[axis] (0, -0.25) -- (2.7, -0.25) node[right, font=\scriptsize] {$D$ (特征)};
    \draw[axis] (2.75, -0.1) -- (3.25, 0.3) node[right, font=\scriptsize] {$T$ (序列)};
    \node[font=\scriptsize, color=ochre] at (1.5, -0.85) {省去均值归约与 $\beta$；LLM 更快更省内存};
  \end{scope}
\end{tikzpicture}}
\caption{BatchNorm、LayerNorm 与 RMSNorm 的张量归一化切片几何对比。阴影区块表示同一次均值/方差（或均方根）归约所覆盖的元素集合。}
\label{fig:norm-slices}
\end{figure}

\subsection{为什么归一化能让训练提速十倍？}

早期文献将 BatchNorm 的成功归因于减少“内部协变量偏移（Internal Covariate Shift）”。然而 Santurkar 等人（2018）的经典实验证明，哪怕在 BN 之后人为注入分布噪声，网络依然收敛得更好！真正深层的数学原因是\textbf{归一化显著平滑了损失函数的 Lipschitz 常数与 $\beta$-光滑性}，同时赋予网络\textbf{权重尺度不变性（Scale Invariance）}：
\[
  \BN\!\left((\alpha W) x\right) = \BN(W x), \qquad \forall \alpha > 0.
\]
对尺度不变的函数求导可知 $\nabla_{\alpha W} \mathcal{L} = \frac{1}{\alpha} \nabla_W \mathcal{L}$：如果某层权重不慎变得过大，它的梯度会自动缩小 $\frac{1}{\alpha}$；反之权重过小则梯度自动放大——这种隐式的\textbf{自调节步长机制}让我们可以安全地使用比无归一化网络大一个数量级的学习率，大幅加快收敛。

此外，在视觉检测/分割任务中当单卡批量被迫降至 $N=1\text{或}2$ 时，BN 的批统计量方差极大，此时可使用 \textbf{GroupNorm}（将通道 $C$ 分成 $G$ 组，在每组 $(C/G, H, W)$ 上归一化），同样彻底摆脱对批量大小 $N$ 的依赖。而在自然语言与大模型中，序列长度可变且自回归推理时逐 token 生成，BN 既不自然又无法在线使用，因此 LayerNorm 与 RMSNorm 成为绝对统治者（RMSNorm 在 GPU 上少一次均值归约与写回，实际可减少约 $10\%\sim 40\%$ 的归一化算子耗时）。

\begin{failBox}{BatchNorm 工程排错头号杀手：\texttt{model.train()} 与 \texttt{model.eval()}}
\begin{enumerate}
  \item \textbf{训练与推理的行为割裂}：训练时 BN 使用当前小批量的实时均值 $\mu_B$ 与方差 $\sigma_B^2$，同时以指数滑动平均更新全局统计量 $\hat{\mu} \leftarrow (1-m)\hat{\mu} + m\mu_B$（PyTorch 默认动量 $m=0.1$）。推理时往往只输入单张图片（$N=1$），此时必须调用 \texttt{model.eval()} 切换为使用冻结的 $(\hat{\mu}, \hat{\sigma}^2)$！\textbf{新手最常见的惨案}是验证或部署时忘记 \texttt{model.eval()}，导致单样本推理时方差为 $0$，输出全部崩溃；或者批量极小时训练统计量噪声剧烈，验证精度莫名暴跌。
  \item \textbf{跨卡多机同步陷阱}：数据并行训练中，若每张 GPU 分到的微批量太小（如 $N_{\text{local}} = 2$），默认 BN 只在单卡内部算均值方差，必须替换为 \texttt{SyncBatchNorm} 或直接换用 GroupNorm/LayerNorm。
  \item \textbf{半精度 \texttt{fp16} 方差溢出}：在混合精度训练中，对 $x^2$ 求和极易突破 \texttt{fp16} 上限 $65504$，因此工业实现（PyTorch/Apex/FlashNorm）一律在归一化内部将中间累加器强制提升至 \texttt{fp32} 计算。
\end{enumerate}
\end{failBox}

\section{实测验证：20 层 ReLU 网络的激活传播实验}
\label{sec:exp-init}

{\rightskip=0pt plus 1fil\relax 为了亲眼见证理论公式~\eqref{eq:linear-var-exp}、\eqref{eq:kaiming-init} 与~\eqref{eq:ln} 的惊人威力，我们运行配套实验脚本 \texttt{code/vol5\_experiments.py} 中的 \texttt{exp2\_init\_norm\_opt}：构造一个 $L = 20$ 层、宽度 $D = 128$ 的纯 ReLU 深度前馈网络，输入批量 $B = 64$ 的标准高斯噪声 $x^{(0)} \sim \mathcal{N}(0, I)$（第 0 层实测标准差 $0.9975$），对比四种配置下各层输出激活的标准差：\par}
\begin{itemize}
  \item \textbf{Small 初始化}（$\sigma_w = 0.5 / \sqrt{D}$）：每层按理论因子 $\frac{1}{2}(0.5)^2 = 0.125$（标准差约 $0.35\times$）塌缩；第 10 层跌至 $2.88 \times 10^{-5}$，第 19 层跌至 $2.14 \times 10^{-9}$，第 20 层实测坍缩至 \textbf{\texttt{6.60e-10}}（绘图 CSV 文件中为对数坐标可读性截断在 \texttt{1.00e-09} 下界）。
  \item \textbf{Large 初始化}（$\sigma_w = 2.2 / \sqrt{D}$）：每层按因子 $\frac{1}{2}(2.2)^2 = 2.42$（标准差约 $1.56\times$）膨胀；第 10 层达 $65.3$，第 20 层实测爆炸至 \textbf{\texttt{6.87e+03}}（\texttt{6.8709e+03}）——与 Small 之间相差整整 \textbf{13 个数量级}！
  \item \textbf{Kaiming 初始化}（$\sigma_w = \sqrt{2/D}$）：完美命中临界点，穿透全部 20 层后既不坍缩也不爆炸，第 20 层标准差稳稳保持在 \textbf{\texttt{0.843}}（\texttt{8.4263e-01}），全程波动落在 $[0.76, 1.00]$ 之间。
  \item \textbf{Kaiming + LayerNorm}：每层 ReLU 后施加 LayerNorm，所有 20 层的激活标准差纹丝不动地锁定在 \textbf{\texttt{0.996}}（\texttt{9.9613e-01}——精确对应 PyTorch 样本标准差贝塞尔修正 $\sqrt{(128-1)/128} \approx 0.9961$！）。
\end{itemize}

图~\ref{fig:init-prop} 直接读取 \texttt{data/init\_prop.csv} 在半对数坐标系中绘出这四条实测曲线，两条倾斜笔直的对数直线生动印证了式~\eqref{eq:linear-var-exp} 的指数定律。

\begin{figure}[htbp]
\centering
\begin{tikzpicture}
  \begin{semilogyaxis}[
    width=0.86\linewidth, height=6.4cm,
    xlabel={网络深度（层号 $l$）},
    ylabel={激活张量标准差 $\operatorname{std}(h^{(l)})$（对数刻度）},
    xmin=-0.4, xmax=20.4, ymin=5e-10, ymax=3e4,
    xtick={0,2,4,6,8,10,12,14,16,18,20},
    grid=both, grid style={tealLine!45, dashed},
    axis line style={slate},
    tick label style={font=\scriptsize},
    label style={font=\footnotesize\sffamily, color=ink},
    legend columns=2,
    legend style={at={(0.5,1.04)}, anchor=south, font=\scriptsize, fill=white, draw=tealLine, legend cell align=left, column sep=8pt}
  ]
    \addplot[ochre, very thick, mark=triangle*, mark size=1.8pt]
      table[x=layer, y=large, col sep=comma]{data/init_prop.csv};
    \addlegendentry{Large ($2.2/\sqrt{D}$)：终值 $6.87\times 10^3$}

    \addplot[navy, very thick, dashed, mark=diamond*, mark size=1.8pt]
      table[x=layer, y=kaiming_ln, col sep=comma]{data/init_prop.csv};
    \addlegendentry{Kaiming + LN：锁定 $0.996$}

    \addplot[teal, very thick, mark=*, mark size=1.7pt]
      table[x=layer, y=kaiming, col sep=comma]{data/init_prop.csv};
    \addlegendentry{Kaiming ($\sqrt{2/D}$)：终值 $0.843$}

    \addplot[plum, very thick, mark=square*, mark size=1.6pt]
      table[x=layer, y=small, col sep=comma]{data/init_prop.csv};
    \addlegendentry{Small ($0.5/\sqrt{D}$)：终值 $6.60\times 10^{-10}$}

    \draw[dotted, slate, thick] (axis cs:0, 1) -- (axis cs:20, 1);
  \end{semilogyaxis}
\end{tikzpicture}
\caption{20 层 ReLU 网络（宽度 $D=128$）前向激活标准差的实测演化（读取自 \texttt{data/init\_prop.csv}）。半对数图上的两条直线清晰暴露了初始化失配时的指数爆炸与指数坍缩，而 Kaiming 初始化与 LayerNorm 将信号牢牢锚定在 $O(1)$ 水平线上。}
\label{fig:init-prop}
\end{figure}

\section{自适应优化器演化：从 SGD 动量到 Adam 与 AdamW}
\label{sec:optimizers}

解决了“信号能穿透深层网络”之后，第二个拦路虎浮出水面：深度神经网络的损失曲面是一个极度病态（ill-conditioned）的高维峡谷——某些参数方向（如大奇异值方向或高频特征）坡度极陡，另一些方向（如低频词嵌入）坡度极其平缓。固定步长的朴素随机梯度下降（SGD）在陡峭方向剧烈震荡、在平坦方向龟速爬行。现代优化器的全部进化史，就是为了解决这一曲率失衡问题。

\subsection{从动量（Momentum）到 AdaGrad 与 RMSProp}

\begin{enumerate}
  \item \textbf{SGD + 动量（Heavy-Ball Momentum）}：为梯度引入一阶指数滑动平均速度 $v_t$：
  \[
    v_t = \beta v_{t-1} + g_t, \qquad \theta_t = \theta_{t-1} - \eta v_t \qquad (\text{典型 }\beta = 0.9).
  \]
  在峡谷两壁来回震荡的高曲率方向，正负梯度在滑动平均中彼此抵消；而在沿谷底一致前进的低曲率方向，速度不断累积加速至 $\frac{1}{1-\beta} = 10\times$。Nesterov 加速梯度（NAG）则先沿动量方向“往前看一步”再算梯度，进一步提升收敛率。

  \item \textbf{AdaGrad（Duchi et al., 2011）}：为每个参数坐标 $i$ 独立累积历史梯度平方和 $G_{t,i} = \sum_{\tau=1}^{t} g_{\tau,i}^2$，并按其平方根倒数缩放步长：
  \[
    \theta_{t,i} = \theta_{t-1,i} - \frac{\eta}{\sqrt{G_{t,i}} + \epsilon}\, g_{t,i}.
  \]
  稀疏/低梯度坐标获得较大有效学习率，频繁/大梯度坐标自动减速。缺点在于 $G_{t,i}$ 单调递增，训练后期有效步长往往过早衰减至零。

  \item \textbf{RMSProp（Hinton, 2012）}：用指数滑动平均替代单调累积，遗忘遥远历史：
  \[
    v_t = \beta_2 v_{t-1} + (1-\beta_2) g_t \had g_t, \qquad \theta_t = \theta_{t-1} - \frac{\eta}{\sqrt{v_t} + \epsilon} \had g_t.
  \]
\end{enumerate}

\subsection{Adam 与偏差校正的严格证明}

\textbf{Adam（Adaptive Moment Estimation；Kingma \& Ba, 2015）} 将“一阶动量（Momentum）”与“二阶自适应步长（RMSProp）”优雅合一，并修补了零初始化带来的冷启动偏差。在每一步 $t = 1, 2, \dots$，给定当前小批量随机梯度 $g_t = \nabla_\theta \mathcal{L}_t(\theta_{t-1})$（默认超参数 $\beta_1 = 0.9,\; \beta_2 = 0.999,\; \epsilon = 10^{-8}$）：
\begin{align}
  m_t &= \beta_1 m_{t-1} + (1 - \beta_1) g_t, & \hat{m}_t &= \frac{m_t}{1 - \beta_1^t}, \label{eq:adam-m} \\
  v_t &= \beta_2 v_{t-1} + (1 - \beta_2) g_t^{\had 2}, & \hat{v}_t &= \frac{v_t}{1 - \beta_2^t}, \label{eq:adam-v} \\
  \theta_t &= \theta_{t-1} - \eta \frac{\hat{m}_t}{\sqrt{\hat{v}_t} + \epsilon}, \label{eq:adam-step}
\end{align}
其中初始状态 $m_0 = 0,\; v_0 = 0$，平方与除法均按元素进行。

\begin{thmBox}
\begin{proposition}[Adam 偏差校正的无偏性证明]
\label{prop:adam-bias}
假设在迭代初期随机梯度序列 $\{g_\tau\}_{\tau=1}^t$ 近似平稳（$\E[g_\tau] \equiv \mu_g$ 且 $\E[g_\tau^{\had 2}] \equiv \mu_{g^2}$）。则由 $m_0 = 0, v_0 = 0$ 递推得到的未经校正估计量严重偏向 $0$，而校正后的 $\hat{m}_t, \hat{v}_t$ 严格无偏：
\[
  \E[\hat{m}_t] = \mu_g, \qquad \E[\hat{v}_t] = \mu_{g^2}.
\]
\end{proposition}
\begin{proof}
反复展开一阶递推式 $m_t = \beta_1 m_{t-1} + (1-\beta_1) g_t$ 直至 $m_0 = 0$：
\[
  m_t = (1 - \beta_1) \sum_{\tau=1}^{t} \beta_1^{t-\tau} g_\tau.
\]
两端取期望，利用平稳假设与等比数列求和公式 $\sum_{k=0}^{t-1} \beta_1^k = \frac{1 - \beta_1^t}{1 - \beta_1}$：
\[
  \E[m_t] = (1 - \beta_1) \left(\sum_{\tau=1}^{t} \beta_1^{t-\tau}\right) \mu_g
  = (1 - \beta_1) \cdot \frac{1 - \beta_1^t}{1 - \beta_1}\,\mu_g
  = (1 - \beta_1^t)\,\mu_g.
\]
因此，在 $t = 1$ 时（若 $\beta_2 = 0.999$），未经校正的 $v_1 = 0.001\, g_1^2$ 竟然比真实二阶矩缩小了整整 $1000$ 倍！若不除以 $(1 - \beta_2^t)$，首步更新就会被 $\sqrt{v_1}$ 放大约 $\sqrt{1000} \approx 31.6$ 倍，引发剧烈震荡。除以精确的归一化因子 $(1-\beta_1^t)$ 与 $(1-\beta_2^t)$ 后，立得 $\E[\hat{m}_t] = \mu_g$ 与 $\E[\hat{v}_t] = \mu_{g^2}$。
\end{proof}
\end{thmBox}

\subsection{AdamW 的核心洞察：为什么 $L_2$ 正则化不等价于权重衰减？}

在经典 SGD 中，在目标函数上添加 $L_2$ 正则项 $\tilde{\mathcal{L}}(\theta) = \mathcal{L}(\theta) + \frac{\lambda}{2}\norm{\theta}_2^2$ 后，更新步为
\[
  \theta_t = \theta_{t-1} - \eta \left(g_t + \lambda \theta_{t-1}\right) = (1 - \eta\lambda)\theta_{t-1} - \eta g_t,
\]
此时“$L_2$ 正则化”与“每步将权重按比例 $(1-\eta\lambda)$ 收缩的\textbf{权重衰减（Weight Decay）}”是\textbf{完全等价}的。正因如此，早期深度学习框架（包括 PyTorch 最初的 \texttt{Adam(weight\_decay=...)}) 直接把 $\lambda \theta_{t-1}$ 粗暴加进梯度 $g_t$ 里！

Loshchilov 与 Hutter（2019）敏锐地指出了这一隐蔽而致命的错误，并提出了如今几乎统治全部大模型预训练的 \textbf{AdamW（Decoupled Weight Decay）}。

\begin{thmBox}
\begin{theorem}[Adam+$L_2$ 与 AdamW 解耦权重衰减的不等价性]
\label{thm:adamw}
设原始梯度为 $g_t = \nabla_\theta \mathcal{L}_t(\theta_{t-1})$，正则系数为 $\lambda > 0$。比较两种更新规则（为清晰起见先忽略 $\beta_1$ 一阶动量效应或令 $\beta_1 = 0$，完整带 $\beta_1$ 情形结论相同）：
\begin{itemize}
  \item \textbf{耦合版（Adam + $L_2$）}：将梯度替换为 $\tilde{g}_t = g_t + \lambda \theta_{t-1}$ 并送入 Adam，得到
  \begin{equation}
    \theta_t^{\text{L2}} = \theta_{t-1} - \eta \frac{\hat{m}_t(\tilde{g})}{\sqrt{\hat{v}_t(\tilde{g})} + \epsilon}
    \approx \theta_{t-1} - \frac{\eta \lambda}{\sqrt{\hat{v}_t(\tilde{g})} + \epsilon} \had \theta_{t-1} - \eta \frac{\hat{m}_t(g)}{\sqrt{\hat{v}_t(\tilde{g})} + \epsilon}.
    \label{eq:adam-l2}
  \end{equation}
  \item \textbf{解耦版（AdamW）}：先按原始梯度 $g_t$ 计算自适应步长，再将权重衰减项直接解耦施加于参数：
  \begin{equation}
    \theta_t^{\text{W}} = (1 - \eta \lambda)\,\theta_{t-1} - \eta \frac{\hat{m}_t(g)}{\sqrt{\hat{v}_t(g)} + \epsilon}.
    \label{eq:adamw}
  \end{equation}
\end{itemize}
在自适应预条件矩阵 $P_t = \diag(\sqrt{\hat{v}_t} + \epsilon)$ 非标量倍单位阵时，二者\textbf{绝不等价}，且 Adam+$L_2$ 会系统性扭曲各参数的正则化强度。
\end{theorem}
\begin{proof}
对比式~\eqref{eq:adam-l2} 与式~\eqref{eq:adamw} 中作用在 $\theta_{t-1}$ 上的收缩因子：
\begin{enumerate}
  \item 在 \textbf{AdamW} 中，无论某坐标 $i$ 的历史梯度大小如何，它在每一步承受的相对衰减率都是严格均匀的 $\eta\lambda$，且衰减项 $\lambda\theta_{t-1}$ \textbf{不会污染}二阶矩估计 $v_t$。
  \item 在 \textbf{Adam + $L_2$} 中，权重衰减项被迫除以自适应分母 $\sqrt{\hat{v}_{t,i}} + \epsilon$，等价于坐标 $i$ 的有效权重衰减系数变成了
  \[
    \lambda_{\text{eff}, i} = \frac{\lambda}{\sqrt{\hat{v}_{t,i}} + \epsilon}.
  \]
  这意味着：\textbf{历史梯度越大的重要权重（$\sqrt{\hat{v}_{t,i}}$ 大），受到的 $L_2$ 惩罚反而越弱}；而历史梯度微小的权重（$\sqrt{\hat{v}_{t,i}}$ 小），反而承受了剧烈放大的收缩！更糟的是，$\tilde{g}_t^2 = (g_t + \lambda\theta_{t-1})^2$ 还会把参数幅值混入二阶矩 $v_t$ 中，使最佳权重衰减系数 $\lambda$ 与学习率 $\eta$ 强耦合，调参极其痛苦。
\end{enumerate}
只有当所有坐标的 $\sqrt{\hat{v}_{t,i}}$ 恒为同一常数（即退化为无自适应缩放的 SGD）时，两者才重合。
\end{proof}
\end{thmBox}

实践中还有一个极其关键的细节：\textbf{不要对偏置项 $b$ 与归一化层的缩放/平移参数（$\gamma, \beta$）施加权重衰减}！把 $\gamma$ 往 $0$ 衰减会直接扼杀归一化层的输出通道，且一维向量参数量极小、几乎不引起过拟合。在 PyTorch 中，标准做法是按参数维度过滤，仅对二维及以上的权重矩阵（\texttt{p.ndim >= 2}）开启 \texttt{weight\_decay}。

\section{Dropout 正则化与学习率调度动力学}
\label{sec:dropout-lr}

\subsection{反向 Dropout（Inverted Dropout）的期望守恒}

在全连接层或注意力投影层中，神经元之间容易形成脆弱的“协同适应（co-adaptation）”，导致过拟合。Srivastava 等人（2014）提出的 \textbf{Dropout} 在每次前向传播时，以失活概率 $p \in [0, 1)$ 独立地将每个隐藏激活随机置零。

{\rightskip=0pt plus 1em\relax 现代框架一律采用\textbf{反向 Dropout（Inverted Dropout）}：在\textbf{训练阶段}独立采样伯努利掩码 $m_j \sim \operatorname{Bernoulli}(1 - p)$，并将存活神经元同步放大 $\frac{1}{1-p}$ 倍：\par}
\begin{equation}
  \tilde{h}_j = \frac{m_j}{1 - p}\, h_j \quad (\text{训练}), \qquad
  \tilde{h}_j = h_j \quad (\text{推理}).
  \label{eq:inv-dropout}
\end{equation}
立即验证期望守恒性：
\[
  \E_m\!\left[\tilde{h}_j\right] = \frac{\E[m_j]}{1 - p}\, h_j = \frac{1 - p}{1 - p}\, h_j = h_j.
\]
由于训练时已经完成了缩放补偿，推理阶段无需任何乘法改动（零额外开销！），同时它可近似看作对 $2^m$ 个共享权重的稀疏子网络执行几何平均集成。

\subsection{线性预热（Warmup）与余弦退火（Cosine Annealing）}

为什么几乎所有现代 Transformer 与深层 ResNet 的训练都以“先升后降”的\textbf{线性预热 + 余弦退火}调度学习率？
\begin{equation}
  \eta_t = \begin{cases}
    \eta_{\max}\,\dfrac{t}{T_{\text{warm}}}, & 1 \le t \le T_{\text{warm}}, \\[6pt]
    \eta_{\min} + \dfrac{1}{2}\left(\eta_{\max} - \eta_{\min}\right)\left(1 + \cos\!\left(\pi \dfrac{t - T_{\text{warm}}}{T_{\text{total}} - T_{\text{warm}}}\right)\right), & T_{\text{warm}} < t \le T_{\text{total}}.
  \end{cases}
  \label{eq:warmup-cosine}
\end{equation}
其背后有清晰的动力学原因：
\begin{itemize}
  \item \textbf{为什么前几百步必须预热？}在初始化瞬间，参数远离最优谷底，梯度幅度大且方向高度随机；更重要的是，对于 Adam/AdamW，尽管有偏差校正，前几步的二阶矩 $\hat{v}_t$ 仅仅基于极少数几个样本估计，统计方差极大！若一开始就施放满额学习率 $\eta_{\max}$，极易将深层权重一脚踢飞到梯度饱和区或破坏归一化尺度。线性预热给 $m_t, v_t$ 留出了稳定估计的窗口。
  \item \textbf{为什么后期采用余弦退火？}训练中段保持较大步长以快速穿越平坦高原并逃离尖锐局部极小；训练末段学习率平滑降至 $\eta_{\min} \approx 0.1\eta_{\max}$，压低小批量随机梯度噪声 $\eta^2 \Sigma_g$，让参数精细沉降到平坦极小值盆地的最深处。
\end{itemize}

下面给出一个简洁可运行的 PyTorch 实战片段，把本章全部最佳实践（Kaiming 初始化、RMSNorm、参数分组 AdamW、线性预热 + 余弦退火）浓缩在一起：

\begin{lstlisting}[caption={现代深层网络训练配置的标准模板：RMSNorm + 参数分组 AdamW + 余弦预热}]
import math, torch
import torch.nn as nn

class RMSNorm(nn.Module):
    def __init__(self, dim: int, eps: float = 1e-6):
        super().__init__()
        self.eps = eps
        self.weight = nn.Parameter(torch.ones(dim))

    def forward(self, x: torch.Tensor) -> torch.Tensor:
        # 在 fp32 下计算均方根，防止半精度平方溢出
        rms = torch.rsqrt(x.float().pow(2).mean(dim=-1, keepdim=True) + self.eps)
        return (x.float() * rms).type_as(x) * self.weight

def build_adamw_groups(model: nn.Module, lr: float = 3e-4, wd: float = 0.1):
    # 仅对矩阵权重(ndim >= 2)启用解耦衰减；偏置与 Norm 缩放参数不衰减
    decay = [p for p in model.parameters() if p.requires_grad and p.ndim >= 2]
    no_decay = [p for p in model.parameters() if p.requires_grad and p.ndim < 2]
    return torch.optim.AdamW([
        {"params": decay, "weight_decay": wd},
        {"params": no_decay, "weight_decay": 0.0},
    ], lr=lr, betas=(0.9, 0.95))

def lr_schedule(step: int, warm: int, total: int, max_lr: float, min_lr: float):
    if step < warm:
        return max_lr * (step + 1) / warm
    prog = (step - warm) / max(1, total - warm)
    return min_lr + 0.5 * (max_lr - min_lr) * (1.0 + math.cos(math.pi * prog))
\end{lstlisting}

\begin{extBox}{损失地形的平坦极小值、稳定性边缘与 $\mu$P 超参数迁移}
为什么拥有数亿参数的过参数化神经网络，在把训练损失压到几乎为零之后，依然能在测试集上泛化良好？现代深度学习理论从损失地形（Loss Landscape）的几何视角给出了三个激动人心的洞见：
\begin{enumerate}
  \item \textbf{平坦极小值 vs 尖锐极小值}：Hessian 最大特征值 $\lambda_{\max}(H)$ 越小，极小值盆地越“平坦”。由于训练损失面与测试损失面之间存在微小的水平漂移，落入尖锐窄缝的解在测试集上损失暴增，而落入宽阔平坦盆地的解对扰动高度鲁棒（Hochreiter \& Schmidhuber, 1997；Keskar et al., 2017）。较小批量与适度大初始学习率带来的随机噪声，恰好起到了“把参数震出尖锐坑洼”的隐式正则化作用。
  \item \textbf{稳定性边缘（Edge of Stability；Cohen et al., 2021）}：经典凸优化要求学习率 $\eta < 2/\lambda_{\max}(H)$ 才能单调下降；但实测表明，神经网络训练往往会持续推高曲率直到 $\lambda_{\max}(H) \approx 2/\eta$，然后贴着这一“稳定性边缘”边震荡边稳步降低损失！
  \item \textbf{最大更新参数化（$\mu$P；Yang et al., 2022）}：如果我们在 40M 参数的小模型上调好了最优学习率，把宽度扩大 $16\times$ 到 10B 模型时，传统标准参数化（SP）下的最优学习率会剧烈偏移，而每次尝试 10B 模型的成本高达数百万美元。$\mu$P 严格推导了权重初始化方差与每层学习率随宽度 $n$ 的缩放幂律（隐藏层权重方差 $\propto 1/n$、Adam 学习率 $\propto 1/n$，而输入嵌入层保持 $\Theta(1)$），使得每层激活的变化量 $\Delta h^{(l)}$ 在宽度趋于无穷时恒为 $\Theta(1)$，从而实现\textbf{零样本超参数迁移}——在小模型上调好的学习率可直接平移至百亿大模型！
\end{enumerate}
\end{extBox}

\section*{本章小结与核心脉络}
\addcontentsline{toc}{section}{本章小结与核心脉络}

本章揭开了深度网络从“算得出梯度”到“真正训得动”的动力学密码：
\begin{itemize}
  \item \textbf{初始化定基调}：Kaiming 初始化 $\Var(w) = 2/n_{\text{in}}$ 精确抵消 ReLU 的半波能量衰减，将 20 层前向标准差稳在 $O(1)$；
  \item \textbf{归一化护航程}：BatchNorm 适合大批量视觉卷积（牢记 \texttt{train/eval} 切换），LayerNorm 与更轻快的 RMSNorm 则是序列与大模型标配；
  \item \textbf{优化器控步伐}：Adam 通过偏差校正消除冷启动偏差，AdamW 将权重衰减从自适应分母中解耦，配合预热与余弦退火构成现代深度学习训练的标准引擎。
\end{itemize}

然而，即便集齐了 Kaiming 初始化、LayerNorm 与 AdamW，如果我们把纯全连接 MLP 直接喂给 $224 \times 224$ 的自然图像或千词长文本，参数量依然会瞬间失控，且 20 层以上纯直筒网络仍然面临退化瓶颈。如何把世界的\textbf{空间平移对称性}、\textbf{跨层恒等高速公路}与\textbf{时间递推结构}直接编码进网络拓扑中？这正是下一章的主题。

\section*{本章分层习题}
\addcontentsline{toc}{section}{本章分层习题}

\exTier{A 层：概念与手算热身}{检验初始化数值、归一化手算与优化器首步行为}

\begin{enumerate}
  \item[\textbf{5.2.1}] \textbf{归一化算子手算对比}：给定一个小批量特征张量 $X \in \R^{2 \times 3}$（$B=2, D=3$）：
  \[
    X = \begin{pmatrix} 1 & 2 & 3 \\ 4 & 8 & 12 \end{pmatrix},
  \]
  取 $\gamma = \mathbf{1},\; \beta = \mathbf{0},\; \epsilon = 0$。分别手算 $\BN(X)$、$\LN(X)$ 与 $\RMSN(X)$ 的精确输出矩阵，并观察第二行恰好是第一行的 $4$ 倍时，LayerNorm 与 RMSNorm 对这两行的输出有何异同？
  \exHint{注意 $\LN$ 与 $\RMSN$ 均具有正尺度不变性 $\operatorname{Norm}(\alpha x) = \operatorname{Norm}(x)$（$\alpha > 0$），故两行输出完全相同！而 $\BN$ 沿列统计，两行结果并不相同。}

  \item[\textbf{5.2.2}] \textbf{Adam 初始步与偏差校正直观}：设一维参数 $\theta_0 = 1.0$，前两步收到的梯度均为常数 $g_1 = g_2 = 0.2$，取 $\beta_1 = 0.9,\; \beta_2 = 0.999,\; \eta = 0.01,\; \epsilon = 0$。分别计算：
  \begin{enumerate}
    \item 带偏差校正的 Adam 第 1 步更新量 $\Delta \theta_1$；
    \item 去掉偏差校正（直接用 $m_1/\sqrt{v_1}$）时的第 1 步更新量，比较两者倍数差异。
  \end{enumerate}
  \exHint{带校正时 $\hat{m}_1 = g_1,\; \hat{v}_1 = g_1^2$，因此无论梯度绝对值大小如何，第 1 步步长恒为 $\pm \eta$！去掉校正后放大了 $\frac{1-\beta_1}{\sqrt{1-\beta_2}} = \frac{0.1}{\sqrt{0.001}} \approx 3.16$ 倍。}
\end{enumerate}

\exTier{B 层：定理推导与严谨证明}{深入方差传播、归一化梯度与优化器等价条件}

\begin{enumerate}
  \item[\textbf{5.2.3}] \textbf{LeakyReLU 与 SELU/GELU 的初始化方差推广}：
  \begin{enumerate}
    \item 设 $\sigma(z) = \max(a z, z)$ 为斜率 $a \in (0, 1)$ 的 LeakyReLU，严格证明在零均值对称预激活假设下 $\E[\sigma(z)^2] = \frac{1 + a^2}{2}\Var(z)$，进而推出 Kaiming 初始化方差 $\Var(w) = \frac{2}{(1+a^2)n_{\text{in}}}$；
    \item 对于标准高斯输入 $z \sim \mathcal{N}(0, 1)$，求激活函数 $\operatorname{GELU}(z) = z \Phi(z)$ 的精确二阶矩 $\E[\operatorname{GELU}(z)^2]$（给出闭式推导与数值结果），并与 ReLU 的二阶矩 $0.5$（增益 $\sqrt{2}$）比较，解释为何在单位方差处 GELU 的精确方差保持增益反而略大于 $\sqrt{2}$。
  \end{enumerate}
  \exHint{对 LeakyReLU 将积分拆为 $(-\infty, 0)$ 与 $(0, +\infty)$ 两段并利用 $p(-z) = p(z)$。}

  \item[\textbf{5.2.4}] \textbf{LayerNorm 的张量 VJP 与尺度不变梯度正交性}：对于单向量无仿射参数的标准归一化 $y = \frac{x - \mu}{\sigma}$（其中 $\mu = \frac{1}{D}\sum_i x_i,\; \sigma = \sqrt{\frac{1}{D}\sum_i (x_i-\mu)^2}$），设上游伴随量为 $\bar{y} \in \R^D$。
  \begin{enumerate}
    \item 证明反向传播公式为
    \[
      \bar{x} = \frac{1}{\sigma}\left(\bar{y} - \frac{1}{D}\sum_{j=1}^D \bar{y}_j - y \cdot \frac{1}{D}\sum_{j=1}^D \bar{y}_j y_j\right);
    \]
    \item 验证 $\sum_i \bar{x}_i = 0$ 且 $\inner{\bar{x}}{x - \mu\mathbf{1}} = 0$，从几何上解释为什么归一化层的反向梯度天然垂直于中心化输入向量！
  \end{enumerate}
  \exHint{由均值平移不变性与正尺度不变性，$y(x + c\mathbf{1}) = y(x)$ 且 $y(\alpha x) = y(x)$，对 $c$ 和 $\alpha$ 求导即可一眼看出正交性！}
\end{enumerate}

\exTier{C 层：代码实现与数值实验}{动手验证反向传播方差与 AdamW vs Adam+$L_2$}

\begin{enumerate}
  \item[\textbf{5.2.5}] \textbf{复现并扩展 20 层反向梯度传播实验}：修改 \texttt{code/vol5\_experiments.py} 中的 \texttt{exp2\_init\_norm\_opt}，在记录前向激活标准差的同时，对第 20 层施加损失 $\mathcal{L} = \norm{h^{(20)}}_2^2 / (BD)$ 并回传梯度，记录四种配置下第 $1 \sim 20$ 层权重梯度 $\norm{\nabla_{W^{(l)}} \mathcal{L}}_F$ 的演化曲线。对比 \texttt{Kaiming} 与 \texttt{Kaiming + RMSNorm} 的反向梯度稳定性。
  \exHint{将各层权重设为 \texttt{requires\_grad=True}，运行 \texttt{loss.backward()} 后读取 \texttt{W.grad.norm().item()}。}

  \item[\textbf{5.2.6}] {\rightskip=0pt plus 1em\relax \textbf{病态二次型上的 Adam+$L_2$ vs AdamW 实证}：构造二维病态二次损失 $\mathcal{L}(\theta_1, \theta_2) = 100\,\theta_1^2 + 0.01\,\theta_2^2$，从 $(\theta_1, \theta_2) = (2.0, 2.0)$ 出发，分别用 Adam（\texttt{weight\_decay=0.1}，即耦合 $L_2$）与 AdamW（\texttt{weight\_decay=0.1}）优化 200 步。打印两者的最终权重与每步有效衰减比例，验证定理~\ref{thm:adamw} 预言的不等价现象。\par}
  \exHint{注意 PyTorch 中 \texttt{torch.optim.Adam(..., weight\_decay=0.1)} 实现的是耦合 $L_2$，而 \texttt{torch.optim.AdamW} 实现的是解耦衰减；在低曲率轴 $\theta_2$ 上两者差异最为悬殊。}
\end{enumerate}

\exTier{D 层：开放探究与前沿迁移}{探索前沿归一化、无归一化训练与优化器设计}

\begin{enumerate}
  \item[\textbf{5.2.7}] \textbf{为什么反向 Dropout 与 BatchNorm 串联时可能出现“方差偏移（Variance Shift）”？}：Li 等人（2019）发现在现代 CNN 中同时串联 Dropout 与 BatchNorm 时测试精度反而下降。设特征 $x$ 经失活率为 $p$ 的反向 Dropout 得到 $\tilde{x} = \frac{m}{1-p} x$，证明虽然 $\E[\tilde{x}] = \E[x]$，但其二阶矩却被放大了 $\frac{1}{1-p}$ 倍（$\Var(\tilde{x}) = \frac{1}{1-p}\Var(x) + \frac{p}{1-p}(\E[x])^2$）。由此解释为什么 BN 在训练时记录的滑动方差 $\hat{\sigma}^2$ 与推理时（关闭 Dropout）的真实输入方差严重失配，以及为何应当把 Dropout 放在 BN 之后（或仅放在最终分类头前）。
  \exHint{利用全方差公式 $\Var(M X) = \E[M^2]\E[X^2] - (\E[M]\E[X])^2$，其中 $M = m/(1-p)$ 满足 $\E[M]=1,\; \E[M^2] = 1/(1-p)$。}

  \item[\textbf{5.2.8}] \textbf{探究低显存自适应优化器（Adafactor / 8-bit Adam / Muon）}：标准 AdamW 需要为每个参数额外保存两份 \texttt{fp32} 状态（一阶矩 $m_t$ 与二阶矩 $v_t$），使优化器显存占用高达模型权重本身的 $2\times$（在半精度训练中更是占到参数显存的 $4\times$）。阅读 Adafactor（用秩-1 矩阵分解 $v_{ij} \approx r_i c_j$ 近似二阶矩）与近期兴起的矩阵正交化优化器 Muon（利用 Newton-Schulz 迭代近似谱正交更新）的核心思想，讨论：为何矩阵权重在结构上与一维向量不同，按奇异值谱而非逐元素缩放步长能带来怎样的收敛收益？
  \exHint{对比逐元素预条件 $\diag(v_t)^{-1/2} m_t$ 与矩阵极分解 $UV^\top$（将所有非零奇异值拉平为 $1$）对矩阵条件数的不同作用。}
\end{enumerate}

\chapter[卷积、残差网络与循环序列建模]{空间与序列归纳偏置：卷积网络、残差连接与循环网络}
\label{ch:cnn-resnet-rnn}

\chapEpigraph{对称，无论你将这个词的意义定义得多么宽泛或狭隘，都是人类数千年来试图借以理解并创造秩序、美与完美的观念。}{Hermann Weyl（赫尔曼·外尔，《对称》，1952）}

\begin{chapIntro}{本章导读：把世界的对称性与高速公路写进架构}
多层感知机（MLP）虽然在理论上具备万能逼近能力，却对数据的几何结构“一无所知”：如果把一张 $224 \times 224 \times 3$ 像素的彩色图像拉直成 $150{,}528$ 维向量送入宽度为 $1024$ 的全连接隐藏层，仅仅第一层就需要 $150{,}528 \times 1024 \approx 1.54 \times 10^8$ 个参数——而且只要把图像中的猫向右平移一个像素，所有权重都得重新学习一遍！

将关于现实世界结构的先验假设（\textbf{归纳偏置，Inductive Bias}）直接铸入网络拓扑，催生了现代深度学习的三座里程碑，也是本章的核心主线：
\begin{enumerate}
  \item \textbf{空间局部性与平移等变（CNN）}：严格推导二维卷积的输出尺寸、感受野递推与参数量，证明卷积的\textbf{平移等变定理}，揭示底层硬件如何用 \texttt{im2col} 把卷积化归为稠密矩阵乘法（GEMM），并分析深度可分离卷积的加速比。
  \item \textbf{深度残差高速公路（ResNet）}：剖析深层直筒网络的退化难题，严格证明\textbf{残差梯度高速公路定理}（含 Pre-Activation 洁净路径分析），并用 24 层实测数据（\texttt{data/resnet\_grad.csv}）亲眼见证 ResNet 将底层梯度从 \texttt{1.78e-10} 拉升至 \texttt{1.01e-01} 的九个数量级跨越！
  \item \textbf{时间递推与门控高速公路（RNN / LSTM / GRU）}：完整推导随时间反向传播（BPTT）的连乘雅可比与奇异值消失/爆炸条件，揭示 LSTM 加法细胞状态如何构成“时间维度上的残差高速公路”，并指出 $O(T)$ 串行依赖瓶颈——为第 4 章 Transformer 自注意力的横空出世埋下决定性伏笔。
\end{enumerate}
\end{chapIntro}

\section{二维卷积的数学、对称性与矩阵本质}
\label{sec:cnn-math}

自然图像天然满足两条强有力的结构对称性：\textbf{局部性（Locality）}（相邻像素构成边缘、纹理等局部模式）与\textbf{平移对称性（Translation Symmetry）}（同一只猫出现在左上角或右下角，识别它所需的局部边缘检测器完全相同）。卷积神经网络（CNN；LeCun et al., 1989; Krizhevsky et al., 2012）将这两条先验编码为两大硬约束：\textbf{局部感受野}与\textbf{空间权重共享}。

\subsection{二维卷积定义、输出尺寸与感受野递推}

\begin{defBox}
\begin{definition}[多通道二维卷积层]
设输入特征图张量为 $X \in \R^{C_{\text{in}} \times H \times W}$，卷积核张量为 $K \in \R^{C_{\text{out}} \times C_{\text{in}} \times k_H \times k_W}$，偏置为 $b \in \R^{C_{\text{out}}}$，零填充（padding）为 $P$，步幅（stride）为 $S$。深度学习中习称的二维卷积（严格数学意义上为互相关，因核参数可学习，翻转与否等价）定义为：
\begin{equation}
  Y[o, i, j] = b[o] + \sum_{c=0}^{C_{\text{in}}-1} \sum_{u=0}^{k_H-1} \sum_{v=0}^{k_W-1}
  K[o, c, u, v]\, \tilde{X}\!\left[c,\, i S + u,\, j S + v\right],
  \label{eq:conv2d}
\end{equation}
其中 $\tilde{X}$ 为四周补零 $P$ 圈后的输入张量。若引入空洞率（dilation）$d \ge 1$，则等效核尺寸变为 $k_{\text{eff}} = k + (k - 1)(d - 1)$。
\end{definition}
\end{defBox}

直接清点滑动窗口的合法起止坐标，立得三个最常用的工程核心公式：
\begin{itemize}
  \item \textbf{输出空间尺寸}（以方形核 $k_H = k_W = k$ 为例）：
  \begin{equation}
    H_{\text{out}} = \left\lfloor \frac{H + 2P - k_{\text{eff}}}{S} \right\rfloor + 1, \qquad
    W_{\text{out}} = \left\lfloor \frac{W + 2P - k_{\text{eff}}}{S} \right\rfloor + 1.
    \label{eq:conv-out-size}
  \end{equation}
  当 $S = 1,\; d = 1$ 且 $k$ 为奇数时，取 $P = (k - 1)/2$ 可精确保持 $H_{\text{out}} = H$（即 \texttt{same} 填充，如 $3\times 3$ 卷积配 $P=1$）。
  \item \textbf{参数量与乘加计算量（FLOPs）}：
  \[
    \text{Params} = C_{\text{out}} C_{\text{in}} k_H k_W + C_{\text{out}}, \qquad
    \text{FLOPs} \approx 2\, H_{\text{out}} W_{\text{out}} C_{\text{out}} C_{\text{in}} k_H k_W.
  \]
  回看开篇的 $224 \times 224 \times 3$ 图像：一个 $C_{\text{out}} = 64$、核尺寸 $3\times 3$ 的卷积层仅需 $64 \times 3 \times 9 + 64 = 1{,}792$ 个参数，比全连接层暴降整整\textbf{五个数量级}，且参数量完全与图像分辨率 $(H, W)$ 无关！
  \item \textbf{感受野（Receptive Field）增长递推}：第 $l$ 层输出的一个像素能“看见”原始输入图像中多大的区域？记第 $l$ 层的核尺寸与步幅为 $k_l, s_l$，输入层初始感受野 $r_0 = 1$，累计下采样步幅（jump）为 $J_l = \prod_{i=1}^{l} s_i$（$J_0 = 1$），则有精确递推：
  \begin{equation}
    r_l = r_{l-1} + (k_l - 1)\prod_{i=1}^{l-1} s_i = r_{l-1} + (k_l - 1) J_{l-1}.
    \label{eq:receptive-field}
  \end{equation}
\end{itemize}

式~\eqref{eq:receptive-field} 立刻揭示了 VGGNet（Simonyan \& Zisserman, 2015）的设计精髓：连续堆叠\textbf{三层}步幅 $S=1$ 的 $3\times 3$ 卷积，感受野依次增长为 $1 \to 3 \to 5 \to 7$，与单层 $7\times 7$ 卷积完全相同；但若通道数恒为 $C$，三层 $3\times 3$ 的参数量仅为 $3 \times (9 C^2) = 27 C^2$，而单层 $7\times 7$ 需要 $49 C^2$（节省 $45\%$ 参数），同时三层之间还夹入了\textbf{三次非线性 ReLU 激活}，表达力大幅增强！

\subsection{平移等变性定理与池化的局部不变性}

为什么卷积比全连接层在视觉任务上具备天然的样本效率？核心在于群论意义下的\textbf{平移等变性（Translation Equivariance）}：把输入图像平移向量 $\delta$ 后再做卷积，与先做卷积再把输出特征图平移 $\delta$ 完全等价。

\begin{thmBox}
\begin{theorem}[卷积的平移等变性]
\label{thm:conv-equivariance}
在无边界截断的无限二维网格 $\mathbb{Z}^2$ 上，定义平移算子 $(T_\delta f)(x) := f(x - \delta)$（$\delta \in \mathbb{Z}^2$）与离散卷积 $(f * g)(x) := \sum_{u \in \mathbb{Z}^2} f(x - u)\, g(u)$（或互相关 $\sum_u f(x + u) g(u)$）。则卷积算子与任意平移算子可交换：
\begin{equation}
  (T_\delta f) * g = T_\delta (f * g).
  \label{eq:equivariance}
\end{equation}
\end{theorem}
\begin{proof}
对任意坐标 $x \in \mathbb{Z}^2$ 直接按定义展开左端：
\begin{align*}
  \big((T_\delta f) * g\big)(x)
  &= \sum_{u \in \mathbb{Z}^2} (T_\delta f)(x - u)\, g(u)
   = \sum_{u \in \mathbb{Z}^2} f\!\left((x - \delta) - u\right) g(u) \\
  &= (f * g)(x - \delta)
   = \big(T_\delta (f * g)\big)(x).
\end{align*}
（对于互相关定义 $(f \star g)(x) = \sum_u f(x+u) g(u)$，同样有 $\sum_u f(x-\delta+u)g(u) = T_\delta(f \star g)(x)$，结论完全成立。）
\end{proof}
\end{thmBox}

务必精准区分\textbf{等变性（Equivariance）}与\textbf{不变性（Invariance）}：
\begin{itemize}
  \item \textbf{卷积提供平移等变性}：目标检测与语义分割需要知道物体\emph{在哪里}，猫向右移 10 像素，对应的特征激活峰值也同步向右移 10 像素。
  \item \textbf{池化（Pooling）与全局平均池化（GAP）提供近似平移不变性}：对分类任务，猫在左侧还是右侧类别标签都是“猫”。局部最大池化（MaxPool）在窗口内抹平微小抖动，而网络末端的全局平均池化 $\bar{y}_c = \frac{1}{HW}\sum_{h,w} Y_{c,h,w}$ 对空间位置求和，从而实现了整体平移不变性。（注意：带步幅 $S > 1$ 的下采样仅保留 $S$ 倍整数网格的离散等变性，且零填充会在图像边缘微弱破坏完美等变性。）
\end{itemize}

\subsection{从嵌套循环到硬件加速：\texttt{im2col} 的 GEMM 等价与深度可分离卷积}

如果按照定义式~\eqref{eq:conv2d} 写出六层嵌套 \texttt{for} 循环，在 GPU 上运行速度会慢上千倍。现代深度学习底层库（cuDNN、oneDNN）实现高性能卷积的经典核心之一，是 \textbf{\texttt{im2col}（Image-to-Column）变换}——它把二维滑动窗口卷积严丝合缝地变成一次大型稠密矩阵乘法（GEMM），从而吃满 GPU 的张量核心（Tensor Cores，详见第 6 卷第 5 章）！

具体构造极其优雅：
\begin{enumerate}
  \item \textbf{展开输入特征图}：对输出特征图的每一个空间位置 $(i, j)$（共有 $L = H_{\text{out}} W_{\text{out}}$ 个位置），取出对应输入感受野里的 $C_{\text{in}} \times k_H \times k_W$ 局部块，拉直成一列（长度为 $K = C_{\text{in}} k_H k_W$）。把所有 $L$ 列并排拼成矩阵 $X_{\text{col}} \in \R^{(C_{\text{in}} k_H k_W) \times (H_{\text{out}} W_{\text{out}})}$。
  \item \textbf{拉平卷积核}：将四维卷积核张量 $K \in \R^{C_{\text{out}} \times C_{\text{in}} \times k_H \times k_W}$ 直接重塑（零拷贝 \texttt{reshape}）为二维权重矩阵 $W_{\text{row}} \in \R^{C_{\text{out}} \times (C_{\text{in}} k_H k_W)}$。
  \item \textbf{一次 GEMM 完结}：矩阵乘积
  \[
    Y_{\text{mat}} = W_{\text{row}} X_{\text{col}} + b \mathbf{1}_L^\top \in \R^{C_{\text{out}} \times (H_{\text{out}} W_{\text{out}})}
  \]
  的每一项恰好就是式~\eqref{eq:conv2d}！最后将 $Y_{\text{mat}}$ 重塑回 $(C_{\text{out}}, H_{\text{out}}, W_{\text{out}})$ 即可。而且它的反向传播 VJP 直接复用第~\ref{ch:mlp-autodiff} 章命题~\ref{prop:linear-vjp}：$\bar{W}_{\text{row}} = \bar{Y}_{\text{mat}} X_{\text{col}}^\top$，以及 $\bar{X} = \operatorname{col2im}(W_{\text{row}}^\top \bar{Y}_{\text{mat}})$（其中 \texttt{col2im} 将重叠像素处的梯度累加，再次印证多分支梯度累加原则！）。
\end{enumerate}

\textbf{深度可分离卷积（Depthwise Separable Convolution；MobileNet）}进一步将标准卷积因式分解为两步：先做\textbf{逐通道深度卷积（Depthwise Conv）}（每个通道独立使用一个 $k \times k$ 核，不混合通道，耗时 $H_{\text{out}} W_{\text{out}} C_{\text{in}} k^2$），再做\textbf{逐点 $1\times 1$ 卷积（Pointwise Conv）}（纯粹在每个像素上做通道线性混合，耗时 $H_{\text{out}} W_{\text{out}} C_{\text{in}} C_{\text{out}}$）。两者总计算量与标准卷积之比为：
\begin{equation}
  \frac{H_{\text{out}} W_{\text{out}} C_{\text{in}} k^2 + H_{\text{out}} W_{\text{out}} C_{\text{in}} C_{\text{out}}}{H_{\text{out}} W_{\text{out}} C_{\text{in}} C_{\text{out}} k^2}
  = \frac{1}{C_{\text{out}}} + \frac{1}{k^2}.
  \label{eq:dw-sep-ratio}
\end{equation}
当 $k = 3,\; C_{\text{out}} = 256$ 时，计算量与参数量骤降至原来的 $\frac{1}{256} + \frac{1}{9} \approx 11.5\%$——足足提速近 \textbf{9 倍}！

\section{ResNet 与残差梯度高速公路定理}
\label{sec:resnet}

手握 Kaiming 初始化、BatchNorm 与 $3\times 3$ 卷积之后，2015 年的研究者满怀期待地把 VGG 式的直筒网络从 20 层堆到了 56 层，却遭遇了一个令人困惑的奇异现象——\textbf{退化问题（Degradation Problem）}：56 层直筒网络不仅\emph{测试误差}更高，连\textbf{训练误差都显著高于 20 层网络}！

这绝不是过拟合（过拟合应当表现为训练误差更低），而是一个纯粹的\textbf{优化难题}：理论上，只要把多出来的 36 层设为恒等映射 $h(x) = x$，56 层网络的训练误差至少应当不差于 20 层网络；但实践证明，让一堆非线性层去逼近“什么都不做的恒等映射”竟然出奇地困难！

He 等人（2016）提出的\textbf{残差网络（Residual Network, ResNet）}用一行看似平淡无奇的加法改写了整个深度学习史：与其让堆叠层直接拟合目标映射 $\mathcal{H}(x)$，不如让它拟合\textbf{残差（Residual）} $\mathcal{F}(x) = \mathcal{H}(x) - x$，再通过一条无参数的\textbf{恒等跳跃连接（Identity Shortcut）}把输入直接加到输出上：
\[
  x_{l+1} = x_l + \mathcal{F}(x_l; W_l).
\]
当最优映射接近恒等映射时，网络只需把残差支路权重 $W_l$ 推向 $0$ 即可轻松实现恒等——而这仅仅是故事的表面。它更惊人的魔力藏在反向传播里。

\begin{thmBox}
\begin{theorem}[残差网络的梯度高速公路定理（He et al., 2016）]
\label{thm:resnet-highway}
考虑由 $L$ 个洁净预激活（Pre-Activation）残差块串联而成的深层网络，其递推满足 $x_{i+1} = x_i + \mathcal{F}(x_i; W_i)$（$i = 0, \dots, L-1$）。对任意浅层 $l$ 与任意更深层 $L > l$：
\begin{enumerate}
  \item \textbf{前向展开为残差和}：
  \begin{equation}
    x_L = x_l + \sum_{i=l}^{L-1} \mathcal{F}(x_i; W_i).
    \label{eq:resnet-fwd}
  \end{equation}
  \item \textbf{反向梯度分解为恒等直达项与残差调制项}：对任意标量损失 $\mathcal{L}$，
  \begin{equation}
    \frac{\partial \mathcal{L}}{\partial x_l}
    = \frac{\partial \mathcal{L}}{\partial x_L}\left(I + \frac{\partial}{\partial x_l}\sum_{i=l}^{L-1} \mathcal{F}(x_i; W_i)\right),
    \label{eq:resnet-bwd}
  \end{equation}
  （这里采用行向量梯度约定 $\frac{\partial \mathcal{L}}{\partial x} \in \R^{1 \times D}$；若采用列向量梯度 $\nabla_x \mathcal{L} \in \R^D$，则等价写作 $\nabla_{x_l}\mathcal{L} = \left(I + \frac{\partial}{\partial x_l}\sum_{i=l}^{L-1}\mathcal{F}(x_i; W_i)\right)^\top \nabla_{x_L}\mathcal{L}$）。
\end{enumerate}
\end{theorem}
\begin{proof}
对 $x_{i+1} - x_i = \mathcal{F}(x_i; W_i)$ 从 $i = l$ 到 $L-1$ 望远镜求和（telescoping sum），所有中间项两两抵消，立得式~\eqref{eq:resnet-fwd}。

对式~\eqref{eq:resnet-fwd} 两端关于 $x_l$ 应用多元链式法则：
\[
  \frac{\partial \mathcal{L}}{\partial x_l}
  = \frac{\partial \mathcal{L}}{\partial x_L} \cdot \frac{\partial x_L}{\partial x_l}
  = \frac{\partial \mathcal{L}}{\partial x_L} \left(\frac{\partial x_l}{\partial x_l} + \frac{\partial}{\partial x_l}\sum_{i=l}^{L-1}\mathcal{F}(x_i; W_i)\right)
  = \frac{\partial \mathcal{L}}{\partial x_L} \left(I + \frac{\partial}{\partial x_l}\sum_{i=l}^{L-1}\mathcal{F}(x_i; W_i)\right).
\]
进一步，若逐层连乘局部雅可比 $\frac{\partial x_L}{\partial x_l} = \prod_{i=l}^{L-1} \left(I + J_{\mathcal{F}_i}\right)$ 并按二项式展开，它正是从第 $L$ 层到第 $l$ 层全部 $2^{L-l}$ 条不同长度路径（从直接跳过的长度 $0$ 恒等路径，到经过 $1, 2, \dots, L-l$ 个残差支路的组合路径）的雅可比之和！
\end{proof}
\end{thmBox}

仔细凝视式~\eqref{eq:resnet-bwd}，它包含两重深刻启示：
\begin{enumerate}
  \item \textbf{梯度绝不会系统性消失}：在普通直筒网络中，梯度被连乘矩阵 $\prod J_i$ 反复挤压；而在 ResNet 中，顶层梯度 $\frac{\partial \mathcal{L}}{\partial x_L}$ 通过恒等项 $I$ \textbf{畅通无阻地直接空降到任意浅层 $l$}，完全绕开了所有中间权重矩阵！即便残差支路的雅可比项非常小，只要它不恰好对整个小批量恒等于 $-I$（这在随机权重与数据变化下几乎不可能发生），浅层梯度就始终保持饱满的 $O(1)$ 量级——这就是“残差梯度高速公路”的真正含义。
  \item \textbf{为什么 Pre-Activation ResNet 与 Pre-LN Transformer 更优越？}最初的 ResNet v1 在加法之后又套了一层 ReLU：$x_{l+1} = \ReLU(x_l + \mathcal{F}(x_l))$，这条 ReLU 像收费站一样挡在了主干道上，破坏了完美的望远镜求和式~\eqref{eq:resnet-fwd}；He 等人在 ResNet v2（Identity Mappings）中把 BN 和 ReLU 挪进残差支路内部（\texttt{BN -> ReLU -> Conv -> BN -> ReLU -> Conv}），让主干道变成纯粹的加法 $x_{l+1} = x_l + \mathcal{F}(x_l)$，使 1000 层网络也能丝滑收敛！你会在第 4 章立刻发现，当今所有大语言模型（GPT-3、LLaMA、DeepSeek）所采用的 \textbf{Pre-Norm 残差流} $x \leftarrow x + \operatorname{Attn}(\operatorname{Norm}(x))$，正是 ResNet v2 洁净高速公路的直系传人！
\end{enumerate}

在工程细节上，ResNet-50/101/152 为了在高通道数下控制计算量，采用了\textbf{瓶颈块（Bottleneck Block）}：$1\times 1$ 降维（如 $256 \to 64$）$\to 3\times 3$ 空间卷积（$64 \to 64$）$\to 1\times 1$ 升维（$64 \to 256$）；当跨阶段下采样导致空间尺寸或通道数变化时，在跳跃连接上放置一个 $1\times 1$ 线性投影 $W_s$ 对齐维度：$x_{l+1} = W_s x_l + \mathcal{F}(x_l)$。

\section{实测验证：24 层直筒网络 vs ResNet 梯度实测}
\label{sec:exp-resnet}

为了定量验证定理~\ref{thm:resnet-highway}，我们运行 \texttt{code/vol5\_experiments.py} 中的 \texttt{exp3\_resnet\_grad\_highway}：构建一个 $L = 24$ 层、宽 $D = 64$、批量 $B = 32$ 的深度网络，每层权重独立采样自 $W_l \sim \mathcal{N}(0, (0.65/\sqrt{D})^2)$（故意模拟略低于临界点的收缩权重，常见于训练中后期或饱和非线性），对比：
\begin{itemize}
  \item \textbf{Plain 直筒网络}：$h_l = \tanh(h_{l-1} W_l)$；
  \item \textbf{ResNet 残差网络}：$h_l = h_{l-1} + 0.35\,\tanh(h_{l-1} W_l)$。
\end{itemize}
在顶层施加均方损失 $\mathcal{L} = \frac{1}{BD}\norm{h_{24}}_F^2$ 并执行一次反向传播，读取每一层激活梯度 $\norm{\nabla_{h_l} \mathcal{L}}_F$。实测文件 \texttt{data/resnet\_grad.csv} 给出了触目惊心的对比（绘制于图~\ref{fig:resnet-grad}）：
\begin{itemize}
  \item {\rightskip=0pt plus 1em\relax \textbf{24 层 Plain 直筒网络}：由于前向激活已先经历 24 层收缩，第 24 层顶层梯度仅为 \textbf{\texttt{9.10e-07}}（\texttt{9.1049e-07}）；沿着反向路径回传 23 层后，梯度按指数律狂泻，到达第 1 层时仅剩 \textbf{\texttt{1.78e-10}}（\texttt{1.7770e-10}）——底层权重几乎完全冻结！\par}
  \item \textbf{24 层 ResNet 残差网络}：第 24 层顶层梯度为 \textbf{\texttt{5.91e-02}}（\texttt{5.9071e-02}），沿着恒等高速公路回传至第 1 层时不仅没有衰减，还稳健维持在 \textbf{\texttt{1.01e-01}}（\texttt{1.0068e-01}）——第 1 层梯度幅值比直筒网络足足高出\textbf{近九个数量级（$5.7 \times 10^8$ 倍）}！
\end{itemize}

\begin{figure}[htbp]
\centering
\begin{tikzpicture}
  \begin{semilogyaxis}[
    width=0.86\linewidth, height=6.5cm,
    xlabel={网络层号 $l$（梯度由第 24 层反向传向第 1 层）},
    ylabel={激活梯度范数 $\norm{\partial \mathcal{L}/\partial h_l}_F$（对数刻度）},
    xmin=0.5, xmax=24.5, ymin=5e-11, ymax=5e-1,
    xtick={1,4,8,12,16,20,24},
    grid=both, grid style={tealLine!45, dashed},
    axis line style={slate},
    tick label style={font=\scriptsize},
    label style={font=\footnotesize\sffamily, color=ink},
    legend style={at={(0.97,0.70)}, anchor=east, font=\scriptsize, fill=white, draw=tealLine, legend cell align=left}
  ]
    \addplot[teal, very thick, mark=*, mark size=1.7pt]
      table[x=layer, y=resnet_grad, col sep=comma]{data/resnet_grad.csv};
    \addlegendentry{ResNet（24 层）：$5.91\times 10^{-2}$ 至 $1.01\times 10^{-1}$}

    \addplot[ochre, very thick, mark=square*, mark size=1.6pt]
      table[x=layer, y=plain_grad, col sep=comma]{data/resnet_grad.csv};
    \addlegendentry{Plain（24 层）：$9.10\times 10^{-7}$ 至 $1.78\times 10^{-10}$}

    % Annotate 9 orders of magnitude gap at layer 1
    \draw[{Stealth[length=2mm]}-{Stealth[length=2mm]}, navy, thick]
      (axis cs:2.0, 2.5e-10) -- (axis cs:2.0, 6.5e-2)
      node[midway, right=2pt, font=\scriptsize\sffamily\bfseries, color=navy, fill=white, inner sep=1.5pt, align=left] {相差约 $5.7\times 10^8$ 倍\\（近 9 个数量级）};
  \end{semilogyaxis}
\end{tikzpicture}
\caption{24 层 Plain 直筒网络与 ResNet 反向激活梯度范数实测对比（读取自 \texttt{data/resnet\_grad.csv}）。直筒网络在半对数坐标中呈指数级滑坡，而 ResNet 依托恒等跳跃连接将梯度完整送抵第 1 层。}
\label{fig:resnet-grad}
\end{figure}

\section{循环神经网络：BPTT、梯度病理与 LSTM 时间高速公路}
\label{sec:rnn-lstm}

自然语言、音频波形与时间序列具有与图像截然不同的归纳偏置：\textbf{序列长度 $T$ 可变}，且信息按时间步 $t = 1, \dots, T$ 因果演化。循环神经网络（RNN）的思路是在\textbf{时间维度上共享同一组权重}，维护一个动态隐藏状态 $h_t \in \R^d$ 来摘要迄今为止的全部历史。

\subsection{Elman RNN 与随时间反向传播（BPTT）的精确推导}

标准单层 Elman RNN 的前向递推为（初始状态 $h_0 = \mathbf{0}$）：
\begin{equation}
  a_t = W_{xh} x_t + W_{hh} h_{t-1} + b_h, \qquad
  h_t = \tanh(a_t), \qquad
  \hat{y}_t = W_{hy} h_t + b_y,
  \label{eq:elman-rnn}
\end{equation}
其中总损失为各时间步损失之和 $\mathcal{L} = \sum_{t=1}^{T} \mathcal{L}_t(\hat{y}_t, y_t)$。由于同一矩阵 $W_{hh}$ 在每一个时间步 $k = 1, \dots, T$ 都被复用，沿时间展开后的计算图就是一个深度为 $T$、权重共享的 DAG！

应用第~\ref{ch:mlp-autodiff} 章的多分支梯度累加公式~\eqref{eq:vjp}，第 $t$ 步的损失 $\mathcal{L}_t$ 对循环权重 $W_{hh}$ 的贡献，必须累加它在每一个更早时间步 $k \le t$ 上的直接作用：
\begin{equation}
  \frac{\partial \mathcal{L}}{\partial W_{hh}}
  = \sum_{t=1}^{T} \sum_{k=1}^{t}
    \frac{\partial \mathcal{L}_t}{\partial h_t}
    \left(\prod_{j=k+1}^{t} \frac{\partial h_j}{\partial h_{j-1}}\right)
    \frac{\partial^+ h_k}{\partial W_{hh}},
  \label{eq:bptt}
\end{equation}
其中 $\frac{\partial^+ h_k}{\partial W_{hh}}$ 表示固定 $h_{k-1}$ 时 $h_k$ 对第 $k$ 步所用 $W_{hh}$ 的“直接局部偏导”。

现在盯紧式~\eqref{eq:bptt} 中间的核心连乘因子——跨时间步 $(t - k)$ 的状态传输雅可比：
\begin{equation}
  \frac{\partial h_j}{\partial h_{j-1}} = \diag\!\left(1 - h_j^2\right) W_{hh},
  \qquad\Longrightarrow\qquad
  \frac{\partial h_T}{\partial h_t} = \prod_{j=t+1}^{T} \diag\!\left(1 - h_j^2\right) W_{hh},
  \label{eq:bptt-jacobian}
\end{equation}
（注意此处 $h_j = \tanh(W_{hh} h_{j-1} + \dots)$ 为列向量，故雅可比矩阵第 $(p,q)$ 项为 $\frac{\partial (h_j)_p}{\partial (h_{j-1})_q} = (1 - (h_j)_p^2)(W_{hh})_{pq}$，即 $\diag(1 - h_j^2) W_{hh}$；在左乘行向量或对列向量伴随量回传时，对应转置 $W_{hh}^\top \diag(1 - h_j^2)$。）

\begin{thmBox}
\begin{proposition}[BPTT 长期依赖的奇异值判定条件（Pascanu et al., 2013）]
\label{prop:bptt-spectral}
设 $\gamma = \sup_a \abs{\tanh'(a)} = 1$，记 $W_{hh}$ 的最大奇异值为 $\sigma_{\max}(W_{hh}) = \norm{W_{hh}}_2$。则跨越 $\Delta t = T - t$ 个时间步的梯度传播范数满足上界：
\begin{equation}
  \norm{\frac{\partial h_T}{\partial h_t}}_2
  \le \prod_{j=t+1}^{T} \norm{\diag(1 - h_j^2)}_2 \norm{W_{hh}}_2
  \le \left(\gamma\,\sigma_{\max}(W_{hh})\right)^{T - t}.
  \label{eq:bptt-bound}
\end{equation}
由此：
\begin{itemize}
  \item 若 $\sigma_{\max}(W_{hh}) < 1$（再叠加 $\tanh'$ 在非零处严格小于 $1$），长期梯度必然以指数速度\textbf{消失}，网络只能学到最近几步的局部关联，彻底遗忘远距离上下文；
  \item 反之，若 $\sigma_{\max}(W_{hh}) > 1$（且特征方向未被 $\tanh'$ 压制），梯度沿时间呈指数\textbf{爆炸}，引发剧烈数值震荡。
\end{itemize}
\end{proposition}
\end{thmBox}

面对梯度的两极病态，工程上有截然不同的对策：
\begin{itemize}
  \item \textbf{对抗梯度爆炸：全局梯度范数裁剪（Gradient Norm Clipping）}。给定阈值 $\tau > 0$（如 $\tau = 1.0$），计算全体参数梯度拼成的总欧氏范数 $\norm{g}_2$，若超标则沿原方向等比缩短：
  \[
    g \leftarrow \frac{\tau}{\max(\tau, \norm{g}_2)}\, g.
  \]
  注意：梯度裁剪保持梯度\textbf{方向不变}而仅限制步长上限，是解决梯度爆炸（以及大模型训练中偶发损失尖刺）的特效药；但对已经跌成 $10^{-10}$ 的梯度消失却无能为力！
  \item \textbf{对抗梯度消失：重构时间拓扑——LSTM 与 GRU 门控高速公路}！
\end{itemize}

\subsection{LSTM 与 GRU：时间维度上的残差高速公路}

早在 ResNet 诞生近二十年前，Hochreiter 与 Schmidhuber（1997）就在\textbf{长短期记忆网络（LSTM）}中独立发现了完全同构的数学思想！LSTM 开辟了一条独立于隐藏状态 $h_t$ 的\textbf{细胞状态主干道（Cell State）} $c_t \in \R^d$，并设置三个受 Sigmoid 控制的软门控（值域 $(0,1)$）：\textbf{遗忘门 $f_t$}、\textbf{输入门 $i_t$} 与\textbf{输出门 $o_t$}：
\begin{align}
  f_t &= \sigma\!\left(W_f [h_{t-1}, x_t] + b_f\right), &
  i_t &= \sigma\!\left(W_i [h_{t-1}, x_t] + b_i\right), \label{eq:lstm-gates1} \\
  \tilde{c}_t &= \tanh\!\left(W_c [h_{t-1}, x_t] + b_c\right), &
  o_t &= \sigma\!\left(W_o [h_{t-1}, x_t] + b_o\right), \label{eq:lstm-gates2} \\
  c_t &= f_t \had c_{t-1} + i_t \had \tilde{c}_t, &
  h_t &= o_t \had \tanh(c_t). \label{eq:lstm-cell}
\end{align}

请盯住式~\eqref{eq:lstm-cell} 中细胞状态的\textbf{加法更新} $c_t = f_t \had c_{t-1} + i_t \had \tilde{c}_t$！计算它相邻时间步的直接导数（忽略经由 $h_{t-1}$ 绕回门控的次级路径）：
\[
  \frac{\partial c_t}{\partial c_{t-1}} = \diag(f_t) + \text{（门控间接项）}.
\]
将它与 ResNet 的 $x_{l+1} = x_l + \mathcal{F}(x_l)$ 并排放在一起（见图~\ref{fig:arch-compare}）！你瞬间就会豁然开朗：\textbf{LSTM 的细胞状态 $c_t$ 正是时间轴上的加性门控残差高速公路}！
\begin{itemize}
  \item 当网络需要长距离记住某个早期主语或关键事实时，只需学到遗忘门 $f_t \approx \mathbf{1}$（因此实践中常将遗忘门偏置 $b_f$ 初始化为 $+1.0$ 甚至 $+2.0$，让初始阶段默认保留记忆！），细胞状态梯度就能以接近 $1$ 的因子沿时间无损回传数十上百步；
  \item 当语义切换到新段落时，网络将 $f_t \approx \mathbf{0}$ 清空旧记忆，同时打开输入门 $i_t$ 写入新内容。
\end{itemize}

Cho 等人（2014）提出的\textbf{门控循环单元（GRU）}把细胞状态与隐藏状态合二为一，用仅两个门（更新门 $z_t$ 与重置门 $r_t$）实现几乎同等的长期记忆效果，参数量减少 $25\%$：
\[
  z_t = \sigma(W_z[h_{t-1}, x_t]), \quad
  r_t = \sigma(W_r[h_{t-1}, x_t]), \quad
  \tilde{h}_t = \tanh(W_h[r_t \had h_{t-1}, x_t]), \quad
  h_t = (1 - z_t) \had h_{t-1} + z_t \had \tilde{h}_t.
\]

\begin{figure}[htbp]
\centering
\resizebox{\linewidth}{!}{%
\begin{tikzpicture}[
  box/.style={rectangle, rounded corners=1.4mm, draw=teal, thick, fill=tealLight,
              minimum width=23mm, minimum height=7.5mm, align=center, font=\sffamily\scriptsize, text=ink},
  op/.style={circle, draw=navy, thick, fill=navyLight, inner sep=1.6pt, font=\scriptsize\bfseries, text=navy},
  gate/.style={rectangle, rounded corners=1.2mm, draw=ochre, thick, fill=ochreLight,
               minimum width=11mm, minimum height=6mm, align=center, font=\scriptsize\bfseries, text=ochre},
  arr/.style={-{Stealth[length=2mm]}, thick, slate},
  hwy/.style={-{Stealth[length=2.3mm]}, very thick, tealDark}
]
  % ===== (a) Plain Block =====
  \begin{scope}[shift={(0,0)}]
    \node[font=\sffamily\bfseries\small, color=slate] at (1.3, 4.9) {(a) Plain 直筒块};
    \node[font=\scriptsize, color=ink] (px0) at (1.3, 0.2) {$x_l$};
    \node[box] (p1) at (1.3, 1.4) {Conv / Linear + BN};
    \node[box, fill=ochreLight, draw=ochre] (p2) at (1.3, 2.55) {ReLU / $\tanh$};
    \node[box] (p3) at (1.3, 3.7) {Conv / Linear + BN};
    \node[font=\scriptsize, color=ink] (px1) at (1.3, 4.45) {$x_{l+1}$};
    \draw[arr] (px0) -- (p1);
    \draw[arr] (p1) -- (p2);
    \draw[arr] (p2) -- (p3);
    \draw[arr] (p3) -- (px1);
    \node[font=\scriptsize, color=ochre, align=center] at (1.3, -0.55) {梯度连乘 $\prod J_i$\\深层指数衰减};
  \end{scope}

  % ===== (b) Pre-Activation ResNet Block =====
  \begin{scope}[shift={(4.4,0)}]
    \node[font=\sffamily\bfseries\small, color=tealDark] at (1.7, 4.9) {(b) Pre-Act ResNet 块};
    \node[font=\scriptsize, color=ink] (rx0) at (0.3, 0.2) {$x_l$};
    \node[box, minimum width=22mm] (r1) at (2.3, 1.3) {BN $\to$ ReLU};
    \node[box, minimum width=22mm] (r2) at (2.3, 2.3) {Conv ($W_1$)};
    \node[box, minimum width=22mm] (r3) at (2.3, 3.3) {BN $\to$ ReLU $\to$ Conv};
    \node[op] (radd) at (0.3, 3.95) {$+$};
    \node[font=\scriptsize, color=ink] (rx1) at (0.3, 4.55) {$x_{l+1}$};

    % Clean identity highway on the left
    \draw[hwy] (rx0) -- (radd);
    \draw[hwy] (radd) -- (rx1);
    \node[rotate=90, font=\tiny\sffamily\bfseries, color=tealDark] at (-0.08, 2.15) {恒等梯度高速公路 ($I$)};

    % Residual branch
    \draw[arr] (0.3, 0.65) -| (r1.south);
    \draw[arr] (r1) -- (r2);
    \draw[arr] (r2) -- (r3);
    \draw[arr] (r3.north) |- (radd.east);
    \node[font=\scriptsize, color=tealDark, align=center] at (1.5, -0.55) {空间深度主干道\\$\frac{\partial x_L}{\partial x_l} = I + \frac{\partial}{\partial x_l}\sum \mathcal{F}_i$};
  \end{scope}

  % ===== (c) LSTM Cell State Highway =====
  \begin{scope}[shift={(9.3,0)}]
    \node[font=\sffamily\bfseries\small, color=navy] at (3.6, 4.9) {(c) LSTM 细胞状态时间高速公路};
    % Top cell state highway
    \node[font=\scriptsize\bfseries, color=tealDark] (ct0) at (0.1, 3.85) {$c_{t-1}$};
    \node[op] (cmul) at (1.7, 3.85) {$\had$};
    \node[op] (cadd) at (4.1, 3.85) {$+$};
    \node[font=\scriptsize\bfseries, color=tealDark] (ct1) at (6.8, 3.85) {$c_t$};

    \draw[hwy] (ct0) -- (cmul);
    \draw[hwy] (cmul) -- (cadd);
    \draw[hwy] (cadd) -- (ct1);
    \node[font=\tiny\sffamily\bfseries, color=tealDark] at (3.4, 4.25) {时间加性高速公路（$f_t \approx 1$ 时梯度直通）};

    % Gates below (generous spacing, zero touching)
    \node[gate] (gf) at (1.7, 1.8) {$f_t$ ($\sigma$)};
    \node[gate] (gi) at (3.2, 1.8) {$i_t$ ($\sigma$)};
    \node[gate, draw=teal, fill=tealLight, text=tealDark] (gc) at (4.8, 1.8) {$\tilde{c}_t$ ($\tanh$)};
    \node[gate] (go) at (6.4, 1.8) {$o_t$ ($\sigma$)};

    \node[op] (imul) at (4.1, 2.8) {$\had$};
    \node[op] (omul) at (6.4, 2.8) {$\had$};
    \node[font=\scriptsize, color=ink] (ht0) at (0.1, 0.75) {$h_{t-1}, x_t$};
    \node[font=\scriptsize\bfseries, color=navy] (ht1) at (7.3, 2.8) {$h_t$};

    \draw[arr] (ht0.east) -- (6.4, 0.75);
    \draw[arr] (1.7, 0.75) -- (gf.south);
    \draw[arr] (3.2, 0.75) -- (gi.south);
    \draw[arr] (4.8, 0.75) -- (gc.south);
    \draw[arr] (6.4, 0.75) -- (go.south);

    \draw[arr] (gf.north) -- (cmul.south);
    \draw[arr] (gi.north) |- (imul.west);
    \draw[arr] (gc.north) |- (imul.east);
    \draw[arr] (imul.north) -- (cadd.south);
    \draw[arr] (go.north) -- (omul.south);
    \draw[arr] (5.55, 3.85) |- node[pos=0.28, right, font=\tiny, color=slate] {$\tanh$} (omul.west);
    \draw[arr] (omul.east) -- (ht1.west);

    \node[font=\scriptsize, color=navy, align=center] at (3.6, -0.55) {时间深度主干道：$c_t = f_t \had c_{t-1} + i_t \had \tilde{c}_t$\\仍受限于 $O(T)$ 串行计算依赖};
  \end{scope}
\end{tikzpicture}}
\caption{深度学习中两条伟大“加性高速公路”的结构同构对比：(a) Plain 直筒块梯度连乘衰减；(b) Pre-Activation ResNet 在\textbf{层深度维度}打开恒等高速公路；(c) LSTM 通过细胞状态 $c_t$ 在\textbf{时间维度}打开门控加性高速公路。}
\label{fig:arch-compare}
\end{figure}

\subsection{循环网络的两个根本瓶颈：为何历史走向了 Transformer？}

尽管 LSTM 与 GRU 在 2014--2017 年统治了机器翻译与语音识别，它们却存在两道难以逾越的物理天花板：
\begin{enumerate}
  \item \textbf{无法沿序列并行训练（$O(T)$ 串行墙）}：在式~\eqref{eq:lstm-cell} 中，计算第 $t$ 步的 $(h_t, c_t)$ 必须严格等待第 $t-1$ 步完成。面对长度 $T = 4096$ 的文本段落，拥有数万个并行核心的现代 GPU 只能被迫一步步串行推进，算力利用率极低，根本无法扩展到万亿 token 的预训练规模。
  \item \textbf{固定维度信息瓶颈与 $O(T)$ 交互路径}：无论前面已经读入了一百个词还是一万个词，RNN/LSTM 必须把全部历史塞进同一个固定维度向量 $(h_t, c_t) \in \R^{2d}$ 中；相距 $T$ 步的两个词之间交互必须历经 $T$ 次非线性递推，信息损耗不可避免。
\end{enumerate}

能否彻底抛弃时间递推，让序列中\textbf{任意两个位置以 $O(1)$ 路径长度直接对话}，并且在训练时将全部 $T$ 个位置打包成一个大型矩阵乘法一次性并行算完？这个灵魂发问，将在下一章（第 4 章）直接引爆改变当代人工智能格局的\textbf{自注意力机制与 Transformer 架构}！

下面给出一段简洁可运行的 PyTorch 实战片段，展示 \texttt{im2col} 与 \texttt{F.conv2d} 的精确数值等价性，以及 Pre-Activation 残差块的标准写法：

\begin{lstlisting}[caption={\texttt{im2col}（\texttt{F.unfold}）与 \texttt{F.conv2d} 数值等价验证及 Pre-Act 残差块}]
import torch
import torch.nn as nn
import torch.nn.functional as F

# 1. 验证 im2col + GEMM 与标准 2D 卷积严格等价
x = torch.randn(2, 3, 16, 16)
w = torch.randn(8, 3, 3, 3)                  # (C_out, C_in, kH, kW)
y_conv = F.conv2d(x, w, stride=2, padding=1) # 输出 (2, 8, 8, 8)

x_col = F.unfold(x, kernel_size=3, padding=1, stride=2)  # (2, 27, 64)
w_row = w.view(8, -1)                                    # (8, 27)
y_gemm = (w_row @ x_col).view(2, 8, 8, 8)
assert torch.allclose(y_conv, y_gemm, atol=1e-5)

# 2. Pre-Activation 残差块：保持恒等主干道完全洁净
class PreActResBlock(nn.Module):
    def __init__(self, ch: int):
        super().__init__()
        self.branch = nn.Sequential(
            nn.BatchNorm2d(ch), nn.ReLU(inplace=True),
            nn.Conv2d(ch, ch, kernel_size=3, padding=1, bias=False),
            nn.BatchNorm2d(ch), nn.ReLU(inplace=True),
            nn.Conv2d(ch, ch, kernel_size=3, padding=1, bias=False),
        )
    def forward(self, x: torch.Tensor) -> torch.Tensor:
        return x + self.branch(x)            # 纯加法高速主干道！
\end{lstlisting}

\begin{failBox}{CNN、ResNet 与 RNN 工程排错三例}
\begin{enumerate}
  \item \textbf{卷积尺寸向下取整与残差分支错位}：使用式~\eqref{eq:conv-out-size} 做 $S=2$ 下采样时，若输入分辨率为奇数（如 $7 \times 7$），带 $P=1$ 的 $3\times 3$ 卷积输出 $\lfloor(7+2-3)/2\rfloor + 1 = 4$，而 $2\times 2$ 池化或不同填充可能得到 $3$，导致残差相加时报维度不匹配。且在 \texttt{Conv2d} 后紧接 \texttt{BatchNorm2d} 时，卷积层务必设 \texttt{bias=False}（否则偏置会被 BN 减去均值完全抵消，纯属冗余参数）。
  \item \textbf{Post-Norm 深度退化陷阱}：把归一化与激活放在残差加法\emph{之后}（如早期的 Post-Act ResNet 或 Post-LN Transformer），在层数加深后往往极难收敛或必须依赖极精细的 Warmup；尽量使用 Pre-Act / Pre-Norm 保持主干道为纯加法流。
  \item \textbf{截断 BPTT 忘记 \texttt{h.detach()} 引发显存泄漏}：在长文本流式训练 RNN/LSTM 时，常按每 $K = 64$ 步执行一次截断反向传播（Truncated BPTT）。如果不把传递给下一段的隐藏状态显式切断计算图（\texttt{h = h.detach()}），PyTorch 会把从第 1 步到第 $10000$ 步的所有历史计算图永久挂在内存中，很快撑爆显存（OOM）！
\end{enumerate}
\end{failBox}

\begin{extBox}{视觉与序列的当代演进：ViT vs ConvNeXt 与选择性状态空间模型（Mamba）}
本章介绍的卷积与循环并未被 Transformer 简单淘汰，反而在过去几年中完成了惊艳的现代复兴与大融合：
\begin{itemize}
  \item \textbf{Vision Transformer（ViT）与 ConvNeXt 的双雄争霸}：ViT（Dosovitskiy et al., 2020）把图像切成 $16\times 16$ 图像块（patch）当成 token 喂给纯 Transformer，在海量数据下突破了 CNN 局部感受野的天花板。但 Liu 等人（2022）随即推出 \textbf{ConvNeXt}：保留纯 CNN 结构，仅借用 Transformer 的现代配方（大 $7\times 7$ 深度可分离卷积、LayerNorm 取代 BatchNorm、GELU、倒置瓶颈与 AdamW + 强数据增强），就在 ImageNet 上与 Swin Transformer 战成平手甚至反超！这揭示了一个发人深省的事实：过去许多归功于“注意力本身”的性能跃升，有相当一部分来自更优的训练动力学与宏观模块设计。
  \item \textbf{状态空间模型（S4 / Mamba；Gu \& Dao, 2023）}：针对 Transformer 的 $O(T^2)$ 注意力复杂度，Mamba 回归了类似 RNN 的线性递归状态方程 $h_t = \bar{A}_t h_{t-1} + \bar{B}_t x_t$，但实现了两项突破：一是通过\textbf{线性对角递推}，利用 GPU 上的\textbf{并行前缀扫描（Parallel Associative Scan）}在 $O(\log T)$ 并行深度内一次性算完全部时间步，彻底打破了非线性 RNN 的串行魔咒；二是让状态转移矩阵 $(\bar{A}_t, \bar{B}_t)$ 随当前输入 $x_t$ 动态变化（称为\textbf{选择性机制}，在数学上与 LSTM 遗忘门高度呼应），从而在推理时保持 $O(1)$ 恒定显存与 $O(T)$ 线性长上下文吞吐！
\end{itemize}
\end{extBox}

\section*{本章小结与核心脉络}
\addcontentsline{toc}{section}{本章小结与核心脉络}

从第 1 章的泛化万用积木 MLP，到第 2 章的稳定训练动力学，再到本章的结构归纳偏置，我们完成了现代深度架构演进的前半程全景图：
\begin{itemize}
  \item \textbf{CNN} 把二维局部性与平移等变性写进共享卷积核，并用 \texttt{im2col} 与深度可分离结构榨干硬件算力；
  \item \textbf{ResNet} 用一条洁净的恒等加法主干道在深度轴上打通了梯度高速公路，将 24 层底层梯度提升了近九个数量级；
  \item \textbf{RNN/LSTM} 将同样的加性高速公路搬上了时间轴，但仍受困于 $O(T)$ 串行计算墙与固定容量隐状态瓶颈。
\end{itemize}
下一章，我们将彻底卸下串行循环的枷锁，亲手推导出缩放点积注意力、多头机制、RoPE 旋转位置编码与 KV Cache，并从零编写一个可在笔记本电脑上学会反转序列的\textbf{迷你 GPT}！

\section*{本章分层习题}
\addcontentsline{toc}{section}{本章分层习题}

\exTier{A 层：概念与手算热身}{熟练掌握卷积尺寸、感受野、参数量与门控直觉}

\begin{enumerate}
  \item[\textbf{5.3.1}] \textbf{多层卷积的特征图尺寸、参数量与感受野手算}：某视觉主干网络接收 $224 \times 224 \times 3$ 图像，依次经过三层卷积：
  \begin{itemize}
    \item Conv1：$C_{\text{out}} = 64,\; k = 7,\; S = 2,\; P = 3$；
    \item MaxPool：$k = 3,\; S = 2,\; P = 1$；
    \item Conv2：$C_{\text{out}} = 128,\; k = 3,\; S = 1,\; P = 1$（带空洞率 $d = 2$）。
  \end{itemize}
  逐层计算输出特征图尺寸 $(C, H, W)$、每层可训练参数量（设 \texttt{bias=False}）以及 Conv2 输出像素在原图上的精确感受野大小。
  \exHint{对 Conv2 先算等效核尺寸 $k_{\text{eff}} = 3 + 2(2-1) = 5$，再代入式~\eqref{eq:conv-out-size} 与式~\eqref{eq:receptive-field}（池化层同样贡献感受野增长）。}

  \item[\textbf{5.3.2}] \textbf{ResNet 瓶颈块 vs 基础块的参数效率对比}：在通道数为 $C = 256$ 时，比较以下两个残差块的参数量与单像素乘加量：
  \begin{enumerate}
    \item 基础块（BasicBlock）：两层 $3\times 3$ 卷积（$256 \to 256 \to 256$）；
    \item 瓶颈块（Bottleneck）：$1\times 1$ 降维（$256 \to 64$）$\to 3\times 3$ 卷积（$64 \to 64$）$\to 1\times 1$ 升维（$64 \to 256$）。
  \end{enumerate}
  \exHint{基础块约 $2 \times 9 \times 256^2 \approx 1.18\text{M}$ 参数；瓶颈块仅 $69{,}632$ 参数——节省了超过 $94\%$！}
\end{enumerate}

\exTier{B 层：定理推导与严谨证明}{严格推导卷积反向传播、残差路径展开与循环梯度}

\begin{enumerate}
  \item[\textbf{5.3.3}] \textbf{一维卷积的反向传播与转置卷积}：考虑步幅 $S=1$、无填充的一维互相关输出 $y_i = \sum_{u=0}^{k-1} w_u x_{i+u}$。给定上游梯度 $\bar{y}$，严格推导出 $\bar{w}_u$ 与 $\bar{x}_m$ 的显式表达式，并证明：对输入的反向传播 $\bar{x}$ 恰好等价于将上游梯度 $\bar{y}$ 两端补零 $k-1$ 格后，与\textbf{空间翻转（flip）后的卷积核}做互相关！
  \exHint{对 $\mathcal{L}(y)$ 求全微分 $d\mathcal{L} = \sum_i \bar{y}_i \sum_u w_u dx_{i+u}$，令 $m = i + u$ 交换求和次序即可一眼看出翻转核 $w_{m-i}$。}

  \item[\textbf{5.3.4}] \textbf{带缩放残差分支的方差增长与 DeepNorm / $\frac{1}{\sqrt{L}}$ 缩放}：考虑 $L$ 个独立初始化的洁净残差块 $x_{l+1} = x_l + \alpha\,\mathcal{F}_l(x_l)$，假设每层输入与残差输出不相关且 $\Var(\mathcal{F}_l(x_l)) = \Var(x_l)$。
  \begin{enumerate}
    \item 证明第 $L$ 层激活方差满足 $\Var(x_L) = (1 + \alpha^2)^L \Var(x_0)$；
    \item 当 $\alpha = 1$ 时，$L=100$ 层的方差将膨胀为 $2^{100}$！证明取残差分支缩放因子 $\alpha = \frac{1}{\sqrt{L}}$（这正是 GPT-2 初始化对残差输出投影层除以 $\sqrt{2 N_{\text{layers}}}$ 的来源！）可使顶层方差稳定收敛于 $e \cdot \Var(x_0)$（与深度 $L$ 无关）。
  \end{enumerate}
  \exHint{利用经典极限 $\lim_{L \to \infty} (1 + 1/L)^L = e$。}
\end{enumerate}

\exTier{C 层：代码实现与数值实验}{手写核心模块并复现深度与时间梯度实测}

\begin{enumerate}
  \item[\textbf{5.3.5}] \textbf{比较 Post-Act 与 Pre-Act ResNet 的 50 层梯度传播}：将实验脚本 \texttt{vol5\_experiments.py} 中的残差梯度实验 \texttt{exp3\_resnet\_grad\_highway} 扩展至 $L = 50$ 层，加入第三条曲线——Post-Activation 残差块 $h_l = \ReLU(h_{l-1} + 0.35\,\ReLU(h_{l-1} W_l))$（Kaiming 初始化）。比较 Plain、Post-Act ResNet 与 Pre-Act ResNet 从第 50 层到第 1 层的梯度范数衰减比率，验证洁净恒等主干道的优势。
  \exHint{注意 Post-Act 中的外层 $\ReLU$ 每层约阻断一半通道的负值，在 50 层累积后同样会造成明显衰减。}

  \item[\textbf{5.3.6}] \textbf{手工实现单步 LSTM 单元并测试长距离记忆复制任务}：不调用 \texttt{nn.LSTM}，仅用 \texttt{nn.Linear}（将四个门合并为单次矩阵乘以提速！）严格按照式~\eqref{eq:lstm-gates1}--\eqref{eq:lstm-cell} 实现 \texttt{MyLSTMCell}，并与 \texttt{nn.LSTMCell} 逐元素精度互验（误差 $< 10^{-6}$）。随后在“首词记忆任务”（输入长度为 $T = 40$ 的随机序列，要求在最后一步输出第 1 步的符号）上对比 Elman RNN 与 \texttt{MyLSTMCell}（分别测试遗忘门偏置 $b_f = 0$ 与 $b_f = +2.0$）的收敛速度。
  \exHint{将 $[h_{t-1}, x_t]$ 拼接后一次性投影到 $4d$ 维，再用 \texttt{.chunk(4, dim=-1)} 切分为 $(f, i, \tilde{c}, o)$。}
\end{enumerate}

\exTier{D 层：开放探究与前沿迁移}{深入群等变网络与并行线性递归}

\begin{enumerate}
  \item[\textbf{5.3.7}] \textbf{超越平移：旋转等变卷积（Group Equivariant CNN）}：标准 2D 卷积只对平移群 $(\mathbb{Z}^2, +)$ 等变，但对医学病理切片或卫星遥感图而言，图像旋转 $90^\circ, 180^\circ, 270^\circ$ 后的语义完全等价。查阅 Cohen 与 Welling（2016）的 $p4$ 群卷积思想：如何通过把单个 $3\times 3$ 卷积核旋转 $4$ 个角度共享参数，将特征图提升到群函数 $F(g)$（$g \in p4$）上，从而严格保证旋转等变性而不增加参数量？
  \exHint{将标准卷积定义中的平移群求和替换为有限群 $G$ 上的群卷积 $(f *_G \psi)(g) = \sum_{h \in G} f(h)\,\psi(g^{-1} h)$。}

  \item[\textbf{5.3.8}] \textbf{线性递归的并行前缀扫描（Parallel Scan）原理}：为什么 Mamba/LRU 等现代线性状态空间模型可以突破 RNN 的 $O(T)$ 串行墙？给定标量线性递推 $h_t = a_t h_{t-1} + b_t$（$t = 1, \dots, T$），定义二元运算符 $\bullet$ 作用于数对 $(a, b)$：
  \[
    (a_2, b_2) \bullet (a_1, b_1) := \left(a_2 a_1,\; a_2 b_1 + b_2\right).
  \]
  证明运算符 $\bullet$ 满足\textbf{结合律}，并解释为何结合律允许 GPU 用树状归约在仅 $O(\log_2 T)$ 个并行步内算出全部 $h_1, \dots, h_T$，而带 $\tanh$ 的非线性 Elman RNN 却无法应用此技巧！
  \exHint{把 $h_t = a_t h_{t-1} + b_t$ 写成 $2\times 2$ 仿射矩阵乘积 $\binom{h_t}{1} = \begin{pmatrix} a_t & b_t \\ 0 & 1 \end{pmatrix}\binom{h_{t-1}}{1}$，矩阵乘法天然满足结合律！}
\end{enumerate}

\chapter[注意力机制与 Transformer 架构]{注意力机制与 Transformer 架构\\（从零手写迷你 GPT）}
\label{ch:attention-transformer}

\chapEpigraph{每个人都知道什么是注意力。它是心智在诸多同时呈现的对象或思绪之中，清晰而鲜明地择取其一。}{William James（威廉·詹姆斯，《心理学原理》，1890）}

\begin{chapIntro}{本章导读：卸下串行枷锁，重构序列计算的底层范式}
上一章结尾我们看到，循环神经网络（RNN/LSTM）受困于两道难以逾越的物理墙：其一，隐藏状态必须按时间步 $h_1 \to h_2 \to \cdots \to h_T$ 严格\textbf{串行}递推，导致 GPU 成千上万个算术核心在序列长度维度上被迫空转；其二，无论前文长达几百还是几万词，所有历史信息都必须压缩进一个固定维度的状态向量 $h_t \in \R^d$，长程信息不可避免地被冲刷殆尽。

2017 年 Vaswani 等人发表的《Attention Is All You Need》以彻底的极简主义掀翻了棋盘：\textbf{完全抛弃循环与卷积，仅靠“自注意力（Self-Attention）”让序列中任意两个位置以 $O(1)$ 的路径长度直接对话！} 本章将以第一性原理层层拆解现代大语言模型的全部核心零件：
\begin{enumerate}
  \item \textbf{缩放点积注意力与 $\sqrt{d_k}$ 缩放定理}：从可微软词典寻址出发，严格证明若不除以 $\sqrt{d_k}$，内积方差将随维度线性膨胀，把 $\softmax$ 推入雅可比矩阵趋零的饱和死区；
  \item \textbf{多头机制、张量变形舞步与因果掩码}：逐行追踪 \texttt{(B, T, C)} 到 \texttt{(B, H, T, d\_h)} 的张量形状流转，严格证明上三角 $-\infty$ 掩码保证了未来词元对当前位置具有\textbf{严格为零的梯度}，并介绍现代 LLM 标配的分组查询注意力（GQA）；
  \item \textbf{位置编码与 RoPE 旋转位置编码}：证明无位置编码的自注意力具有置换等变性，完整推导二维复平面与旋转矩阵下的 RoPE 相对位置定理 $R_\theta(m)^\top R_\theta(n) = R_\theta(n - m)$；
  \item \textbf{Pre-LN 架构、SwiGLU 与 KV Cache 显存账本}：衔接到第 3 章的残差高速公路定理，精确核算预填充（Prefill）与自回归解码（Decode）的计算量及 KV Cache 显存公式；
  \item \textbf{可运行实测（\texttt{data/transformer\_train.csv}）}：用纯 PyTorch 从零手写 2 层因果\textbf{迷你 GPT}，亲眼见证它在 60--80 步内将 6 词序列反转任务的整句精确匹配率从 $0\%$ 拉升至 $100.00\%$！
\end{enumerate}
\end{chapIntro}

\section{从软词典寻址到缩放点积注意力}
\label{sec:sdpa}

\subsection{注意力即“可微分的软词典查询”}

在传统程序设计中，哈希词典（dictionary）存储着一组键值对 $\{(k_1, v_1), \dots, (k_T, v_T)\}$。当你提交一个查询 $q$ 时，硬词典执行精确匹配：若 $q = k_j$ 则原样返回 $v_j$，否则返回空值。这种“非黑即白”的离散寻址在数学上是不可微的，无法通过第 1 章的反向传播端到端学习。

\textbf{注意力机制（Attention Mechanism）}的天才之处，是将硬词典软化为连续可微的概率凸组合：
\begin{enumerate}
  \item 用查询向量 $q \in \R^{d_k}$ 与所有键向量 $k_1, \dots, k_T \in \R^{d_k}$ 逐一计算相似度打分 $s_j = \operatorname{score}(q, k_j)$；
  \item 用 $\softmax$ 将打分归一化为概率分布 $\alpha_j = \frac{\exp(s_j)}{\sum_{r=1}^T \exp(s_r)}$（满足 $\alpha_j > 0,\; \sum_{j=1}^T \alpha_j = 1$）；
  \item 以权重 $\alpha_j$ 对所有值向量 $v_1, \dots, v_T \in \R^{d_v}$ 做加权平均，输出 $o = \sum_{j=1}^T \alpha_j v_j \in \R^{d_v}$。
\end{enumerate}

在\textbf{自注意力（Self-Attention）}中，查询、键与值都源自同一输入序列 $X \in \R^{T \times d_{\text{model}}}$（每个词元 $x_i$ 既是“发问者”，也是“被检索的索引”与“被提取的内容”），仅通过三组独立线性投影区分角色：
\begin{equation}
  Q = X W_Q \in \R^{T \times d_k}, \qquad
  K = X W_K \in \R^{T \times d_k}, \qquad
  V = X W_V \in \R^{T \times d_v}.
  \label{eq:qkv-proj}
\end{equation}
于是序列中任意两个相隔甚远的词元 $i$ 与 $j$（无论相距 $5$ 步还是 $5000$ 步），只需一次内积 $q_i^\top k_j$ 即可直接交互——最大信号传播路径长度从 RNN 的 $O(T)$ 一举压缩为\textbf{常数 $O(1)$}，且全部 $T$ 个位置的计算化归为高度并行的矩阵乘法！

\begin{defBox}
\begin{definition}[缩放点积注意力，Scaled Dot-Product Attention]
给定查询矩阵 $Q \in \R^{T_q \times d_k}$、键矩阵 $K \in \R^{T_k \times d_k}$ 与值矩阵 $V \in \R^{T_k \times d_v}$，缩放点积注意力定义为：
\begin{equation}
  \Attn(Q, K, V) = \softmax\!\left(\frac{Q K^\top}{\sqrt{d_k}}\right) V \in \R^{T_q \times d_v},
  \label{eq:sdpa}
\end{equation}
其中 $\softmax$ 沿矩阵的\textbf{最后一维（每一行，即键维度 $T_k$）}独立归一化。记 $Z = \frac{Q K^\top}{\sqrt{d_k}} \in \R^{T_q \times T_k}$ 为注意力对数几率（logits），$P = \softmax(Z) \in \R^{T_q \times T_k}$ 为注意力权重矩阵，则输出 $O = P V$。
\end{definition}
\end{defBox}

\begin{example}[手算迷你示例：$T=3,\; d_k=d_v=2$ 的软词典检索]
\label{ex:sdpa-hand}
设三个词元对应的键与值分别为 $k_1 = (1, 0)^\top,\; k_2 = (0, 1)^\top,\; k_3 = (-1, 0)^\top$，以及 $v_1 = (10, 0)^\top,\; v_2 = (0, 10)^\top,\; v_3 = (-10, 0)^\top$。现给定一个偏向键 $k_1$ 的查询 $q = (2, 0.5)^\top$：
\begin{enumerate}
  \item \textbf{原始点积与缩放}：$q^\top K^\top = (2,\; 0.5,\; -2)$。除以 $\sqrt{d_k} = \sqrt{2} \approx 1.4142$ 得缩放对数几率
  \[
    z = (1.4142,\; 0.3536,\; -1.4142).
  \]
  \item \textbf{行 $\softmax$ 权重}：$\exp(z) \approx (4.113,\; 1.424,\; 0.243)$，求和为 $5.780$，故
  \[
    \alpha = \softmax(z) \approx (0.712,\; 0.246,\; 0.042).
  \]
  \item \textbf{加权读取值向量}：
  \[
    o = 0.712 (10, 0)^\top + 0.246 (0, 10)^\top + 0.042 (-10, 0)^\top \approx (6.70,\; 2.46)^\top.
  \]
\end{enumerate}
可见注意力自动将 $71.2\%$ 的权重分配给最匹配的键 $k_1$，同时保留了 $24.6\%$ 给正交邻近项 $k_2$，几乎屏蔽了反向项 $k_3$——并且整个过程对 $q, k_j, v_j$ 处处光滑可微！
\end{example}

\subsection{为什么必须除以 $\sqrt{d_k}$？——方差膨胀与 Softmax 饱和死区定理}

初学者常疑惑：式~\eqref{eq:sdpa} 中的分母 $\sqrt{d_k}$ 究竟是经验技巧还是数学必然？下面给出完整的概率与微积分证明。

\begin{thmBox}
\begin{theorem}[$\sqrt{d_k}$ 缩放定理与 Softmax 雅可比梯度消失]
\label{thm:sqrt-dk}
设查询向量 $q \in \R^{d_k}$ 与键向量 $k \in \R^{d_k}$ 的各分量相互独立，且满足零均值、单位方差：$\E[q_r] = \E[k_r] = 0,\; \Var(q_r) = \Var(k_r) = 1$（$r = 1, \dots, d_k$）。则：
\begin{enumerate}
  \item 未缩放点积 $s = q^\top k$ 满足 $\E[s] = 0$ 且 $\Var(s) = d_k$；因此其典型量级按 $\Theta(\sqrt{d_k})$ 增长。令 $z = s / \sqrt{d_k}$ 可精确恢复单位方差 $\Var(z) = 1$（与维度 $d_k$ 无关）。
  \item 记 $p = \softmax(z) \in \Delta^{T-1}$，则 $\softmax$ 的雅可比矩阵为：
  \begin{equation}
    J_{\softmax}(z) = \frac{\partial p}{\partial z} = \diag(p) - p p^\top.
    \label{eq:softmax-jacobian}
  \end{equation}
  当未缩放 logits 随 $\sqrt{d_k}$ 放大时，设最大与次大未缩放分差为 $\Delta = s_{(1)} - s_{(2)} = c\sqrt{d_k} > 0$，则 $p$ 指数级坍缩至单热顶点 $e_{i^*}$，且雅可比矩阵的 Frobenius 范数与最大奇异值均按 $O(e^{-c\sqrt{d_k}})$ 趋近于零，导致经由注意力权重回传给 $Q, K$ 的梯度彻底消失。
\end{enumerate}
\end{theorem}
\begin{proof}
\textbf{第一部分（内积方差）}：因各分量独立且零均值，
\[
  \E[q^\top k] = \sum_{r=1}^{d_k} \E[q_r k_r] = \sum_{r=1}^{d_k} \E[q_r]\E[k_r] = 0.
\]
对每一项乘积 $q_r k_r$，由独立性有 $\Var(q_r k_r) = \E[q_r^2 k_r^2] - (\E[q_r k_r])^2 = \E[q_r^2]\E[k_r^2] - 0 = 1 \cdot 1 = 1$。由于不同坐标 $r \ne r'$ 相互独立，方差直接相加：
\[
  \Var(q^\top k) = \sum_{r=1}^{d_k} \Var(q_r k_r) = d_k.
\]
于是除以 $\sqrt{d_k}$ 后，$\Var\!\left(\frac{q^\top k}{\sqrt{d_k}}\right) = \frac{1}{d_k}\Var(q^\top k) = 1$。

\textbf{第二部分（Softmax 雅可比与饱和坍缩）}：对 $p_i = \frac{e^{z_i}}{\sum_m e^{z_m}}$ 求偏导 $\frac{\partial p_i}{\partial z_j}$，利用商的求导法则：
\begin{itemize}
  \item 当 $i = j$ 时：
  \[
    \frac{\partial p_i}{\partial z_i}
    = \frac{e^{z_i}\sum_m e^{z_m} - e^{z_i} e^{z_i}}{\left(\sum_m e^{z_m}\right)^2}
    = p_i - p_i^2 = p_i(1 - p_i).
  \]
  \item 当 $i \ne j$ 时：
  \[
    \frac{\partial p_i}{\partial z_j}
    = \frac{0 - e^{z_i} e^{z_j}}{\left(\sum_m e^{z_m}\right)^2}
    = -p_i p_j.
  \]
\end{itemize}
合并两式即得 $\frac{\partial p_i}{\partial z_j} = p_i(\delta_{ij} - p_j)$，写成矩阵形式正是式~\eqref{eq:softmax-jacobian}：$J_{\softmax}(z) = \diag(p) - p p^\top$。

现在考察未缩放情形 $z = s$，设 $i^* = \argmax_m s_m$ 且与其余分量至少相差 $\Delta = c\sqrt{d_k} > 0$。令 $\epsilon := 1 - p_{i^*}$，则：
\[
  \epsilon = \frac{\sum_{m \ne i^*} e^{s_m - s_{i^*}}}{1 + \sum_{m \ne i^*} e^{s_m - s_{i^*}}}
  \le (T - 1) e^{-c\sqrt{d_k}}.
\]
由于对任意非最大项 $j \ne i^*$ 均有 $p_j \le \epsilon$，而最大项 $p_{i^*}(1 - p_{i^*}) \le \epsilon$，因此雅可比矩阵的每一元素满足：
\[
  \abs{\big(J_{\softmax}(s)\big)_{ij}}
  \le \begin{cases}
    p_{i^*}(1 - p_{i^*}) \le \epsilon, & i = j = i^*,\\
    p_i \le \epsilon, & \text{其余一切 }(i, j).
  \end{cases}
\]
由 Gershgorin 圆盘定理或直接对角求和可知，雅可比矩阵的谱范数满足 $\norm{J_{\softmax}(s)}_2 \le 2\epsilon \le 2(T-1) e^{-c\sqrt{d_k}}$。当 $d_k = 64$ 或 $128$ 时，典型未缩放分差可达 $10 \sim 20$，此时 $e^{-15} \approx 3 \times 10^{-7}$！上游传给注意力概率的梯度 $\bar{P} = \bar{O} V^\top$ 在穿过式~\eqref{eq:softmax-jacobian} 化为 logits 梯度时：
\begin{equation}
  \bar{Z}_{i,:} = P_{i,:} \had \left(\bar{P}_{i,:} - \inner{P_{i,:}}{\bar{P}_{i,:}}\mathbf{1}^\top\right),
  \label{eq:attn-softmax-vjp}
\end{equation}
由于 $P_{i,:}$ 几乎是单热向量 $e_{i^*}$，$\bar{P}_{i, i^*} - \inner{P_{i,:}}{\bar{P}_{i,:}} \approx 0$ 且其余位置被 $P_{i, j} \approx 0$ 彻底乘零，于是 $\bar{Q} = \frac{1}{\sqrt{d_k}}\bar{Z} K$ 与 $\bar{K} = \frac{1}{\sqrt{d_k}}\bar{Z}^\top Q$ 全面熄火。
\end{proof}
\end{thmBox}

式~\eqref{eq:attn-softmax-vjp} 同样给出了缩放点积注意力反向传播（VJP）的核心公式——顺带验证第 1 章的反向模式自动微分：给定上游梯度 $\bar{O} \in \R^{T \times d_v}$，依次可得 $\bar{V} = P^\top \bar{O}$，$\bar{P} = \bar{O} V^\top$，按式~\eqref{eq:attn-softmax-vjp} 算出 $\bar{Z}$，最后 $\bar{Q} = d_k^{-1/2}\bar{Z} K$ 与 $\bar{K} = d_k^{-1/2}\bar{Z}^\top Q$。

\section{多头注意力、因果掩码与分组查询（GQA）}
\label{sec:mha-mask-gqa}

\subsection{多头注意力（MHA）与张量形状变形舞步}

单一注意力头在每个位置上只能生成一张权重分布 $P \in \R^{T \times T}$，但自然语言中的依赖关系是多维交织的：同一个代词“它”可能既需要在\emph{句法指代子空间}回看十词之前的主语，又需要在\emph{局部修饰子空间}紧盯相邻动词。

\textbf{多头注意力（Multi-Head Attention, MHA）}将总通道维度 $C = d_{\text{model}}$ 均匀切分为 $H$ 个独立子空间，每个头的维度为 $d_h = C / H$：
\begin{equation}
  \operatorname{MHA}(X) = \operatorname{Concat}\!\left(\operatorname{head}_1, \dots, \operatorname{head}_H\right) W^O,
  \qquad
  \operatorname{head}_h = \Attn(X W_h^Q,\; X W_h^K,\; X W_h^V),
  \label{eq:mha}
\end{equation}
其中 $W_h^Q, W_h^K, W_h^V \in \R^{C \times d_h}$，$W^O \in \R^{C \times C}$。注意到把 $H$ 个头的参数量加总，总投影参数为 $3 H (C d_h) + C^2 = 4 C^2$——\textbf{与单头满维注意力参数量完全一致}，却换来了 $H$ 组可同时关注的独立表示子空间！

工程实现绝不会用 Python \texttt{for} 循环逐头计算，而是把所有头的 $Q, K, V$ 融合成单次大矩阵乘法 $W_{QKV} \in \R^{C \times 3C}$，再借由精准的\textbf{张量形状舞步（Tensor Shape Choreography）}交给批量矩阵乘（Batched GEMM）一次性并行扫过批次 $B$ 与头数 $H$：

\begin{center}
\footnotesize
\begin{tabularx}{\linewidth}{@{}p{26mm}p{34mm}X@{}}
\toprule
\textbf{操作步骤} & \textbf{张量形状} & \textbf{数学与内存含义} \\
\midrule
输入序列 $X$ & \texttt{(B, T, C)} & $B$ 条样本，序列长 $T$，模型宽度 $C = H \cdot d_h$ \\
融合投影 + 切分 & $3\times$\texttt{(B, T, C)} & 单次线性层 \texttt{qkv(x)} 后按 \texttt{dim=-1} 均切出 $Q, K, V$ \\
拆头与转置 & \texttt{(B, H, T, d\_h)} & \textbf{先} \texttt{view(B,T,H,d\_h)}，\textbf{再} \texttt{transpose(1,2)} 将 $H$ 移入批维 \\
缩放内积与掩码 & \texttt{(B, H, T, T)} & \texttt{(q @ k.transpose(-2,-1)) / sqrt(d\_h)} 加掩码 \\
归一化与取值 & \texttt{(B, H, T, d\_h)} & \texttt{softmax(att, dim=-1) @ v} 完成全部 $B\cdot H$ 组软检索 \\
合并多头 + 输出 & \texttt{(B, T, C)} & \texttt{transpose(1,2).contiguous().view(B,T,C)} 接 $W^O$ \\
\bottomrule
\end{tabularx}
\end{center}

请特别留意最后一步里的 \texttt{.contiguous()}：\texttt{transpose(1, 2)} 只交换了张量的步长（stride）元数据而未搬运底层内存，若不先整理为连续内存就直接调用 \texttt{.view(B, T, C)}，PyTorch 会立即抛出运行时错误。

\subsection{因果掩码（Causal Masking）：零信息泄露的严格保证}

在自回归语言模型（GPT、Llama 等解码器架构）训练中，我们把整段文本 $(x_1, \dots, x_T)$ 一次性并行输入网络，要求位置 $i$ 仅利用前缀 $x_{\le i}$ 来预测下一词 $x_{i+1}$。如果位置 $i$ 在自注意力中“偷看”了 $j > i$ 的未来词元，训练损失将瞬间作弊归零，但在推理生成时未来尚不存在，模型便彻底瘫痪。

解决方案是在 $\softmax$ \textbf{之前}，给注意力 logits 矩阵加上一个上三角负无穷掩码 $M \in \R^{T \times T}$：
\begin{equation}
  M_{ij} = \begin{cases}
    0, & j \le i,\\
    -\infty, & j > i,
  \end{cases}
  \qquad
  P = \softmax\!\left(\frac{Q K^\top}{\sqrt{d_k}} + M\right).
  \label{eq:causal-mask}
\end{equation}

\begin{thmBox}
\begin{proposition}[因果掩码的零未来梯度定理]
\label{prop:causal-zero-grad}
在式~\eqref{eq:causal-mask} 的因果自注意力中，对任意位置 $i$ 与任意未来位置 $j > i$，不仅前向注意力权重严格满足 $P_{ij} = 0$，而且任意只依赖位置 $i$ 输出 $o_i$ 的损失 $\mathcal{L}(o_i)$ 对未来键、值向量 $(k_j, v_j)$ 的反向梯度也\textbf{恒严格为零}：
\[
  \frac{\partial o_i}{\partial k_j} = \mathbf{0}, \qquad \frac{\partial o_i}{\partial v_j} = \mathbf{0} \qquad (\forall\, j > i).
\]
\end{proposition}
\begin{proof}
前向显然：对 $j > i$，$Z_{ij} = -\infty$，故 $P_{ij} = \frac{e^{-\infty}}{\sum_{r \le i} e^{Z_{ir}}} = 0$，从而 $o_i = \sum_{r=1}^i P_{ir} v_r$ 与 $v_j$ 无关，立得 $\frac{\partial o_i}{\partial v_j} = P_{ij} I = \mathbf{0}$。

对键向量 $k_j$（$j > i$），由链式法则和式~\eqref{eq:softmax-jacobian}，对任意列 $m \in \{1, \dots, T\}$：
\[
  \frac{\partial P_{im}}{\partial Z_{ij}} = P_{im}(\delta_{mj} - P_{ij}).
\]
当 $m = j$ 时 $P_{ij} = 0$，故 $P_{ij}(1 - P_{ij}) = 0$；当 $m \ne j$ 时同样有 $-P_{im} P_{ij} = 0$。因此整行导数 $\frac{\partial P_{i,:}}{\partial Z_{ij}} \equiv \mathbf{0}$，从而 $\frac{\partial o_i}{\partial k_j} = \sum_m \frac{\partial o_i}{\partial P_{im}}\frac{\partial P_{im}}{\partial Z_{ij}}\frac{\partial Z_{ij}}{\partial k_j} = \mathbf{0}$。未来词元在正反向两个方向上都被数学隔离！
\end{proof}
\end{thmBox}

\begin{insightBox}{为什么不能在 $\softmax$ 之后再把上三角置零？}
如果先对完整 $Q K^\top/\sqrt{d_k}$ 做 $\softmax$ 再把 $j > i$ 的权重乘零，会有两重灾难：第一，剩余下三角权重之和不再等于 $1$（且分母 $Z_i = \sum_{r=1}^T e^{Z_{ir}}$ 已经偷偷混入了未来词元 $k_{j>i}$ 的信息！）；第二，由分母泄露导致的 $\partial P_{im}/\partial Z_{ij} = -P_{im}P_{ij} \ne 0$ 会把未来词元的梯度传给前文！\textbf{掩码必须在 $\softmax$ 之前以加法 $-\infty$ 施加。}
\end{insightBox}

\subsection{分组查询注意力（GQA）：现代大模型的显存瘦身术}

标准 MHA 为每个查询头都配备专属的键头与值头（$n_{\text{kv}} = n_q = H$）。但在后文第~\ref{sec:block-kvcache} 节我们将看到，大模型长文本推理的头号显存瓶颈是保存历史键值的 \textbf{KV Cache}。
\begin{itemize}
  \item \textbf{多查询注意力（MQA；Shazeer, 2019）}：令所有 $n_q$ 个查询头共享\textbf{同一组}键与值（$n_{\text{kv}} = 1$），KV 显存骤降 $H$ 倍，但有时会轻微损及质量与训练稳定性；
  \item \textbf{分组查询注意力（GQA；Ainslie et al., 2023）}：取两全其美的折中——把 $n_q$ 个查询头均分为 $n_{\text{kv}}$ 组（例如 Llama-3 8B 中 $n_q = 32,\; n_{\text{kv}} = 8$，每 4 个查询头共享 1 组 $K, V$）。GQA 几乎保留了 MHA 的全部表达精度，同时把 KV Cache 显存和解码访存带宽直接砍去 $75\%$（$4\times$ 压缩），已成为当今主流开源大模型的默认标准。
\end{itemize}

\section{位置编码与旋转位置编码（RoPE）}
\label{sec:pe-rope}

\subsection{自注意力的置换等变性：没有位置编码的“词袋”}

与 RNN 按时间步推进、CNN 按空间坐标滑窗不同，原始自注意力把输入视为无序集合！

\begin{thmBox}
\begin{proposition}[无位置编码自注意力的置换等变性]
\label{prop:attn-perm}
设 $\Pi \in \{0, 1\}^{T \times T}$ 为任意 $T \times T$ 置换矩阵（满足 $\Pi^\top \Pi = \Pi \Pi^\top = I$）。对无因果掩码且无位置编码的自注意力层 $F(X) = \Attn(X W_Q, X W_K, X W_V)$，恒有：
\begin{equation}
  F(\Pi X) = \Pi\, F(X).
  \label{eq:attn-perm}
\end{equation}
\end{proposition}
\begin{proof}
代入 $\Pi X$ 得 $Q' = \Pi Q,\; K' = \Pi K,\; V' = \Pi V$，从而内积矩阵变为：
\[
  \frac{Q' (K')^\top}{\sqrt{d_k}} = \frac{\Pi Q K^\top \Pi^\top}{\sqrt{d_k}}.
\]
因为 $\softmax$ 对每一行独立运算且关于行内元素次序对称，行置换与列置换均与 $\softmax$ 交换：$\softmax(\Pi Z \Pi^\top) = \Pi\,\softmax(Z)\,\Pi^\top$。最后右乘 $V' = \Pi V$ 并利用正交性 $\Pi^\top \Pi = I$：
\[
  F(\Pi X) = \Pi\,\softmax\!\left(\frac{Q K^\top}{\sqrt{d_k}}\right)\Pi^\top \Pi V = \Pi\, F(X).
\]
\end{proof}
\end{thmBox}

这意味着如果不注入位置信息，“猫追狗”与“狗追猫”里每个词拿到的上下文表征完全一样——模型退化成了一个高级词袋！因此，我们必须为 Transformer 装上“空间坐标仪”。

\subsection{正弦绝对位置编码的灵感与局限}

原始 Transformer（Vaswani et al., 2017）将确定性正弦位置向量 $\mathrm{PE}_m \in \R^d$ 直接\emph{加}在词嵌入上：$x_m \leftarrow e_m + \mathrm{PE}_m$，其第 $2i$ 与 $2i+1$ 维定义为：
\begin{equation}
  \mathrm{PE}_{(m, 2i)} = \sin(m \theta_i), \qquad
  \mathrm{PE}_{(m, 2i+1)} = \cos(m \theta_i), \qquad
  \theta_i = 10000^{-2i/d}.
  \label{eq:sin-pe}
\end{equation}
不同维度构成从高频（$i=0,\; \theta_0 = 1$）到低频（$i=d/2-1,\; \theta \approx 10^{-4}$）的连续“二进制时钟”。由三角和角公式，$\mathrm{PE}_{m+\Delta}$ 可由 $\mathrm{PE}_m$ 经仅依赖偏移 $\Delta$ 的线性旋转矩阵得到。

但加法式绝对位置编码有一个隐蔽缺陷：展开内积 $(e_m + \mathrm{PE}_m)^\top W_Q^\top W_K (e_n + \mathrm{PE}_n)$ 时，会产生内容--位置交叉项与 $\mathrm{PE}_m^\top W_Q^\top W_K \mathrm{PE}_n$；投影矩阵 $W_Q^\top W_K$ 会破坏完美旋转结构，使得模型在训练长度之外的外推能力较弱。

\subsection{旋转位置编码（RoPE）：复平面上的相对位置奇迹}

苏剑林等人（Su et al., 2021）提出的\textbf{旋转位置编码（Rotary Positional Embedding, RoPE）}提出了一个惊艳的几何洞见：\textbf{不要把位置向量“加”到词嵌入里，而是把查询 $q_m$ 与键 $k_n$ 按各自的位置索引 $(m, n)$ 在二维子平面上直接“旋转”对应角度！}

先看最核心的二维情形（$d_h = 2$）。将二维向量 $q = (q_1, q_2)^\top \in \R^2$ 视为复数 $\tilde{q} = q_1 + \mathrm{i} q_2 \in \C$，位置 $m$ 对应的角度旋转操作就是乘以单位复数 $e^{\mathrm{i} m \theta}$，等价于左乘二维正交旋转矩阵：
\begin{equation}
  R_\theta(m) = \begin{pmatrix}
    \cos(m\theta) & -\sin(m\theta) \\
    \sin(m\theta) & \cos(m\theta)
  \end{pmatrix} \in \mathrm{SO}(2).
  \label{eq:rope-2d}
\end{equation}
对一般偶数维度 $d_h = 2K$，将向量两两配对为 $K$ 个二维子平面，采用多尺度角频率 $\theta_r = \text{base}^{-2(r-1)/d_h}$（经典 $\text{base} = 10000$，现代长上下文模型如 Llama-3 扩大至 $500{,}000$），构成 $d_h \times d_h$ 分块对角旋转矩阵：
\begin{equation}
  R_{\Theta, m}^{d_h} = \diag\!\left(R_{\theta_1}(m),\; R_{\theta_2}(m),\; \dots,\; R_{\theta_{d_h/2}}(m)\right).
  \label{eq:rope-nd}
\end{equation}

\begin{thmBox}
\begin{theorem}[RoPE 的相对位置内积不变性与保范定理]
\label{thm:rope}
对任意位置索引 $m, n \in \mathbb{Z}$、角频率 $\theta \in \R$ 及向量 $q, k \in \R^{d_h}$，旋转矩阵严格满足群乘法关系：
\begin{equation}
  R_\theta(m)^\top R_\theta(n) = R_\theta(n - m),
  \label{eq:rope-group}
\end{equation}
由此推出经 RoPE 旋转后的查询与键内积\textbf{仅依赖于相对位置偏移 $n - m$}，且旋转不改变任何向量欧氏范数：
\begin{equation}
  \big(R_{\Theta, m}^{d_h} q\big)^\top \big(R_{\Theta, n}^{d_h} k\big)
  = q^\top R_{\Theta, n - m}^{d_h} k,
  \qquad
  \norm{R_{\Theta, m}^{d_h} q}_2 = \norm{q}_2.
  \label{eq:rope-rel}
\end{equation}
\end{theorem}
\begin{proof}
\textbf{方法一（二维旋转矩阵直接展开）}：对单个二维子块，取转置 $R_\theta(m)^\top$ 并与 $R_\theta(n)$ 相乘：
\begin{align*}
  R_\theta(m)^\top R_\theta(n)
  &= \begin{pmatrix}
    \cos m\theta & \sin m\theta \\
    -\sin m\theta & \cos m\theta
  \end{pmatrix}
  \begin{pmatrix}
    \cos n\theta & -\sin n\theta \\
    \sin n\theta & \cos n\theta
  \end{pmatrix} \\
  &= \begin{pmatrix}
    \cos m\theta \cos n\theta + \sin m\theta \sin n\theta &
    -\cos m\theta \sin n\theta + \sin m\theta \cos n\theta \\
    -\sin m\theta \cos n\theta + \cos m\theta \sin n\theta &
    \sin m\theta \sin n\theta + \cos m\theta \cos n\theta
  \end{pmatrix} \\
  &= \begin{pmatrix}
    \cos\big((n - m)\theta\big) & -\sin\big((n - m)\theta\big) \\
    \sin\big((n - m)\theta\big) & \cos\big((n - m)\theta\big)
  \end{pmatrix}
  = R_\theta(n - m).
\end{align*}
\textbf{方法二（复平面一行极简证）}：将 $q, k \in \R^2$ 记为复数 $\tilde{q}, \tilde{k} \in \C$，则实内积满足 $u^\top v = \operatorname{Re}(\tilde{u}^* \tilde{v})$（其中 $*$ 表示复共轭）。于是：
\[
  \big(R_\theta(m) q\big)^\top \big(R_\theta(n) k\big)
  = \operatorname{Re}\!\left[\big(\tilde{q} e^{\mathrm{i}m\theta}\big)^* \big(\tilde{k} e^{\mathrm{i}n\theta}\big)\right]
  = \operatorname{Re}\!\left[\tilde{q}^* \tilde{k}\, e^{\mathrm{i}(n-m)\theta}\right]
  = q^\top R_\theta(n - m) k.
\]
对分块对角矩阵 $R_{\Theta, m}^{d_h}$，逐块应用上述恒等式即得 $(R_{\Theta, m}^{d_h})^\top R_{\Theta, n}^{d_h} = R_{\Theta, n-m}^{d_h}$；特别地，令 $n = m$ 得 $(R_{\Theta, m}^{d_h})^\top R_{\Theta, m}^{d_h} = R_{\Theta, 0}^{d_h} = I$，故 $R_{\Theta, m}^{d_h}$ 为正交矩阵，严格保持 $\norm{R_{\Theta, m}^{d_h} q}_2 = \norm{q}_2$。
\end{proof}
\end{thmBox}

RoPE 的工程实现也极为轻巧：利用分块稀疏性，根本不需要显式构造 $d_h \times d_h$ 矩阵，仅需 $O(d_h)$ 逐元素乘加即可完成：
\[
  R_{\Theta, m}^{d_h} q
  = \begin{pmatrix} q_1 \\ q_2 \\ \vdots \\ q_{d_h-1} \\ q_{d_h} \end{pmatrix}
  \had \begin{pmatrix} \cos m\theta_1 \\ \cos m\theta_1 \\ \vdots \\ \cos m\theta_K \\ \cos m\theta_K \end{pmatrix}
  + \begin{pmatrix} -q_2 \\ q_1 \\ \vdots \\ -q_{d_h} \\ q_{d_h-1} \end{pmatrix}
  \had \begin{pmatrix} \sin m\theta_1 \\ \sin m\theta_1 \\ \vdots \\ \sin m\theta_K \\ \sin m\theta_K \end{pmatrix}.
\]
同时，当多个几何级数频率 $\theta_r$ 叠加时，由 Abel 分部求和可证明内积上界随相对距离 $|n - m|$ 增大而呈现自然的\textbf{远程衰减性}——既保留绝对位置的 $O(1)$ 计算速度，又精确实现了相对位置建模，因此成为 Llama、Qwen、DeepSeek、Mistral 等几乎全部现代开源 LLM 的首选位置编码。

\section{Transformer 块、SwiGLU 与 KV Cache 复杂度分析}
\label{sec:block-kvcache}

\subsection{Pre-LN vs Post-LN 与 SwiGLU 前馈网络}

单个 Transformer 解码器块由“通信层（多头自注意力，沿序列维 $T$ 混合信息）”与“计算层（逐位置前馈网络 FFN，沿通道维 $C$ 提炼特征）”交替叠加而成。如何把归一化层放在残差连接里，直接决定了百层大模型能否稳定训练：
\begin{itemize}
  \item \textbf{Post-LN（2017 原始论文）}：$x' = \LN\big(x + \mathrm{MHA}(x)\big)$。正如第 3 章残差梯度高速公路定理所警示的，外层 $\LN$ 直接横亘在恒等主干道上，深层堆叠后每层都会对残差流反复重缩放，若无极精细的学习率 Warmup，训练极易在初始阶段发散。
  \item \textbf{Pre-LN（GPT-2 / Llama 现代标配）}：
  \begin{equation}
    u_l = x_l + \mathrm{MHA}\!\big(\LN(x_l)\big), \qquad
    x_{l+1} = u_l + \mathrm{FFN}\!\big(\LN(u_l)\big).
    \label{eq:pre-ln}
  \end{equation}
  从第 1 层输入直到最后一层输出之前，存在一条\textbf{完全不被任何非线性或归一化阻断的洁净加法主干道}（仅在顶层分类头之前补加一层终末 $\LN_f$），因此梯度可直通底层！现代架构通常进一步用第 2 章介绍的 \textbf{RMSNorm} 取代 $\LN$，省去减均值开销而提速 $10\% \sim 20\%$。
\end{itemize}

在逐位置 FFN 中，原始 Transformer 采用两层全连接 $\mathrm{FFN}(x) = \ReLU(x W_1 + b_1) W_2 + b_2$（中间升维至 $4d$，总参数 $8d^2$）。现代大模型（PaLM、Llama、Qwen）普遍升级为带门控的 \textbf{SwiGLU}（Shazeer, 2020）：
\begin{equation}
  \mathrm{SwiGLU}(x) = \Big(\operatorname{SiLU}(x W_{\text{gate}}) \had (x W_{\text{up}})\Big) W_{\text{down}},
  \qquad
  \operatorname{SiLU}(z) = z \cdot \sigma(z).
  \label{eq:swiglu}
\end{equation}
因为引入了第三个矩阵 $W_{\text{gate}}$，为了保持总参数量与传统 FFN 的 $8d^2$ 持平，工程上将中间隐藏宽度调整为 $d_{\text{ff}} = \frac{8}{3} d_{\text{model}}$（并常向上圆整至 $256$ 的倍数以对齐 GPU 张量核心），此时参数量为 $3 \cdot d_{\text{model}} \cdot \frac{8}{3} d_{\text{model}} = 8 d_{\text{model}}^2$，在同算力下稳定带来更低的验证困惑度。

图~\ref{fig:transformer-arch} 直观展示了 Pre-LN Transformer 块的洁净残差流与因果掩码矩阵结构。

\begin{figure}[htbp]
\centering
\resizebox{0.97\linewidth}{!}{%
\begin{tikzpicture}[
  box/.style={rectangle, rounded corners=1.5mm, draw=#1, thick, fill=#1!10, minimum width=30mm, minimum height=7.2mm, align=center, font=\sffamily\footnotesize, text=ink},
  add/.style={circle, draw=tealDark, very thick, fill=tealLight, inner sep=1.2pt, font=\bfseries\small, text=tealDark},
  arr/.style={-{Stealth[length=2mm]}, thick, ink!80},
  hwy/.style={-{Stealth[length=2.2mm]}, very thick, tealDark}
]
  % ==================== Panel (a): Pre-LN Transformer Block ====================
  \node[font=\sffamily\bfseries\small, color=tealDark] at (1.0, 6.75) {(a) Pre-LN Transformer 块与洁净残差流};

  \node[font=\sffamily\footnotesize, color=slate] (xin) at (0, -0.3) {输入隐藏状态 $x_l \in \R^{T\times C}$};
  \coordinate (s1) at (0, 0.35);
  \node[box=navy] (ln1) at (0, 1.15) {RMSNorm / LayerNorm};
  \node[box=teal] (mha) at (0, 2.2) {因果多头自注意力（+ RoPE）};
  \node[add] (a1) at (0, 3.15) {$+$};

  \coordinate (s2) at (0, 3.65);
  \node[box=navy] (ln2) at (0, 4.35) {RMSNorm / LayerNorm};
  \node[box=ochre] (ffn) at (0, 5.35) {SwiGLU 前馈网络（FFN）};
  \node[add] (a2) at (2.8, 5.35) {$+$};
  \node[font=\sffamily\footnotesize, color=tealDark] (xout) at (2.8, 6.18) {输出 $x_{l+1}$};

  % Wait: let's route residual cleanly on the right!
  \draw[hwy] (xin) -- (s1);
  \draw[arr] (s1) -- (ln1);
  \draw[arr] (ln1) -- (mha);
  \draw[arr] (mha) -- (a1);
  \draw[hwy] (s1) -- (2.3, 0.35) |- (a1) node[pos=0.25, right=2pt, font=\scriptsize\sffamily, color=tealDark] {恒等主干};

  \draw[hwy] (a1) -- (s2);
  \draw[arr] (s2) -- (ln2);
  \draw[arr] (ln2) -- (ffn);
  \draw[arr] (ffn) -- (a2);
  \draw[hwy] (s2) -| (a2) node[pos=0.72, right=2pt, font=\scriptsize\sffamily, color=tealDark] {恒等主干};
  \draw[hwy] (a2) -- (xout);

  % ==================== Panel (b): Causal Mask Matrix ====================
  \begin{scope}[shift={(6.4, 0.55)}]
    \node[font=\sffamily\bfseries\small, color=tealDark] at (2.3, 6.20) {(b) $5\times 5$ 因果掩码矩阵（$\softmax$ 前）};

    \foreach \i in {0,...,4} {
      \foreach \j in {0,...,4} {
        \pgfmathsetmacro{\x}{\j * 0.92}
        \pgfmathsetmacro{\y}{(4 - \i) * 0.92}
        \ifnum\j>\i
          \fill[ochreLight] (\x, \y) rectangle (\x+0.88, \y+0.88);
          \draw[ochre!55] (\x, \y) rectangle (\x+0.88, \y+0.88);
          \node[font=\scriptsize, color=ochre] at (\x+0.44, \y+0.44) {$-\infty$};
        \else
          \fill[tealLight] (\x, \y) rectangle (\x+0.88, \y+0.88);
          \draw[teal!60] (\x, \y) rectangle (\x+0.88, \y+0.88);
          \node[font=\scriptsize, color=tealDark] at (\x+0.44, \y+0.44) {$z_{\i\j}$};
        \fi
      }
    }
    % Axis labels
    \node[font=\scriptsize\sffamily, color=slate] at (2.3, 4.92) {被关注键位置 $j$（Keys）};
    \node[font=\scriptsize\sffamily, color=slate, rotate=90] at (-0.42, 2.3) {查询位置 $i$（Queries）};
    \foreach \k in {0,...,4} {
      \node[font=\tiny\ttfamily, color=slate] at ({\k*0.92+0.44}, 4.65) {t\k};
      \node[font=\tiny\ttfamily, color=slate] at (-0.18, {(4-\k)*0.92+0.44}) {t\k};
    }
    \node[font=\scriptsize\sffamily, color=ochre, align=center] at (2.3, -0.45) {上三角置 $-\infty \implies P_{ij}=0$ 且 $\partial o_i / \partial k_j = 0$};
  \end{scope}
\end{tikzpicture}}
\caption{Pre-LN Transformer 解码器块与因果掩码原理示意。(a) 两条残差主干道保持纯加法流通，保证深层梯度直达底层；(b) $\softmax$ 前对未来位置施加 $-\infty$ 掩码，从前向与反向同时杜绝未来信息泄露。}
\label{fig:transformer-arch}
\end{figure}

\subsection{预填充 vs 自回归解码：KV Cache 复杂度与显存账本}

训练阶段我们已知全部 $T$ 个词元，整条序列并行算完一次前向即可。但在\textbf{推理生成}时，大模型必须逐词自回归输出，分为两个物理特性迥异的阶段：

\begin{enumerate}
  \item \textbf{预填充阶段（Prefill）}：一次性处理长度为 $T_p$ 的用户输入提示词（Prompt）。此时 $Q, K, V \in \R^{T_p \times C}$ 全部并行计算，注意力矩阵乘法耗时 $O(T_p^2 d_{\text{model}})$，由于矩阵宽大，GPU 处于\textbf{算力受限（Compute-Bound）}的高效状态。同时将每一层算出的所有历史键张量 $K_{1:T_p}$ 与值张量 $V_{1:T_p}$ 缓存在显存中——这便是 \textbf{KV Cache}。
  \item \textbf{逐词解码阶段（Autoregressive Decode）}：生成第 $t$ 个新词元时，由于因果掩码，\textbf{过去 $1, \dots, t-1$ 个词元的键与值向量完全不会随新词的到来而改变}！
  \begin{itemize}
    \item \textbf{若不使用 KV Cache}：每吐出一个新词都要把长度为 $t$ 的完整前缀从头算一遍，单步注意力耗费 $O(t^2 d)$，生成 $T$ 个词总复杂度飙升至惊人的 $\sum_{t=1}^T O(t^2 d) = O(T^3 d)$！
    \item \textbf{启用 KV Cache 后}：只需对\textbf{当前单个新词元} $x_t$ 投影出 $q_t, k_t, v_t \in \R^{1 \times d}$，把 $(k_t, v_t)$ 追加进缓存，再用单向量 $q_t$ 去点乘缓存矩阵 $K_{1:t} \in \R^{t \times d}$ 并加权求和 $V_{1:t}$。单步注意力计算量暴降至 $O(t d)$，生成 $T$ 个词的总计算量回到 $O(T^2 d)$！
  \end{itemize}
\end{enumerate}

但天下没有免费的午餐：KV Cache 用\textbf{显存空间换取计算时间}。对一个 $L$ 层、批大小 $B$、序列长度 $T$、KV 头数 $n_{\text{kv}}$、单头维度 $d_h$、每标量占 $b_{\text{elem}}$ 字节（FP16/BF16 为 $2$ 字节，INT8 为 $1$ 字节）的模型，KV Cache 占用的精确显存为：
\begin{equation}
  \text{Mem}_{\mathrm{KV}} = 2 \cdot L \cdot B \cdot T \cdot n_{\text{kv}} \cdot d_h \cdot b_{\text{elem}},
  \label{eq:kv-cache-mem}
\end{equation}
其中最前面的因子 $2$ 对应 $K$ 与 $V$ 两份缓存。

\begin{example}[Llama-3 8B 规模模型的 KV Cache 显存精算]
设 $L = 32,\; n_q = 32,\; d_h = 128$，使用 BF16（$b_{\text{elem}} = 2$ 字节），单条请求（$B = 1$）上下文长度 $T = 8192$：
\begin{itemize}
  \item 若采用标准 \textbf{MHA}（$n_{\text{kv}} = 32$）：
  \[
    \text{Mem}_{\mathrm{KV}}^{\mathrm{MHA}} = 2 \times 32 \times 1 \times 8192 \times 32 \times 128 \times 2 \text{ B} = 4{,}294{,}967{,}296 \text{ B} = 4.00\text{ GiB} \approx 4.29\text{ GB}.
  \]
  当并发批大小升至 $B = 32$ 时，仅 KV Cache 就要吞掉 \textbf{128 GiB} 显存，远远超过 8B 模型权重本身（约 16 GB）！
  \item 采用 \textbf{GQA}（$n_{\text{kv}} = 8$）后：单条请求的 KV Cache 直接缩减 $4$ 倍至 \textbf{1.00 GiB}（$1.07\text{ GB}$），使并发吞吐瞬间翻四番！这也是第 6 卷第 5--6 章将深入剖析的 PagedAttention（vLLM）与 KV 量化优化的直接起点。
\end{itemize}
\end{example}

\section{实测实验：从零手写迷你 GPT 学会序列反转}
\label{sec:exp-minigpt}

纸上谈兵终觉浅。下面我们直接用不到 45 行纯 PyTorch 代码，不调用 \texttt{nn.Transformer}，亲手组装一个具备完整因果多头自注意力与 Pre-LN 残差流的 \textbf{Mini-GPT}（来自配套脚本 \texttt{code/vol5\_experiments.py} 中的 \texttt{exp4\_mini\_gpt\_reverse}），并在经典的算法推理任务——\textbf{6 词序列反转}上检验它的真本领！

\textbf{任务构造}：词表大小为 $12$（数字词元 \texttt{0..9}、分隔符 \texttt{10} 记作 \texttt{|}、填充符 \texttt{11}）。每条样本随机生成 6 个数字并拼接其倒序，例如：
\[
  \texttt{3 8 1 0 7 4 | 4 7 0 1 8 3}
\]
总长度 $13$，自回归输入 $x = s_{0:11}$、目标 $y = s_{1:12}$（长度 $12$）。\textbf{交叉熵损失仅在分隔符之后的 6 个反转答案位置上计算}——因为前 6 个数字是完全随机生成的，不可预测。要答对这 6 个位置，模型必须学会用注意力跨越不定距离（从间隔 $1$ 步到间隔 $11$ 步）精准定位对称位置并搬运词元！

\begin{lstlisting}[caption={纯 PyTorch 手写因果多头自注意力与 Pre-LN 迷你 GPT（与 \texttt{vol5\_experiments.py} 一致）},label={lst:minigpt}]
import math, torch, torch.nn as nn, torch.nn.functional as F

class CausalSelfAttn(nn.Module):
    def __init__(self, d_model: int, n_head: int, max_len: int):
        super().__init__()
        self.n_head, self.d_head = n_head, d_model // n_head
        self.qkv = nn.Linear(d_model, 3 * d_model, bias=False)
        self.proj = nn.Linear(d_model, d_model, bias=False)
        mask = torch.tril(torch.ones(max_len, max_len)).view(1, 1, max_len, max_len)
        self.register_buffer("mask", mask)

    def forward(self, x: torch.Tensor) -> torch.Tensor:
        B, T, C = x.shape
        q, k, v = self.qkv(x).chunk(3, dim=-1)
        q = q.view(B, T, self.n_head, self.d_head).transpose(1, 2)  # (B,H,T,d_h)
        k = k.view(B, T, self.n_head, self.d_head).transpose(1, 2)
        v = v.view(B, T, self.n_head, self.d_head).transpose(1, 2)
        att = (q @ k.transpose(-2, -1)) / math.sqrt(self.d_head)    # 缩放点积
        att = att.masked_fill(self.mask[:, :, :T, :T] == 0, float("-inf"))
        att = F.softmax(att, dim=-1)                                # 沿键维归一化
        y = (att @ v).transpose(1, 2).contiguous().view(B, T, C)
        return self.proj(y)

class MiniGPT(nn.Module):
    def __init__(self, vocab=12, seq_len=13, d_model=48, n_head=4, n_layer=2):
        super().__init__()
        self.tok = nn.Embedding(vocab, d_model)
        self.pos = nn.Embedding(seq_len, d_model)
        self.blocks = nn.ModuleList([nn.ModuleDict({
            "ln1": nn.LayerNorm(d_model),
            "attn": CausalSelfAttn(d_model, n_head, seq_len),
            "ln2": nn.LayerNorm(d_model),
            "ffn": nn.Sequential(nn.Linear(d_model, 4*d_model),
                                 nn.GELU(), nn.Linear(4*d_model, d_model)),
        }) for _ in range(n_layer)])
        self.ln_f = nn.LayerNorm(d_model)
        self.head = nn.Linear(d_model, vocab, bias=False)

    def forward(self, idx: torch.Tensor) -> torch.Tensor:
        B, T = idx.shape
        x = self.tok(idx) + self.pos(torch.arange(T, device=idx.device))
        for blk in self.blocks:
            x = x + blk["attn"](blk["ln1"](x))    # Pre-LN 残差注意力
            x = x + blk["ffn"](blk["ln2"](x))     # Pre-LN 残差前馈
        return self.head(self.ln_f(x))
\end{lstlisting}

我们配置 $d_{\text{model}} = 48,\; H = 4$（单头维度 $d_h = 12$）、$L = 2$ 层，采用 AdamW（$\eta = 2\times 10^{-3},\; \lambda = 10^{-2}$）、批大小 $64$ 在线采样训练 $320$ 步，并在 $256$ 条独立验证序列上追踪两项严格区分的指标：\textbf{教师强制（teacher-forcing）逐词准确率}（\texttt{val\_tok\_acc}）与\textbf{真实自回归贪心生成的整句精确匹配率}（\texttt{val\_seq\_acc}——仅向模型喂入原序列与分隔符 \texttt{[src, sep]}，由模型自己逐词自回归解码 6 步，全程绝不泄露任何真实答案前缀！）。实测数据 \texttt{data/transformer\_train.csv} 记录了极具戏剧性的“顿悟式”收敛（见图~\ref{fig:minigpt-train}）：
\begin{itemize}
  \item \textbf{第 0 步（随机初始化）}：交叉熵损失 \texttt{2.6166}（接近 $\ln 12 \approx 2.485$ 的盲猜基线），逐词准确率仅 \texttt{8.33\%}，自回归整句精确率 \texttt{0.00\%}；
  \item \textbf{第 20 步}：损失降至 \texttt{2.0962}，逐词准确率升至 \texttt{34.38\%}，但自回归 6 词全对率仍为 \texttt{0.00\%}（模型还在摸索镜像位置索引）；
  \item \textbf{第 40 步（相变跃迁！）}：一旦注意力头锁定了“位置 $6+k$ 对齐位置 $5-k$”的复制电路，损失瞬间暴跌至 \texttt{0.2933}，逐词准确率飙升至 \texttt{99.87\%}，真实自回归整句精确率一举跃上 \texttt{99.22\%}；
  \item \textbf{第 60--80 步}：损失进一步跌破 \texttt{0.0321} 与 \texttt{0.0149}，教师强制逐词准确率与真实自回归整句精确率双双锁死在 \textbf{\texttt{100.00\%}}；
  \item \textbf{第 320 步}：交叉熵损失收敛至 \texttt{0.0018}，自回归整句准确率稳稳保持 \textbf{\texttt{100.00\%}}——一个参数量不足 6 万的纯手写迷你 GPT，彻底掌握了任意未见 6 位数序列的独立逐词倒序生成！
\end{itemize}

\begin{figure}[htbp]
\centering
\begin{tikzpicture}
  % Left plot: Cross-Entropy Loss
  \begin{axis}[
    name=lossplot,
    width=0.47\linewidth, height=5.8cm,
    title={\sffamily\bfseries\footnotesize (a) 训练交叉熵损失收敛曲线},
    xlabel={训练步数（Step）},
    ylabel={交叉熵损失 $\mathcal{L}_{\mathrm{CE}}$},
    xmin=-6, xmax=326, ymin=-0.08, ymax=2.85,
    xtick={0,80,160,240,320},
    grid=both, grid style={tealLine!45, dashed},
    axis line style={slate},
    tick label style={font=\scriptsize},
    label style={font=\footnotesize\sffamily, color=ink}
  ]
    \addplot[ochre, very thick, mark=*, mark size=1.6pt]
      table[x=step, y=loss, col sep=comma]{data/transformer_train.csv};
    \node[font=\scriptsize\sffamily, color=ochre, anchor=west] at (axis cs:55, 2.35) {step=0: 2.6166};
    \node[font=\scriptsize\sffamily, color=tealDark, anchor=south west] at (axis cs:92, 0.22) {step=80: 0.0149};
    \draw[-{Stealth[length=1.8mm]}, tealDark, thin] (axis cs:105, 0.22) -- (axis cs:82, 0.05);
  \end{axis}

  % Right plot: Val Token Acc & Exact Seq Acc
  \begin{axis}[
    at={(lossplot.east)}, anchor=west, xshift=1.45cm,
    width=0.47\linewidth, height=5.8cm,
    title={\sffamily\bfseries\footnotesize (b) 验证集逐词与整句精确匹配率},
    xlabel={训练步数（Step）},
    ylabel={验证准确率（\%）},
    xmin=-6, xmax=326, ymin=-3, ymax=108,
    xtick={0,80,160,240,320},
    ytick={0,25,50,75,100},
    grid=both, grid style={tealLine!45, dashed},
    axis line style={slate},
    tick label style={font=\scriptsize},
    label style={font=\footnotesize\sffamily, color=ink},
    legend style={at={(0.97,0.22)}, anchor=east, font=\scriptsize, fill=white, draw=tealLine, legend cell align=left}
  ]
    \addplot[teal, very thick, mark=*, mark size=1.6pt]
      table[x=step, y=val_tok_acc, col sep=comma]{data/transformer_train.csv};
    \addlegendentry{逐词准确率（Tok Acc）}

    \addplot[navy, very thick, dashed, mark=square*, mark size=1.5pt]
      table[x=step, y=val_seq_acc, col sep=comma]{data/transformer_train.csv};
    \addlegendentry{整句精确率（Seq Acc）}

    \node[font=\scriptsize\sffamily\bfseries, color=tealDark, fill=white, inner sep=1pt]
      at (axis cs:195, 86) {step $\ge 60$: 100.00\%};
  \end{axis}
\end{tikzpicture}
\caption{纯 PyTorch 2 层迷你 GPT 在 6 词序列反转任务上的实测训练曲线（读取自 \texttt{data/transformer\_train.csv}）。第 20--40 步间模型完成复制注意力电路的“顿悟”相变，80 步后损失降至 \texttt{0.0149}，整句准确率锁定 \texttt{100.00\%}。}
\label{fig:minigpt-train}
\end{figure}

\begin{failBox}{手写 Transformer 最致命的三个隐蔽 Bug}
\begin{enumerate}
  \item \textbf{忘记因果掩码导致训练作弊与推理崩盘}：若漏掉 \texttt{masked\_fill}，位置 $t$ 会直接抄写位置 $t+1$ 的答案词元，训练损失三步跌至 $0$；可一旦进入自回归采样循环，模型从未学过从前文推理，只能胡言乱语。自检金标准：随机扰动输入张量第 $t+1$ 步之后的所有元素，检查位置 $t$ 的前向输出差值是否严格为 \texttt{0.0}！
  \item \textbf{多头拆分的错误 \texttt{view} 顺序静默打乱序列}：新手常把 \texttt{q.view(B, T, H, d\_h).transpose(1, 2)} 误写成一行 \texttt{q.view(B, H, T, d\_h)}。后者代码\textbf{不会报错}，却会把不同时间步的特征强行切片拼凑进错误的头里，彻底摧毁语义对齐！
  \item \textbf{把 $\softmax$ 算在错误的维度上}：对形状为 \texttt{(B, H, T, T)} 的注意力矩阵，必须执行 \texttt{F.softmax(att, dim=-1)}（沿键维度使每行之和为 $1$）。若错写成 \texttt{dim=-2}（每列之和为 $1$），不仅每个查询的权重之和不再归一，还会把未来行的数值揉进当前行的分母里，双重破坏因果性！
\end{enumerate}
\end{failBox}

\begin{extBox}{打破 $O(T^2)$ 内存墙：FlashAttention 预告与线性注意力 / 状态空间对偶}
标准注意力实现必须把 $T \times T$ 的注意力矩阵 $P$ 写入显存（HBM）再读回，当上下文长达 $T = 128\text{K}$ 时，仅单头注意力矩阵就需要 $128\text{K}^2 \times 2\text{ B} = 32\text{ GB}$ 显存！现代系统与算法正从两条路线攻克此墙：
\begin{itemize}
  \item \textbf{FlashAttention（Dao et al., 2022/2023；精确注意力，零近似）}：利用 GPU 片上极高速 SRAM 实施\textbf{分块平铺（Tiling）}，结合第 1 章提到的在线安全 $\softmax$ 递推，边读取小块 $Q_i, K_j, V_j$ 边动态维护运行最大值 $m_i$ 与归一分母 $\ell_i$，\textbf{全程从不将 $O(T^2)$ 矩阵落盘至 HBM}！反向传播时仅重算局部块，使显存复杂度从 $O(T^2)$ 骤降至 $O(T)$，实测提速 $2\sim 4\times$（详见第 6 卷第 5 章）。
  \item \textbf{线性注意力与结构化状态空间对偶（SSD / Mamba-2）}：若将 $\exp(q_i^\top k_j)$ 近似为特征映射内积 $\phi(q_i)^\top \phi(k_j)$，由矩阵乘法结合律：
  \[
    \sum_{j=1}^i \big(\phi(q_i)^\top \phi(k_j)\big) v_j
    = \phi(q_i)^\top \underbrace{\sum_{j=1}^i \phi(k_j) v_j^\top}_{S_i \in \R^{d \times d}},
  \]
  只需维护一个固定大小的状态矩阵 $S_i = S_{i-1} + \phi(k_i) v_i^\top$，复杂度一举化为线性 $O(T d^2)$！Dao 与 Gu（2024）进一步揭示：带衰减因子的线性注意力在矩阵形式下与第 3 章的 Mamba 选择性状态空间模型存在严格的\textbf{数学对偶性}——注意力与循环递归在半个世纪后终于在统一框架下合流！
\end{itemize}
\end{extBox}

\section*{本章小结与核心脉络}
\addcontentsline{toc}{section}{本章小结与核心脉络}

本章彻底拆解了现代基础模型的心脏——Transformer：
\begin{itemize}
  \item \textbf{缩放点积注意力}以 $O(1)$ 路径长度打通全局交互，分母 $\sqrt{d_k}$ 精确抵消高维内积方差膨胀，避免 $\softmax$ 雅可比矩阵 $\diag(p) - pp^\top$ 陷入饱和零梯度区；
  \item \textbf{多头、因果掩码与 GQA} 分别解决了多子空间表征、自回归零未来梯度泄露与长文本显存瘦身三大工程命题；
  \item \textbf{RoPE} 借由正交旋转群恒等式 $R_\theta(m)^\top R_\theta(n) = R_\theta(n-m)$，以 $O(d)$ 代价在绝对位置框架内精确实现了相对位置内积；
  \item \textbf{Pre-LN、SwiGLU 与 KV Cache} 将训练稳定性与推理吞吐推向极致，而 45 行纯 PyTorch 手写迷你 GPT 在 80 步内斩获 $100\%$ 序列反转精确率，印证了这一架构的惊人威力。
\end{itemize}
有了如此强大的通用架构之后，下一个核心问题是：\textbf{没有海量人工标签时，如何喂养它成长为通识大模型，又该如何最优分配模型规模与数据算力？} 这正是第 5 章的主题。

\section*{本章分层习题}
\addcontentsline{toc}{section}{本章分层习题}

\exTier{A 层：概念与手算热身}{掌握缩放注意力、参数量估算与 KV Cache 显存账本}

\begin{enumerate}
  \item[\textbf{5.4.1}] \textbf{因果缩放点积注意力完整手算}：给定长度 $T=3$、维度 $d_k = d_v = 2$ 的序列矩阵：
  \[
    Q = \begin{pmatrix} 1 & 0 \\ 0 & 1 \\ 1 & 1 \end{pmatrix}, \qquad
    K = \begin{pmatrix} 1 & 0 \\ 1 & 1 \\ 0 & 1 \end{pmatrix}, \qquad
    V = \begin{pmatrix} 2 & 0 \\ 0 & 2 \\ 1 & 1 \end{pmatrix}.
  \]
  分别计算\textbf{无因果掩码}与\textbf{带因果掩码}（式~\eqref{eq:causal-mask}）时的缩放 logits 矩阵 $Z$、权重矩阵 $P$ 及最终输出矩阵 $O$（保留三位小数）。对比第 1 行输出的变化。
  \exHint{带因果掩码时，第 1 行只能看见 $j=1$，故 $P_{1,:} = (1, 0, 0)$，输出恒等于 $v_1 = (2, 0)$！}

  \item[\textbf{5.4.2}] \textbf{70B 大模型的参数结构与 KV Cache 精算}：某现代解码器大模型配置为 $L = 80$ 层、$d_{\text{model}} = 8192$、查询头 $n_q = 64$、KV 头 $n_{\text{kv}} = 8$（$d_h = 128$）、SwiGLU 中间维 $d_{\text{ff}} = 28672$、词表 $|V| = 128{,}000$（不共享输入/输出嵌入）。
  \begin{enumerate}
    \item 估算该模型的总可训练参数量（忽略 RMSNorm 极少量参数）；
    \item 若以 BF16（2 字节）部署，上下文长度 $T = 32{,}768$、并发批大小 $B = 16$，依据式~\eqref{eq:kv-cache-mem} 精确计算 KV Cache 显存占用（以 GiB 为单位），并与若改用标准 MHA（$n_{\text{kv}} = 64$）时的显存对比。
  \end{enumerate}
  \exHint{每层 GQA 投影参数为 $W_Q(d^2) + W_K(d \cdot n_{\text{kv}} d_h) + W_V(d \cdot n_{\text{kv}} d_h) + W_O(d^2)$；SwiGLU 为 $3 d \cdot d_{\text{ff}}$。}
\end{enumerate}

\exTier{B 层：定理推导与严谨证明}{深入注意力梯度流、RoPE 远程衰减与跨层信息容量}

\begin{enumerate}
  \item[\textbf{5.4.3}] \textbf{Softmax 雅可比半正定性与缩放点积注意力完整 VJP}：
  \begin{enumerate}
    \item 利用 Cauchy--Schwarz 不等式证明：对任意概率分布 $p \in \Delta^{T-1}$ 及任意向量 $u \in \R^T$，式~\eqref{eq:softmax-jacobian} 给出的矩阵 $J_{\softmax} = \diag(p) - p p^\top$ 是对称半正定矩阵（$u^\top J_{\softmax} u \ge 0$），且其零空间恰好由全一向量 $\mathbf{1}$ 张成；
    \item 由此解释为何给注意力 logits 任意一行同时加上标量常数 $c$（如为了数值稳定而减去行最大值 $\max_j Z_{ij}$），不仅前向概率保持不变，反向梯度 $\bar{Z}$ 也严格满足 $\sum_{j=1}^T \bar{Z}_{ij} = 0$。
  \end{enumerate}
  \exHint{把 $u^\top (\diag(p) - pp^\top) u$ 改写成以 $p$ 为概率测度的方差 $\Var_{j \sim p}(u_j) = \E_p[u^2] - (\E_p[u])^2 \ge 0$！}

  \item[\textbf{5.4.4}] \textbf{RoPE 编码的相对距离远程衰减上界}：设 $d_h = 2K$，角频率 $\theta_r = 10000^{-2(r-1)/d_h}$，记 $\tilde{q}_r, \tilde{k}_r \in \C$ 为第 $r$ 个二维子平面上的复数表示，相对距离为 $\Delta = n - m$。定义 $h_r := \tilde{q}_r^* \tilde{k}_r$ 及部分和 $S_r := \sum_{j=1}^r e^{\mathrm{i} \Delta \theta_j}$（$S_0 = 0$）。利用 Abel 分部求和恒等式：
  \[
    \sum_{r=1}^{K} h_r e^{\mathrm{i} \Delta \theta_r}
    = h_K S_K - \sum_{r=1}^{K-1} (h_{r+1} - h_r) S_r,
  \]
  证明 RoPE 注意力内积绝对值受控于 $\max_r |h_{r+1} - h_r| \sum_{r=1}^{K} |S_r|$（外加边界项），并说明为何当频率 $\theta_r$ 呈几何级数分布时，$\sum_r |S_r|$ 随相对距离 $|\Delta|$ 增大呈现振荡衰减趋势。
  \exHint{高、中、低频单位复数 $e^{\mathrm{i}\Delta\theta_r}$ 随 $|\Delta|$ 增大相位错开，相消干涉导致部分和 $|S_r|$ 显著小于 $r$。}
\end{enumerate}

\exTier{C 层：代码实现与数值实验}{手写 RoPE、无缩放消融对比与 KV Cache 增量解码验证}

\begin{enumerate}
  \item[\textbf{5.4.5}] \textbf{无 $\sqrt{d_k}$ 缩放消融实验}：复制代码~\ref{lst:minigpt} 中的 \texttt{MiniGPT}，将模型宽度放大至 $d_{\text{model}} = 256,\; H = 2$（即单头维度 $d_h = 128$），分别在开启与关闭 \texttt{/ math.sqrt(self.d\_head)} 两种设置下训练序列反转任务。统计前 100 步中注意力权重矩阵的\textbf{平均信息熵} $H(P) = -\frac{1}{T}\sum_{i,j} P_{ij}\ln P_{ij}$ 以及 $W_Q$ 梯度范数，定量验证定理~\ref{thm:sqrt-dk} 的预测。
  \exHint{未缩放时初始注意力熵接近 $0$（退化为单热尖峰），梯度范数较缩放组小几个数量级。}

  \item[\textbf{5.4.6}] \textbf{为手写 \texttt{MiniGPT} 添加 RoPE 与增量 KV Cache 推理}：
  \begin{enumerate}
    \item 去掉 \texttt{self.pos} 绝对位置嵌入，用 $O(d_h)$ 的 \texttt{cos/sin} 逐元素旋转为 \texttt{CausalSelfAttn} 接入 RoPE，并验证相对位置内积不变性（误差 $< 10^{-6}$）；
    \item 给 \texttt{CausalSelfAttn.forward} 新增可选参数 \texttt{kv\_cache}，支持单词增量解码。对比“每次重算全前缀”与“启用 KV Cache 逐步生成”输出的最后一步 logits，确保两者最大绝对误差 $< 10^{-5}$！
  \end{enumerate}
  \exHint{增量解码第 $t$ 步时，不要忘记把 RoPE 的位置索引设为当前绝对步号 $t$（而非 $0$），并且单步查询 $q_t$ 无需再加下三角掩码（它合法可见缓存中的全部历史键 $K_{0:t}$）。}
\end{enumerate}

\exTier{D 层：开放探究与前沿迁移}{深入长上下文外推与注意力回路的机械可解释性}

\begin{enumerate}
  \item[\textbf{5.4.7}] \textbf{长上下文外推：位置内插（PI）与 NTK-Aware / YaRN 缩放}：一个在长度 $T_{\text{train}} = 4096$ 上用 RoPE 训练的 LLM，若直接推理 $T_{\text{test}} = 16384$，模型困惑度会瞬间崩溃。查阅 Chen 等人（2023）的\textbf{位置内插（Position Interpolation）}（将 $m$ 替换为 $m / 4$）与 \textbf{YaRN}（按频率高低分区缩放：高频维度保持不变以保留局部词序分辨率，低频维度做内插）思想，从定理~\ref{thm:rope} 的多频时钟视角剖析为何“高频不插、低频内插”远优于均匀内插。
  \exHint{均匀缩小 $m \to m/4$ 会把相邻词元的微小高频相位差压缩 $4$ 倍，导致模型分不清紧邻两个词的先后关系！}

  \item[\textbf{5.4.8}] \textbf{归纳头（Induction Heads）与上下文学习回路}：Olsson 等人（2022）在研究 Transformer 机械可解释性时发现，两层及以上自注意力模型会自发涌现一种由两层协同构成的\textbf{归纳头（Induction Head）}电路：第 1 层的“前词头”把位置 $i-1$ 的词元复制到位置 $i$ 的残差流中；第 2 层的“匹配头”看到当前词元 \texttt{A} 时，在历史中寻找“前一个词是 \texttt{A}”的位置（即曾紧跟在 \texttt{A} 后的词元 \texttt{B}），从而在上下文中直接完成模式补全 \texttt{[A][B] ... [A] -> [B]}。试结合本章迷你 GPT 的实验现象（第 20--40 步的相变跃迁），分析为什么单层无位置协同的注意力难以实现复杂多跳复制，而两层残差流组合如何构成可编程的符号电路？
  \exHint{将残差流视作所有注意力头共同读写的“共享内存总线”，第 1 层写入位置/邻接标签，第 2 层按该标签发起第二跳查询。}
\end{enumerate}

\chapter[自监督预训练、对比学习与缩放律]{自监督表示学习、预训练范式与缩放定律}
\label{ch:pretrain-scaling}

\chapEpigraph{观其伴，知其词。}{J. R. Firth（约翰·鲁珀特·弗斯，1957）}

\begin{chapIntro}{本章导读：从“人工标注红利”迈向“无监督规模红利”}
在深度学习的前十年（2012--2018），几乎所有耀眼的突破都建立在\textbf{有监督学习}之上：人类花费数年时间为百万张图像贴上“波斯猫”“消防车”标签，神经网络再通过交叉熵拟合这些离散标签。但有监督范式很快触及天花板——人工标注既昂贵又充满偏见，且标签里的信息量极其稀疏（一张百万像素图像往往只配了一个 $10$ 比特的类别编号）。

真正的范式大迁徙始于\textbf{自监督预训练（Self-Supervised Pretraining）}：既然互联网上流淌着数以万亿计的原始文本、代码与图文对，我们完全可以把\textbf{数据本身的一部分遮盖起来，让模型用另一部分去预测它}——无需任何人工标注，便能把整个人类文明的语言、逻辑与视觉常识压缩进神经网络的权重中！本章将系统建立这一当代基础模型范式的四大支柱：
\begin{enumerate}
  \item \textbf{词元化与 BPE 算法}：在微型语料上手工推演字节对编码（BPE）的逐轮合并过程，剖析词表大小 $|V|$ 与序列长度 $T$ 之间的系统权衡；
  \item \textbf{语言预训练双轨}：深入剖析双向掩码去噪（BERT / MLM）的训练--推理失配难题，以及因果自回归（GPT / CLM）的概率链式分解与困惑度（Perplexity）信息论本质；
  \item \textbf{对比学习与 InfoNCE 互信息下界定理}：从第一性原理给出\textbf{零近似的有限样本严格证明} $I(X; Y) \ge \ln N - \mathcal{L}_N^{\mathrm{InfoNCE}}$，阐明超球面 $\mathbb{S}^{d-1}$ 上的对齐--均匀性平衡与温度 $\tau$ 的硬负样本聚焦效应，并用双模态 CLIP 风格实测（\texttt{data/infonce\_train.csv}）验证 Top-1 检索率从 $0\%$ 跃升至 $100.00\%$；
  \item \textbf{神经缩放定律与 Chinchilla 最优算力分配}：调用第 4 卷第 2 章的拉格朗日乘数法，严格推导出在固定算力预算 $C \approx 6ND$ 下的最优参数量与数据量幂律指数 $a = \frac{\beta}{\alpha+\beta},\; b = \frac{\alpha}{\alpha+\beta}$ 与 $D^* \approx 20 N^*$ 黄金法则。
\end{enumerate}
\end{chapIntro}

\section{从监督标注到词元化：BPE 算法与词表权衡}
\label{sec:bpe}

神经网络只能吞吐整数索引序列与浮点数张量，无法直接阅读原始字符串。如何把连续文本切分为离散词元（Token）——即\textbf{词元化（Tokenization）}——是大模型输入管道的第一道精密阀门。

两种朴素方案都存在致命短板：
\begin{itemize}
  \item \textbf{词级切分（Word-level）}：按空格和标点切词。英语屈折变化（\texttt{play/played/playing}）、拼写错误、代码标识符与专有名词会让词表轻松膨胀至数百万，不仅撑爆嵌入矩阵 $E \in \R^{|V|\times d}$，还会对一切未登录词（OOV）退化为丧失信息的 \texttt{<UNK>}；
  \item \textbf{字符级切分（Char-level）}：词表极小（如 ASCII 仅 128 个），绝无 OOV，但一个词被拆成 $5\sim 10$ 个字符，序列长度 $T$ 暴增数倍——对第~\ref{ch:attention-transformer} 章注意力 $O(T^2)$ 复杂度而言无异于性能灾难。
\end{itemize}

现代大语言模型（GPT-2/3/4、Llama、Qwen、DeepSeek）一致采用\textbf{子词（Subword）词元化}，其中最核心的算法是由 Gage（1994）发明、Sennrich 等人（2016）引入 NLP 并由 GPT-2 推广至字节级的\textbf{字节对编码（Byte Pair Encoding, BPE）}。

\begin{defBox}
\begin{definition}[字节对编码（BPE）贪心合并算法]
给定训练语料（预先按空白/正则粗切为带频次的词列表）及目标词表大小 $|V_{\text{target}}|$：
\begin{enumerate}
  \item \textbf{初始化}：将每个词拆分为基本符号序列（字符或 UTF-8 字节，共 256 个基础字节，确保任何语言与表情符号都\textbf{永不出现 \texttt{<UNK>}}），并在词尾附加边界标记 \texttt{</w>}（或 GPT-2 的词首空格符号），形成初始词表 $V_0$；
  \item \textbf{贪心迭代}：在当前语料中统计所有\textbf{相邻符号对} $(u, v)$ 的出现总频次，选出频次最高的一对 $(u^*, v^*)$，将其合并为新符号 $w = u^* v^*$ 加入词表，并把语料中所有相邻的 $(u^*, v^*)$ 替换为 $w$；
  \item \textbf{终止}：重复步骤 2 直至 $|V| = |V_{\text{target}}|$，并按顺序保存全部合并规则表（Merge Table）。推理编码新文本时，严格按合并规则表的优先级依次合并即可。
\end{enumerate}
\end{definition}
\end{defBox}

\begin{example}[手算全过程：微型语料上的 5 轮 BPE 合并追踪]
\label{ex:bpe-trace}
设微型语料包含 4 个词及其频次：\texttt{low}（5 次）、\texttt{lower}（2 次）、\texttt{newest}（6 次）、\texttt{widest}（3 次）。添加词尾标记 \texttt{</w>} 后，初始字符词表为 $\{d, e, i, l, n, o, r, s, t, w, \texttt{</w>}\}$（共 11 个符号），语料表示为：
\begin{center}
\footnotesize
\begin{tabular}{@{}lclc@{}}
\toprule
\textbf{初始切分} & \textbf{频次} & \textbf{5 轮合并后切分} & \textbf{频次} \\
\midrule
\texttt{l o w </w>} & 5 & \texttt{low </w>} & 5 \\
\texttt{l o w e r </w>} & 2 & \texttt{low e r </w>} & 2 \\
\texttt{n e w e s t </w>} & 6 & \texttt{n e w est</w>} & 6 \\
\texttt{w i d e s t </w>} & 3 & \texttt{w i d est</w>} & 3 \\
\bottomrule
\end{tabular}
\end{center}
逐轮统计最高频相邻对并合并：
\begin{enumerate}
  \item \textbf{第 1 轮}：\texttt{(e, s)} 在 \texttt{newest}（6）与 \texttt{widest}（3）中共现 $6+3 = 9$ 次（全场最高）$\implies$ 合并为 \texttt{es}；
  \item \textbf{第 2 轮}：\texttt{(es, t)} 共现 $6+3 = 9$ 次 $\implies$ 合并为 \texttt{est}；
  \item \textbf{第 3 轮}：\texttt{(est, </w>)} 共现 $6+3 = 9$ 次 $\implies$ 合并为 \texttt{est</w>}（自动习得了英语最高级后缀！）；
  \item \textbf{第 4 轮}：\texttt{(l, o)} 在 \texttt{low}（5）与 \texttt{lower}（2）中共现 $5+2 = 7$ 次 $\implies$ 合并为 \texttt{lo}；
  \item \textbf{第 5 轮}：\texttt{(lo, w)} 共现 $5+2 = 7$ 次 $\implies$ 合并为 \texttt{low}（自动习得了词根！）。
\end{enumerate}
经过仅仅 5 轮合并，词表从 11 个符号温和增至 16 个，但整部语料的总词元数从 $5\times 4 + 2\times 6 + 6\times 7 + 3\times 7 = 95$ 骤降至 $5\times 2 + 2\times 4 + 6\times 4 + 3\times 4 = 54$（压缩率达 $43\%$）！更妙的是，即使测试时遇到从未见过的生词 \texttt{lowest}，该规则表也会优雅地切分为 \texttt{[low, est</w>]}——精准复用了“词根 + 最高级后缀”的形态学结构！
\end{example}

\begin{insightBox}{词表大小 $|V|$ 与序列长度 $T$ 的现代工程权衡}
为什么早期 GPT-2 用 $|V| \approx 50\text{K}$，而现代 Llama-3 与 Qwen-2.5 纷纷将词表扩大至 $128\text{K} \sim 152\text{K}$？
\begin{itemize}
  \item \textbf{大词表的收益}：更低的\textbf{词元生育率（Fertility，每词或每汉字平均所需词元数）}，大幅缩短同等文本的序列长度 $T$——尤其对中文、代码和多语言文本，能将推理生成速度与等效上下文窗口提升 $30\% \sim 50\%$；
  \item \textbf{大词表的代价}：输入嵌入与顶层 LM Head 矩阵各占 $|V| \cdot d_{\text{model}}$ 参数（若 $d_{\text{model}}=4096,\; |V|=128\text{K}$，仅这两张表就占超过 $10$ 亿参数，且末层 $\softmax$ 需在 $128\text{K}$ 维上计算）；此外，过于庞大的词表会制造大量低频“幽灵词元”，在有限预训练数据下得不到充分梯度更新。
\end{itemize}
\end{insightBox}

\section{语言预训练双雄：掩码去噪（BERT）与因果自回归（GPT）}
\label{sec:mlm-clm}

语言学家 Firth 在 1957 年写下名言：“观其伴，知其词（You shall know a word by the company it keeps）。” 现代大语言模型的全部魔力，正是把这句分布式语义假说翻译成了两种可微分的自监督损失函数。

\subsection{掩码语言建模（BERT / MLM）：双向去噪与训练--推理失配}

BERT（Devlin et al., 2019）采用\textbf{仅编码器（Encoder-only）}双向 Transformer（无因果掩码，每个位置都能同时看见左侧与右侧全部上下文）。训练时，从输入序列 $x = (x_1, \dots, x_T)$ 中随机挑选 $15\%$ 的位置集合 $\mathcal{M}$，构造损坏输入 $\tilde{x}$，最大化被遮盖词元的条件对数似然：
\begin{equation}
  \mathcal{L}_{\mathrm{MLM}}(\theta) = -\frac{1}{|\mathcal{M}|} \sum_{i \in \mathcal{M}} \ln p_\theta\!\left(x_i \mid \tilde{x}\right).
  \label{eq:mlm-loss}
\end{equation}
MLM 的双向语境使其在文本分类、命名实体识别与语义检索（句子嵌入）上极为出色，但它有两项与生俱来的局限：
\begin{enumerate}
  \item \textbf{\texttt{[MASK]} 训练--推理失配}：预训练时输入里充斥着人工占位符 \texttt{[MASK]}，可下游微调和真实推理时句子里根本不会出现 \texttt{[MASK]}！为此 BERT 发明了著名的 \textbf{80--10--10 补丁}（被选中的 $15\%$ 词元中，仅 $80\%$ 替换为 \texttt{[MASK]}，$10\%$ 替换为随机词，$10\%$ 保持原词不变），但也仅能缓解而无法根除失配，且无法直接按联合概率逐词生成流畅长文。
  \item \textbf{样本计算效率偏低}：每跑完一整趟前向传播，只有 $15\%$ 的位置产生损失与梯度信号，其余 $85\%$ 位置的隐藏输出未被直接用于预测监督。
\end{enumerate}

\subsection{因果自回归语言建模（GPT / CLM）与困惑度的信息论本质}

GPT 系列（Radford et al., 2018--2020）回归概率论最纯粹的\textbf{乘法链式法则}：任何长度为 $T$ 的文本序列的联合概率都可以\textbf{零近似地精确分解}为一系列下一词条件概率之积：
\begin{equation}
  p_\theta(x_1, x_2, \dots, x_T) = \prod_{t=1}^{T} p_\theta(x_t \mid x_{<t})
  \;\iff\;
  \mathcal{L}_{\mathrm{CLM}}(\theta) = -\frac{1}{T}\sum_{t=1}^{T} \ln p_\theta(x_t \mid x_{<t}).
  \label{eq:clm-loss}
\end{equation}
配合第~\ref{ch:attention-transformer} 章的因果掩码，单次并行前向传播就能同时对序列中\textbf{全部 $100\%$ 词元}施加监督信号，而且不存在任何人工占位符，预训练目标与推理自回归生成完全一致！

评估语言模型最核心的指标是\textbf{困惑度（Perplexity, PPL）}，它就是每词元平均负对数似然（交叉熵损失）的自然指数：
\begin{equation}
  \mathrm{PPL}(\theta) = \exp\!\left(\mathcal{L}_{\mathrm{CLM}}(\theta)\right)
  = \left(\prod_{t=1}^{T} p_\theta(x_t \mid x_{<t})\right)^{-1/T}.
  \label{eq:ppl}
\end{equation}
困惑度有一个极直观的物理含义——\textbf{等效均匀分支因子}：如果模型在每一步都像投掷一枚 $K$ 面均匀骰子（$p = 1/K$），那么 $\mathcal{L}_{\mathrm{CLM}} = \ln K$，从而 $\mathrm{PPL} = K$！因此 $\mathrm{PPL} = 8$ 意味着模型在预测下一个词时，平均只在 $8$ 个等可能候选中犹豫；而在信息论视角下，期望交叉熵恰好分解为：
\[
  \E_{x \sim p_{\text{data}}}\!\left[-\ln p_\theta(x)\right]
  = \underbrace{H(p_{\text{data}})}_{\text{自然语言固有熵（常数）}} + \underbrace{\DKL\!\left(p_{\text{data}} \,\|\, p_\theta\right)}_{\ge 0},
\]
因此最小化下一词预测损失严格等价于压缩真实人类语言分布与模型分布之间的 KL 散度！

\begin{center}
\footnotesize
\begin{tabularx}{\linewidth}{@{}p{24mm}p{26mm}p{24mm}X@{}}
\toprule
\textbf{预训练范式} & \textbf{代表架构与掩码} & \textbf{单次监督比例} & \textbf{核心优势与适用场景} \\
\midrule
掩码去噪（MLM） & BERT / RoBERTa（双向无掩码） & $15\%$（仅被遮盖位） & 双向上下文表征强，适合检索嵌入、分类与序列标注；无法直接自回归生成 \\
因果自回归（CLM） & GPT / Llama / Qwen（下三角因果） & $100\%$（全序列） & 训练与生成完全统一，极易规模化，具备上下文学习与通用推理涌现能力 \\
编解码去噪（Seq2Seq） & T5 / BART（编码双向 + 解码因果） & 损坏跨度（约 $15\%\sim 30\%$） & 兼顾双向上下文理解与自由生成，常见于翻译、摘要与多模态条件生成 \\
\bottomrule
\end{tabularx}
\end{center}

\section{对比学习、CLIP 与 InfoNCE 互信息下界定理}
\label{sec:contrastive-infonce}

语言天然由离散词元构成，可以轻松计算词表上的 $\softmax$；但图像、音频和传感器信号是高维连续像素——直接在像素空间做“下一像素预测”会把大量算力浪费在高频纹理噪声上。怎样在不做像素重建的情况下，迫使神经网络学出高层语义表征？

\textbf{对比学习（Contrastive Learning）}给出了一个极简而震撼的答案：\textbf{拉近语义一致的“正样本对”，推远无关的“负样本对”！}
\begin{itemize}
  \item \textbf{单模态自监督（SimCLR；Chen et al., 2020）}：把同一张图像经两次随机裁剪与色彩抖动得到的两个视图视为正样本对 $(x_i, x_i^+)$，把同批次中其余所有图像视为负样本；
  \item \textbf{跨模态图文对齐（CLIP；Radford et al., 2021）}：从互联网抓取 4 亿组（图像 $x_i^A$, 配对文本 $x_i^B$），用视觉编码器 $f_A$ 和文本编码器 $f_B$ 将它们分别映射到同一个单位超球面 $\mathbb{S}^{d-1}$ 上（即 $L_2$ 归一化 $\hat{z}_i^A = f_A(x_i^A)/\norm{f_A(x_i^A)}_2,\; \hat{z}_i^B = f_B(x_i^B)/\norm{f_B(x_i^B)}_2$），并在 $N \times N$ 余弦相似度矩阵 $S_{ij} = \inner{\hat{z}_i^A}{\hat{z}_j^B}$ 上执行对称对比分类！
\end{itemize}

\begin{defBox}
\begin{definition}[对称温度缩放 InfoNCE 对比损失]
给定一个包含 $N$ 个配对样本的批次 $\{(\hat{z}_i^A, \hat{z}_i^B)\}_{i=1}^N \subset \mathbb{S}^{d-1} \times \mathbb{S}^{d-1}$ 与温度超参数 $\tau > 0$，记余弦相似度矩阵 $S_{ij} = \inner{\hat{z}_i^A}{\hat{z}_j^B} \in [-1, 1]$。按行（以 $A$ 检索 $B$）与按列（以 $B$ 检索 $A$）的对称 InfoNCE 损失定义为：
\begin{equation}
  \mathcal{L}_{A \to B} = -\frac{1}{N}\sum_{i=1}^{N} \ln \frac{\exp(S_{ii}/\tau)}{\sum_{j=1}^{N} \exp(S_{ij}/\tau)},
  \qquad
  \mathcal{L}_{B \to A} = -\frac{1}{N}\sum_{j=1}^{N} \ln \frac{\exp(S_{jj}/\tau)}{\sum_{i=1}^{N} \exp(S_{ij}/\tau)},
  \label{eq:infonce-dir}
\end{equation}
\begin{equation}
  \mathcal{L}_{\mathrm{InfoNCE}} = \frac{1}{2}\left(\mathcal{L}_{A \to B} + \mathcal{L}_{B \to A}\right).
  \label{eq:infonce-sym}
\end{equation}
\end{definition}
\end{defBox}

从直观上看，式~\eqref{eq:infonce-dir} 不过是把 $N \times N$ 相似度矩阵 $S/\tau$ 的每一行和每一列都当成一个 $N$ 分类问题，正确标签恰好落在主对角线 $(i, i)$ 上！但它在信息论深处蕴藏着一条极其优美的核心定理（Oord et al., 2018; Poole et al., 2019）：\textbf{最小化 InfoNCE 损失，严格等价于最大化两个模态表征之间互信息 $I(X; Y)$ 的变分下界！} 下面我们给出一个完全无需大 $N$ 渐近近似、对任意有限批大小 $N \ge 2$ 皆精确成立的“天书级”证明。

\begin{thmBox}
\begin{theorem}[InfoNCE 互信息下界定理]
\label{thm:infonce-mi}
设随机变量 $(X, Y) \sim p(x, y)$ 具有联合分布与边缘分布 $p(x), p(y)$。考虑如下 $N$ 选 $1$ 对比采样过程：先均匀抽取正样本槽位 $K \sim \operatorname{Unif}\{1, \dots, N\}$，采样 $X \sim p(x)$，令 $Y_K \sim p(y \mid X)$，而其余 $N-1$ 个负样本 $Y_j \sim p(y)$（$j \ne K$）相互独立且独立于 $(X, Y_K)$。对任意可测打分函数 $f: \mathcal{X} \times \mathcal{Y} \to \R$，定义总体 InfoNCE 损失：
\begin{equation}
  \mathcal{L}_N^{\mathrm{InfoNCE}}(f)
  := \E_{K, X, Y_{1:N}}\!\left[-\ln \frac{\exp\big(f(X, Y_K)\big)}{\sum_{j=1}^{N} \exp\big(f(X, Y_j)\big)}\right].
  \label{eq:infonce-pop}
\end{equation}
则对任意打分函数 $f$ 与任意有限批大小 $N \ge 2$，恒有严格互信息下界：
\begin{equation}
  I(X; Y) \ge \ln N - \mathcal{L}_N^{\mathrm{InfoNCE}}(f).
  \label{eq:infonce-mi-bound}
\end{equation}
并且当且仅当打分函数满足 $e^{f(x, y)} \propto \frac{p(y \mid x)}{p(y)} = \frac{p(x, y)}{p(x)p(y)}$（即 $f(x, y) = \ln \frac{p(x, y)}{p(x)p(y)} + c(x)$，点互信息加任意仅依赖 $x$ 的常数）时，变分交叉熵间隙归零。
\end{theorem}
\begin{proof}
记 $q(k \mid x, y_{1:N}) := \frac{\exp(f(x, y_k))}{\sum_{j=1}^N \exp(f(x, y_j))}$ 为由打分函数诱导的关于槽位 $K$ 的变分后验，而 $p(k \mid x, y_{1:N})$ 为真实的贝叶斯后验。由式~\eqref{eq:infonce-pop}，期望交叉熵可精确拆分为条件熵加非负 KL 散度：
\begin{align}
  \mathcal{L}_N^{\mathrm{InfoNCE}}(f)
  &= \E_{X, Y_{1:N}}\!\left[\sum_{k=1}^{N} p(k \mid X, Y_{1:N})\big(-\ln q(k \mid X, Y_{1:N})\big)\right] \nonumber \\
  &= H(K \mid X, Y_{1:N}) + \E_{X, Y_{1:N}}\!\Big[\DKL\!\big(p(K \mid X, Y_{1:N}) \,\|\, q(K \mid X, Y_{1:N})\big)\Big] \nonumber \\
  &\ge H(K \mid X, Y_{1:N}), \label{eq:infonce-step1}
\end{align}
等号成立当且仅当 $q(k \mid x, y_{1:N}) = p(k \mid x, y_{1:N})$。利用贝叶斯公式算出真实后验：
\[
  p(k \mid x, y_{1:N})
  = \frac{\frac{1}{N} p(y_k \mid x) \prod_{j \ne k} p(y_j)}{\sum_{m=1}^{N} \frac{1}{N} p(y_m \mid x) \prod_{j \ne m} p(y_j)}
  = \frac{p(y_k \mid x)/p(y_k)}{\sum_{m=1}^{N} p(y_m \mid x)/p(y_m)},
\]
可见取 $e^{f(x, y)} \propto \frac{p(y\mid x)}{p(y)}$ 时变分后验与真实后验完全重合，KL 间隙归零！

接下来将 $H(K \mid X, Y_{1:N})$ 与 $I(X; Y)$ 精确相连。由互信息定义与 $K \sim \operatorname{Unif}\{1,\dots,N\}$：
\begin{equation}
  H(K \mid X, Y_{1:N}) = H(K) - I(K; X, Y_{1:N}) = \ln N - I(K; X, Y_{1:N}).
  \label{eq:infonce-step2}
\end{equation}
关键一步在于界定 $I(K; X, Y_{1:N})$。注意两项精确独立性：
\begin{itemize}
  \item 当 $X$ 被边缘积分掉后，给定 $K = k$ 时 $Y_k$ 的边缘分布仍是 $p(y)$，故整个元组 $Y_{1:N}$ 的边缘联合分布恒为 $\prod_{j=1}^N p(y_j)$，完全不依赖 $k$！因此 $K \perp\!\!\!\perp Y_{1:N}$，即 $I(K; Y_{1:N}) = 0$。
  \item 生成过程中 $K$ 与 $X$ 独立抽取，故 $I(X; K) = 0$。
\end{itemize}
现在用互信息链式法则以两种顺序展开 $I(X; K, Y_{1:N})$：
\begin{align*}
  \text{顺序一：}\quad I(X; K, Y_{1:N})
  &= \underbrace{I(X; K)}_{=\,0} + I(X; Y_{1:N} \mid K) \\
  &= \frac{1}{N}\sum_{k=1}^{N} I(X; Y_k, Y_{-k} \mid K = k) = I(X; Y),
\end{align*}
最后一步是因为给定 $K=k$ 时，负样本 $Y_{-k}$ 独立于 $(X, Y_k)$，不携带关于 $X$ 的任何额外信息。
\begin{align*}
  \text{顺序二：}\quad I(X; K, Y_{1:N})
  &= \underbrace{I(X; Y_{1:N})}_{\ge\, 0} + I(X; K \mid Y_{1:N})
   \ge I(X; K \mid Y_{1:N}).
\end{align*}
同时对 $I(K; X, Y_{1:N})$ 使用链式法则并代入 $I(K; Y_{1:N}) = 0$：
\[
  I(K; X, Y_{1:N}) = \underbrace{I(K; Y_{1:N})}_{=\,0} + I(K; X \mid Y_{1:N}) = I(X; K \mid Y_{1:N}) \le I(X; Y).
\]
将 $I(K; X, Y_{1:N}) \le I(X; Y)$ 代回式~\eqref{eq:infonce-step2} 与式~\eqref{eq:infonce-step1}，立得：
\[
  \mathcal{L}_N^{\mathrm{InfoNCE}}(f) \ge \ln N - I(X; Y)
  \;\iff\;
  I(X; Y) \ge \ln N - \mathcal{L}_N^{\mathrm{InfoNCE}}(f).
\]
\end{proof}
\end{thmBox}

定理~\ref{thm:infonce-mi} 立刻解开了一个困扰工程界多年的核心直觉：\textbf{为什么 CLIP 和 SimCLR 执着于把批大小 $N$ 推到 $4{,}096$ 甚至 $32{,}768$？} 因为式~\eqref{eq:infonce-mi-bound} 的右端被 $\ln N$ 死死封顶！即使损失降到 $0$，批大小 $N = 8$ 最多只能保证 $\ln 8 \approx 2.08$ nats（$3$ 比特）的互信息下界；只有把负样本数 $N$ 做大，$\ln N$ 上限才会随之抬高，从而把互信息下界收紧到真正高语义容量区间！

\subsection{超球面上的对齐与均匀性，以及温度 $\tau$ 的硬负样本聚焦}

把式~\eqref{eq:infonce-dir} 的单样本损失改写展开（利用 $\ln(\exp(a) + \sum \exp(b)) = \ln\sum \exp$）：
\begin{equation}
  \ell_i = \underbrace{-\frac{1}{\tau} \inner{\hat{z}_i^A}{\hat{z}_i^B}}_{\text{对齐项（Alignment）：拉近正对}}
  + \underbrace{\ln \sum_{j=1}^{N} \exp\!\left(\frac{\inner{\hat{z}_i^A}{\hat{z}_j^B}}{\tau}\right)}_{\text{均匀项（Uniformity）：超球面上排斥负对}}.
  \label{eq:align-uniform}
\end{equation}
Wang 与 Isola（2020）揭示了对比学习的优雅几何本质：在单位超球面 $\mathbb{S}^{d-1}$ 上，InfoNCE 同时优化两个相互制衡的几何目标——\textbf{对齐性（Alignment）}将语义相同的正对推向同一点（$\cos \to +1$），而\textbf{均匀性（Uniformity）}如电荷排斥般把所有样本近似均匀地铺满整个超球面（平均负对余弦 $\cos \approx 0$），从而最大化信息熵并彻底杜绝表征坍缩！

再对 $\ell_i$ 求关于锚点向量 $\hat{z}_i^A$ 的梯度（记 $p_{ij} = \frac{\exp(S_{ij}/\tau)}{\sum_m \exp(S_{im}/\tau)}$）：
\begin{equation}
  \frac{\partial \ell_i}{\partial \hat{z}_i^A}
  = \frac{1}{\tau}\left[-(1 - p_{ii}) \hat{z}_i^B + \sum_{j \ne i} p_{ij}\, \hat{z}_j^B\right].
  \label{eq:infonce-grad}
\end{equation}
式~\eqref{eq:infonce-grad} 一眼揭露了温度 $\tau$ 的物理作用：
\begin{itemize}
  \item \textbf{较小温度 $\tau$（如 $0.05 \sim 0.15$）}：$\softmax$ 权重 $p_{ij}$ 急剧向那些余弦相似度较高、最容易混淆的\textbf{硬负样本（Hard Negatives）}集中，产生精准的局部斥力；
  \item \textbf{过大温度 $\tau$（如 $\tau \gg 1$）}：$p_{ij} \approx 1/N$ 趋于平坦，模型对早已推远的简单负样本和难分伯仲的硬负样本“一视同仁”，梯度辨别力大幅稀释；
  \item \textbf{过小温度 $\tau$（如 $\tau \to 0.005$）}：梯度幅值乘上 $1/\tau$ 剧烈放大，且容易把批次中偶然语义相近的“假负样本（如另一张也拍着金毛犬的照片）”强行拆散，破坏全局语义拓扑！因此 CLIP 将 $\ln(1/\tau)$ 设为可学习参数，并严格截断上界 $1/\tau \le 100$（即 $\tau \ge 0.01$）。
\end{itemize}

图~\ref{fig:clip-matrix} 直观描绘了 CLIP 双编码器、$N\times N$ 余弦相似度矩阵与超球面 $\mathbb{S}^{d-1}$ 对齐--均匀几何。

\begin{figure}[htbp]
\centering
\resizebox{0.97\linewidth}{!}{%
\begin{tikzpicture}[
  enc/.style={rectangle, rounded corners=1.5mm, draw=#1, thick, fill=#1!10, minimum width=24mm, minimum height=9mm, align=center, font=\sffamily\footnotesize, text=ink},
  arr/.style={-{Stealth[length=2mm]}, thick, slate}
]
  % Left: Modality A & B encoders
  \node[font=\sffamily\bfseries\small, color=tealDark] at (1.5, 4.9) {(a) 双模态编码器与 $N\times N$ 相似度矩阵};

  \node[enc=teal] (ea) at (-0.6, 3.5) {模态 A 编码器\\$f_A \to L_2\text{ 归一}$};
  \node[enc=navy] (eb) at (-0.6, 1.3) {模态 B 编码器\\$f_B \to L_2\text{ 归一}$};

  \node[font=\scriptsize\sffamily, color=tealDark] at (-0.6, 4.3) {$\hat{z}_1^A, \dots, \hat{z}_N^A \in \mathbb{S}^{d-1}$};
  \node[font=\scriptsize\sffamily, color=navy] at (-0.6, 0.5) {$\hat{z}_1^B, \dots, \hat{z}_N^B \in \mathbb{S}^{d-1}$};

  % 4x4 similarity matrix
  \begin{scope}[shift={(1.7, 0.9)}]
    \foreach \i in {0,...,3} {
      \foreach \j in {0,...,3} {
        \pgfmathsetmacro{\x}{\j * 0.78}
        \pgfmathsetmacro{\y}{(3 - \i) * 0.78}
        \ifnum\i=\j
          \fill[teal!28] (\x, \y) rectangle (\x+0.74, \y+0.74);
          \draw[tealDark, thick] (\x, \y) rectangle (\x+0.74, \y+0.74);
          \node[font=\scriptsize\bfseries, color=tealDark] at (\x+0.37, \y+0.37) {$+$};
        \else
          \fill[ochreLight] (\x, \y) rectangle (\x+0.74, \y+0.74);
          \draw[ochre!45] (\x, \y) rectangle (\x+0.74, \y+0.74);
          \node[font=\scriptsize, color=ochre!80] at (\x+0.37, \y+0.37) {$-$};
        \fi
      }
    }
    \node[font=\scriptsize\sffamily, color=navy] at (1.56, 3.38) {$\hat{z}_1^B \;\;\; \hat{z}_2^B \;\;\; \hat{z}_3^B \;\;\; \hat{z}_4^B$};
    \node[font=\scriptsize\sffamily, color=tealDark, rotate=90] at (-0.30, 1.56) {$\hat{z}_1^A \;\; \cdots \;\; \hat{z}_4^A$};
    \node[font=\scriptsize\sffamily, color=ink, align=center] at (1.56, -0.45) {主对角线拉近（$\cos \to +1$）\\非对角线排斥（$\cos \to 0$）};
  \end{scope}

  \draw[arr, teal] (ea.east) -- (1.35, 3.5);
  \draw[arr, navy] (eb.east) -- (1.35, 1.3);

  % Right: Hypersphere S^{d-1} Alignment & Uniformity
  \begin{scope}[shift={(8.4, 2.45)}]
    \node[font=\sffamily\bfseries\small, color=tealDark] at (0, 2.45) {(b) 超球面 $\mathbb{S}^{d-1}$ 上的对齐与均匀};
    \draw[tealLine, very thick, fill=tealLight!35] (0,0) circle (1.75cm);
    \draw[tealLine, dashed] (-1.75,0) arc (180:360:1.75 and 0.45);
    \draw[tealLine, dotted] (1.75,0) arc (0:180:1.75 and 0.45);

    % 3 positive pairs well separated on sphere
    \coordinate (a1) at (38:1.75);
    \coordinate (b1) at (48:1.75);
    \coordinate (a2) at (158:1.75);
    \coordinate (b2) at (168:1.75);
    \coordinate (a3) at (276:1.75);
    \coordinate (b3) at (286:1.75);

    \foreach \p in {a1, a2, a3} {\fill[teal] (\p) circle (2.6pt);}
    \foreach \p in {b1, b2, b3} {\fill[navy] (\p) circle (2.6pt);}

    % Alignment arrows (pull together)
    \draw[{Stealth[length=1.6mm]}-{Stealth[length=1.6mm]}, tealDark, thick] (a1) -- (b1);
    \draw[{Stealth[length=1.6mm]}-{Stealth[length=1.6mm]}, tealDark, thick] (a2) -- (b2);
    \draw[{Stealth[length=1.6mm]}-{Stealth[length=1.6mm]}, tealDark, thick] (a3) -- (b3);

    % Uniformity repulsion dashed arcs inside
    \draw[{Stealth[length=1.6mm]}-{Stealth[length=1.6mm]}, ochre, dashed, thick] (43:1.25) to[bend right=15] (163:1.25);
    \draw[{Stealth[length=1.6mm]}-{Stealth[length=1.6mm]}, ochre, dashed, thick] (163:1.25) to[bend right=15] (281:1.25);
    \draw[{Stealth[length=1.6mm]}-{Stealth[length=1.6mm]}, ochre, dashed, thick] (281:1.25) to[bend right=15] (43:1.25);

    \node[font=\scriptsize\sffamily, color=tealDark, anchor=west] at (43:1.9) {对齐正对};
    \node[font=\scriptsize\sffamily, color=ochre] at (0, -0.05) {均匀斥力};
    \node[font=\scriptsize\sffamily, color=slate] at (0, -2.15) {$\bullet$ 模态 A（青） \quad $\bullet$ 模态 B（蓝）};
  \end{scope}
\end{tikzpicture}}
\caption{CLIP 风格跨模态对比学习几何机制。(a) 两个独立编码器将异构模态投射到共享超球面并构造 $N \times N$ 相似度矩阵，通过对称交叉熵拉升主对角线；(b) 超球面上正对对齐靠拢、不同语义概念均匀铺开。}
\label{fig:clip-matrix}
\end{figure}

\section{实测实验：双模态 CLIP 风格 InfoNCE 对比对齐}
\label{sec:exp-infonce}

为了亲手验证式~\eqref{eq:infonce-sym} 与定理~\ref{thm:infonce-mi}，我们运行 \texttt{code/vol5\_experiments.py} 中的 \texttt{exp5\_infonce\_contrastive}：合成 $N = 128$ 个共享 $8$ 维潜在概念 $Z \in \R^{128 \times 8}$，分别经由\textbf{两种截然不同的非线性观测模态}生成输入数据（维度 $D_{\text{in}} = 16$）：
\[
  X_A = \tanh(Z W_A) + 0.08\,\varepsilon_A, \qquad
  X_B = \sin(Z W_B) + 0.08\,\varepsilon_B.
\]
原始特征空间 $X_A$ 与 $X_B$ 的坐标毫无对齐关系（初始平均余弦相似度为负！）。我们用两个独立的 2 层 ReLU MLP 编码器 \texttt{enc\_a}, \texttt{enc\_b}（$16 \to 32 \to 16$）将它们映射到共享单位超球面 $\mathbb{S}^{15}$，取温度 $\tau = 0.15$，以 AdamW（$\eta = 5\times 10^{-3}$）优化对称 InfoNCE 损失 200 轮：

\begin{lstlisting}[caption={纯 PyTorch 对称 CLIP 风格 InfoNCE 核心训练步（与 \texttt{vol5\_experiments.py} 一致）},label={lst:infonce}]
tau = 0.15
labels = torch.arange(N)                         # 正确配对全在主对角线 0..N-1
for ep in range(201):
    za = F.normalize(enc_a(Xa), dim=-1)          # 投影至单位超球面 S^{d-1}
    zb = F.normalize(enc_b(Xb), dim=-1)
    sim = za @ zb.T                              # (N, N) 余弦相似度矩阵
    logits = sim / tau
    loss = 0.5 * (F.cross_entropy(logits, labels)
                + F.cross_entropy(logits.T, labels))
    opt.zero_grad(); loss.backward(); opt.step()
\end{lstlisting}

实测记录文件 \texttt{data/infonce\_train.csv} 清晰刻画了超球面几何的自发成型（绘制于图~\ref{fig:infonce-train}）：
\begin{itemize}
  \item \textbf{第 0 轮（未对齐随机投影）}：InfoNCE 损失高达 \texttt{5.5298}（高于均匀随机基线 $\ln 128 \approx 4.8520$），跨模态 Top-1 检索准确率为 \texttt{0.00\%}，正对平均余弦 \texttt{pos\_cos = -0.204} 与负对平均余弦 \texttt{neg\_cos = -0.190} 几乎完全混在一起；
  \item \textbf{第 50 轮}：损失骤降至 \texttt{1.3946}，Top-1 跨模态检索准确率飙升至 \texttt{90.62\%}，正对余弦拉升至 \texttt{+0.840}，而负对余弦被精准压制在正交零点附近（\texttt{-0.006}）；
  \item \textbf{第 90--100 轮}：损失进一步跌至 \texttt{0.8177} 与 \texttt{0.7531}，Top-1 检索准确率触顶 \textbf{\texttt{100.00\%}}（正对余弦 \texttt{0.931}）；
  \item \textbf{第 200 轮}：损失降至 \texttt{0.5099}，Top-1 保持 \textbf{\texttt{100.00\%}}，\texttt{pos\_cos = +0.967}，\texttt{neg\_cos = -0.008}！代入定理~\ref{thm:infonce-mi}，此时学得的双模态互信息下界达到 $\ln 128 - 0.5099 = 4.342\text{ nats}$（约 $6.26$ 比特，已逼近 $128$ 样本所能承载的理论信息上限 $\log_2 128 = 7$ 比特）！
\end{itemize}

\begin{figure}[htbp]
\centering
\begin{tikzpicture}
  % Left plot: Loss and Top-1 Retrieval Acc (two axes cleanly or separate panels -> separate panels!)
  \begin{axis}[
    name=leftp,
    width=0.47\linewidth, height=5.8cm,
    title={\sffamily\bfseries\footnotesize (a) InfoNCE 损失与跨模态 Top-1 检索率},
    xlabel={训练轮数（Epoch）},
    ylabel={Top-1 检索准确率（\%）},
    xmin=-4, xmax=204, ymin=-4, ymax=108,
    xtick={0,50,100,150,200},
    ytick={0,25,50,75,100},
    grid=both, grid style={tealLine!45, dashed},
    axis line style={slate},
    tick label style={font=\scriptsize},
    label style={font=\footnotesize\sffamily, color=ink},
    legend style={at={(0.97,0.22)}, anchor=east, font=\scriptsize, fill=white, draw=tealLine, legend cell align=left}
  ]
    \addplot[teal, very thick, mark=*, mark size=1.5pt]
      table[x=epoch, y=top1_acc, col sep=comma]{data/infonce_train.csv};
    \addlegendentry{Top-1 检索率（\%）}

    % Overlay scaled loss? Better: show loss right-axis OR annotate clearly! Let's use clean right y-axis on panel (a).
    \node[font=\scriptsize\sffamily, color=tealDark, fill=white, inner sep=1pt]
      at (axis cs:142, 86) {ep $\ge 90$: 100.00\%};
    \node[font=\scriptsize\sffamily, color=ochre, anchor=west]
      at (axis cs:18, 12) {Loss: 5.53 $\to$ 0.51};
  \end{axis}

  % Right plot: Alignment (pos_cos) vs Uniformity (neg_cos)
  \begin{axis}[
    at={(leftp.east)}, anchor=west, xshift=1.45cm,
    width=0.47\linewidth, height=5.8cm,
    title={\sffamily\bfseries\footnotesize (b) 超球面余弦：正对对齐 vs 负对均匀},
    xlabel={训练轮数（Epoch）},
    ylabel={平均余弦相似度 $\cos(\hat{z}^A, \hat{z}^B)$},
    xmin=-4, xmax=204, ymin=-0.32, ymax=1.08,
    xtick={0,50,100,150,200},
    ytick={-0.2,0,0.2,0.4,0.6,0.8,1.0},
    grid=both, grid style={tealLine!45, dashed},
    axis line style={slate},
    tick label style={font=\scriptsize},
    label style={font=\footnotesize\sffamily, color=ink},
    legend style={at={(0.97,0.48)}, anchor=east, font=\scriptsize, fill=white, draw=tealLine, legend cell align=left}
  ]
    \addplot[navy, very thick, mark=*, mark size=1.5pt]
      table[x=epoch, y=pos_cos, col sep=comma]{data/infonce_train.csv};
    \addlegendentry{正对 \texttt{pos\_cos}（对齐）}

    \addplot[ochre, very thick, dashed, mark=square*, mark size=1.4pt]
      table[x=epoch, y=neg_cos, col sep=comma]{data/infonce_train.csv};
    \addlegendentry{负对 \texttt{neg\_cos}（均匀）}

    \draw[slate, dotted, thick] (axis cs:0, 0) -- (axis cs:200, 0);
    \node[font=\scriptsize\sffamily, color=navy, anchor=east] at (axis cs:195, 0.84) {ep=200: +0.967};
    \node[font=\scriptsize\sffamily, color=ochre, anchor=east] at (axis cs:195, -0.14) {ep=200: -0.008};
  \end{axis}
\end{tikzpicture}
\caption{纯 PyTorch 双模态对称 InfoNCE 对比对齐实测结果（读取自 \texttt{data/infonce\_train.csv}，$N=128,\; \tau=0.15$）。两个非线性失配模态在 100 轮内实现 \texttt{100.00\%} 跨模态检索，正对余弦升至 \texttt{+0.967}，负对余弦收敛于超球面正交基线 \texttt{-0.008}。}
\label{fig:infonce-train}
\end{figure}

\section{神经缩放定律与 Chinchilla 算力最优分配}
\label{sec:scaling-laws}

当架构固定为 Pre-LN Transformer、目标固定为自回归下一词预测后，深度学习界迎来了一个堪比统计物理热力学定律的惊人经验发现——\textbf{神经缩放定律（Neural Scaling Laws）}（Kaplan et al., 2020; Hoffmann et al., 2022）：大语言模型的测试交叉熵损失并不随架构微调剧烈抖动，而是关于模型非嵌入参数量 $N$、训练词元总数 $D$ 与总训练算力 $C$ 的极其平滑的\textbf{幂律函数}！

\subsection{算力预算 $C \approx 6ND$ 的第一性原理核算}

训练一个拥有 $N$ 个参数、吃进 $D$ 个词元的稠密 Transformer 需要多少浮点运算（FLOPs）？我们用第 1 章的矩阵乘法规则直接心算：
\begin{itemize}
  \item \textbf{前向传播}：每个参数权重在每个词元上恰好参与一次乘法与一次加法（即 $1\text{ MAC} = 2\text{ FLOPs}$），故单个词元的前向主计算量约为 $2N$ FLOPs（当序列长 $T$ 不算极端庞大时，线性投影主导于注意力内积）；
  \item \textbf{反向传播}：回顾第 1 章线性层 VJP，每个矩阵层需算两次同等规模矩阵乘——一次传给激活输入 $\bar{X} = \bar{Y} W^\top$（$2N$ FLOPs），一次累积权重梯度 $\bar{W} = X^\top \bar{Y}$（$2N$ FLOPs），合计 $4N$ FLOPs；
  \item \textbf{训练总算力}：前向加反向每词元耗费 $2N + 4N = 6N$ FLOPs，乘以总词元数 $D$，即得著名近似公式：
  \begin{equation}
    C \approx 6 N D.
    \label{eq:compute-6nd}
  \end{equation}
\end{itemize}

\subsection{Chinchilla 算力最优分配定理（拉格朗日乘数法严格推导）}

DeepMind 的 Hoffmann 等人（2022）拟合了跨五个算力数量级的 400 余个语言模型实验，提出如下参数化损失模型：
\begin{equation}
  L(N, D) = E + \frac{A}{N^\alpha} + \frac{B}{D^\beta},
  \label{eq:chinchilla-loss}
\end{equation}
其中 $E > 0$ 代表自然语言的不可约熵下限（即使无限大模型与无限数据也无法低于 $E$），第二项 $A N^{-\alpha}$ 刻画模型容量不足的逼近误差，第三项 $B D^{-\beta}$ 刻画有限训练数据的统计估计误差。

现在假设你手握一笔固定且昂贵的 GPU 集群总算力预算 $C$（例如耗资数百万美元的 $C = 5 \times 10^{23}$ FLOPs）：\textbf{究竟应该训一个“巨无霸参数但数据少”的模型，还是“参数适中但充分喂饱海量词元”的模型？} 这正是第 4 卷第 2 章拉格朗日乘数法的完美舞台！

\begin{thmBox}
\begin{theorem}[Chinchilla 算力最优缩放定理]
\label{thm:chinchilla}
在固定总训练算力约束 $C = 6 N D$ 下，最小化参数化损失函数 $L(N, D) = E + A N^{-\alpha} + B D^{-\beta}$（$A, B, \alpha, \beta > 0$），其唯一全局最优的模型参数量 $N^*(C)$ 与训练词元数 $D^*(C)$ 严格服从如下幂律：
\begin{equation}
  N^*(C) = G \left(\frac{C}{6}\right)^a,
  \qquad
  D^*(C) = G^{-1} \left(\frac{C}{6}\right)^b,
  \label{eq:chinchilla-opt}
\end{equation}
其中最优缩放指数与比例常数精确为：
\begin{equation}
  a = \frac{\beta}{\alpha + \beta}, \qquad
  b = \frac{\alpha}{\alpha + \beta}, \qquad
  G = \left(\frac{\alpha A}{\beta B}\right)^{\frac{1}{\alpha + \beta}}.
  \label{eq:chinchilla-exp}
\end{equation}
\end{theorem}
\begin{proof}
构造带等式约束的拉格朗日函数（为简化代数，记 $M := C / 6$，约束为 $N D = M$）：
\[
  \Lambda(N, D, \lambda) = E + A N^{-\alpha} + B D^{-\beta} + \lambda (N D - M).
\]
对 $N$ 与 $D$ 分别求一阶驻点条件 $\frac{\partial \Lambda}{\partial N} = 0$ 与 $\frac{\partial \Lambda}{\partial D} = 0$：
\begin{align*}
  -\alpha A N^{-\alpha - 1} + \lambda D &= 0 \implies \alpha A N^{-\alpha} = \lambda N D = \lambda M, \\
  -\beta B D^{-\beta - 1} + \lambda N &= 0 \implies \beta B D^{-\beta} = \lambda N D = \lambda M.
\end{align*}
两式右端相等，立得容量误差与数据误差的\textbf{最优均衡方程}：
\begin{equation}
  \alpha A N^{-\alpha} = \beta B D^{-\beta}
  \implies D^\beta = \frac{\beta B}{\alpha A} N^\alpha.
  \label{eq:chinchilla-balance}
\end{equation}
将算力约束 $D = M / N$ 代入式~\eqref{eq:chinchilla-balance}：
\[
  \left(\frac{M}{N}\right)^\beta = \frac{\beta B}{\alpha A} N^\alpha
  \implies N^{\alpha + \beta} = \frac{\alpha A}{\beta B}\, M^\beta.
\]
两端同时开 $(\alpha + \beta)$ 次方，立即解出：
\[
  N^*(C) = \left(\frac{\alpha A}{\beta B}\right)^{\frac{1}{\alpha + \beta}} M^{\frac{\beta}{\alpha + \beta}}
  = G \left(\frac{C}{6}\right)^a, \qquad a = \frac{\beta}{\alpha + \beta}.
\]
再代回 $D^*(C) = M / N^*(C)$，注意 $1 - a = \frac{\alpha}{\alpha + \beta} = b$，得：
\[
  D^*(C) = G^{-1} \left(\frac{C}{6}\right)^b, \qquad b = \frac{\alpha}{\alpha + \beta}.
\]
最后验证二阶充分性：消去 $D = M/N$ 后，单变量目标 $\phi(N) = E + A N^{-\alpha} + B M^{-\beta} N^{\beta}$ 的二阶导数为
\[
  \phi''(N) = \alpha(\alpha + 1) A N^{-\alpha - 2} + \beta(\beta - 1) B M^{-\beta} N^{\beta - 2}.
\]
在驻点处由式~\eqref{eq:chinchilla-balance} 有 $B M^{-\beta} N^{\beta} = \frac{\alpha}{\beta} A N^{-\alpha}$，代入得
\[
  \phi''(N^*) = \alpha A (N^*)^{-\alpha - 2}\Big[(\alpha + 1) + (\beta - 1)\Big]
  = \alpha(\alpha + \beta) A (N^*)^{-\alpha - 2} > 0,
\]
故 $(N^*, D^*)$ 为严格全局极小值点！
\end{proof}
\end{thmBox}

式~\eqref{eq:chinchilla-exp} 具有极深的直觉美感：指数 $a$ 和 $b$ 恰好与 $\alpha, \beta$ \textbf{交叉对应}，且 $a + b = 1$！如果损失对参数更敏感（$\alpha > \beta$），最优策略反而应把更多比例的增量算力倾注在相对“落后”的数据维度 $D$ 上（$b > a$），使两项边际收益始终保持式~\eqref{eq:chinchilla-balance} 的动态平衡。

代入 Hoffmann 等人（2022）的实测拟合值：
\[
  E \approx 1.69, \quad A \approx 406.4, \quad B \approx 410.7, \quad \alpha \approx 0.34, \quad \beta \approx 0.28,
\]
立得：
\[
  a = \frac{0.28}{0.34 + 0.28} \approx 0.452, \qquad
  b = \frac{0.34}{0.34 + 0.28} \approx 0.548.
\]
这彻底推翻了早期 Kaplan 等人（2020）“参数应远快于数据扩张（$a \approx 0.73, b \approx 0.27$）”的旧认知，揭示出\textbf{参数规模与训练词元数应当近乎等比例（$\approx 1:1$）随 $\sqrt{C}$ 同步扩张}！在百亿至千亿参数量级上，最优词元与参数比恰好落在著名的 \textbf{Chinchilla 黄金法则}附近：
\[
  D^* \approx 20\, N^*.
\]

\begin{example}[Gopher 280B vs Chinchilla 70B：同等算力的戏剧性逆转]
在总算力预算均为 $C \approx 5.76 \times 10^{23}$ FLOPs 下：
\begin{itemize}
  \item \textbf{旧路线 Gopher}：堆出了 $N = 280\text{B}$ 参数，但受算力所限只喂了 $D = 300\text{B}$ 词元（$D/N \approx 1.07$，严重“饥饿欠训”）；
  \item \textbf{Chinchilla 最优配比}：将参数缩小 4 倍至 $N = 70\text{B}$，把省下的算力全部用来喂饱 $D = 1.4\text{T}$ 词元（$D/N = 20$）。结果 $70\text{B}$ 的 Chinchilla 不仅在几乎全部基准上碾压 $4$ 倍大的 Gopher，而且部署推理成本也骤降为四分之一！
\end{itemize}
\textbf{当代延伸（推理最优的“过训练”范式）}：Chinchilla 定律优化的是\emph{纯训练算力}。但在实际工业部署中，一个开源模型（如 Llama-3 8B）会被全球用户推理数十亿次！把推理成本摊入生命周期总成本后，最划算的策略是有意\textbf{超越 Chinchilla 比例做“过训练（Over-training）”}——例如 Llama-3 8B 喂入了高达 $15\text{T}$ 词元（$D/N \approx 1875$），用一次性的超额预训练算力换取了一个极小、极快、极强的推理利器。
\end{example}

\begin{failBox}{自监督预训练与对比学习三大高频踩坑}
\begin{enumerate}
  \item \textbf{对比学习没有负样本导致“表征坍缩（Representation Collapse）”}：若只保留式~\eqref{eq:align-uniform} 中的正对对齐项 $-\inner{\hat{z}_i^A}{\hat{z}_i^B}$ 而删去负样本分母，优化器会在三步之内找到一个完美零损失的作弊解——把所有输入统统映射为同一个常向量 $\hat{z} \equiv c$！所有样本余弦相似度全为 $+1$，表征彻底报废。
  \item \textbf{温度 $\tau$ 设定极端导致训练停滞或发散}：在单位超球面上 $\cos \in [-1, 1]$，若忘记除以温度（即默认 $\tau = 1$），$\exp(\pm 1)$ 仅相差 $7.4$ 倍，哪怕有 128 个负样本，损失也很难压低、梯度软弱无力；反之若把 $\tau$ 设得过小（如 $0.002$），logits 高达 $500$，瞬间触发 $\exp$ 上溢 \texttt{NaN} 或梯度爆炸。经验甜蜜区通常在 $\tau \in [0.05, 0.2]$。
  \item \textbf{盲目追求大参数量而忽视数据配比（算力次优陷阱）}：在个人或实验室有限算力下，强行训练一个数十亿参数却只看了一点数据的“欠训巨兽”，不仅验证困惑度远逊于同算力充分训练的小模型，还会消耗成倍的显存与推理延迟。先按式~\eqref{eq:chinchilla-opt} 核算 $C \approx 6ND$ 与数据预算，是当代训练大模型的第一铁律。
\end{enumerate}
\end{failBox}

\begin{extBox}{视觉自监督（MAE / DINO）与大模型“涌现能力”之辩}
\begin{itemize}
  \item \textbf{掩码自编码器（MAE；He et al., 2022）}：为何把 BERT 的 $15\%$ 掩码直接照搬到图像上效果平平？因为图像像素具有极高的空间冗余（遮掉 $15\%$ 靠周边插值就能猜出，无需理解物体）。MAE 将掩码比例大幅拉高至惊人的 \textbf{$75\%$}，且仅将可见的 $25\%$ 图像块送入重型 ViT 编码器（提速 $3\sim 4\times$），再配一个极轻量解码器重建被遮盖像素，彻底激发了视觉高层语义！
  \item \textbf{无需负样本也能防坍缩的 DINO（Caron et al., 2021）}：通过\textbf{动量教师（EMA Teacher）自蒸馏}与两大对称破缺技巧——\textbf{中心化（Centering，防模型坍缩到单一类别）}与\textbf{锐化（Sharpening，防输出退化为均匀分布）}，在完全不使用负样本的情况下让 ViT 自发学会了像素级前景物体分割！
  \item \textbf{“涌现能力（Emergent Abilities）”是物理相变还是度量幻觉？}：随着模型参数与算力跨越临界点，算术推理与代码生成准确率常呈现“先躺平为零、后突然飙升”的陡峭曲线。Schaeffer 等人（2023）提出尖锐反驳：底层交叉熵损失其实一直沿着式~\eqref{eq:chinchilla-loss} 平滑幂律下降！所谓的“陡崖涌现”，很大程度上源于\textbf{整句精确匹配（Exact Match）等非线性指标的指数放大效应}——回忆第~\ref{ch:attention-transformer} 章图~\ref{fig:minigpt-train} 中的 6 词序列反转实验：当逐词准确率 $p$ 从 $0.4$ 升至 $0.99$ 时，整句 6 词全对概率 $p^6$ 会从 $0.004$ 突然跳变到 $0.94$！平滑的逐词对数似然改进，经过连乘度量后就会折射出惊艳的宏观相变。
\end{itemize}
\end{extBox}

\section*{本章小结与核心脉络}
\addcontentsline{toc}{section}{本章小结与核心脉络}

本章打通了现代基础模型从“原始无标注数据”成长为“通用智能底座”的完整数学链条：
\begin{itemize}
  \item \textbf{BPE 子词词元化}以贪心高频合并平衡了词表规模与序列长度，同时彻底消灭了 OOV；
  \item \textbf{MLM 与 CLM} 揭示了自监督学习的两种核心构造，而因果下一词预测凭借 $100\%$ 词元监督效率与精确概率链式分解，奠定了现代大语言模型的主流地位；
  \item \textbf{对比学习 InfoNCE} 在信息论上严格给出互信息下界 $I(X; Y) \ge \ln N - \mathcal{L}_N^{\mathrm{InfoNCE}}$，在几何上统一了超球面 $\mathbb{S}^{d-1}$ 的正对对齐与负对均匀排斥，双模态实测以 10 行核心代码取得 $100\%$ 跨模态检索；
  \item \textbf{Chinchilla 缩放定律}用拉格朗日乘数法证明了最优算力分配幂律 $a = \frac{\beta}{\alpha+\beta},\; b = \frac{\alpha}{\alpha+\beta}$，将大模型研发从“盲目炼丹”带入了可预测的定量工程科学。
\end{itemize}
至此，我们已经能用深度网络做判别、表征对齐与离散文本自回归生成。但对于连续世界里绚丽多彩的图像、音频与分子三维构象，该如何从纯噪声中凭空“画”出栩栩如生的样本？下一章（本卷压轴章），我们将从第 4 卷 EM 算法的 ELBO 启程，一举征服 \textbf{VAE、扩散模型（DDPM）与流匹配（Flow Matching）}！

\section*{本章分层习题}
\addcontentsline{toc}{section}{本章分层习题}

\exTier{A 层：概念与手算热身}{熟练掌握 BPE 合并追踪、困惑度计算与算力配比估算}

\begin{enumerate}
  \item[\textbf{5.5.1}] \textbf{BPE 分词手算与新词切分}：给定微型语料及词频：\texttt{hug}（10）、\texttt{pug}（5）、\texttt{pun}（12）、\texttt{bun}（4）、\texttt{hugs}（5），均在词尾加 \texttt{</w>}。
  \begin{enumerate}
    \item 列出初始字符词表，并手工追踪前 4 轮 BPE 合并（给出每轮最高频符号对、频次与新符号）；
    \item 使用这 4 条合并规则，切分未登录测试词 \texttt{bugs</w>} 与 \texttt{puns</w>}。
  \end{enumerate}
  \exHint{第 1 轮统计 \texttt{(u, g)} 出现 $10+5+5 = 20$ 次，\texttt{(u, n)} 出现 $12+4 = 16$ 次，故先合并 \texttt{ug}！}

  \item[\textbf{5.5.2}] \textbf{困惑度、交叉熵与 Chinchilla 算力预算换算}：
  \begin{enumerate}
    \item 某中文大模型词表 $|V| = 100{,}000$。若模型在某测试集上每词元平均交叉熵损失为 $\mathcal{L}_{\mathrm{CE}} = 2.3026$ nats，其困惑度 $\mathrm{PPL}$ 是多少？若采用每词元平均信息量（单位：比特/词元，即 bits-per-token），该模型相比完全盲猜压缩了多少倍？
    \item 假设你拥有 $64$ 张峰值 BF16 算力为 $1000\text{ TFLOPS}$、实际算力利用率（MFU）为 $40\%$ 的 GPU，连续训练 $14$ 天。计算总有效算力 $C$（单位 FLOPs），并按 Chinchilla 法则 $C \approx 6ND$ 与 $D^* \approx 20 N^*$，算出你应当训练的最优模型参数量 $N^*$ 与词元数 $D^*$。
  \end{enumerate}
  \exHint{14 天 $= 14 \times 86{,}400 = 1{,}209{,}600$ 秒；由 $C = 6 N (20 N) = 120 N^2$ 立得 $N^* = \sqrt{C / 120}$。}
\end{enumerate}

\exTier{B 层：定理推导与严谨证明}{严格推导 InfoNCE 梯度动力学与缩放定律约束极值}

\begin{enumerate}
  \item[\textbf{5.5.3}] \textbf{InfoNCE 的硬负样本聚焦定理}：给定单位球面上的锚点 $\hat{z}_i^A$ 与负样本集合 $\{\hat{z}_j^B\}_{j\ne i}$，记余弦相似度 $s_j = \inner{\hat{z}_i^A}{\hat{z}_j^B}$。根据式~\eqref{eq:infonce-grad}，负样本 $j$ 贡献的相对梯度权重为 $w_j(\tau) = \frac{p_{ij}}{\sum_{m \ne i} p_{im}} = \frac{\exp(s_j / \tau)}{\sum_{m \ne i} \exp(s_m / \tau)}$。
  \begin{enumerate}
    \item 证明当 $\tau \to 0^+$ 时，$w_j(\tau)$ 趋近于仅在最难负样本 $j^* = \argmax_{j \ne i} s_j$ 上取 $1$ 的单热分布（假设最大值唯一），此时 InfoNCE 退化为经典的\textbf{三元组硬挖掘损失（Hard-Negative Triplet Loss）}；
    \item 证明当 $\tau \to +\infty$ 时，$w_j(\tau) \to \frac{1}{N-1}$，且一阶泰勒展开显示所有负样本按相等权重贡献排斥力。
  \end{enumerate}
  \exHint{分子分母同除以 $\exp(s_{j^*} / \tau)$ 即可立得 $\tau \to 0^+$ 极限。}

  \item[\textbf{5.5.4}] \textbf{含推理摊销成本的广义缩放最优解}：真实工业大模型不仅要付训练算力 $C_{\text{train}} = 6 N D$，上线后还要服务总计 $D_{\text{inf}}$ 个推理词元（每个生成词元的前向主计算量约为 $2 N$ FLOPs），故全生命周期总算力为：
  \[
    C_{\text{total}} = 6 N D + 2 N D_{\text{inf}}.
  \]
  在固定总预算 $C_{\text{total}}$ 下最小化式~\eqref{eq:chinchilla-loss} 的测试损失 $L(N, D)$。用拉格朗日乘数法推导新的均衡方程，并证明：只要 $D_{\text{inf}} > 0$，最优词元参数比 $D^* / N^*$ 严格\textbf{大于}纯训练最优的 Chinchilla 比值（解释为何 Llama-3 8B 要在 15T 词元上“过训练”！）。
  \exHint{新约束对 $N$ 求导多出 $2\lambda D_{\text{inf}}$，均衡条件变为 $\alpha A N^{-\alpha} = \beta B D^{-\beta}\!\left(1 + \frac{D_{\text{inf}}}{3D}\right) > \beta B D^{-\beta}$。}
\end{enumerate}

\exTier{C 层：代码实现与数值实验}{亲手编写 BPE 分词器与温度/批大小消融实验}

\begin{enumerate}
  \item[\textbf{5.5.5}] \textbf{30 行纯 Python 从零实现微型 BPE 分词器}：不依赖任何第三方库，编写 \texttt{train\_bpe(corpus, num\_merges)} 与 \texttt{encode(word, merges)} 两个函数，在例~\ref{ex:bpe-trace} 的语料与《天书中的智能》一段 500 词英文技术段落上运行 50 轮合并。打印前 15 条学到的合并规则以及压缩前后的总词元数对比。
  \exHint{用 \texttt{collections.Counter} 统计相邻二元组频次，合并时线性扫描替换即可。}

  \item[\textbf{5.5.6}] \textbf{InfoNCE 温度 $\tau$ 与批大小 $N$ 消融实验}：基于代码~\ref{lst:infonce} 与实验脚本 \texttt{vol5\_experiments.py} 中的对比函数 \texttt{exp5\_\allowbreak infonce\_\allowbreak contrastive}：
  \begin{enumerate}
    \item 固定 $N = 128$，对比 $\tau \in \{0.02,\; 0.15,\; 1.0,\; 5.0\}$ 四组温度在 200 轮后的 Top-1 检索率、\texttt{pos\_cos} 与 \texttt{neg\_cos}；
    \item 固定 $\tau = 0.15$，分别用小批对比（每次从 128 样本中随机抽子批 $B \in \{8, 32, 128\}$）训练 200 轮，在全部 128 样本上评测全局 Top-1 检索率，定量验证定理~\ref{thm:infonce-mi} 中批大小对互信息下界的收紧作用！
  \end{enumerate}
  \exHint{注意在 $\tau = 1.0$ 和 $5.0$ 时观察 \texttt{neg\_cos} 与收敛速度是否显著恶化。}
\end{enumerate}

\exTier{D 层：开放探究与前沿迁移}{探索 SigLIP、数据质量缩放律与合成数据坍缩}

\begin{enumerate}
  \item[\textbf{5.5.7}] \textbf{从 Softmax InfoNCE 到 Sigmoid 逐对对比（SigLIP）}：标准 CLIP 的 $\softmax$ 分母需要跨所有 GPU 执行全局 \texttt{all-gather} 同步才能拼出完整的 $N \times N$ 矩阵，成为超大规模分布式训练的通信瓶颈。Zhai 等人（2023）提出的 \textbf{SigLIP} 彻底放弃了 $\softmax$，把 $N \times N$ 矩阵的每一格当作独立的二元逻辑回归（正对标签 $+1$，负对标签 $-1$，配可学习偏置 $b$ 与温度 $t$）：
  \[
    \mathcal{L}_{\mathrm{SigLIP}} = -\frac{1}{N}\sum_{i=1}^{N}\sum_{j=1}^{N} \ln \frac{1}{1 + \exp\!\big(-y_{ij}(t\, S_{ij} + b)\big)}.
  \]
  分析 SigLIP 为何在分布式显存和通信上极具优势，且为何必须引入负初始化偏置 $b \approx -10$ 才能防止初期压倒性多数的负样本淹没正样本梯度？
  \exHint{当 $N = 32768$ 时正负样本比例高达 $1 : 32767$，初始 $b \approx -\ln N$ 恰好把先验赔率拉平！}

  \item[\textbf{5.5.8}] \textbf{高质量数据饱和与“模型坍缩（Model Collapse）”}：随着全网高质量公开人类文本接近被消费殆尽（约 $10^{13}\sim 10^{14}$ 词元），Chinchilla 公式~\eqref{eq:chinchilla-loss} 隐含的“每词元独立同分布且永不重复”假设面临挑战。查阅 Muennighoff 等人（2023）的\textbf{数据受限缩放定律（Data-Constrained Scaling Laws）}（重复训练至多 4 轮几乎无损，超过后边际收益指数衰减）以及 Shumailov 等人（2024）关于用 AI 生成数据递归训练下一代模型引发\textbf{分布尾部退化（Model Collapse）}的理论，探讨为何“可验证的合成数据（如数学证明验证器与代码单元测试筛选）”有望打破这一诅咒？
  \exHint{纯无筛选自采样会因有限样本方差逐代抹除长尾概率质量；而外部验证器（如 Lean 证明检查器或编译器）注入了真实的新信息（负熵）！}
\end{enumerate}

\chapter[深度生成：VAE、扩散与流匹配]{深度生成模型：VAE、扩散模型（DDPM）与流匹配}
\label{ch:vae-diffusion-flow}

\chapEpigraph{我们为自己构筑事实的图像。}{Ludwig Wittgenstein（路德维希·维特根斯坦，《逻辑哲学论》，2.1）}

\begin{chapIntro}{本章导读：从“理解数据”到“从纯噪声中凭空创造世界”}
在前五章里，我们已经掌握了深度神经网络的完整工具箱：能用反向模式自动微分计算任意深度计算图的精确梯度，能用 Kaiming 初始化与残差连接驯服百层网络的训练动力学，也能用 Transformer 与自监督预训练对离散文本序列建立极高精度的条件概率 $p_\theta(x_t \mid x_{<t})$。然而，面对高维\textbf{连续世界}——一张 $1024 \times 1024 \times 3$ 的照片、一段高保真声波，或一个由数百个原子三维坐标构成的蛋白质构象——逐像素做离散自回归既缓慢又违背空间整体性。我们真正渴望的是学会一个显式或隐式的连续生成分布 $p_\theta(x)$，能够从最平庸的各向同性高斯白噪声 $\varepsilon \sim \mathcal{N}(0, I)$ 出发，凭空“雕刻”出逼真且多样的新样本。

本章作为第 5 卷的压轴篇章，将从\textbf{第 4 卷第 6 章 EM 算法的证据下界（ELBO）}这一熟悉桥梁启程，用同一套变分推断与微分方程语言，贯穿当代深度生成模型的三大正统范式：
\begin{enumerate}
  \item \textbf{变分自编码器（VAE）}：将 EM 的逐样本变分后验升级为神经网络摊销编码器 $q_\phi(z \mid x)$，严格证明对角高斯 KL 散度的闭式公式，剖析\textbf{重参数化技巧（Reparameterization Trick）}为何能穿透随机节点并把梯度方差压缩数百倍，以及后验坍缩（Posterior Collapse）的机理与对策；
  \item \textbf{去噪扩散概率模型（DDPM）与分数匹配}：将单步隐变量拓展为 $T$ 步层次化马尔可夫加噪链，用数学归纳法证明前向任意步闭式边缘分布 $q(x_t \mid x_0)$，通过贝叶斯定理与配方法\textbf{零跳步}导出逆向后验均值 $\tilde{\mu}_t$ 与方差 $\tilde{\beta}_t$，揭示简洁的 $\varepsilon$-预测均方误差与\textbf{去噪分数匹配（Denoising Score Matching）}的严格等价性，并以纯 PyTorch 二维四模态环分布实测（\texttt{data/ddpm\_train.csv}）直观验证；
  \item \textbf{流匹配（Flow Matching / Rectified Flow）}：突破扩散模型的弯曲随机轨迹，直接在噪声 $x_0$ 与真实数据 $x_1$ 之间架设笔直的常微分方程（ODE）传输路径 $x_t = (1-t)x_0 + t x_1$，严格证明条件流匹配（CFM）与边际速度场回归的梯度等价性与连续性方程守恒，阐明它何以成为 Stable Diffusion 3、Flux 等新一代文生图/文生视频引擎的共同底座。
\end{enumerate}
\end{chapIntro}

\section{从 EM 算法的 ELBO 到变分自编码器（VAE）}
\label{sec:vae}

\begin{recallBox}{第 4 卷第 6 章：隐变量模型的 ELBO 恒等分解}
给定观测数据 $x$ 与不可观测的连续隐变量 $z$，若联合生成模型为 $p_\theta(x, z) = p_\theta(x \mid z)\, p(z)$，则对\textbf{任意}辅助变分分布 $q(z)$，边际对数似然都严格满足如下恒等式：
\begin{equation}
  \ln p_\theta(x)
  = \underbrace{\E_{z \sim q(z)}\!\left[\ln \frac{p_\theta(x, z)}{q(z)}\right]}_{\mathcal{L}_{\mathrm{ELBO}}(q,\, \theta;\, x)}
  + \underbrace{\DKL\!\big(q(z) \,\|\, p_\theta(z \mid x)\big)}_{\ge\, 0}.
  \label{eq:elbo-recall}
\end{equation}
由于 KL 散度非负，最大化证据下界 $\mathcal{L}_{\mathrm{ELBO}}$ 既能抬高真实的边际对数似然 $\ln p_\theta(x)$，又能迫使变分分布 $q(z)$ 贴近真实后验 $p_\theta(z \mid x)$。
\end{recallBox}

在第 4 卷的高斯混合模型（GMM）里，隐变量 $z$ 只是 $K$ 个离散簇编号，真实后验 $p_\theta(z \mid x)$ 可由贝叶斯公式解析算出（E 步），因此 EM 算法畅通无阻。但当我们把解码器 $p_\theta(x \mid z)$ 升级为一个多层深度神经网络（例如把 64 维隐向量 $z$ 映射为百万像素高清图像的卷积或 Transformer 解码器）时，经典 EM 瞬间撞上两堵铜墙铁壁：
\begin{enumerate}
  \item \textbf{真实后验彻底不可积分}：$p_\theta(z \mid x) = \frac{p_\theta(x \mid z)p(z)}{\int p_\theta(x \mid z')p(z')\,dz'}$ 的分母涉及深度非线性网络在连续高维空间上的积分，不存在任何闭式解，E 步无从下笔；
  \item \textbf{逐样本变分参数不可扩展}：若按传统平均场变分推断，为每张样本 $x^{(i)}$ 单独维护参数 $(\mu^{(i)}, \sigma^{(i)})$，面对 $N = 10^7$ 张图像不仅显存爆炸，且测试遇到新图片时仍需从头运行迭代优化。
\end{enumerate}

Kingma 与 Welling（2014）以及 Rezende 等人（2014）提出的\textbf{变分自编码器（Variational Autoencoder, VAE）}用一记惊艳的\textbf{摊销变分推断（Amortized Variational Inference）}同时粉碎了这两道壁垒：不再为每个样本单独优化变分参数，而是训练一个共享参数的\textbf{神经编码器（识别网络）} $q_\phi(z \mid x)$，直接把输入样本 $x$ 映射为近似后验的参数 $(\mu_\phi(x),\, \sigma_\phi(x))$！

\subsection{VAE 的概率结构与 ELBO 目标}

\begin{defBox}
\begin{definition}[变分自编码器（VAE）及其证据下界]
\label{def:vae}
设观测样本 $x \in \R^D$，隐变量 $z \in \R^d$（通常 $d \ll D$）。标准高斯 VAE 由以下两部分耦合而成：
\begin{enumerate}
  \item \textbf{生成模型（解码器）}：先验取标准各向同性高斯 $p(z) = \mathcal{N}(0, I_d)$；给定 $z$，由参数为 $\theta$ 的解码网络给出观测似然 $p_\theta(x \mid z)$（对连续像素常取 $\mathcal{N}(\mu_\theta(z),\, \sigma_{\text{dec}}^2 I_D)$，对应均方误差重建；对二值图像取逐像素伯努利分布，对应二元交叉熵）；
  \item \textbf{摊销推断模型（编码器）}：由参数为 $\phi$ 的编码网络输出均值向量 $\mu_\phi(x) \in \R^d$ 与对角标准差向量 $\sigma_\phi(x) \in \R_+^d$，定义因式分解高斯变分后验：
  \begin{equation}
    q_\phi(z \mid x) = \mathcal{N}\!\big(z;\, \mu_\phi(x),\, \diag(\sigma_\phi^2(x))\big)
    = \prod_{j=1}^{d} \mathcal{N}\!\big(z_j;\, \mu_{\phi, j}(x),\, \sigma_{\phi, j}^2(x)\big).
    \label{eq:vae-encoder}
  \end{equation}
\end{enumerate}
将式~\eqref{eq:elbo-recall} 中的联合密度拆为 $p_\theta(x, z) = p_\theta(x \mid z)\,p(z)$，立得单样本的 \textbf{VAE 证据下界（ELBO）}：
\begin{equation}
  \mathcal{L}_{\mathrm{VAE}}(\theta, \phi;\, x)
  = \underbrace{\E_{z \sim q_\phi(z \mid x)}\!\left[\ln p_\theta(x \mid z)\right]}_{\text{重建保真项（Reconstruction）}}
  \;-\; \underbrace{\DKL\!\big(q_\phi(z \mid x) \,\|\, p(z)\big)}_{\text{先验正则项（KL Regularization）}}
  \;\le\; \ln p_\theta(x).
  \label{eq:vae-elbo}
\end{equation}
训练时最小化整个数据集上的负 ELBO：$\min_{\theta,\phi}\; -\frac{1}{N}\sum_{i=1}^{N} \mathcal{L}_{\mathrm{VAE}}(\theta, \phi;\, x^{(i)})$；生成新样本时则彻底丢弃编码器，直接采 $z \sim \mathcal{N}(0, I_d)$ 并前向经过解码器即可！
\end{definition}
\end{defBox}

观察式~\eqref{eq:vae-elbo} 令人拍案叫绝的对称张力：
\begin{itemize}
  \item 第一项\textbf{重建保真项}要求编码器把 $x$ 压缩进 $z$ 后，解码器仍能把 $x$ 原封不动还原出来；若只有这一项，模型就退化为普通确定性自编码器——它只会把训练样本零散塞进隐空间的偏僻角落，各点之间全是“无人区”，随机采样 $z \sim \mathcal{N}(0, I_d)$ 解码出来只会是一团乱码；
  \item 第二项\textbf{KL 正则项}强行把每个样本对应的后验云团 $q_\phi(z \mid x)$ 拉向同一个标准高斯球心 $\mathcal{N}(0, I_d)$ 并撑开其方差，迫使不同样本的后验云团彼此平滑重叠、填满整个先验支撑集，从而保证隐空间处处可插值、处处可生成！
\end{itemize}

\subsection{对角高斯 KL 散度的精确闭式推导}

式~\eqref{eq:vae-elbo} 中的第二项甚至不需要任何蒙特卡洛采样——对于对角高斯编码器与标准正态先验，它拥有极其简洁的精确解析公式！

\begin{thmBox}
\begin{theorem}[对角高斯与标准正态分布的闭式 KL 散度]
\label{thm:gaussian-kl}
设 $q(z) = \mathcal{N}\!\big(\mu,\, \diag(\sigma_1^2, \dots, \sigma_d^2)\big)$（其中 $\sigma_j > 0$）且 $p(z) = \mathcal{N}(0, I_d)$，则两者之间的 KL 散度精确等于：
\begin{equation}
  \DKL\!\big(q(z) \,\|\, p(z)\big)
  = \frac{1}{2} \sum_{j=1}^{d} \left(\mu_j^2 + \sigma_j^2 - \ln \sigma_j^2 - 1\right).
  \label{eq:gaussian-kl-closed}
\end{equation}
而且该式严格非负，等号成立当且仅当对所有维度 $j = 1, \dots, d$ 均有 $\mu_j = 0$ 且 $\sigma_j^2 = 1$。
\end{theorem}
\begin{proof}
由于 $q(z) = \prod_{j=1}^d q_j(z_j)$ 与 $p(z) = \prod_{j=1}^d p_j(z_j)$ 均在各坐标上独立因式分解，对数密度比自然拆为逐维求和：
\begin{align*}
  \ln q(z) &= -\frac{d}{2}\ln(2\pi) - \frac{1}{2}\sum_{j=1}^{d} \ln\sigma_j^2 - \frac{1}{2}\sum_{j=1}^{d} \frac{(z_j - \mu_j)^2}{\sigma_j^2}, \\
  \ln p(z) &= -\frac{d}{2}\ln(2\pi) - \frac{1}{2}\sum_{j=1}^{d} z_j^2.
\end{align*}
两式相减，常数项 $-\frac{d}{2}\ln(2\pi)$ 精确抵消：
\[
  \ln \frac{q(z)}{p(z)} = \frac{1}{2} \sum_{j=1}^{d} \left(z_j^2 - \frac{(z_j - \mu_j)^2}{\sigma_j^2} - \ln\sigma_j^2\right).
\]
在 $z \sim q(z)$ 下对上式逐项取数学期望。由高斯分布的一阶与二阶矩定义，对每个坐标 $j$ 均有：
\[
  \E_{q}\!\left[\frac{(z_j - \mu_j)^2}{\sigma_j^2}\right] = \frac{\Var_{q}(z_j)}{\sigma_j^2} = \frac{\sigma_j^2}{\sigma_j^2} = 1,
  \qquad
  \E_{q}\!\left[z_j^2\right] = \Var_{q}(z_j) + \big(\E_{q}[z_j]\big)^2 = \sigma_j^2 + \mu_j^2.
\]
代入期望表达式，立即得到式~\eqref{eq:gaussian-kl-closed}：
\[
  \DKL(q \,\|\, p)
  = \E_{z \sim q}\!\left[\ln \frac{q(z)}{p(z)}\right]
  = \frac{1}{2}\sum_{j=1}^{d}\left(\mu_j^2 + \sigma_j^2 - \ln\sigma_j^2 - 1\right).
\]
最后验证非负性与取等条件：令 $g(u) = u - \ln u - 1$（$u > 0$），其导数 $g'(u) = 1 - \frac{1}{u}$ 在 $u \in (0, 1)$ 严格为负、在 $u \in (1, +\infty)$ 严格为正，故 $g(u) \ge g(1) = 0$ 且等号仅在 $u = 1$ 取得。结合 $\mu_j^2 \ge 0$，可见式~\eqref{eq:gaussian-kl-closed} 每一项均非负，且总和为零当且仅当 $\mu_j = 0$ 与 $\sigma_j^2 = 1$ 对全部 $j$ 同时成立。
\end{proof}
\end{thmBox}

\begin{example}[手算验收：对角高斯 KL 散度的直观惩罚]
设隐空间维度 $d = 2$，编码器对某输入样本输出 $\mu = (1.0,\, -0.5)$ 与标准差 $\sigma = (0.5,\, 2.0)$（方差 $\sigma^2 = (0.25,\, 4.0)$）。代入式~\eqref{eq:gaussian-kl-closed}：
\begin{align*}
  \text{第 1 维} &: \frac{1}{2}\big(1.0^2 + 0.25 - \ln 0.25 - 1\big) = \frac{1}{2}(0.25 + 1.3863) \approx 0.8181\text{ nats}, \\
  \text{第 2 维} &: \frac{1}{2}\big((-0.5)^2 + 4.0 - \ln 4.0 - 1\big) = \frac{1}{2}(3.25 - 1.3863) \approx 0.9319\text{ nats}.
\end{align*}
合计 $\DKL(q \,\|\, p) \approx 1.750\text{ nats}$。注意 $-\ln \sigma_j^2$ 项的精妙之处：无论编码器试图把方差缩成确定性点质量（$\sigma_j \to 0^+ \implies -\ln\sigma_j^2 \to +\infty$），还是无限拉宽（$\sigma_j \to +\infty \implies \sigma_j^2 \to +\infty$），KL 散度都会施加无穷大惩罚！工程代码中，编码器末层几乎总是直接输出\textbf{对数方差} $s_j = \ln \sigma_j^2 \in \R$（无范围约束），此时式~\eqref{eq:gaussian-kl-closed} 化作一行数值极稳定的代码：\texttt{-0.5 * torch.sum(1 + logvar - mu.pow(2) - logvar.exp(), dim=-1)}。
\end{example}

\subsection{重参数化技巧（Reparameterization Trick）与低方差梯度证明}
\label{sec:reparam}

搞定了 KL 散度项后，还剩下式~\eqref{eq:vae-elbo} 的第一项重建期望 $\E_{z \sim q_\phi(z \mid x)}[f_\theta(z)]$（其中 $f_\theta(z) := \ln p_\theta(x \mid z)$）。对解码器参数 $\theta$ 求梯度毫无难度，因为期望分布不依赖 $\theta$，梯度可直接穿进期望号内。真正的拦路虎在于对\textbf{编码器参数 $\phi$} 求梯度：
\[
  \nabla_\phi \E_{z \sim q_\phi(z \mid x)}\!\left[f_\theta(z)\right] \;=\; ?
\]
在计算图中，$z \sim q_\phi(z \mid x)$ 是一个\textbf{随机采样节点}——第 1 章的反向模式自动微分（链式法则）无法直接穿透一个依赖参数 $\phi$ 的随机抽样门！历史上主要有两条出路：

\begin{enumerate}
  \item \textbf{分数函数估计器（REINFORCE / Log-Derivative Trick）}：利用恒等式 $\nabla_\phi q_\phi(z) = q_\phi(z)\nabla_\phi \ln q_\phi(z)$，将梯度改写为：
  \begin{equation}
    \nabla_\phi \E_{z \sim q_\phi}\!\left[f(z)\right]
    = \int f(z)\,\nabla_\phi q_\phi(z)\,dz
    = \E_{z \sim q_\phi}\!\left[f(z)\,\nabla_\phi \ln q_\phi(z)\right].
    \label{eq:reinforce}
  \end{equation}
  虽然式~\eqref{eq:reinforce} 严格无偏且适用于离散隐变量（它在第 6 卷强化学习策略梯度中是核心基石），但它在连续生成模型中存在致命短板：\textbf{完全无视 $f(z)$ 自身可微的精确梯度信息 $\nabla_z f(z)$}，仅靠黑箱标量打分 $f(z)$ 盲目试错，导致单样本蒙特卡洛估计的方差高得惊人，训练极难收敛！
  \item \textbf{重参数化技巧（Reparameterization / Pathwise Estimator）}：将随机性从参数 $\phi$ 上\textbf{剥离}到一个固定无参的外部噪声源 $\varepsilon \sim \mathcal{N}(0, I_d)$，把采样改写为关于 $(\mu_\phi, \sigma_\phi, \varepsilon)$ 的确定性可微映射：
  \begin{equation}
    z = g_\phi(\varepsilon, x) = \mu_\phi(x) + \sigma_\phi(x) \had \varepsilon,
    \qquad \varepsilon \sim \mathcal{N}(0, I_d).
    \label{eq:reparam}
  \end{equation}
  此时随机变量 $g_\phi(\varepsilon, x)$ 的概率分布与 $q_\phi(z \mid x)$ \textbf{完全同分布}，且期望的积分测度不再依赖 $\phi$！于是梯度运算符可以合法穿过期望号，沿着确定性路径直接调用第 1 章的链式法则：
  \begin{equation}
    \nabla_\phi \E_{z \sim q_\phi(z \mid x)}\!\left[f_\theta(z)\right]
    = \E_{\varepsilon \sim \mathcal{N}(0, I_d)}\!\left[\nabla_\phi f_\theta\!\big(\mu_\phi(x) + \sigma_\phi(x)\had \varepsilon\big)\right]
    = \E_{\varepsilon}\!\left[\big(\nabla_z f_\theta(z)\big)^\top \frac{\partial g_\phi(\varepsilon, x)}{\partial \phi}\right].
    \label{eq:reparam-grad}
  \end{equation}
\end{enumerate}

重参数化到底把梯度方差降低了多少？我们用一道极干净的解析基准给出精确回答！

\begin{thmBox}
\begin{proposition}[重参数化与 REINFORCE 的解析方差鸿沟]
\label{prop:reparam-var}
考虑一维二次目标 $f(z) = z^2$ 与变分分布 $q_\mu(z) = \mathcal{N}(\mu, \sigma^2)$（$\sigma > 0$ 固定），其真实均值梯度为 $\frac{\partial}{\partial \mu}\E_{q_\mu}[z^2] = \frac{\partial}{\partial \mu}(\mu^2 + \sigma^2) = 2\mu$。设 $\hat{g}_{\mathrm{reparam}}$ 与 $\hat{g}_{\mathrm{RF}}$ 分别为重参数化与 REINFORCE 的单样本无偏梯度估计器，则两者方差精确为：
\begin{equation}
  \Var(\hat{g}_{\mathrm{reparam}}) = 4\sigma^2,
  \qquad
  \Var(\hat{g}_{\mathrm{RF}}) = \frac{\mu^4}{\sigma^2} + 14\mu^2 + 15\sigma^2.
  \label{eq:reparam-vs-rf-var}
\end{equation}
\end{proposition}
\begin{proof}
记标准正态变量 $\varepsilon \sim \mathcal{N}(0, 1)$，利用标准正态高阶中心矩 $\E[\varepsilon] = \E[\varepsilon^3] = \E[\varepsilon^5] = 0$、$\E[\varepsilon^2] = 1$、$\E[\varepsilon^4] = 3$、$\E[\varepsilon^6] = 15$：
\begin{enumerate}
  \item \textbf{重参数化估计器}：令 $z = \mu + \sigma \varepsilon$，直接求导得 $\hat{g}_{\mathrm{reparam}} = \frac{\partial}{\partial \mu}(\mu + \sigma\varepsilon)^2 = 2\mu + 2\sigma\varepsilon$。显然 $\E[\hat{g}_{\mathrm{reparam}}] = 2\mu$，其方差为
  \[
    \Var(\hat{g}_{\mathrm{reparam}}) = \Var(2\sigma\varepsilon) = 4\sigma^2,
  \]
  完全不随均值位置 $\mu$ 膨胀，且随不确定度 $\sigma \to 0$ 趋于零。
  \item \textbf{REINFORCE 估计器}：由 $\frac{\partial}{\partial \mu}\ln q_\mu(z) = \frac{z - \mu}{\sigma^2} = \frac{\varepsilon}{\sigma}$，得
  \[
    \hat{g}_{\mathrm{RF}} = z^2 \frac{z - \mu}{\sigma^2}
    = (\mu + \sigma\varepsilon)^2 \frac{\varepsilon}{\sigma}
    = \frac{\mu^2}{\sigma}\varepsilon + 2\mu\varepsilon^2 + \sigma\varepsilon^3.
  \]
  易见 $\E[\hat{g}_{\mathrm{RF}}] = 2\mu\,\E[\varepsilon^2] = 2\mu$（同样无偏）。计算其二阶矩（奇数次项期望全为零）：
  \begin{align*}
    \E[\hat{g}_{\mathrm{RF}}^2]
    &= \frac{\mu^4}{\sigma^2}\E[\varepsilon^2] + 4\mu^2\E[\varepsilon^4] + \sigma^2\E[\varepsilon^6] + 2\left(\frac{\mu^2}{\sigma}\right)\sigma\,\E[\varepsilon^4] \\
    &= \frac{\mu^4}{\sigma^2} + 12\mu^2 + 15\sigma^2 + 6\mu^2
    = \frac{\mu^4}{\sigma^2} + 18\mu^2 + 15\sigma^2.
  \end{align*}
  减去均值平方 $(2\mu)^2 = 4\mu^2$，即得 $\Var(\hat{g}_{\mathrm{RF}}) = \frac{\mu^4}{\sigma^2} + 14\mu^2 + 15\sigma^2$。
\end{enumerate}
\end{proof}
\end{thmBox}

代入一组最平常的数值 $\mu = 3,\; \sigma = 0.5$：重参数化的梯度方差仅为 $4 \times 0.25 = 1.0$，而 REINFORCE 的方差竟高达 $324 + 126 + 3.75 = \mathbf{453.75}$！更致命的是，当编码器学得越来越精准（$\sigma \to 0$）时，重参数化方差收敛到 $0$，REINFORCE 方差却按 $O(\mu^4 / \sigma^2)$ \textbf{发散爆炸}！正是重参数化带来的极低方差，使得 VAE 在每个训练步只需采样 $L = 1$ 个噪声 $\varepsilon$ 即可配合 AdamW 丝滑收敛。

\begin{failBox}{VAE 的两大经典顽疾：后验坍缩（Posterior Collapse）与图像模糊}
\begin{enumerate}
  \item \textbf{后验坍缩（Posterior Collapse）}：当解码器 $p_\theta(x \mid z)$ 过于强大（例如采用自回归 LSTM、PixelCNN 或 Transformer 解码器）时，即使完全无视输入隐码 $z$，解码器单凭像素或词元之间的局部相关性也能取得很低的负对数似然 $-\ln p_\theta(x \mid z)$。此时优化器为了白白赚取 KL 项的零损失，会迅速把所有样本的编码器推向恒等先验：
  \[
    q_\phi(z \mid x) \equiv \mathcal{N}(0, I_d) \quad (\forall x) \implies \DKL\!\big(q_\phi(z \mid x) \,\|\, p(z)\big) \to 0,\quad I_q(X; Z) \to 0!
  \]
  这就是令人头疼的\textbf{后验坍缩}——编码器彻底摆烂，隐变量 $z$ 与输入 $x$ 互信息归零，模型失去了所有可控表征能力。\textbf{三大工业急救法}：① \textbf{KL 退火（KL Annealing）}：在目标函数中引入系数 $\beta_t$，训练前期令 $\beta_t$ 从 $0$ 线性热身至 $1$，先让编码器把有用信息写入 $z$ 再逐步施加先验约束；\mbox{② \textbf{自由比特（Free Bits）}}：把逐维 KL 改为 $\max(\lambda_{\text{free}},\, \DKL^{(j)})$，允许每个维度免费携带至少 $\lambda_{\text{free}}$ nats 信息不受压制；③ \textbf{离散码本（VQ-VAE）}：将连续隐变量量化到离散码本，彻底规避高斯 KL 坍缩（这也是 Stable Diffusion 早期潜空间压缩器的先驱）。
  \item \textbf{高斯似然导致的“平均化模糊”}：若解码器假定像素服从各向同性高斯 $p_\theta(x \mid z) = \mathcal{N}(\mu_\theta(z), \sigma^2 I)$，最大化对数似然等价于最小化像素 $L_2$ 误差 $\|x - \mu_\theta(z)\|_2^2$。当同一隐码 $z$ 对应多种合理解码细节（如头发丝向左或向右偏 1 像素）时，最小化 $L_2$ 风险的最优预测是\textbf{所有可能结果的算术平均}——在视觉上便呈现出一层挥之不去的“磨砂模糊感”。
\end{enumerate}
\end{failBox}

\section{扩散模型（DDPM）：层次化加噪、逆向后验与分数匹配}
\label{sec:ddpm}

为什么单步 VAE 在高分辨率图像生成上总显吃力？本质原因在于它试图让解码器\textbf{一步登天}：一次性把低维正态噪声直接映射为百万像素的精细流形。既然一步跨不过天堑，一个天才而返璞归真的想法呼之欲出：\textbf{把“从纯噪声画出一幅画”拆解成 $T$ 个微小的、几乎线性的渐进式去噪小步}！

这正是 Sohl-Dickstein 等人（2015）奠基、Ho 等人（2020）以极简形式引爆全球的\textbf{去噪扩散概率模型（Denoising Diffusion Probabilistic Models, DDPM）}。在数学上，DDPM 可以精确视作一个\textbf{与数据同维的 $T$ 层层次化 VAE}，其精妙之处在于：\textbf{前向编码器（加噪过程）完全固定、无需任何可训练参数}，从而将全部神经网络容量倾注于学习逆向去噪！

\subsection{前向马尔可夫加噪与任意步闭式采样定理}

给定数据分布样本 $x_0 \sim q(x_0)$（$x_0 \in \R^D$）及预先设定的单调递增小方差表（Noise Schedule）$0 < \beta_1 < \beta_2 < \dots < \beta_T < 1$，定义长度为 $T$ 的前向马尔可夫加噪链：
\begin{equation}
  q(x_{1:T} \mid x_0) = \prod_{t=1}^{T} q(x_t \mid x_{t-1}),
  \qquad
  q(x_t \mid x_{t-1}) = \mathcal{N}\!\big(x_t;\, \sqrt{1 - \beta_t}\, x_{t-1},\; \beta_t I\big).
  \label{eq:ddpm-forward-step}
\end{equation}
注意系数 $\sqrt{1 - \beta_t}$ 的用意：它恰好保持方差量级守恒——若 $\Var(x_{t-1}) = I$，则 $\Var(\sqrt{1-\beta_t}\,x_{t-1} + \sqrt{\beta_t}\,\varepsilon) = (1-\beta_t)I + \beta_t I = I$！随着步数 $t$ 推进，原始信号不断按 $\sqrt{1-\beta_t}$ 衰减、高斯白噪声不断注入，当 $T$ 足够大时，$x_T$ 的分布便与纯标准高斯 $\mathcal{N}(0, I)$ 几乎无法区分。

为方便书写，引入两个贯穿全章的标准记号：
\begin{equation}
  \alpha_t := 1 - \beta_t, \qquad \bar{\alpha}_t := \prod_{s=1}^{t} \alpha_s
  \qquad (\text{约定 } \bar{\alpha}_0 := 1).
  \label{eq:ddpm-alpha}
\end{equation}
训练时，我们绝不需要真的执行 $t$ 次串行加噪才能拿到 $x_t$——高斯分布的线性可加性赋予了我们“一步直达任意时刻 $t$”的闭式时空穿梭机！

\begin{thmBox}
\begin{theorem}[DDPM 前向扩散的任意步闭式边缘分布]
\label{thm:ddpm-forward}
在前向马尔可夫过程~\eqref{eq:ddpm-forward-step} 下，对任意步数 $t \in \{1, 2, \dots, T\}$，给定初始数据 $x_0$ 时的条件边缘分布精确为：
\begin{equation}
  q(x_t \mid x_0) = \mathcal{N}\!\big(x_t;\, \sqrt{\bar{\alpha}_t}\, x_0,\; (1 - \bar{\alpha}_t) I\big),
  \label{eq:ddpm-xt-closed}
\end{equation}
等价地，可用单个标准高斯噪声 $\varepsilon \sim \mathcal{N}(0, I)$ 一步重参数化采样：
\begin{equation}
  x_t = \sqrt{\bar{\alpha}_t}\, x_0 + \sqrt{1 - \bar{\alpha}_t}\,\varepsilon.
  \label{eq:ddpm-xt-reparam}
\end{equation}
\end{theorem}
\begin{proof}
对时间步 $t$ 施用数学归纳法。
\begin{itemize}
  \item \textbf{奠基步骤（$t = 1$）}：由于 $\bar{\alpha}_1 = \alpha_1 = 1 - \beta_1$，式~\eqref{eq:ddpm-forward-step} 直接给出
  \[
    q(x_1 \mid x_0) = \mathcal{N}\!\big(\sqrt{\alpha_1}\, x_0,\, (1 - \alpha_1) I\big)
    = \mathcal{N}\!\big(\sqrt{\bar{\alpha}_1}\, x_0,\, (1 - \bar{\alpha}_1) I\big),
  \]
  命题在 $t = 1$ 成立。
  \item \textbf{归纳假设与递推（$t-1 \implies t$）}：假设对 $t - 1$ 成立 $x_{t-1} = \sqrt{\bar{\alpha}_{t-1}}\, x_0 + \sqrt{1 - \bar{\alpha}_{t-1}}\, \bar{\varepsilon}_{t-1}$，其中 $\bar{\varepsilon}_{t-1} \sim \mathcal{N}(0, I)$。按第 $t$ 步转移规则~\eqref{eq:ddpm-forward-step} 引入独立标准正态噪声 $\varepsilon_t \sim \mathcal{N}(0, I)$（与 $\bar{\varepsilon}_{t-1}$ 独立）：
  \begin{align*}
    x_t &= \sqrt{\alpha_t}\, x_{t-1} + \sqrt{1 - \alpha_t}\,\varepsilon_t \\
        &= \sqrt{\alpha_t}\!\left(\sqrt{\bar{\alpha}_{t-1}}\, x_0 + \sqrt{1 - \bar{\alpha}_{t-1}}\,\bar{\varepsilon}_{t-1}\right) + \sqrt{1 - \alpha_t}\,\varepsilon_t \\
        &= \sqrt{\alpha_t \bar{\alpha}_{t-1}}\, x_0
           + \underbrace{\sqrt{\alpha_t(1 - \bar{\alpha}_{t-1})}\,\bar{\varepsilon}_{t-1} + \sqrt{1 - \alpha_t}\,\varepsilon_t}_{\text{两个独立零均值高斯随机向量之和}}.
  \end{align*}
  由于独立高斯随机向量相加仍为高斯分布且方差直接相加，上式第二部分服从零均值各向同性高斯分布，其标量方差精确化简为：
  \[
    \alpha_t(1 - \bar{\alpha}_{t-1}) + (1 - \alpha_t)
    = \alpha_t - \alpha_t\bar{\alpha}_{t-1} + 1 - \alpha_t
    = 1 - \bar{\alpha}_t.
  \]
  又 $\sqrt{\alpha_t \bar{\alpha}_{t-1}} = \sqrt{\bar{\alpha}_t}$，故 $x_t$ 与 $\sqrt{\bar{\alpha}_t}\,x_0 + \sqrt{1 - \bar{\alpha}_t}\,\varepsilon$（$\varepsilon \sim \mathcal{N}(0, I)$）严格同分布。由数学归纳法原理，式~\eqref{eq:ddpm-xt-closed} 对所有 $t = 1, \dots, T$ 成立。
\end{itemize}
\end{proof}
\end{thmBox}

\subsection{变分下界分解与逆向条件后验的精确推导}

现在定义参数化逆向去噪马尔可夫链：从纯高斯先验 $p(x_T) = \mathcal{N}(0, I)$ 起步，每步由参数为 $\theta$ 的神经网络预测高斯去噪转移核：
\begin{equation}
  p_\theta(x_{0:T}) = p(x_T)\prod_{t=1}^{T} p_\theta(x_{t-1} \mid x_t),
  \qquad
  p_\theta(x_{t-1} \mid x_t) = \mathcal{N}\!\big(x_{t-1};\, \mu_\theta(x_t, t),\; \Sigma_\theta(x_t, t)\big).
  \label{eq:ddpm-reverse}
\end{equation}
沿用式~\eqref{eq:elbo-recall} 的 ELBO 思想（将整条轨迹 $x_{1:T}$ 视作高维隐变量 $z$），负对数似然受控于变分上界（VLB）：
\[
  \E_{q(x_0)}\!\left[-\ln p_\theta(x_0)\right]
  \le \E_{q(x_{0:T})}\!\left[\ln \frac{q(x_{1:T} \mid x_0)}{p_\theta(x_{0:T})}\right]
  = \E_{q}\!\left[-\ln p(x_T) + \sum_{t=2}^{T} \ln \frac{q(x_t \mid x_{t-1})}{p_\theta(x_{t-1} \mid x_t)} + \ln\frac{q(x_1 \mid x_0)}{p_\theta(x_0 \mid x_1)}\right].
\]
此处若直接做蒙特卡洛估计方差极大，因为分子 $q(x_t \mid x_{t-1})$ 与分母 $p_\theta(x_{t-1} \mid x_t)$ 的条件方向正好相反！关键一招是借助马尔可夫性与贝叶斯定理，给每个前向步添上 $x_0$ 条件并翻转方向：
\[
  q(x_t \mid x_{t-1}) = q(x_t \mid x_{t-1}, x_0)
  = \frac{q(x_{t-1} \mid x_t, x_0)\, q(x_t \mid x_0)}{q(x_{t-1} \mid x_0)} \qquad (t \ge 2).
\]
代入求和中，相邻项的 $\ln q(x_t \mid x_0)$ 与 $-\ln q(x_{t-1} \mid x_0)$ 发生华丽的\textbf{望远镜对消（Telescoping Cancellation）}，仅余首尾两项 $\ln q(x_T \mid x_0) - \ln q(x_1 \mid x_0)$，而 $-\ln q(x_1 \mid x_0)$ 又恰好与末项 $+\ln q(x_1 \mid x_0)$ 抵消！整理后，整条扩散链的变分上界优雅分解为一系列完全可解析计算的 KL 散度之和：
\begin{equation}
  \mathcal{L}_{\mathrm{VLB}}(\theta)
  = \E_{q}\!\Bigg[
    \underbrace{\DKL\!\big(q(x_T \mid x_0) \,\|\, p(x_T)\big)}_{L_T\text{（无参数常数）}}
    + \sum_{t=2}^{T} \underbrace{\DKL\!\big(q(x_{t-1} \mid x_t, x_0) \,\|\, p_\theta(x_{t-1} \mid x_t)\big)}_{L_{t-1}\text{（逐层高斯去噪匹配）}}
    \underbrace{-\ln p_\theta(x_0 \mid x_1)}_{L_0\text{（末步重建）}}
  \Bigg].
  \label{eq:ddpm-vlb}
\end{equation}

要让 $p_\theta(x_{t-1} \mid x_t)$ 去精准贴近目标分布 $q(x_{t-1} \mid x_t, x_0)$，我们必须先把后者\textbf{彻底算出来}！令人欣喜的是，虽然无条件逆向分布 $q(x_{t-1} \mid x_t)$ 不可知（需要知道整个数据分布），但只要\textbf{额外给定起点 $x_0$}，$q(x_{t-1} \mid x_t, x_0)$ 竟然是一个完全解析的高斯分布！

\begin{thmBox}
\begin{theorem}[DDPM 逆向条件后验的精确闭式公式]
\label{thm:ddpm-posterior}
对任意 $t \in \{2, \dots, T\}$，给定当前加噪样本 $x_t$ 与原始干净样本 $x_0$ 时，前一步 $x_{t-1}$ 的真实条件后验分布严格服从各向同性高斯分布：
\begin{equation}
  q(x_{t-1} \mid x_t, x_0) = \mathcal{N}\!\big(x_{t-1};\, \tilde{\mu}_t(x_t, x_0),\; \tilde{\beta}_t I\big),
  \label{eq:ddpm-post-gauss}
\end{equation}
其精确方差与均值分别为：
\begin{equation}
  \tilde{\beta}_t = \frac{1 - \bar{\alpha}_{t-1}}{1 - \bar{\alpha}_t}\,\beta_t,
  \qquad
  \tilde{\mu}_t(x_t, x_0) = \frac{\sqrt{\bar{\alpha}_{t-1}}\,\beta_t}{1 - \bar{\alpha}_t}\, x_0
  + \frac{\sqrt{\alpha_t}(1 - \bar{\alpha}_{t-1})}{1 - \bar{\alpha}_t}\, x_t.
  \label{eq:ddpm-post-mean-x0}
\end{equation}
进一步，利用式~\eqref{eq:ddpm-xt-reparam} 将 $x_0 = \frac{1}{\sqrt{\bar{\alpha}_t}}\big(x_t - \sqrt{1 - \bar{\alpha}_t}\,\varepsilon\big)$ 代入，后验均值可化为极其简洁的\textbf{噪声参数化形式}：
\begin{equation}
  \tilde{\mu}_t(x_t, \varepsilon) = \frac{1}{\sqrt{\alpha_t}}\!\left(x_t - \frac{1 - \alpha_t}{\sqrt{1 - \bar{\alpha}_t}}\,\varepsilon\right).
  \label{eq:ddpm-post-mean-eps}
\end{equation}
\end{theorem}
\begin{proof}
\textbf{第一步：贝叶斯定理展开指数二次型。} 由贝叶斯公式与前向马尔可夫性质 $q(x_t \mid x_{t-1}, x_0) = q(x_t \mid x_{t-1})$，将 $q(x_{t-1} \mid x_t, x_0)$ 视为关于变量 $x_{t-1}$ 的密度（吸收一切仅依赖 $x_t, x_0$ 的归一化常数）：
\[
  q(x_{t-1} \mid x_t, x_0)
  = \frac{q(x_t \mid x_{t-1})\, q(x_{t-1} \mid x_0)}{q(x_t \mid x_0)}
  \propto q(x_t \mid x_{t-1})\, q(x_{t-1} \mid x_0).
\]
将式~\eqref{eq:ddpm-forward-step} 与定理~\ref{thm:ddpm-forward} 中的两个高斯密度代入，提取指数括号内的全部 $x_{t-1}$ 项：
\begin{align*}
  q(x_{t-1} \mid x_t, x_0)
  &\propto \exp\!\left(
    -\frac{1}{2}\!\left[
      \frac{\norm{x_t - \sqrt{\alpha_t}\, x_{t-1}}_2^2}{\beta_t}
      + \frac{\norm{x_{t-1} - \sqrt{\bar{\alpha}_{t-1}}\, x_0}_2^2}{1 - \bar{\alpha}_{t-1}}
    \right]
  \right) \\
  &\propto \exp\!\left(
    -\frac{1}{2}\!\left[
      \underbrace{\left(\frac{\alpha_t}{\beta_t} + \frac{1}{1 - \bar{\alpha}_{t-1}}\right)}_{A_t} \norm{x_{t-1}}_2^2
      - 2 \underbrace{\left(\frac{\sqrt{\alpha_t}}{\beta_t}\, x_t + \frac{\sqrt{\bar{\alpha}_{t-1}}}{1 - \bar{\alpha}_{t-1}}\, x_0\right)^\top}_{b_t^\top} x_{t-1}
    \right]
  \right).
\end{align*}

\textbf{第二步：配方法精确读出方差 $\tilde{\beta}_t$ 与均值 $\tilde{\mu}_t(x_t, x_0)$。} 对比标准高斯指数 $-\frac{1}{2\tilde{\beta}_t}\norm{x_{t-1} - \tilde{\mu}_t}_2^2 = -\frac{1}{2}\big(\tilde{\beta}_t^{-1}\norm{x_{t-1}}_2^2 - 2\tilde{\beta}_t^{-1}\tilde{\mu}_t^\top x_{t-1} + \text{const}\big)$，立即有 $\tilde{\beta}_t = A_t^{-1}$ 与 $\tilde{\mu}_t = \tilde{\beta}_t b_t$。先通分化简精度系数 $A_t$（注意 $\beta_t = 1 - \alpha_t$ 且 $\alpha_t \bar{\alpha}_{t-1} = \bar{\alpha}_t$）：
\[
  A_t = \frac{\alpha_t}{\beta_t} + \frac{1}{1 - \bar{\alpha}_{t-1}}
      = \frac{\alpha_t(1 - \bar{\alpha}_{t-1}) + \beta_t}{\beta_t(1 - \bar{\alpha}_{t-1})}
      = \frac{\alpha_t - \bar{\alpha}_t + 1 - \alpha_t}{\beta_t(1 - \bar{\alpha}_{t-1})}
      = \frac{1 - \bar{\alpha}_t}{\beta_t(1 - \bar{\alpha}_{t-1})}.
\]
取倒数即得后验方差：
\[
  \tilde{\beta}_t = A_t^{-1} = \frac{1 - \bar{\alpha}_{t-1}}{1 - \bar{\alpha}_t}\,\beta_t.
\]
将其乘以线性系数向量 $b_t$，精确约分即得式~\eqref{eq:ddpm-post-mean-x0}：
\[
  \tilde{\mu}_t(x_t, x_0)
  = \frac{1 - \bar{\alpha}_{t-1}}{1 - \bar{\alpha}_t}\,\beta_t
    \left(\frac{\sqrt{\alpha_t}}{\beta_t}\, x_t + \frac{\sqrt{\bar{\alpha}_{t-1}}}{1 - \bar{\alpha}_{t-1}}\, x_0\right)
  = \frac{\sqrt{\alpha_t}(1 - \bar{\alpha}_{t-1})}{1 - \bar{\alpha}_t}\, x_t
  + \frac{\sqrt{\bar{\alpha}_{t-1}}\,\beta_t}{1 - \bar{\alpha}_t}\, x_0.
\]

\textbf{第三步：代换 $x_0$ 化为纯噪声表达 $\tilde{\mu}_t(x_t, \varepsilon)$。} 由式~\eqref{eq:ddpm-xt-reparam} 反解 $x_0 = \frac{1}{\sqrt{\bar{\alpha}_t}}\big(x_t - \sqrt{1 - \bar{\alpha}_t}\,\varepsilon\big)$，并利用恒等式 $\frac{\sqrt{\bar{\alpha}_{t-1}}}{\sqrt{\bar{\alpha}_t}} = \frac{1}{\sqrt{\alpha_t}}$，代入上式：
\begin{align*}
  \tilde{\mu}_t(x_t, \varepsilon)
  &= \frac{\sqrt{\alpha_t}(1 - \bar{\alpha}_{t-1})}{1 - \bar{\alpha}_t}\, x_t
   + \frac{\beta_t}{\sqrt{\alpha_t}(1 - \bar{\alpha}_t)}\!\left(x_t - \sqrt{1 - \bar{\alpha}_t}\,\varepsilon\right) \\
  &= \frac{\alpha_t(1 - \bar{\alpha}_{t-1}) + (1 - \alpha_t)}{\sqrt{\alpha_t}(1 - \bar{\alpha}_t)}\, x_t
   - \frac{1 - \alpha_t}{\sqrt{\alpha_t}\sqrt{1 - \bar{\alpha}_t}}\,\varepsilon.
\end{align*}
注意第一项的分子恰为 $\alpha_t - \bar{\alpha}_t + 1 - \alpha_t = 1 - \bar{\alpha}_t$，与分母中的 $(1 - \bar{\alpha}_t)$ 奇迹般完全约去！于是提出公因子 $\frac{1}{\sqrt{\alpha_t}}$，即得式~\eqref{eq:ddpm-post-mean-eps}：
\[
  \tilde{\mu}_t(x_t, \varepsilon) = \frac{1}{\sqrt{\alpha_t}}\!\left(x_t - \frac{1 - \alpha_t}{\sqrt{1 - \bar{\alpha}_t}}\,\varepsilon\right).
\]
\end{proof}
\end{thmBox}

\subsection{极简 $\varepsilon$-预测目标与去噪分数匹配（Score Matching）的深层统一}

现在回到式~\eqref{eq:ddpm-vlb} 中的核心训练项 $L_{t-1} = \DKL\big(q(x_{t-1} \mid x_t, x_0) \,\|\, p_\theta(x_{t-1} \mid x_t)\big)$。Ho 等人（2020）将逆向方差固定为与输入无关的常数 $\Sigma_\theta(x_t, t) = \sigma_t^2 I$（实测取 $\sigma_t^2 = \beta_t$ 或 $\sigma_t^2 = \tilde{\beta}_t$ 表现相近），并依照式~\eqref{eq:ddpm-post-mean-eps} 的解析形式，用神经网络 $\varepsilon_\theta(x_t, t)$ 预测被掺入的噪声 $\varepsilon$，从而参数化逆向均值：
\begin{equation}
  \mu_\theta(x_t, t) := \frac{1}{\sqrt{\alpha_t}}\!\left(x_t - \frac{1 - \alpha_t}{\sqrt{1 - \bar{\alpha}_t}}\,\varepsilon_\theta(x_t, t)\right).
  \label{eq:ddpm-mu-theta}
\end{equation}
由两个同方差各向同性高斯分布的 KL 散度公式（只剩均值欧氏距离平方），立得：
\[
  L_{t-1}(\theta)
  = \E_{x_0, \varepsilon}\!\left[\frac{1}{2\sigma_t^2}\norm{\tilde{\mu}_t(x_t, \varepsilon) - \mu_\theta(x_t, t)}_2^2\right]
  = \E_{x_0, \varepsilon}\!\left[
      \underbrace{\frac{(1 - \alpha_t)^2}{2\sigma_t^2 \alpha_t (1 - \bar{\alpha}_t)}}_{w_t}
      \norm{\varepsilon - \varepsilon_\theta(\sqrt{\bar{\alpha}_t}\,x_0 + \sqrt{1 - \bar{\alpha}_t}\,\varepsilon,\, t)}_2^2
    \right].
\]
Ho 等人发现了一件极富启发性的工程事实：直接\textbf{舍弃}随时间步剧烈变化的前置权重 $w_t$（它把过重权重压在噪声极小的 $t \approx 1$ 细枝末节上，反而不利于学习全局结构），让所有时间步 $t \in \{1, \dots, T\}$ 等权平均，不仅把整个变分下界压缩成了令人震撼的一行均方误差，反而大幅提升了真实图像生成质量！这就是著名的 \textbf{DDPM 简化训练目标}：
\begin{equation}
  \mathcal{L}_{\mathrm{simple}}(\theta)
  = \E_{t \sim \mathcal{U}\{1,\dots,T\},\; x_0 \sim q(x_0),\; \varepsilon \sim \mathcal{N}(0, I)}
    \!\left[\norm{\varepsilon - \varepsilon_\theta\!\left(\sqrt{\bar{\alpha}_t}\, x_0 + \sqrt{1 - \bar{\alpha}_t}\,\varepsilon,\; t\right)}_2^2\right].
  \label{eq:ddpm-lsimple}
\end{equation}
训练时只需随机抽一个数据 $x_0$、随机抽一个步号 $t$、随机撒一撮白噪声 $\varepsilon$，让网络看着被污染的 $x_t$ 猜出刚才撒进去的 $\varepsilon$ 即可！推理采样时则从 $x_T \sim \mathcal{N}(0, I)$ 起步，沿 $t = T, T-1, \dots, 1$ 逐轮减去预测噪声并补入微小扰动 $z \sim \mathcal{N}(0, I)$（末步 $t=1$ 不加噪）：
\begin{equation}
  x_{t-1} = \frac{1}{\sqrt{\alpha_t}}\!\left(x_t - \frac{1 - \alpha_t}{\sqrt{1 - \bar{\alpha}_t}}\,\varepsilon_\theta(x_t, t)\right) + \sigma_t z \qquad (t > 1).
  \label{eq:ddpm-sample-step}
\end{equation}

更令人惊叹的是，DDPM 与另一条截然不同的物理几何脉络——Song 与 Ermon（2019, 2020）的\textbf{基于分数的生成模型（Score-Based Generative Modeling）}——在数学深处完全合流！

\begin{thmBox}
\begin{proposition}[噪声预测与去噪分数匹配（Denoising Score Matching）的严格等价]
\label{prop:dsm}
给定条件加噪核 $q(x_t \mid x_0) = \mathcal{N}\!\big(\sqrt{\bar{\alpha}_t}\,x_0,\, (1 - \bar{\alpha}_t) I\big)$ 及关系 $x_t = \sqrt{\bar{\alpha}_t}\,x_0 + \sqrt{1 - \bar{\alpha}_t}\,\varepsilon$，其关于加噪样本 $x_t$ 的\textbf{条件对数密度梯度（Stein 分数函数）}精确满足：
\begin{equation}
  \nabla_{x_t} \ln q(x_t \mid x_0)
  = -\frac{x_t - \sqrt{\bar{\alpha}_t}\, x_0}{1 - \bar{\alpha}_t}
  = -\frac{\varepsilon}{\sqrt{1 - \bar{\alpha}_t}}.
  \label{eq:score-eps-link}
\end{equation}
并且，在全部可能的函数 $s(x_t, t)$ 中，最小化去噪分数匹配目标 $\E_{x_0, x_t}\big[\norm{s(x_t, t) - \nabla_{x_t}\ln q(x_t \mid x_0)}_2^2\big]$ 的唯一理论最优解恰好是\textbf{边际数据分布的真实分数函数} $s^*(x_t, t) = \nabla_{x_t} \ln q_t(x_t)$！
\end{proposition}
\begin{proof}
对 $\ln q(x_t \mid x_0) = -\frac{D}{2}\ln\big(2\pi(1-\bar{\alpha}_t)\big) - \frac{\norm{x_t - \sqrt{\bar{\alpha}_t}\, x_0}_2^2}{2(1-\bar{\alpha}_t)}$ 关于 $x_t$ 直接求梯度，立得式~\eqref{eq:score-eps-link}：
\[
  \nabla_{x_t}\ln q(x_t \mid x_0)
  = -\frac{x_t - \sqrt{\bar{\alpha}_t}\, x_0}{1 - \bar{\alpha}_t}
  = -\frac{\sqrt{1 - \bar{\alpha}_t}\,\varepsilon}{1 - \bar{\alpha}_t}
  = -\frac{\varepsilon}{\sqrt{1 - \bar{\alpha}_t}}.
\]
再证最优解为边际分数：由于 $L_2$ 风险的逐点极小值由条件期望给出（Vincent, 2011），对任意固定 $x_t$，最优预测器满足
\begin{align*}
  s^*(x_t, t)
  &= \E_{x_0 \sim q(x_0 \mid x_t)}\!\left[\nabla_{x_t} \ln q(x_t \mid x_0)\right]
   = \int \frac{q(x_0)\, q(x_t \mid x_0)}{q_t(x_t)} \cdot \frac{\nabla_{x_t} q(x_t \mid x_0)}{q(x_t \mid x_0)}\, dx_0 \\
  &= \frac{1}{q_t(x_t)} \nabla_{x_t} \int q(x_0)\, q(x_t \mid x_0)\, dx_0
   = \frac{\nabla_{x_t} q_t(x_t)}{q_t(x_t)}
   = \nabla_{x_t} \ln q_t(x_t).
\end{align*}
\end{proof}
\end{thmBox}

\begin{insightBox}{预测噪声 $\varepsilon$，本质上是在高维空间画出一张指向数据流形的引力场}
把式~\eqref{eq:score-eps-link} 代回逆向采样更新~\eqref{eq:ddpm-sample-step}，就会看到极其动人的物理图景：
\[
  -\frac{1 - \alpha_t}{\sqrt{1 - \bar{\alpha}_t}}\,\varepsilon_\theta(x_t, t)
  \;=\; \beta_t\, \underbrace{\left(-\frac{\varepsilon_\theta(x_t, t)}{\sqrt{1 - \bar{\alpha}_t}}\right)}_{s_\theta(x_t, t) \approx \nabla_{x_t} \ln q_t(x_t)}!
\]
原来，每一步“减去预测噪声 $\varepsilon_\theta$”，在几何上就是\textbf{沿着加噪数据分布的对数概率梯度（分数向量场）$\nabla_{x_t} \ln q_t(x_t)$ 向上爬坡}！而末尾补入的微小随机项 $\sigma_t z$ 则扮演退火朗之万动力学（Annealed Langevin Dynamics）中的热扰动，帮助粒子跃出低概率浅坑、在多模态数据峰群间按真实概率权重均匀落脚！在连续时间极限 $T \to \infty$ 下，Song 等人（2021）证明整个 DDPM 正是一条由 Anderson（1982）逆时随机微分方程（Reverse-Time SDE）刻画的扩散物理过程。
\end{insightBox}

\section{实战验证：30 行纯 PyTorch 二维环形分布 DDPM 实测}
\label{sec:ddpm-exp}

为了亲眼见证纯高斯白噪声如何在神经网络的几十次微步牵引下坍缩到复杂的多模态几何流形，我们在 \texttt{code/vol5\_experiments.py} 的 \texttt{exp6\_ddpm\_2d\_ring} 中，用区区 30 行核心 PyTorch 代码手写了一个二维 DDPM（$T = 40$，线性噪声表 $\beta_t \in [10^{-3}, 0.08]$）。目标数据分布设定为半径 $r = 2.0$ 上的\textbf{四模态高斯环}：四个簇心分别位于 $(\pm 2, 0)$ 与 $(0, \pm 2)$，各簇标准差 $\sigma = 0.18$。去噪网络 $\varepsilon_\theta(x_t, t)$ 仅由一个 3 层 \texttt{SiLU} 多层感知机（输入拼接 $(x_t, t/T) \in \R^3$，隐藏宽度 $64$）构成，完全不依赖任何现成扩散库。

\begin{lstlisting}[caption={纯 PyTorch 二维四模态环 DDPM 核心训练与逆向采样代码（与 \texttt{code/vol5\_experiments.py} 完全一致）},label={lst:ddpm-2d}]
import math, torch
import torch.nn as nn
import torch.nn.functional as F

T = 40
betas = torch.linspace(1e-3, 0.08, T)
alphas = 1.0 - betas
abar = torch.cumprod(alphas, dim=0)

centers = torch.tensor([[2.0, 0.0], [0.0, 2.0], [-2.0, 0.0], [0.0, -2.0]])
def sample_data(n):
    idx = torch.randint(0, 4, (n,))
    return centers[idx] + 0.18 * torch.randn(n, 2)

eps_net = nn.Sequential(
    nn.Linear(3, 64), nn.SiLU(),
    nn.Linear(64, 64), nn.SiLU(),
    nn.Linear(64, 2)
)
opt = torch.optim.AdamW(eps_net.parameters(), lr=4e-3)

# ---- 1. DDPM Training Step (Eq. ddpm-lsimple) ----
for step in range(601):
    x0 = sample_data(256)
    t = torch.randint(0, T, (256,))
    eps = torch.randn_like(x0)
    ab = abar[t].unsqueeze(-1)
    xt = torch.sqrt(ab) * x0 + torch.sqrt(1.0 - ab) * eps
    t_in = ((t.float() + 1.0) / T).unsqueeze(-1)
    loss = F.mse_loss(eps_net(torch.cat([xt, t_in], dim=-1)), eps)
    opt.zero_grad(); loss.backward(); opt.step()

# ---- 2. DDPM Reverse Sampling Loop (Eq. ddpm-sample-step) ----
@torch.no_grad()
def sample_ddpm(n=400):
    x = torch.randn(n, 2)
    for t in reversed(range(T)):
        t_in = torch.full((n, 1), (t + 1) / T)
        eps_pred = eps_net(torch.cat([x, t_in], dim=-1))
        c1 = 1.0 / math.sqrt(alphas[t].item())
        c2 = (1.0 - alphas[t].item()) / math.sqrt(1.0 - abar[t].item())
        mean = c1 * (x - c2 * eps_pred)
        x = mean + math.sqrt(betas[t].item()) * torch.randn_like(x) if t > 0 else mean
    return x
\end{lstlisting}

为了量化生成样本的质量，我们在训练过程中每隔 $60$ 步从纯标准正态噪声独立采样 $400$ 个点，计算它们到四个真实簇心最近距离的均值（\texttt{gen\_err}），并与真实数据自身的最近簇心距离下限（\texttt{real\_err}）直接对照。值得顺手做一次精妙的解析验算：对于二维各向同性高斯簇 $c + \sigma \varepsilon$（$\sigma = 0.18$），由于四簇间距 $2\sqrt{2} \approx 2.83 \gg \sigma$，样本到最近簇心的欧氏距离服从缩放瑞利分布（Rayleigh Distribution）$\sigma\sqrt{\varepsilon_1^2 + \varepsilon_2^2}$，其理论期望精确为：
\[
  \E[\text{dist}_{\text{real}}] = \sigma \sqrt{\frac{\pi}{2}} \approx 0.18 \times 1.2533 = 0.2256,
\]
与脚本实测的 $400$ 个真实样本均值 \texttt{real\_err = 0.2232} 吻合误差不足 $1\%$！

实测数据（完整保存在 \texttt{data/ddpm\_train.csv}）呈现了清晰的收敛进程（见图~\ref{fig:ddpm-train}）：
\begin{itemize}
  \item \textbf{\texttt{step = 0}（未训练初始化）}：预测噪声几乎为零，每维噪声预测均方误差 \texttt{eps\_mse = 0.9883}（紧贴标准正态方差理论值 $1.0$）；此时逆向采样的 400 个点杂乱散落，平均距最近簇心高达 \texttt{gen\_err = 2.522}；
  \item \textbf{\texttt{step = 180}}：模型迅速刻画出四个簇的全局引力盆地，\texttt{eps\_mse} 降至 \texttt{0.5192}，生成样本距离骤降 $80\%$ 至 \texttt{0.506}；
  \item \textbf{\texttt{step = 360}}：分布形态进一步收紧，\texttt{eps\_mse} 降至 \texttt{0.4360}，\texttt{gen\_err = 0.372}；
  \item \textbf{\texttt{step = 600}}：生成样本已稳稳汇聚于四座高斯孤岛上，\texttt{gen\_err = 0.297}（第 480 步已达 \texttt{0.292}），逼近真实分布的瑞利物理下限 \texttt{0.223}！
\end{itemize}

\begin{figure}[htbp]
\centering
\begin{tikzpicture}
  % Left panel: Generation distance to nearest true mode
  \begin{axis}[
    width=0.47\linewidth,
    height=5.6cm,
    xlabel={训练步数（Step）},
    ylabel={到最近真实模态的平均距离},
    xmin=-12, xmax=612,
    ymin=0, ymax=2.8,
    grid=major,
    grid style={dashed, gray!25},
    legend pos=north east,
    legend style={font=\scriptsize, fill=white, fill opacity=0.9, draw=tealLine},
    tick label style={font=\scriptsize},
    label style={font=\footnotesize\sffamily},
    title={\small\sffamily\bfseries (a) 生成样本收敛至四模态真实流形},
    title style={color=tealDark}
  ]
    \addplot[teal, very thick, mark=*, mark size=1.7pt]
      table[col sep=comma, x=step, y=gen_err] {data/ddpm_train.csv};
    \addlegendentry{DDPM 生成距离 \texttt{gen\_err}}
    \addplot[ochre, thick, dashed]
      table[col sep=comma, x=step, y=real_err] {data/ddpm_train.csv};
    \addlegendentry{真实分布下限 \texttt{0.223}}
    \node[font=\scriptsize\sffamily, color=tealDark, anchor=west] at (axis cs:76, 1.73) {step=0: 2.522};
    \node[font=\scriptsize\sffamily, color=tealDark, anchor=south west] at (axis cs:185, 0.58) {step=180: 0.506};
    \node[font=\scriptsize\sffamily, color=tealDark, anchor=south east] at (axis cs:592, 0.36) {step=600: 0.297};
  \end{axis}

  % Right panel: Epsilon prediction MSE
  \begin{axis}[
    xshift=0.51\linewidth,
    width=0.47\linewidth,
    height=5.6cm,
    xlabel={训练步数（Step）},
    ylabel={噪声预测均方误差（$\varepsilon$-MSE）},
    xmin=-12, xmax=612,
    ymin=0.2, ymax=1.1,
    grid=major,
    grid style={dashed, gray!25},
    legend pos=north east,
    legend style={font=\scriptsize, fill=white, fill opacity=0.9, draw=navy!30},
    tick label style={font=\scriptsize},
    label style={font=\footnotesize\sffamily},
    title={\small\sffamily\bfseries (b) $\varepsilon$-预测 MSE 与不可约噪声底噪},
    title style={color=navy}
  ]
    \addplot[navy, very thick, mark=square*, mark size=1.5pt]
      table[col sep=comma, x=step, y=eps_mse] {data/ddpm_train.csv};
    \addlegendentry{噪声预测 \texttt{eps\_mse}}
    \draw[slate, densely dotted, thick] (axis cs:0, 0.42) -- (axis cs:600, 0.42);
    \node[font=\scriptsize\sffamily, color=slate, anchor=north east] at (axis cs:590, 0.31) {大 $t$ 不可约条件方差底噪};
    \node[font=\scriptsize\sffamily, color=navy, anchor=west] at (axis cs:22, 0.99) {step=0: 0.988};
  \end{axis}
\end{tikzpicture}
\caption{纯 PyTorch 二维四模态高斯环 DDPM 实测曲线（直接读取自 \texttt{data/ddpm\_train.csv}，$T = 40$，半径 $r=2.0$，$\sigma=0.18$）。(a) 仅需 600 步微调，400 个纯噪声采样点经逆向去噪后的最近簇心距离从 \texttt{2.522} 锐减至 \texttt{0.297}，逼近真实分布下限 \texttt{0.223}；(b) $\varepsilon$-MSE 从 \texttt{0.988} 迅速跌落并稳定在 $0.36\sim 0.48$ 区间，呈现扩散模型特有的“高噪声时间步不可约底噪”现象。}
\label{fig:ddpm-train}
\end{figure}

\begin{remark}[为什么图~\ref{fig:ddpm-train}(b) 中的 $\varepsilon$-MSE 永远不会降到 $0$？]
初学扩散模型的读者看到训练损失停在 $0.40$ 左右且小幅波动，常误以为模型“欠拟合”或“没训好”。其实恰恰相反，这是\textbf{Bayes 最优回归的必然特征}！当随机抽到大时间步（如 $t = 38$）时，原始信号 $\sqrt{\bar{\alpha}_t}\,x_0$ 已被严重淹没，同一观测点 $x_t$ 可能以相近后验概率来自四个不同模态 $x_0$，对应四种截然不同的真实噪声 $\varepsilon$。哪怕拥有上帝视角的最优预测器 $\varepsilon^*(x_t, t) = \E[\varepsilon \mid x_t]$，其最小可能均方误差仍严格等于条件方差 $\E[\Var(\varepsilon \mid x_t)] > 0$（且每个小批随机抽取的 $t$ 分布略有起伏，造成单批损失自然抖动）。因此监控扩散模型是否收敛，绝不能只盯单步噪声 MSE，而要定期运行逆向采样检查真实生成分布指标（如本文的 \texttt{gen\_err} 或图像领域的 FID）！
\end{remark}

\section{流匹配（Flow Matching / Rectified Flow）：拉直生成轨迹}
\label{sec:flow-matching}

DDPM 彻底解决了 VAE 的图像模糊与 GAN 的训练崩溃，但它留下了一个昂贵的工程遗产：\textbf{采样太慢}。因为前向加噪 $\sqrt{\bar{\alpha}_t} x_0 + \sqrt{1-\bar{\alpha}_t}\varepsilon$ 与逆向马尔可夫链构成的是高曲率的随机折线轨迹（即使换成对应的确定性概率流 ODE 依然显著弯曲），为了压制数值离散化误差，经典 DDPM 往往需要串行调用超大神经网络 $50 \sim 1000$ 次才能画出一张图！

能否抛开复杂的马尔可夫方差表 $\beta_t, \alpha_t, \bar{\alpha}_t$，直接回答一个返璞归真的几何问题：\textbf{如果强行规定每个噪声点 $x_0 \sim \mathcal{N}(0, I)$ 必须沿着两点间最短的直线奔向真实数据点 $x_1 \sim p_{\text{data}}$，神经网络只需学会预测这条直线上的速度场，行不行？}

这正是 Lipman 等人（2023）的\textbf{流匹配（Flow Matching）}与 Liu 等人（2023）的\textbf{校正流（Rectified Flow）}所给出的惊艳答案——它不仅在数学上比 DDPM 简洁十倍，更直接成为了 Stable Diffusion 3、Flux.1 与 Meta Movie Gen 等新一代旗舰多模态生成模型的统一引擎！

\subsection{连续时间流与时间坐标约定}

\begin{remark}[重要记号约定：DDPM 与流匹配的时间方向互反！]
在阅读两类文献时务必留意主流时间方向约定：
\begin{itemize}
  \item \textbf{DDPM 习惯}：$t = 0$ 是真实数据，$t = T$ 是纯高斯噪声（前向加噪、逆向生成）；
  \item \textbf{流匹配（Flow Matching）习惯}：$t = 0$ 是纯高斯先验 $x_0 \sim p_0 = \mathcal{N}(0, I)$，$t = 1$ 是真实数据 $x_1 \sim p_1 = p_{\text{data}}$，生成过程就是沿着常微分方程从 $t = 0$ 正向积分至 $t = 1$！
\end{itemize}
本节严格遵循流匹配的标准记号（$t \in [0, 1]$：$x_0$ 为噪声，$x_1$ 为数据）。
\end{remark}

我们希望训练一个由神经网络参数化的时变速度场 $v_\theta(x, t): \R^D \times [0, 1] \to \R^D$，使得粒子从 $x_0 \sim p_0$ 出发，按常微分方程（ODE）：
\begin{equation}
  \frac{dx_t}{dt} = v_\theta(x_t, t), \qquad t \in [0, 1]
  \label{eq:fm-ode}
\end{equation}
演化至终时刻 $t = 1$ 时，粒子群的边缘分布 $p_1$ 恰好吻合真实数据分布 $p_{\text{data}}$。由流体力学中的质量守恒原理，速度场 $u_t(x)$ 驱动的密度演化 $p_t(x)$ 必须满足\textbf{连续性方程（Continuity Equation）}：
\begin{equation}
  \frac{\partial p_t(x)}{\partial t} + \nabla_x \cdot \!\big(p_t(x)\, u_t(x)\big) = 0.
  \label{eq:continuity}
\end{equation}
可是我们根本不知道能把 $p_0$ 搬运到 $p_1$ 的全局边际速度场 $u_t(x)$ 长什么样！流匹配最震撼的天才构想在于：\textbf{先给定一对具体的起点 $x_0$ 与终点 $x_1$，构造一条显式的直线条件路径，再把无数对直线路径的速度场“叠合”起来！}

\subsection{直线插值路径与条件流匹配（CFM）无偏等价定理}

任取独立配对的噪声 $x_0 \sim p_0 = \mathcal{N}(0, I)$ 与真实样本 $x_1 \sim p_1 = p_{\text{data}}$，在它们之间连结一条最简单的\textbf{匀速直线插值路径}：
\begin{equation}
  x_t = (1 - t)\, x_0 + t\, x_1, \qquad t \in [0, 1].
  \label{eq:rf-path}
\end{equation}
在这条直线上，粒子的瞬时条件速度是一个极其清爽的常向量：
\begin{equation}
  u_t(x_t \mid x_0, x_1) = \frac{dx_t}{dt} = x_1 - x_0.
  \label{eq:rf-cond-vel}
\end{equation}
现在直接用神经网络 $v_\theta(x_t, t)$ 对该条件速度做最朴素的均方误差回归，定义\textbf{条件流匹配 / 校正流损失（CFM / Rectified Flow Loss）}：
\begin{equation}
  \mathcal{L}_{\mathrm{CFM}}(\theta)
  = \E_{t \sim \mathcal{U}[0,1],\; x_0 \sim \mathcal{N}(0, I),\; x_1 \sim p_{\text{data}}}
    \!\left[\norm{v_\theta\!\big((1 - t)\, x_0 + t\, x_1,\; t\big) - (x_1 - x_0)}_2^2\right].
  \label{eq:cfm-loss}
\end{equation}
初学者一定会产生一个直击灵魂的疑问：不同样本对 $(x_0, x_1)$ 画出的直线会在中间时刻 $t$ 彼此交叉相撞！同一个点 $x_t$ 处可能对应着截然不同的 $(x_1 - x_0)$ 速度箭头，神经网络去拟合这些互相打架的目标，最终学出来的 ODE 真的能精确复现边际数据分布 $p_1$ 吗？下面这条优美的定理给出斩钉截铁的肯定回答！

\begin{thmBox}
\begin{theorem}[条件流匹配与边际流匹配的梯度等价性及分布守恒定理]
\label{thm:cfm}
在直线路径 $x_t = (1-t)x_0 + t x_1$ 下，记 $p_t(x)$ 为时刻 $t$ 的真实插值边缘密度，并定义该点处所有过路直线的条件平均速度场（边际速度场）：
\begin{equation}
  u_t(x) := \E\!\left[x_1 - x_0 \;\middle|\; x_t = x\right].
  \label{eq:marginal-vel}
\end{equation}
那么以下两条核心性质同时严格成立：
\begin{enumerate}
  \item \textbf{精确分布传输（连续性方程守恒）}：由边际速度场 $u_t(x)$ 驱动的确定性 ODE $\frac{dx}{dt} = u_t(x)$ 在所有时刻 $t \in [0, 1]$ 严格满足连续性方程~\eqref{eq:continuity}，从而将 $p_0 = \mathcal{N}(0, I)$ 精确搬运至 $p_1 = p_{\text{data}}$；
  \item \textbf{梯度零偏差等价}：定义不可直接计算的边际流匹配损失 $\mathcal{L}_{\mathrm{FM}}(\theta) := \E_{t,\, x_t \sim p_t}\!\left[\norm{v_\theta(x_t, t) - u_t(x_t)}_2^2\right]$，则 $\mathcal{L}_{\mathrm{CFM}}(\theta)$ 与 $\mathcal{L}_{\mathrm{FM}}(\theta)$ 仅差一个与参数 $\theta$ 无关的常数 $C$，故两者拥有\textbf{完全相同的精确梯度与全局极小解}：
  \begin{equation}
    \mathcal{L}_{\mathrm{CFM}}(\theta) = \mathcal{L}_{\mathrm{FM}}(\theta) + C
    \;\implies\;
    \nabla_\theta \mathcal{L}_{\mathrm{CFM}}(\theta) = \nabla_\theta \mathcal{L}_{\mathrm{FM}}(\theta),
    \quad
    v^*(x, t) = u_t(x).
    \label{eq:cfm-fm-equiv}
  \end{equation}
\end{enumerate}
\end{theorem}
\begin{proof}
\textbf{第一部分（弱形式验证连续性方程）。} 任取光滑紧支测试函数 $\varphi \in C_c^\infty(\R^D)$。沿随机直线路径 $x_t = (1-t)x_0 + t x_1$，按链式法则对期望 $\E[\varphi(x_t)]$ 关于时间 $t$ 求导：
\[
  \frac{d}{dt}\int_{\R^D} \varphi(x)\, p_t(x)\, dx
  = \frac{d}{dt}\E_{x_0, x_1}\!\left[\varphi(x_t)\right]
  = \E_{x_0, x_1}\!\left[\inner{\nabla \varphi(x_t)}{x_1 - x_0}\right].
\]
对右端应用全期望公式（条件期望的塔性质），先给定 $x_t$ 对 $(x_0, x_1)$ 求内层期望：
\begin{align*}
  \E_{x_0, x_1}\!\left[\inner{\nabla \varphi(x_t)}{x_1 - x_0}\right]
  &= \E_{x_t}\!\left[\inner{\nabla \varphi(x_t)}{\E\!\left[x_1 - x_0 \mid x_t\right]}\right] \\
  &= \int_{\R^D} \inner{\nabla \varphi(x)}{u_t(x)}\, p_t(x)\, dx
   = -\int_{\R^D} \varphi(x)\,\nabla_x \cdot\!\big(p_t(x)\, u_t(x)\big)\, dx,
\end{align*}
最后一步使用了高维分部积分（散度定理）以及 $\varphi$ 紧支集在无穷远处为零的性质。由于上式对一切测试函数 $\varphi$ 恒成立，故 $\partial_t p_t + \nabla_x \cdot (p_t u_t) = 0$ 在分布意义下严格成立！

\textbf{第二部分（展开平方并消去交叉项）。} 将 $\mathcal{L}_{\mathrm{CFM}}(\theta)$ 与 $\mathcal{L}_{\mathrm{FM}}(\theta)$ 展开：
\begin{align*}
  \mathcal{L}_{\mathrm{CFM}}(\theta)
  &= \E_{t, x_0, x_1}\!\left[\norm{v_\theta(x_t, t)}_2^2 - 2\inner{v_\theta(x_t, t)}{x_1 - x_0} + \norm{x_1 - x_0}_2^2\right], \\
  \mathcal{L}_{\mathrm{FM}}(\theta)
  &= \E_{t, x_t}\!\left[\norm{v_\theta(x_t, t)}_2^2 - 2\inner{v_\theta(x_t, t)}{u_t(x_t)} + \norm{u_t(x_t)}_2^2\right].
\end{align*}
两式首项 $\E[\norm{v_\theta(x_t, t)}_2^2]$ 因 $(t, x_t)$ 边缘分布完全一致而严格相等。对于中间交叉项，再次施用全期望公式：
\[
  \E_{x_0, x_1}\!\left[\inner{v_\theta(x_t, t)}{x_1 - x_0}\right]
  = \E_{x_t}\!\left[\inner{v_\theta(x_t, t)}{\E\!\left[x_1 - x_0 \mid x_t\right]}\right]
  = \E_{x_t}\!\left[\inner{v_\theta(x_t, t)}{u_t(x_t)}\right].
\]
于是两损失之差严格等于
\[
  C := \mathcal{L}_{\mathrm{CFM}}(\theta) - \mathcal{L}_{\mathrm{FM}}(\theta)
  = \E_{t, x_0, x_1}\!\left[\norm{x_1 - x_0}_2^2\right] - \E_{t, x_t}\!\left[\norm{u_t(x_t)}_2^2\right],
\]
它完全不含参数 $\theta$（恰好等于条件方差的期望）！两端对 $\theta$ 求梯度立得 $\nabla_\theta \mathcal{L}_{\mathrm{CFM}}(\theta) = \nabla_\theta \mathcal{L}_{\mathrm{FM}}(\theta)$，故最小化极其易算的 $\mathcal{L}_{\mathrm{CFM}}(\theta)$ 严格等价于逼近真实的边际速度场 $u_t(x)$！
\end{proof}
\end{thmBox}

定理~\ref{thm:cfm} 揭示了一个美妙的物理现象——\textbf{轨迹自动“解交叉”（Uncrossing）}！虽然人工随机配对的直线 $(x_0, x_1)$ 在中途相交，但当确定性 ODE 沿着交点处的\emph{平均速度} $u_t(x)$ 行进时，两条相撞的轨迹会在交点处自然“拐弯交换目的地”，绝不真正穿越彼此，而且平均传输距离严格变短！因此流匹配生成的 ODE 轨迹天生远比 DDPM 笔直得多，只需最朴素的显式欧拉步进：
\[
  x_{t + \Delta t} = x_t + \Delta t \cdot v_\theta(x_t, t), \qquad \Delta t = \frac{1}{K},
\]
往往仅需 $K = 10 \sim 25$ 步（配合一次“再拉直” Reflow 甚至只需 $1 \sim 4$ 步！）就能输出惊艳的高清图像！

\subsection{三大深度生成范式的全景几何对比}

至此，VAE、DDPM 与 Flow Matching 三大范式的核心机理已全盘托出。图~\ref{fig:gen-compare} 从概率流形与轨迹几何的角度将三者并排呈现，表~\ref{tab:gen-comparison} 则进一步将它们与自回归模型（第 4--5 章）和生成对抗网络（GAN）放入同一坐标系下做系统对照。

\begin{figure}[htbp]
\centering
\resizebox{\linewidth}{!}{%
\begin{tikzpicture}[
  box/.style={draw, rounded corners=2mm, minimum height=0.82cm, align=center, font=\footnotesize\sffamily},
  arr/.style={-{Stealth[length=2.3mm]}, thick},
  lbl/.style={font=\scriptsize\sffamily, color=slate}
]
  % ==================== Row 1: VAE ====================
  \node[font=\small\sffamily\bfseries, color=tealDark, anchor=west] at (-0.2, 7.85) {(1) 变分自编码器（VAE）：单步摊销瓶颈与重参数化};
  \draw[tealLine, rounded corners=2mm, fill=tealLight!45] (-0.2, 5.43) rectangle (15.6, 7.48);

  \node[box, fill=white, draw=teal, minimum width=1.45cm] (vx) at (0.85, 6.42) {真实数据\\$x$};
  \node[box, fill=tealLight, draw=teal, minimum width=2.1cm] (venc) at (3.65, 6.42) {编码器 $q_\phi(z|x)$\\$(\mu_\phi,\, \sigma_\phi)$};
  \node[box, fill=ochreLight, draw=ochre, minimum width=2.7cm] (vz) at (7.45, 6.42) {重参数化采样\\$z = \mu_\phi + \sigma_\phi \had \varepsilon$};
  \node[box, fill=tealLight, draw=teal, minimum width=2.0cm] (vdec) at (11.15, 6.42) {解码器 $p_\theta(x|z)$\\一步映射};
  \node[box, fill=white, draw=teal, minimum width=1.55cm] (vxhat) at (14.4, 6.42) {重建样本\\$\hat{x}$};

  \draw[arr, teal] (vx) -- (venc);
  \draw[arr, teal] (venc) -- (vz);
  \draw[arr, teal] (vz) -- (vdec);
  \draw[arr, teal] (vdec) -- (vxhat);
  \node[lbl, anchor=south] at (7.45, 6.86) {$\varepsilon \sim \mathcal{N}(0, I)$ 外部注入，反向梯度可穿透};
  \node[lbl, anchor=north, color=ochre] at (7.45, 5.98) {$\DKL(q_\phi \,\|\, \mathcal{N}(0,I))$ 拉齐先验 $\cdot$ 推理仅需 1 步（NFE=1）};

  % ==================== Row 2: DDPM ====================
  \node[font=\small\sffamily\bfseries, color=navy, anchor=west] at (-0.2, 5.05) {(2) 扩散模型（DDPM）：多步随机马尔可夫加噪与逆向去噪};
  \draw[navy!35, rounded corners=2mm, fill=navyLight!55] (-0.2, 2.25) rectangle (15.6, 4.75);

  \node[box, fill=white, draw=navy, minimum width=1.35cm] (dx0) at (0.85, 3.50) {数据\\$x_0$};
  \node[box, fill=white, draw=navy!70, minimum width=1.35cm] (dx1) at (3.55, 3.50) {微噪\\$x_1$};
  \node[font=\bfseries, color=navy] (ddots) at (5.65, 3.50) {$\cdots$};
  \node[box, fill=white, draw=navy!70, minimum width=1.6cm] (dxt) at (7.95, 3.50) {中间态\\$x_t$};
  \node[font=\bfseries, color=navy] (ddots2) at (10.3, 3.50) {$\cdots$};
  \node[box, fill=navyLight, draw=navy, minimum width=2.25cm] (dxT) at (13.6, 3.50) {纯高斯噪声\\$x_T \sim \mathcal{N}(0, I)$};

  % Forward arrows (top, staying inside box)
  \draw[arr, slate, dashed] (dx0.north east) to[bend left=18] node[above=1pt, font=\scriptsize\sffamily, color=slate] {$q(x_1|x_0)$} (dx1.north west);
  \draw[arr, slate, dashed] (dx1.north east) to[bend left=12] node[above=1pt, font=\scriptsize\sffamily, color=slate] {固定加噪 $\beta_t$} (dxt.north west);
  \draw[arr, slate, dashed] (dxt.north east) to[bend left=12] node[above=1pt, font=\scriptsize\sffamily, color=slate] {闭式直达 $\sqrt{\bar{\alpha}_t}x_0 + \sqrt{1-\bar{\alpha}_t}\varepsilon$} (dxT.north west);

  % Reverse arrows (bottom, staying inside box)
  \draw[arr, ochre] (dxT.south west) to[bend left=12] node[below=1pt, font=\scriptsize\sffamily, color=ochre] {逆向去噪 $p_\theta(x_{t-1}|x_t)$（预测 $\varepsilon_\theta$ / 分数爬坡）} (dxt.south east);
  \draw[arr, ochre] (dxt.south west) to[bend left=12] (dx1.south east);
  \draw[arr, ochre] (dx1.south west) to[bend left=18] node[below=1pt, font=\scriptsize\sffamily, color=ochre] {$50\!\sim\!1000$ 步} (dx0.south east);

  % ==================== Row 3: Flow Matching ====================
  \node[font=\small\sffamily\bfseries, color=plum, anchor=west] at (-0.2, 1.75) {(3) 流匹配（Flow Matching / Rectified Flow）：拉直确定性 ODE 传输轨迹};
  \draw[plum!35, rounded corners=2mm, fill=plumLight!60] (-0.2, -0.35) rectangle (15.6, 1.45);

  \node[box, fill=plumLight, draw=plum, minimum width=2.35cm] (fx0) at (1.35, 0.58) {纯高斯噪声\\$x_0 \sim \mathcal{N}(0, I)$};
  \node[box, fill=white, draw=plum, minimum width=3.15cm] (fxt) at (7.70, 0.58) {直线插值态\\$x_t = (1-t)x_0 + t x_1$};
  \node[box, fill=white, draw=plum, minimum width=2.25cm] (fx1) at (14.1, 0.58) {目标数据\\$x_1 \sim p_{\text{data}}$};

  \draw[arr, plum, very thick] (fx0) -- node[above=2pt, font=\scriptsize\sffamily, color=plum, align=center] {回归速度场\\$v_\theta(x_t, t) \approx x_1 - x_0$} (fxt);
  \draw[arr, plum, very thick] (fxt) -- node[above=2pt, font=\scriptsize\sffamily, color=plum, align=center] {确定性 ODE 欧拉步进\\（仅需 $10\!\sim\!25$ 步）} (fx1);
  \node[lbl, anchor=north, color=plum] at (7.70, 0.10) {轨迹平直无冗余随机抖动 $\cdot$ 与现代 DiT 架构无缝契合};
\end{tikzpicture}}
\caption{VAE、DDPM 与流匹配（Flow Matching）三大深度生成范式的结构与传输路径几何对照。从“单步压缩瓶颈”（VAE），经由“多步弯曲随机去噪”（DDPM），演化至“直线常微分方程速度场回归”（Flow Matching），深度生成模型逐步实现了训练稳定性、样本保真度与推理速度的全面突破。}
\label{fig:gen-compare}
\end{figure}

\begin{table}[htbp]
\centering
\caption{现代深度生成模型五大范式的系统对比}
\label{tab:gen-comparison}
{\footnotesize
\begin{tabularx}{\linewidth}{@{}p{21mm}p{29mm}p{17mm}p{18mm}p{20mm}X@{}}
\toprule
\textbf{生成范式} & \textbf{核心训练目标} & \textbf{推理步数 (NFE)} & \textbf{训练稳定性} & \textbf{样本保真与多样性} & \textbf{代表系统与现代角色} \\
\midrule
\textbf{VAE} & ELBO（重建 $+$ 高斯 KL） & 极快（1 步） & 极高 & 多样但易模糊 & 现代潜扩散模型（SD / Sora）的\textbf{潜空间压缩底座} \\
\textbf{GAN} & 极小极大对抗博弈 & 极快（1 步） & 低（易模式坍缩） & 极锐利但覆盖不全 & StyleGAN、VAE 感知/对抗微调损失 \\
\textbf{自回归 (AR)} & 精确下一词交叉熵 & 逐词元串行（$T$ 步） & 极高 & 离散模态之王 & GPT-4、Llama-3、Qwen、DeepSeek \\
\textbf{DDPM} & 噪声 $\varepsilon$-MSE / 分数匹配 & 较慢（$50\sim 1000$） & 极高 & 极高保真且全覆盖 & DDPM、Stable Diffusion 1.5/XL、AlphaFold 3 \\
\textbf{流匹配 (FM)} & 直线速度场 CFM MSE & 较快（$4\sim 25$） & 极高 & 顶级保真与平滑可控 & Stable Diffusion 3、Flux.1、Movie Gen \\
\bottomrule
\end{tabularx}}
\end{table}

\begin{extBox}{现代图像与视频生成的三大加速器：CFG、DDIM 与 DiT}
将本卷的理论推向工业级文生图（Stable Diffusion 3 / Flux）与文生视频（Sora）只需补上三块至关重要的积木：
\begin{enumerate}
  \item \textbf{无分类器引导（Classifier-Free Guidance, CFG；Ho \& Salimans, 2022）}：如何让扩散或流模型严丝合缝地听从文本提示词 $c$（如“一只戴草帽的柴犬”）？在训练时，以 $10\%\sim 20\%$ 的概率随机把文本条件 $c$ 置空为 $\varnothing$，令同一网络同时学会有条件预测 $\varepsilon_\theta(x_t, c)$ 与无条件预测 $\varepsilon_\theta(x_t, \varnothing)$。由贝叶斯公式 $\ln p(c \mid x_t) = \ln p(x_t \mid c) - \ln p(x_t) + \text{const}$，两侧取梯度并代入式~\eqref{eq:score-eps-link} 可知，隐式分类器梯度恰好正比于 $-\big(\varepsilon_\theta(x_t, c) - \varepsilon_\theta(x_t, \varnothing)\big)$！推理时引入引导强度 $w > 1$（常见 $w \in [5, 8]$），构造外推预测：
  \[
    \tilde{\varepsilon}_\theta(x_t, c)
    = \varepsilon_\theta(x_t, \varnothing) + w\,\big(\varepsilon_\theta(x_t, c) - \varepsilon_\theta(x_t, \varnothing)\big),
  \]
  在数学上等价于从锐化后的倾斜分布 $p(x_t)\,p(c \mid x_t)^w$ 中采样——牺牲少量边缘多样性，换取图文对齐度与构图锐度的质的飞跃！
  \item \textbf{去噪扩散隐式模型（DDIM；Song et al., 2021）}：注意到 DDPM 的简化目标~\eqref{eq:ddpm-lsimple} \textbf{只依赖前向边缘分布 $q(x_t \mid x_0)$}，与前向过程是否马尔可夫毫无关系！Song 等人构造了一族共享完全相同边缘分布的非马尔可夫过程，令随机噪声系数 $\eta = 0$ 时逆向过程化为完全确定性的 ODE 离散化，无需重新训练任何权重，即可直接跨步跳跃采样（例如从 $1000$ 步缩减为 $20\sim 50$ 步），速度立提 $20\sim 50$ 倍！
  \item \textbf{潜空间扩散与扩散 Transformer（LDM \& DiT；Rombach et al., 2022; Peebles \& Xie, 2023）}：在 $1024 \times 1024 \times 3$ 原始像素上跑扩散计算量过于恐怖。现代系统采取了\textbf{“VAE $+$ 扩散/流匹配”的绝妙联姻}——先用预训练 VAE 将像素图 $8\times$ 下采样压缩进紧致连续潜空间 $z \in \R^{128 \times 128 \times 16}$，再在潜空间内运行扩散或流匹配，最后由 VAE 解码器一步还原高清像素！而主干去噪网络也由早期的卷积 U-Net 全面进化为 \textbf{DiT（Diffusion Transformer）}：把潜空间网格切块（Patchify）送入第 4 章的双向 Transformer，通过\textbf{自适应层归一化（adaLN-Zero）}把时间步 $t$ 与文本条件 $c$ 动态注入每一层的缩放平移参数 $(\gamma, \beta, \alpha)$ 中——至此，第 4 章的 Transformer 与第 6 章的扩散/流模型完美合璧，视觉与视频生成也正式迈入第 5 章的\textbf{神经缩放定律时代}！
\end{enumerate}
\end{extBox}

\section*{本章小结与第 5 卷收官脉络}
\addcontentsline{toc}{section}{本章小结与第 5 卷收官脉络}

本章将现代深度生成模型的三条核心数学主脉打通为一座结构严密的大厦：
\begin{itemize}
  \item \textbf{变分自编码器（VAE）}承接第 4 卷 EM 算法的 ELBO，借由摊销神经编码器、对角高斯闭式 KL 散度 $\frac{1}{2}\sum_j(\mu_j^2 + \sigma_j^2 - \ln\sigma_j^2 - 1)$ 与低方差重参数化技巧 $z = \mu_\phi + \sigma_\phi \had \varepsilon$，建立了端到端可微的连续隐变量生成框架，至今仍是潜空间扩散模型的压缩基石；
  \item \textbf{去噪扩散模型（DDPM）}用固定前向马尔可夫加噪 $q(x_t \mid x_0) = \mathcal{N}(\sqrt{\bar{\alpha}_t}x_0, (1-\bar{\alpha}_t)I)$ 拆解生成难题，借贝叶斯与配方法导出精确逆向后验均值 $\tilde{\mu}_t = \frac{1}{\sqrt{\alpha_t}}\big(x_t - \frac{1-\alpha_t}{\sqrt{1-\bar{\alpha}_t}}\varepsilon\big)$，最终沉淀为一行的 $\varepsilon$-MSE 目标，并与去噪分数匹配 $\nabla_{x_t}\ln q(x_t \mid x_0) = -\varepsilon / \sqrt{1-\bar{\alpha}_t}$ 严格等价；纯 PyTorch 2D 四模态实测印证了生成误差从 \texttt{2.522} 逼近真实瑞利下限 \texttt{0.223} 的全过程；
  \item \textbf{流匹配（Flow Matching / Rectified Flow）}进一步甩掉繁琐的噪声调度与弯曲随机路径，在噪声与数据间直连线段 $x_t = (1-t)x_0 + t x_1$，证明了条件速度场回归与边际连续性方程的精确等价性，以更少步数、更简代码驱动着当代视觉与视频生成前沿。
\end{itemize}

回顾整部\textbf{第 5 卷《深度学习、自动微分与现代架构》}，我们从第 1 章 60 行手写反向模式自动微分出发，历经第 2 章的初始化与优化动力学、第 3 章的卷积与残差高速公路、第 4 章的缩放注意力与迷你 GPT、第 5 章的自监督预训练与 Chinchilla 缩放律，直至本章的 VAE、扩散与流匹配——现代神经网络的所有核心零件已在你眼前彻底透明。然而，一个预训练好的“通识底座”尚不能自主对齐人类意图、在复杂环境中做长程序贯决策，也还需突破 GPU 显存与集群分布式通信的物理极限。这正是下一卷——\textbf{第 6 卷《强化学习、大模型后训练与 AI 系统》}——等待我们征服的崭新疆域！

\section*{本章分层习题}
\addcontentsline{toc}{section}{本章分层习题}

\exTier{A 层：概念与手算热身}{熟练掌握高斯 KL 手算、重参数化求导与 DDPM 单步噪声表演算}

\begin{enumerate}
  \item[\textbf{5.6.1}] \textbf{VAE 单样本损失与重参数化梯度完整手算}：设隐空间维度 $d = 1$，某输入 $x = 2.0$ 经编码器输出均值 $\mu = 1.0$ 与对数方差 $s := \ln \sigma^2 = 0.0$（即 $\sigma = 1.0$）。设本次采样抽到外部噪声 $\varepsilon = -0.5$，解码器简化为恒等映射 $\hat{x} = z$，重建损失取 $\mathcal{L}_{\text{rec}} = \frac{1}{2}(x - \hat{x})^2$。
  \begin{enumerate}
    \item 计算重参数化隐码 $z$、重建损失 $\mathcal{L}_{\text{rec}}$、式~\eqref{eq:gaussian-kl-closed} 给出的 KL 散度以及总损失 $\mathcal{L}_{\text{tot}} = \mathcal{L}_{\text{rec}} + \DKL$；
    \item 手工计算总损失对编码器输出参数 $\mu$ 与 $s = \ln \sigma^2$ 的精确偏导数 $\frac{\partial \mathcal{L}_{\text{tot}}}{\partial \mu}$ 与 $\frac{\partial \mathcal{L}_{\text{tot}}}{\partial s}$。
  \end{enumerate}
  \exHint{利用 $\sigma = e^{s/2}$，故 $\frac{\partial z}{\partial s} = \frac{1}{2}\sigma\varepsilon$，而 $\frac{\partial \DKL}{\partial s} = \frac{1}{2}(e^s - 1)$。}

  \item[\textbf{5.6.2}] \textbf{微型 $T = 3$ 步 DDPM 系数表与逆向去噪手算}：给定 3 步线性噪声表 $\beta_1 = 0.1,\; \beta_2 = 0.2,\; \beta_3 = 0.3$。
  \begin{enumerate}
    \item 逐一计算 $\alpha_t$、累积乘积 $\bar{\alpha}_t$、每步信噪比 $\mathrm{SNR}(t) := \frac{\bar{\alpha}_t}{1 - \bar{\alpha}_t}$，以及式~\eqref{eq:ddpm-post-mean-x0} 给出的真实后验方差 $\tilde{\beta}_2$ 与 $\tilde{\beta}_3$（保留四位小数）；
    \item 设一维逆向采样行至 $t = 2$ 时状态为 $x_2 = 1.50$，神经网络预测噪声 $\varepsilon_\theta(x_2, 2) = 0.40$，计算去噪均值 $\mu_\theta(x_2, 2)$ 以及该点对应的估计分数 $s_\theta(x_2, 2)$。
  \end{enumerate}
  \exHint{按定义 $\bar{\alpha}_1 = 0.9,\; \bar{\alpha}_2 = 0.72,\; \bar{\alpha}_3 = 0.504$；验证 $\tilde{\beta}_t < \beta_t$ 恒成立。}
\end{enumerate}

\exTier{B 层：定理推导与严谨证明}{深入变分界收紧、多参数化等价性与扩散--流匹配数学同构}

\begin{enumerate}
  \item[\textbf{5.6.3}] \textbf{一般多元高斯 KL 散度公式与标准正态特例}：将定理~\ref{thm:gaussian-kl} 推广到全协方差情形：证明两个 $d$ 维非退化高斯分布 $q = \mathcal{N}(\mu_q, \Sigma_q)$ 与 $p = \mathcal{N}(\mu_p, \Sigma_p)$ 之间的精确 KL 散度为：
  \[
    \DKL(q \,\|\, p) = \frac{1}{2}\!\left[\tr(\Sigma_p^{-1}\Sigma_q) + (\mu_p - \mu_q)^\top \Sigma_p^{-1}(\mu_p - \mu_q) - d + \ln \frac{\det \Sigma_p}{\det \Sigma_q}\right],
  \]
  并验证当 $\mu_p = 0,\; \Sigma_p = I_d,\; \Sigma_q = \diag(\sigma_1^2, \dots, \sigma_d^2)$ 时，该式精确退化为式~\eqref{eq:gaussian-kl-closed}。
  \exHint{利用迹的循环不变性 $\E_q[(z - \mu_p)^\top \Sigma_p^{-1}(z - \mu_p)] = \tr\!\big(\Sigma_p^{-1}\E_q[(z-\mu_p)(z-\mu_p)^\top]\big)$！}

  \item[\textbf{5.6.4}] \textbf{DDPM 的三种等价参数化（$x_0 / \varepsilon / v$-预测）与流匹配的时间坐标同构}：
  \begin{enumerate}
    \item 在 DDPM 中，除了预测噪声 $\varepsilon$，还可令网络直接预测干净数据 $\hat{x}_0(x_t, t)$，或预测 Salimans 与 Ho（2022）提出的角速度目标 $v_t := \sqrt{\bar{\alpha}_t}\,\varepsilon - \sqrt{1 - \bar{\alpha}_t}\,x_0$。写出三者之间精确的双向代数换算公式，并说明为何在极高噪声步（$\bar{\alpha}_t \to 0$）时直接预测 $x_0$ 极不稳定而 $v$-预测始终保持单位方差；
    \item 对流匹配直线插值 $x_s = (1 - s)x_0 + s x_1$（$s \in (0, 1)$），令 $\tilde{x}_s := \frac{x_s}{\sqrt{(1-s)^2 + s^2}}$。证明经此尺度重归一化后，$\tilde{x}_s$ 与 DDPM 前向分布~\eqref{eq:ddpm-xt-reparam} 在满足 $\sqrt{\bar{\alpha}_t} = \frac{s}{\sqrt{(1-s)^2 + s^2}}$ 的时间重参数化下\textbf{完全同分布}！
  \end{enumerate}
  \exHint{对 $v_t$ 验证 $\Var(v_t) = \bar{\alpha}_t \Var(\varepsilon) + (1-\bar{\alpha}_t)\Var(x_0) = 1$（当 $\Var(x_0) = 1$），且 $x_0 = \sqrt{\bar{\alpha}_t}\,x_t - \sqrt{1-\bar{\alpha}_t}\,v_t$。}
\end{enumerate}

\exTier{C 层：代码实现与数值实验}{亲手编写二维流匹配并对比少步采样，复现并治愈 VAE 后验坍缩}

\begin{enumerate}
  \item[\textbf{5.6.5}] \textbf{同一网络上的正面对决：流匹配 vs DDPM 少步采样实验}：复制代码~\ref{lst:ddpm-2d} 的 4 模态高斯环与完全相同的 3 层 \texttt{SiLU} 网络（\texttt{Linear(3,64)-SiLU-Linear(64,64)-SiLU-Linear(64,2)}），仅改动训练循环与采样器，按式~\eqref{eq:cfm-loss} 训练 600 步\textbf{校正流（Rectified Flow）}模型。随后分别以采样步数 $K \in \{5,\; 10,\; 20,\; 40\}$ 运行确定性欧拉积分 $x_{t+\Delta t} = x_t + \Delta t \cdot v_\theta(x_t, t)$，记录各步数下的 \texttt{gen\_err} 并与 DDPM 对比。
  \exHint{流匹配在仅用 $K = 10$ 步欧拉积分时通常即可达到与 40 步 DDPM 相当甚至更低的 \texttt{gen\_err}！}

  \item[\textbf{5.6.6}] \textbf{可控复现并修复 VAE 后验坍缩}：在二维环形数据或 MNIST 子集上编写一个极简 PyTorch VAE（含 \texttt{logvar} 编码与重参数化），设计三组消融实验：
  \begin{enumerate}
    \item 固定 $\beta = 1.0$ 标准训练；
    \item 将 KL 权重刻意加大至 $\beta = 10.0$，观察所有维度的 $\DKL^{(j)}$ 是否迅速坍缩至 $< 0.01$ 且解码重建退化为数据集均图；
    \item 采用线性 KL 退火（前 $30\%$ 轮次将 $\beta$ 从 $0$ 线性增至 $1$），绘制总损失、重建误差与活跃隐维度数量随轮次的变化曲线。
  \end{enumerate}
  \exHint{监测“活跃隐维度”的标准指标是编码器均值在验证集上的方差 $\Var_{x \sim p_{\text{val}}}(\mu_{\phi, j}(x)) > 0.01$。}
\end{enumerate}

\exTier{D 层：开放探究与前沿迁移}{探索校正流迭代拉直（Reflow）与离散扩散大语言模型}

\begin{enumerate}
  \item[\textbf{5.6.7}] \textbf{Reflow（二次拉直）与一步生成（Consistency Models）}：尽管 1-Rectified Flow 显著拉直了轨迹，但由于初始 $(x_0, x_1)$ 独立配对时直线在空间中大量交叉，平均速度场 $u_t(x)$ 仍有残余弯曲。Liu 等人（2023）提出 \textbf{Reflow}：用训练好的 1-Rectified Flow 从噪声 $x_0$ 积分生成配对终点 $\hat{x}_1 = \mathrm{ODE}(x_0)$，再以新配对 $(x_0, \hat{x}_1)$ 重复执行式~\eqref{eq:cfm-loss} 训练 2-Rectified Flow。
  \begin{enumerate}
    \item 利用 Jensen 不等式与连续性方程，证明每次 Reflow 不仅严格保持边缘分布 $p_0, p_1$ 不变，还严格\textbf{单调不增}二次最优传输代价 $\E[\norm{\hat{x}_1 - x_0}_2^2] \le \E[\norm{x_1 - x_0}_2^2]$；
    \item 对比 Reflow 与 Song 等人（2023）的\textbf{一致性模型（Consistency Models）}（直接约束同一条 ODE 轨迹上任意两点 $x_t, x_{t'}$ 映射到同一终点 $x_1$），分析两者在实现 1--4 步极速生成时的核心权衡。
  \end{enumerate}
  \exHint{把 $\hat{x}_1 - x_0 = \int_0^1 u_t(x_t)\,dt$ 代入，对时间积分施用 Jensen 不等式 $\norm{\int_0^1 u_t\,dt}_2^2 \le \int_0^1 \norm{u_t}_2^2\,dt$，再利用条件期望收缩方差 $\E[\norm{u_t(x_t)}_2^2] \le \E[\norm{x_1 - x_0}_2^2]$！}

  \item[\textbf{5.6.8}] \textbf{离散扩散语言模型（MDLM / LLaDA）与统一多模态之争}：本章的高斯扩散与流匹配天然依赖连续欧氏空间中的加性噪声与可微分数 $\nabla_x \ln p_t(x)$，难以直接作用于离散文本词元。近年兴起的\textbf{掩码离散扩散（Masked Discrete Diffusion；Sahoo et al., 2024; Nie et al., 2025）}将前向高斯加噪替换为连续时间吸收态马尔可夫链：每个词元随时间 $t \in [0, 1]$ 以递增概率独立坍缩为 \texttt{[MASK]}，逆向过程则在任意掩码率下迭代并行去噪解码。结合第 5 章的 BERT（MLM）与第 4 章的自回归 GPT，剖析：
  \begin{enumerate}
    \item 掩码离散扩散的变分下界为何可化简为“对所有掩码比例 $t \in (0, 1]$ 加权积分的 BERT 掩码语言建模损失”？
    \item 在推理吞吐（并行解码多词元）、KV Cache 复用与长程逻辑一致性上，离散扩散大模型相比传统因果自回归 LLM 分别具备哪些优势与硬伤？
  \end{enumerate}
  \exHint{BERT 只在固定 $15\%$ 掩码率下训练，无法从 $100\%$ 全掩码白板出发生成；而离散扩散在 $(0, 1]$ 全区间训练，使其能够由全 \texttt{[MASK]} 逐步并行“显影”出整篇文本！}
\end{enumerate}


% ==================== Appendix: Complete Solutions ====================
\appendix
\ctexset{chapter={name={附录\ ,}, number=\Alph{chapter}}}
\renewcommand{\chapWatermarkNum}{\thechapter}
% !TeX root = ../main.tex
\renewcommand{\chapWatermarkNum}{\thechapter}
\ctexset{chapter={name={附录\ ,}, number=\Alph{chapter}}}
\renewcommand{\chaptermark}[1]{\markboth{附录 \thechapter\quad #1}{}}
\chapter[全卷 48 道分层习题完整精解]{全卷 48 道分层习题完整精解}
\label{app:solutions}
\chapEpigraph{分析机编织代数模式，一如提花机编织花朵与枝叶。}{Ada Lovelace（埃达·洛夫莱斯，1843）}

\begin{chapIntro}{本附录使用说明}
本附录收录第 1--6 章全部 48 道分层习题（编号 \textbf{5.1.1--5.6.8}）的完整、自洽精解。全部手算算术、张量 VJP 推导、不等式常数与实验基准均经独立 Python / PyTorch 数值脚本交叉核验：
\begin{itemize}
  \item \textbf{A 层（概念与手算热身）}：完整展开计算图、归一化矩阵、卷积感受野、注意力权重、BPE 合并与 VAE/DDPM 单步手算，同时给出精确表达式与小数结果；
  \item \textbf{B 层（定理推导与严谨证明）}：提供零跳步的矩阵微积分、方差传播、Softmax 半正定性、RoPE 衰减、缩放定律与多参数化等价证明；
  \item \textbf{C 层（代码实现与数值实验）}：给出核心实现架构、常见排错要点与固定随机种子下的实测基准数据；
  \item \textbf{D 层（开放探究与前沿迁移）}：严格证明隐函数微分、方差偏移、并行前缀扫描结合律与 Reflow 传输代价单调性，并系统梳理前沿架构权衡。
\end{itemize}
\end{chapIntro}

% =========================================================================
\section*{第 1 章习题解答（共 8 题）}
\addcontentsline{toc}{section}{第 1 章习题解答（共 8 题）}
% =========================================================================

\subsection*{习题 5.1.1\quad 三输入计算图手算}
引入中间节点 $u = a b$、$v = u + c$、$s = \sigma(v)$ 与输出 $f = s \cdot a$。在点 $(a, b, c) = (2.0, -0.5, 1.0)$ 处：
\begin{enumerate}[label=(\alph*), leftmargin=1.8em]
  \item \textbf{前向传播与计算 DAG}：依次求值：
  \[
    u = 2.0 \times (-0.5) = -1.0, \quad
    v = -1.0 + 1.0 = 0.0, \quad
    s = \sigma(0.0) = \frac{1}{1+e^0} = 0.5, \quad
    f = 0.5 \times 2.0 = 1.0.
  \]
  计算有向无环图（含节点 $a$ 的两条出边）如下：
  \begin{center}
  \begin{tikzpicture}[
    node distance=12mm and 16mm,
    var/.style={rectangle, rounded corners=2pt, draw=teal, fill=tealLight, font=\small, inner sep=4pt},
    op/.style={circle, draw=navy, fill=navyLight, font=\small, inner sep=2.5pt, minimum size=7mm},
    edge/.style={-{Stealth[length=2.2mm]}, thick, color=ink!80}
  ]
    \node[var] (a) {$a = 2.0$};
    \node[var, below=7mm of a] (b) {$b = -0.5$};
    \node[var, below=7mm of b] (c) {$c = 1.0$};
    \node[op, right=15mm of b, yshift=5mm] (mul1) {$\times$};
    \node[op, right=16mm of mul1, yshift=-4mm] (add1) {$+$};
    \node[op, right=16mm of add1] (sig) {$\sigma$};
    \node[op, right=16mm of sig, yshift=7mm] (mul2) {$\times$};
    \node[var, right=13mm of mul2] (out) {$f = 1.0$};

    \draw[edge] (a) -- (mul1);
    \draw[edge] (b) -- (mul1);
    \draw[edge] (mul1) -- node[above, font=\scriptsize] {$u=-1.0$} (add1);
    \draw[edge] (c) -- (add1);
    \draw[edge] (add1) -- node[above, font=\scriptsize] {$v=0.0$} (sig);
    \draw[edge] (sig) -- node[below right, font=\scriptsize] {$s=0.5$} (mul2);
    \draw[edge, color=ochre] (a.east) to[out=12, in=155] node[above, font=\scriptsize, color=ochre] {分支 1：$a=2.0$} (mul2.north west);
    \draw[edge] (mul2) -- (out);
  \end{tikzpicture}
  \end{center}

  \item \textbf{逆拓扑序手算全部伴随量}：初始化 $\bar{f} = 1$，按 $f \to s \to v \to (u, c) \to (a, b)$ 反向累加：
  \begin{align*}
    \bar{s} &= \bar{f} \cdot \frac{\partial f}{\partial s} = 1 \cdot a = 2.0, \\
    \bar{v} &= \bar{s} \cdot \sigma'(v) = \bar{s} \cdot s(1-s) = 2.0 \times 0.5 \times 0.5 = 0.5, \\
    \bar{u} &= \bar{v} \cdot 1 = 0.5, \qquad \bar{c} = \bar{v} \cdot 1 = 0.5, \\
    \bar{b} &= \bar{u} \cdot \frac{\partial u}{\partial b} = 0.5 \times 2.0 = 1.0.
  \end{align*}
  变量 $a$ 出度为 $2$，必须将外层乘法分支 $f = s \cdot a$ 与内层分支 $u = a \cdot b$ 的贡献相加：
  \[
    \bar{a}
    = \underbrace{\bar{f} \cdot s}_{\text{分支 1}} + \underbrace{\bar{u} \cdot b}_{\text{分支 2}}
    = 1 \times 0.5 + 0.5 \times (-0.5)
    = 0.5 - 0.25 = 0.25.
  \]
  综上，$(\bar{a}, \bar{b}, \bar{c}) = (0.25,\; 1.0,\; 0.5)$。直接解析检验 $\frac{\partial f}{\partial a} = \sigma(ab+c) + a b\,\sigma'(ab+c) = 0.5 + (-1)(0.25) = 0.25$，完全一致。
\end{enumerate}

\subsection*{习题 5.1.2\quad ReLU 折纸几何与 XOR 变体}
\textbf{构造（仅需 4 个隐藏 ReLU 神经元）}：先构造一维区间 $[-1, 1]$ 上的对称梯形平台函数
\[
  p(t) := 1 - \ReLU(t - 1) - \ReLU(-t - 1).
\]
当 $|t| \le 1$ 时两个 ReLU 均不激活，故 $p(t) = 1$；当 $|t| > 1$ 时恰有一个激活，故 $p(t) = 2 - |t|$，特别地当 $|t| > 2$ 时 $p(t) < 0$。

将 $x_1$ 与 $x_2$ 方向的梯形函数相加并减去阈值 $1$，定义双层 ReLU 网络：
\[
  f(x_1, x_2) = p(x_1) + p(x_2) - 1
  = 1 - \ReLU(x_1 - 1) - \ReLU(-x_1 - 1) - \ReLU(x_2 - 1) - \ReLU(-x_2 - 1).
\]
其权重矩阵与偏置为：
\[
  W_1 = \begin{pmatrix} 1 & 0 \\ -1 & 0 \\ 0 & 1 \\ 0 & -1 \end{pmatrix}, \quad
  b_1 = \begin{pmatrix} -1 \\ -1 \\ -1 \\ -1 \end{pmatrix}, \quad
  w_2 = (-1,\; -1,\; -1,\; -1)^\top, \quad
  b_2 = 1.
\]
\textbf{正确性验证}：
\begin{itemize}
  \item 若 $(x_1, x_2) \in [-1, 1]^2$，则四项 ReLU 均为 $0$，恒有 $f(x_1, x_2) = 1 > 0$；
  \item 若 $(x_1, x_2) \notin [-2, 2]^2$，则必有 $|x_1| > 2$ 或 $|x_2| > 2$，因此至少有一项 $\ReLU(\pm x_i - 1) > 1$，而其余各项非负，故 $f(x_1, x_2) < 1 - 1 = 0$。
\end{itemize}
\textbf{几何意义与最少神经元说明}：四个隐藏神经元分别对应四条折痕超平面 $x_1 = 1,\; x_1 = -1,\; x_2 = 1,\; x_2 = -1$，共同围出 $[-1, 1]^2$ 上的平顶高原；离开该正方形后函数以斜率 $-1$ 沿四个方向向下倾斜，在正集 $\{f > 0\}$ 的边界构成包含 $[-1, 1]^2$ 且被 $[-2, 2]^2$ 包围的八边形过渡带（在四个顶点处 $|x_1|+|x_2| < 3$）。若仅有 $3$ 个负权 ReLU 神经元，其凹上水平集 $\{f > 0\}$ 至多为三角形；但任何内接于 $[-2, 2]^2$ 的三角形面积不超过 $\frac{1}{2}\cdot 4 \cdot 4 = 8$，且易验证不存在完全落入 $[-2, 2]^2$ 同时包含四个顶点 $(\pm 1, \pm 1)$ 的三角形（底边落于 $x_2=-2$ 时，过 $(\pm 1, 1)$ 要求顶点的横坐标同时满足 $t \le -\frac{2}{3}$ 与 $t \ge \frac{2}{3}$，矛盾）。故 4 个隐藏神经元是凸平台构造下的最简解。

\subsection*{习题 5.1.3\quad 双线性形式与二次型的张量 VJP}
\begin{enumerate}[label=(\arabic*), leftmargin=1.8em]
  \item \textbf{二次型 $f(x, A) = x^\top A x$}：对 $f$ 取一阶微分：
  \[
    df = (dx)^\top A x + x^\top (dA) x + x^\top A (dx).
  \]
  由于 $(dx)^\top A x$ 是标量，取转置不变，即 $(dx)^\top A x = x^\top A^\top dx$。合并 $dx$ 项并对 $dA$ 项使用迹的循环不变性：
  \[
    df = x^\top (A^\top + A)\, dx + \tr\!\big((x x^\top)^\top dA\big).
  \]
  对照 $df = \inner{\bar{x}}{dx} + \inner{\bar{A}}{dA}_F$（上游标量 $\bar{f} = 1$），立得：
  \[
    \bar{x} = \nabla_x f = (A + A^\top) x, \qquad
    \bar{A} = \nabla_A f = x x^\top.
  \]
  \textbf{非对称时的辨析}：当 $A \neq A^\top$ 时，$\nabla_x (x^\top A x)$ \textbf{不等于} $2 A x$。显式反例：取 $A = \begin{psmallmatrix} 1 & 3 \\ 0 & 2 \end{psmallmatrix},\; x = \begin{psmallmatrix} 1 \\ 2 \end{psmallmatrix}$，则 $f(x) = x_1^2 + 3 x_1 x_2 + 2 x_2^2$，精确梯度为 $(A + A^\top)x = \begin{psmallmatrix} 2 & 3 \\ 3 & 4 \end{psmallmatrix}\begin{psmallmatrix} 1 \\ 2 \end{psmallmatrix} = \begin{psmallmatrix} 8 \\ 11 \end{psmallmatrix}$，而 $2 A x = \begin{psmallmatrix} 14 \\ 8 \end{psmallmatrix} \neq \begin{psmallmatrix} 8 \\ 11 \end{psmallmatrix}$。仅当 $A$ 对称时才有 $A + A^\top = 2A$。

  \item \textbf{连乘矩阵 $Z = X W_1 W_2$}：利用全微分 $dZ = (dX) W_1 W_2 + X (dW_1) W_2 + X W_1 (dW_2)$ 与 Frobenius 内积：
  \begin{align*}
    d\mathcal{L} &= \tr(\bar{Z}^\top dZ) \\
    &= \tr(\bar{Z}^\top dX\, W_1 W_2) + \tr(\bar{Z}^\top X\, dW_1\, W_2) + \tr(\bar{Z}^\top X W_1\, dW_2) \\
    &= \tr\!\big((W_1 W_2 \bar{Z}^\top)\, dX\big)
     + \tr\!\big((W_2 \bar{Z}^\top X)\, dW_1\big)
     + \tr\!\big((\bar{Z}^\top X W_1)\, dW_2\big).
  \end{align*}
  由 $\tr(M^\top dU) = \inner{M}{dU}_F$ 对每项取转置即得全部 VJP：
  \[
    \bar{X} = \bar{Z} W_2^\top W_1^\top, \qquad
    \bar{W}_1 = X^\top \bar{Z} W_2^\top, \qquad
    \bar{W}_2 = W_1^\top X^\top \bar{Z}.
  \]
\end{enumerate}

\subsection*{习题 5.1.4\quad 带标签平滑的融合交叉熵 VJP}
\begin{proof}
记 $p = \softmax(z)$，其中 $\ln p_k = z_k - \ln \sum_{j=1}^C e^{z_j}$。首先检验平滑目标 $\tilde{y}_k = (1-\varepsilon)y_k + \frac{\varepsilon}{C}$ 仍构成合法概率分布：
\[
  \sum_{k=1}^C \tilde{y}_k = (1-\varepsilon)\sum_{k=1}^C y_k + \varepsilon \cdot \frac{C}{C} = (1-\varepsilon) + \varepsilon = 1.
\]
\begin{enumerate}[label=(\arabic*), leftmargin=1.8em]
  \item \textbf{梯度推导}：展开融合损失 $\ell_{\text{ls}}(z, y) = -\sum_{k=1}^C \tilde{y}_k \ln p_k$：
  \[
    \ell_{\text{ls}}(z, y)
    = -\sum_{k=1}^C \tilde{y}_k z_k + \left(\sum_{k=1}^C \tilde{y}_k\right) \ln \sum_{j=1}^C e^{z_j}
    = -\sum_{k=1}^C \tilde{y}_k z_k + \ln \sum_{j=1}^C e^{z_j}.
  \]
  对第 $m$ 个 logit 求偏导：
  \[
    \frac{\partial \ell_{\text{ls}}}{\partial z_m}
    = -\tilde{y}_m + \frac{e^{z_m}}{\sum_{j=1}^C e^{z_j}}
    = p_m - \tilde{y}_m,
  \]
  即 $\nabla_z \ell_{\text{ls}}(z, y) = p - \tilde{y}$。这说明只需把命题~\ref{prop:ce-vjp} 中的目标替换为 $\tilde{y}$，无需任何额外数值开销。

  \item \textbf{与均匀分布 KL 正则的等价性}：记均匀分布 $u_k = 1/C$，则
  \begin{align*}
    \ell_{\text{ls}}(z, y)
    &= (1-\varepsilon)\!\left(-\sum_{k=1}^C y_k \ln p_k\right)
       + \varepsilon\!\left(-\sum_{k=1}^C u_k \ln p_k\right) \\
    &= (1-\varepsilon)\,\ell_{\text{ce}}(z, y)
       + \varepsilon \sum_{k=1}^C u_k \ln \frac{u_k}{p_k}
       - \varepsilon \sum_{k=1}^C u_k \ln u_k \\
    &= (1-\varepsilon)\,\ell_{\text{ce}}(z, y) + \varepsilon\,\DKL(u \,\|\, p) + \varepsilon \ln C.
  \end{align*}
  由于 $\varepsilon \ln C$ 是与参数无关的常数，最小化标签平滑交叉熵严格等价于在交叉熵损失上附加 KL 惩罚 $\varepsilon \DKL(u \,\|\, p)$，从而阻止 logits 趋向 $\pm\infty$、抑制过度自信并提升模型校准度。
\end{enumerate}
\end{proof}

\subsection*{习题 5.1.5\quad 扩展 \texttt{Value} 引擎与自动梯度检验}
\textbf{核心算子实现}：在原有 \texttt{Value} 类中补充标量幂、除法、指数、对数与 ReLU，每个节点仅保存自身前向值与闭包的局部导数乘子：
\begin{lstlisting}
def __pow__(self, p):
    assert isinstance(p, (int, float))
    out = Value(self.data ** p, (self,))
    def _b(): self.grad += p * (self.data ** (p - 1)) * out.grad
    out._backward = _b; return out

def __truediv__(self, other):
    other = other if isinstance(other, Value) else Value(other)
    return self * (other ** -1.0)

def exp(self):
    e = math.exp(self.data); out = Value(e, (self,))
    def _b(): self.grad += e * out.grad
    out._backward = _b; return out

def log(self):
    out = Value(math.log(self.data), (self,))
    def _b(): self.grad += (1.0 / self.data) * out.grad
    out._backward = _b; return out

def relu(self):
    out = Value(max(0.0, self.data), (self,))
    def _b(): self.grad += (1.0 if self.data > 0 else 0.0) * out.grad
    out._backward = _b; return out
\end{lstlisting}
\textbf{通用中心差分梯度检验}：对每个参数 $x_i$ 施加 $\pm h$ 扰动，计算 $g_{\text{num},i} = \frac{f(x + h e_i) - f(x - h e_i)}{2h}$，相对误差按
\[
  \mathrm{rel\_err}_i = \frac{\abs{g_{\text{ad},i} - g_{\text{num},i}}}{\max(\abs{g_{\text{ad},i}} + \abs{g_{\text{num},i}},\; 10^{-8})}
\]
度量。注意在检验含 \texttt{relu()} 的网络时，被测参数点应避开隐藏激活恰好等于 $0$ 的不可微折点（否则左、右差分跨越折点会产生 $O(1)$ 虚假偏差）。在 $h = 10^{-6}$ 下，对由 \texttt{relu}、\texttt{exp}、\texttt{log}、除法与幂构成的 2 层标量 MLP 实测，最大相对误差为 $1.82 \times 10^{-9} < 10^{-8}$。

\subsection*{习题 5.1.6\quad 双月数据集上的深度 vs 宽度消融}
\textbf{参数量账本}：含偏置项时，单隐藏层网络 $2 \to 25 \to 1$ 共有 $2\times 25 + 25 + 25\times 1 + 1 = 101$ 个参数；双隐藏层网络 $2 \to 9 \to 8 \to 1$ 共有 $2\times 9 + 9 + 9\times 8 + 8 + 8\times 1 + 1 = 116$ 个参数，二者处于同一百参数量级（若进一步调整为 $2 \to 8 \to 7 \to 1$ 则为 $95$ 参数）。

\textbf{实测对比与几何解释}（固定 Kaiming 初始化、Adam \texttt{lr=0.05} 训练 300 步）：
\begin{center}
\small
\begin{tabular}{@{}lcccc@{}}
\toprule
网络结构（参数量） & 双月训练损失 & 双月准确率 & 同心圆训练损失 & 同心圆准确率 \\
\midrule
浅宽网络 $2 \to 25 \to 1$（101） & $0.00352$ & $100.0\%$ & $0.00039$ & $100.0\%$ \\
深窄网络 $2 \to 9 \to 8 \to 1$（116） & $0.00034$ & $100.0\%$ & $0.00008$ & $100.0\%$ \\
\bottomrule
\end{tabular}
\end{center}
两者最终均可达到 $100\%$ 分类精度，但双隐藏层网络的交叉熵损失低近一个数量级、前期收敛更陡峭。原因在于浅层网络的 $25$ 个 ReLU 只能在输入平面画出至多 $25$ 条互不嵌套的线性折痕（对同心圆只能用多边形逐段逼近）；而双隐藏层网络的第二层可对第一层输出的 $9$ 个分段线性特征做\textbf{复合折纸}，用极少神经元即构造出更紧致光滑的封闭凸包与 S 形弯折决策边界。

\subsection*{习题 5.1.7\quad Pearlmutter 技巧：免构造 Hessian 的 HVP}
\begin{proof}
设梯度映射为 $g(\theta) := \nabla_\theta \mathcal{L}(\theta) \in \R^P$，则其雅可比矩阵恰为对称的 Hessian 矩阵 $J_g(\theta) = \nabla^2_\theta \mathcal{L}(\theta) = H(\theta)$。
\begin{enumerate}[label=(\arabic*), leftmargin=1.8em]
  \item \textbf{Forward-over-Reverse（前向套反向）}：对给定方向向量 $v \in \R^P$，按定义有
  \[
    H(\theta)\, v = J_g(\theta)\, v = \left.\frac{d}{dt}\, g(\theta + t v)\right|_{t=0}.
  \]
  具体计算分两步：（i）先对标量损失 $\mathcal{L}(\theta)$ 执行一次反向模式自动微分（VJP），得到计算 $g(\theta)$ 的计算图，时间与内存均为 $O(\operatorname{Cost}(\mathcal{L}))$；（ii）把 $g(\theta)$ 视为程序映射 $\R^P \to \R^P$，沿切向量 $v$ 执行一次前向模式自动微分（JVP，即对偶数字 $a + \epsilon \dot{a}$ 传播）。前向模式沿单个方向 $v$ 的时间和内存开销仅为原计算图的常数倍（约 $2\times$），故总复杂度为 $O(\operatorname{Cost}(\mathcal{L}))$，完全无需构造 $O(P^2)$ 规模的矩阵 $H$。

  \item \textbf{Reverse-over-Reverse（反向套反向）与常数对比}：构造标量函数 $\phi(\theta) := g(\theta)^\top v = \sum_{i=1}^P g_i(\theta) v_i$（其中 $v$ 视为与 $\theta$ 无关的常数向量），对其再执行一次反向模式微分：
  \[
    \nabla_\theta \phi(\theta) = \nabla_\theta \big(g(\theta)^\top v\big) = J_g(\theta)^\top v = H(\theta)^\top v = H(\theta)\, v,
  \]
  最后一步利用了二阶偏导对称性 $H^\top = H$。二者的渐近复杂度同为 $O(\operatorname{Cost}(\mathcal{L}))$，但常数开销有明显差异：Reverse-over-Reverse 必须为第一次反向图保留全部高阶中间激活以便做第二次反向遍历（PyTorch 中需设 \texttt{create\_graph=True}），显存约增至单次反向的 $2\sim 3$ 倍；而 Forward-over-Reverse 的外层 JVP 无需缓存反向磁带，显存占用更低、算子融合更友好，是计算大型网络 HVP 的首选模式。
\end{enumerate}
\end{proof}

\subsection*{习题 5.1.8\quad 隐函数微分与深度平衡模型}
\begin{proof}
\begin{enumerate}[label=(\arabic*), leftmargin=1.8em]
  \item \textbf{隐函数求导公式}：将收敛不动点视为参数 $\theta$ 的光滑隐函数 $z^*(\theta)$，满足恒等式 $z^*(\theta) = f_\theta\big(z^*(\theta), x\big)$。两端关于 $\theta$ 求全导数（链式法则）：
  \[
    \frac{\partial z^*}{\partial \theta}
    = \frac{\partial f_\theta(z^*, x)}{\partial z^*} \frac{\partial z^*}{\partial \theta}
    + \frac{\partial f_\theta(z^*, x)}{\partial \theta}.
  \]
  记 $J_z := \frac{\partial f_\theta(z^*, x)}{\partial z^*}$ 与 $J_\theta := \frac{\partial f_\theta(z^*, x)}{\partial \theta}$。移项合并得 $(I - J_z)\frac{\partial z^*}{\partial \theta} = J_\theta$。若不动点迭代局部压缩（例如谱半径 $\rho(J_z) < 1$），则 $I - J_z$ 可逆，左乘 $(I - J_z)^{-1}$ 即证：
  \[
    \frac{\partial z^*}{\partial \theta} = (I - J_z)^{-1} J_\theta.
  \]

  \item \textbf{反向传播降为 $O(1)$ 显存的实现}：设下游损失传回的伴随量为 $\bar{z}^* = (\nabla_{z^*} \mathcal{L})^\top$（写成列向量），我们需要的 VJP 为：
  \[
    \bar{\theta} = \left(\frac{\partial z^*}{\partial \theta}\right)^\top \bar{z}^*
    = J_\theta^\top (I - J_z^\top)^{-1} \bar{z}^*.
  \]
  引入辅助向量 $u := (I - J_z^\top)^{-1} \bar{z}^*$，则计算分为两步：
  \begin{enumerate}[label=(\roman*), leftmargin=1.6em]
    \item \textbf{求解线性方程组} $(I - J_z^\top) u = \bar{z}^*$，等价于不动点方程 $u = J_z^\top u + \bar{z}^*$。利用单层单步反向自动微分即可无矩阵地计算 VJP $v \mapsto J_z^\top v$，结合共轭梯度、GMRES 或 Anderson 加速迭代求出 $u$；
    \item \textbf{一次单步 VJP}：计算 $\bar{\theta} = J_\theta^\top u$，即把 $u$ 作为上游梯度对单步映射 $f_\theta(z^*, x)$ 执行一次常规反向传播。
  \end{enumerate}
  整个反向过程\textbf{只依赖最终不动点 $z^*$ 本身}，无论前向迭代 100 步还是 1000 步，均可在无梯度模式（\texttt{torch.\allowbreak no\_grad()}）下完成，彻底免除中间激活轨迹缓存，显存占用恒为 $O(1)$。
\end{enumerate}
\end{proof}

% =========================================================================
\section*{第 2 章习题解答（共 8 题）}
\addcontentsline{toc}{section}{第 2 章习题解答（共 8 题）}
% =========================================================================

\subsection*{习题 5.2.1\quad 归一化算子手算对比}
给定 $X = \begin{psmallmatrix} 1 & 2 & 3 \\ 4 & 8 & 12 \end{psmallmatrix} \in \R^{2\times 3}$，且 $\gamma = \mathbf{1},\; \beta = \mathbf{0},\; \epsilon = 0$。
\begin{enumerate}[label=(\arabic*), leftmargin=1.8em]
  \item \textbf{BatchNorm（沿批维按列统计）}：三列的均值与标准差分别为：
  \begin{itemize}
    \item 第 1 列 $(1, 4)^\top$：$\mu_1 = 2.5,\; \sigma_1 = \sqrt{\frac{(-1.5)^2 + 1.5^2}{2}} = 1.5$；
    \item 第 2 列 $(2, 8)^\top$：$\mu_2 = 5.0,\; \sigma_2 = 3.0$；
    \item 第 3 列 $(3, 12)^\top$：$\mu_3 = 7.5,\; \sigma_3 = 4.5$。
  \end{itemize}
  每列恰好只有两个不相等的数，标准化后必为 $-1$ 与 $+1$：
  \[
    \BN(X) = \begin{pmatrix} -1 & -1 & -1 \\ 1 & 1 & 1 \end{pmatrix}.
  \]

  \item \textbf{LayerNorm（沿特征维按行统计）}：
  \begin{itemize}
    \item 第 1 行 $(1, 2, 3)$：$\mu^{(1)} = 2,\; (\sigma^{(1)})^2 = \frac{(-1)^2 + 0^2 + 1^2}{3} = \frac{2}{3},\; \sigma^{(1)} = \sqrt{\frac{2}{3}}$；
    \item 第 2 行 $(4, 8, 12)$：$\mu^{(2)} = 8,\; \sigma^{(2)} = 4\sqrt{\frac{2}{3}}$。
  \end{itemize}
  归一化后两行\textbf{完全相同}：
  \[
    \LN(X) = \begin{pmatrix}
      -\sqrt{\tfrac{3}{2}} & 0 & \sqrt{\tfrac{3}{2}} \\[3pt]
      -\sqrt{\tfrac{3}{2}} & 0 & \sqrt{\tfrac{3}{2}}
    \end{pmatrix}
    \approx \begin{pmatrix}
      -1.2247 & 0 & 1.2247 \\
      -1.2247 & 0 & 1.2247
    \end{pmatrix}.
  \]

  \item \textbf{RMSNorm（沿行除以均方根，不去均值）}：
  \begin{itemize}
    \item 第 1 行：$\mathrm{RMS}_1 = \sqrt{\frac{1^2 + 2^2 + 3^2}{3}} = \sqrt{\frac{14}{3}}$；
    \item 第 2 行：$\mathrm{RMS}_2 = \sqrt{\frac{4^2 + 8^2 + 12^2}{3}} = 4\sqrt{\frac{14}{3}}$。
  \end{itemize}
  归一化后两行同样\textbf{完全相同}：
  \[
    \RMSN(X) = \sqrt{\frac{3}{14}}\begin{pmatrix} 1 & 2 & 3 \\ 1 & 2 & 3 \end{pmatrix}
    \approx \begin{pmatrix}
      0.4629 & 0.9258 & 1.3887 \\
      0.4629 & 0.9258 & 1.3887
    \end{pmatrix}.
  \]
\end{enumerate}
\textbf{对比结论}：由于第二行 $x^{(2)} = 4 x^{(1)}$，而 $\LN$ 与 $\RMSN$ 均满足正尺度不变性 $\operatorname{Norm}(\alpha x) = \operatorname{Norm}(x)\;(\alpha > 0)$，故二者输出的两行完全一致；区别在于 $\LN$ 额外减去了行均值（具备平移不变性，输出均值为 $0$），而 $\RMSN$ 保留了正偏置方向。相反，$\BN$ 跨样本做列统计，比较的是两个样本之间的相对大小，两行输出恰好相反。

\subsection*{习题 5.2.2\quad Adam 初始步与偏差校正直观}
已知 $m_0 = v_0 = 0,\; g_1 = g_2 = 0.2,\; \beta_1 = 0.9,\; \beta_2 = 0.999,\; \eta = 0.01,\; \epsilon = 0$。第 1 步原始矩估计为：
\[
  m_1 = (1 - \beta_1) g_1 = 0.1 \times 0.2 = 0.02, \qquad
  v_1 = (1 - \beta_2) g_1^2 = 0.001 \times 0.04 = 4 \times 10^{-5}.
\]
\begin{enumerate}[label=(\alph*), leftmargin=1.8em]
  \item \textbf{带偏差校正}：$\hat{m}_1 = \frac{m_1}{1-\beta_1} = g_1 = 0.2$，$\hat{v}_1 = \frac{v_1}{1-\beta_2} = g_1^2 = 0.04$，因此：
  \[
    \Delta \theta_1^{\text{corr}} = -\eta \frac{\hat{m}_1}{\sqrt{\hat{v}_1}} = -0.01 \times \frac{0.2}{0.2} = -0.01, \qquad
    \theta_1^{\text{corr}} = 1.0 - 0.01 = 0.99.
  \]
  （顺带可得第 2 步仍满足 $\hat{m}_2 = 0.2,\; \hat{v}_2 = 0.04$，故 $\Delta \theta_2^{\text{corr}} = -0.01$。）

  \item \textbf{去掉偏差校正}：直接以 $m_1 / \sqrt{v_1}$ 更新：
  \[
    \Delta \theta_1^{\text{nocorr}}
    = -\eta \frac{m_1}{\sqrt{v_1}}
    = -0.01 \times \frac{0.02}{\sqrt{4 \times 10^{-5}}}
    = -0.01 \times \frac{0.1}{\sqrt{0.001}}
    = -0.01\sqrt{10} \approx -0.031623.
  \]
  \textbf{倍数比较}：去掉偏差校正后，首步步长被无端放大了
  \[
    \frac{\Delta \theta_1^{\text{nocorr}}}{\Delta \theta_1^{\text{corr}}}
    = \frac{1 - \beta_1}{\sqrt{1 - \beta_2}}
    = \frac{0.1}{\sqrt{0.001}} = \sqrt{10} \approx 3.1623\text{ 倍}.
  \]
  原因在于 $\beta_2$ 极接近 $1$，未校正的 $v_1$ 被零初始化严重拉低了 $1000$ 倍（开方后约 $31.6$ 倍），而 $m_1$ 仅被拉低 $10$ 倍，分子分母缩水失衡导致训练初期步长剧烈过冲。
\end{enumerate}

\subsection*{习题 5.2.3\quad LeakyReLU 与 GELU 的初始化方差推广}
\begin{enumerate}[label=(\alph*), leftmargin=1.8em]
  \item \textbf{LeakyReLU 初始化方差证明}：设预激活 $z$ 的概率密度 $p(z)$ 偶对称（$p(-z) = p(z)$）且 $\E[z] = 0$，则 $\Var(z) = \E[z^2] = 2\int_0^\infty z^2 p(z)\,dz$。将 $\sigma(z)^2$ 的期望按正负半轴拆分：
  \begin{align*}
    \E[\sigma(z)^2]
    &= \int_{-\infty}^0 (a z)^2 p(z)\,dz + \int_0^\infty z^2 p(z)\,dz \\
    &= a^2 \int_0^\infty u^2 p(-u)\,du + \int_0^\infty z^2 p(z)\,dz
     = (1 + a^2) \int_0^\infty z^2 p(z)\,dz
     = \frac{1 + a^2}{2}\,\Var(z).
  \end{align*}
  代入下一层方差递推 $\Var(z^{(l)}) = n_{\text{in}} \Var(w^{(l)}) \E[\sigma(z^{(l-1)})^2] = n_{\text{in}} \Var(w^{(l)}) \frac{1+a^2}{2}\Var(z^{(l-1)})$，令方差保持因子为 $1$ 即得推广 Kaiming 初始化：
  \[
    \Var(w) = \frac{2}{(1 + a^2)\, n_{\text{in}}}.
  \]
  当 $a = 0$ 时退化为 ReLU 的 $\frac{2}{n_{\text{in}}}$；当 $a = 1$ 时退化为线性层的 $\frac{1}{n_{\text{in}}}$。

  \item \textbf{GELU 精确二阶矩闭式推导与比较}：设 $z \sim \mathcal{N}(0, 1)$，记标准正态密度 $\phi(z) = \frac{1}{\sqrt{2\pi}}e^{-z^2/2}$ 与分布函数 $\Phi(z)$，注意到 $\phi'(z) = -z\phi(z)$。对二阶矩作分部积分（取 $u = z\Phi(z)^2,\; dv = z\phi(z)\,dz = -d\phi(z)$）：
  \[
    \E[\operatorname{GELU}(z)^2]
    = \int_{-\infty}^\infty z^2 \Phi(z)^2 \phi(z)\,dz
    = \int_{-\infty}^\infty \Phi(z)^2 \phi(z)\,dz + 2\int_{-\infty}^\infty z \Phi(z) \phi(z)^2\,dz.
  \]
  分别计算两项积分：
  \begin{itemize}
    \item 第一项：令 $t = \Phi(z) \in (0, 1)$，则 $dt = \phi(z)\,dz$，故 $\int_{-\infty}^\infty \Phi(z)^2 \phi(z)\,dz = \int_0^1 t^2\,dt = \frac{1}{3}$；
    \item 第二项：由于 $\phi(z)^2 = \frac{1}{2\pi} e^{-z^2}$ 且 $\frac{d}{dz}\big(-\frac{1}{2}e^{-z^2}\big) = z e^{-z^2}$，再次分部积分得
    \begin{align*}
      2\int_{-\infty}^\infty z \Phi(z) \phi(z)^2\,dz
      &= \frac{1}{\pi}\int_{-\infty}^\infty \Phi(z)\, d\!\left(-\frac{1}{2}e^{-z^2}\right)
       = \frac{1}{2\pi}\int_{-\infty}^\infty e^{-z^2}\phi(z)\,dz \\
      &= \frac{1}{2\pi\sqrt{2\pi}}\int_{-\infty}^\infty e^{-\frac{3}{2}z^2}\,dz
       = \frac{1}{2\pi\sqrt{2\pi}}\sqrt{\frac{2\pi}{3}}
       = \frac{1}{2\pi\sqrt{3}}.
    \end{align*}
  \end{itemize}
  合并两项，得到惊人简洁的精确闭式解：
  \[
    \E[\operatorname{GELU}(z)^2] = \frac{1}{3} + \frac{1}{2\pi\sqrt{3}} \approx 0.33333 + 0.09189 = 0.42522.
  \]
  \textbf{与 ReLU 对比}：$\E[\operatorname{GELU}(z)^2] \approx 0.4252 < 0.5 = \E[\ReLU(z)^2]$！因此单位方差处精确方差保持增益为 $g_{\text{GELU}} = (0.42522)^{-1/2} \approx 1.5335$，反而略大于 ReLU 的 $\sqrt{2} \approx 1.4142$。几何原因在于：在 $z \in (0, 2)$ 的高概率正区间上，$\Phi(z) < 1$ 使 $z\Phi(z)$ 严格小于 $z$，这一正半轴能量衰减超过了负半轴微弱下凹（$\max_{z<0}|z\Phi(z)| \approx 0.17$）所贡献的极小平方项；仅当输入尺度远大于 $1$ 时 GELU 才渐近趋近于 ReLU（增益趋于 $\sqrt{2}$）。
\end{enumerate}

\subsection*{习题 5.2.4\quad LayerNorm 的张量 VJP 与尺度不变梯度正交性}
\begin{proof}
记中心化向量 $\tilde{x}_i := x_i - \mu$，则 $y_i = \tilde{x}_i / \sigma$，且 $\sum_{j=1}^D y_j = 0$、$\frac{1}{D}\sum_{j=1}^D y_j^2 = 1$。
\begin{enumerate}[label=(\alph*), leftmargin=1.8em]
  \item \textbf{反向传播公式推导}：由 $\sigma^2 = \frac{1}{D}\sum_j \tilde{x}_j^2$ 取微分得
  \[
    d\mu = \frac{1}{D}\sum_{j=1}^D dx_j, \qquad
    d\sigma = \frac{1}{D\sigma}\sum_{j=1}^D \tilde{x}_j (dx_j - d\mu) = \frac{1}{D}\sum_{j=1}^D y_j\, dx_j,
  \]
  其中利用了 $\sum_j y_j = 0$ 消去 $d\mu$ 项。接着对 $y_i = \tilde{x}_i / \sigma$ 求全微分：
  \[
    dy_i = \frac{dx_i - d\mu}{\sigma} - \frac{\tilde{x}_i}{\sigma^2}\, d\sigma
         = \frac{1}{\sigma}\left(dx_i - \frac{1}{D}\sum_{j=1}^D dx_j - y_i \cdot \frac{1}{D}\sum_{j=1}^D y_j\, dx_j\right).
  \]
  代入 $d\mathcal{L} = \sum_{i=1}^D \bar{y}_i\, dy_i$ 并交换求和次序：
  \begin{align*}
    d\mathcal{L}
    &= \frac{1}{\sigma}\sum_{i=1}^D \bar{y}_i\, dx_i
     - \frac{1}{\sigma}\left(\frac{1}{D}\sum_{i=1}^D \bar{y}_i\right)\sum_{j=1}^D dx_j
     - \frac{1}{\sigma}\left(\frac{1}{D}\sum_{i=1}^D \bar{y}_i y_i\right)\sum_{j=1}^D y_j\, dx_j \\
    &= \sum_{j=1}^D \frac{1}{\sigma}\left(\bar{y}_j - \frac{1}{D}\sum_{i=1}^D \bar{y}_i - y_j \cdot \frac{1}{D}\sum_{i=1}^D \bar{y}_i y_i\right) dx_j.
  \end{align*}
  读取 $dx_j$ 的系数即得向量形式：
  \[
    \bar{x} = \frac{1}{\sigma}\left(\bar{y} - \frac{1}{D}\sum_{j=1}^D \bar{y}_j\,\mathbf{1} - y \cdot \frac{1}{D}\sum_{j=1}^D \bar{y}_j y_j\right).
  \]

  \item \textbf{两条正交性的验证与几何解释}：记 $m_1 = \frac{1}{D}\sum_j \bar{y}_j$ 与 $m_2 = \frac{1}{D}\sum_j \bar{y}_j y_j$。直接求和与作内积：
  \begin{align*}
    \sum_{i=1}^D \bar{x}_i
    &= \frac{1}{\sigma}\left(\sum_i \bar{y}_i - D m_1 - m_2 \sum_i y_i\right)
     = \frac{1}{\sigma}(D m_1 - D m_1 - 0) = 0, \\
    \inner{\bar{x}}{x - \mu\mathbf{1}}
    &= \sigma \inner{\bar{x}}{y}
     = \sum_i \bar{y}_i y_i - m_1 \sum_i y_i - m_2 \sum_i y_i^2
     = D m_2 - 0 - D m_2 = 0.
  \end{align*}
  \textbf{几何解释}：对任意平移常数 $c$ 与正尺度因子 $\alpha > 0$，函数 $y(x)$ 满足 $y(x + c\mathbf{1}) \equiv y(x)$ 与 $y(\alpha x) \equiv y(x)$，故复合损失 $\mathcal{L}(y(x))$ 沿全一方向 $\mathbf{1}$ 与径向 $x - \mu\mathbf{1}$ 的方向导数恒为零！这意味着 LayerNorm 在反向传播中相当于一个\textbf{正交投影算子}，自动剥离了梯度中企图改变输入均值或放缩输入范数的冗余分量，使参数更新只沿单位超球面切空间旋转，从而天然抑制梯度爆炸并稳定训练。
\end{enumerate}
\end{proof}

\subsection*{习题 5.2.5\quad 20 层反向梯度传播实验分析}
在 \texttt{exp2\_init\_norm\_opt} 的 20 层网络（$B=64, D=128, L=20$）中为各层权重设置 \texttt{requires\_grad=True}，施加损失 $\mathcal{L} = \frac{1}{BD}\norm{h^{(20)}}_2^2$ 后反向传播，四组配置的权重梯度范数 $\norm{\nabla_{W^{(l)}}\mathcal{L}}_F$ 实测如下：
\begin{center}
\small
\begin{tabular}{@{}lccc@{}}
\toprule
初始化与归一化配置 & 第 1 层 $\norm{\nabla_{W^{(1)}}\mathcal{L}}_F$ & 第 20 层 $\norm{\nabla_{W^{(20)}}\mathcal{L}}_F$ & 现象诊断 \\
\midrule
过小初始化（$\sigma_w = 0.05$） & $4.66 \times 10^{-16}$ & $9.09 \times 10^{-16}$ & 前后向双重指数坍缩 \\
过大初始化（$\sigma_w = 0.20$） & $1.41 \times 10^{8}$ & $2.75 \times 10^{8}$ & 指数爆炸至上溢边缘 \\
Kaiming（$\sqrt{2/D}$） & $1.54$ & $3.01$ & 全 20 层保持 $O(1)$ 稳定 \\
Kaiming + RMSNorm & $1.28 \times 10^{-7}$ & $1.77 \times 10^{-6}$ & 见下方关键排错分析 \\
\bottomrule
\end{tabular}
\end{center}
\textbf{深刻洞察与排错要点}：细心的读者会发现上表最后一行 \texttt{Kaiming + RMSNorm} 的梯度反而比纯 \texttt{Kaiming} 更小！为什么？因为这里施加的损失是 $\mathcal{L} = \frac{1}{BD}\norm{h^{(20)}}_2^2$，而最后一层 RMSNorm 恰好将每个样本的 $\frac{1}{D}\norm{h^{(20)}_b}_2^2$ 强制归一化为常数 $1$（仅靠 $\epsilon = 10^{-6}$ 才留下一丝非零导数，正如习题 5.2.4 所证，归一化梯度与径向严格正交）！若将损失改为对未归一化线性读出头 $z_{\text{out}} = h^{(20)} W_{\text{head}}$ 的交叉熵或 MSE，\texttt{Kaiming + RMSNorm} 的第 1--20 层权重梯度即恢复平稳的 $O(1)$ 分布，且对初始化尺度失配具有极强鲁棒性。

\subsection*{习题 5.2.6\quad 病态二次型上的 Adam+$L_2$ vs AdamW 实证}
考虑 $\mathcal{L}(\theta_1, \theta_2) = 100\,\theta_1^2 + 0.01\,\theta_2^2$，初始点 $(2.0, 2.0)$，取 \texttt{lr=0.01}、\texttt{weight\_decay=0.1} 迭代 200 步：
\begin{itemize}
  \item \textbf{耦合 $L_2$ 的 \texttt{Adam}}：最终权重 $(\theta_1, \theta_2) = (0.4964,\; 0.4964)$；
  \item \textbf{解耦衰减的 \texttt{AdamW}}：最终权重 $(\theta_1, \theta_2) = (0.3727,\; 0.3727)$。
\end{itemize}
\textbf{机理剖析（与定理~\ref{thm:adamw} 对照）}：
\begin{enumerate}[label=(\arabic*), leftmargin=1.8em]
  \item 在每步有效衰减比例上，\texttt{Adam+$L_2$} 把 $\lambda \theta_i$ 混入梯度 $g_i$ 后一起除以 $\sqrt{\hat{v}_{t,i}}$。在初始点附近，$g_1 = 200\theta_1 \approx 400$ 而 $g_2 = 0.02\theta_2 \approx 0.04$，于是高曲率轴 $\theta_1$ 的等效衰减因子 $\frac{\lambda}{\sqrt{\hat{v}_1}} \approx \frac{0.1}{400} = 2.5\times 10^{-4}$ 几乎消失；而低曲率轴 $\theta_2$ 的等效衰减因子高达 $\frac{\lambda}{\sqrt{\hat{v}_2}} \approx \frac{0.1}{0.24} \approx 0.42$，两者相差三个数量级！
  \item 更有趣的是：在无噪声单调下降阶段（梯度恒正），\texttt{Adam+$L_2$} 的分子与分母同号抵消（$\hat{m}_t / \sqrt{\hat{v}_t} \approx +1$），导致耦合 $L_2$ 项在单调段被彻底“吃掉”，步长退化为恒定 $-\eta$；而 \texttt{AdamW} 额外施加与梯度尺度无关的真实乘法收缩 $(1 - \eta\lambda)\theta_{t-1}$，因此 200 步后权重额外按 $e^{-200 \times 0.01 \times 0.1} = e^{-0.2} \approx 0.819$ 比例收缩（实测 $0.3727$ 远低于 $0.4964$），严格印证了二者的不等价性。
\end{enumerate}

\subsection*{习题 5.2.7\quad Dropout 与 BatchNorm 串联的方差偏移}
\begin{proof}
设随机掩码 $m \sim \operatorname{Bernoulli}(1-p)$ 与特征 $x$ 独立，定义缩放因子 $M := \frac{m}{1-p}$。首先计算 $M$ 的一、二阶矩：
\[
  \E[M] = \frac{1 - p}{1 - p} = 1, \qquad
  \E[M^2] = \frac{\E[m^2]}{(1-p)^2} = \frac{1 - p}{(1-p)^2} = \frac{1}{1 - p}.
\]
因此对 $\tilde{x} = M x$，由独立性有 $\E[\tilde{x}] = \E[M]\E[x] = \E[x]$（均值无偏），但其方差为：
\begin{align*}
  \Var(\tilde{x})
  &= \E[M^2]\E[x^2] - (\E[M]\E[x])^2 \\
  &= \frac{1}{1 - p}\big(\Var(x) + (\E[x])^2\big) - (\E[x])^2 \\
  &= \frac{1}{1 - p}\,\Var(x) + \frac{p}{1 - p}\,(\E[x])^2.
\end{align*}
例如当 $p = 0.5,\; \E[x] = 1,\; \Var(x) = 1$ 时，$\Var(\tilde{x}) = 2 \times 1 + 1 \times 1 = 3$，方差被放大整整 $3$ 倍！
\end{proof}
\textbf{对 BN 的破坏与工程对策}：若将 Dropout 置于 BN 之前，训练阶段 BN 累计的滑动方差 $\hat{\sigma}_{\text{run}}^2$ 会收敛到被放大后的 $\Var(\tilde{x})$；但推理阶段关闭 Dropout 后，喂给 BN 的真实特征方差突然跌回 $\Var(x)$，导致 BN 输出幅度系统性萎缩 $\sqrt{\Var(x)/\Var(\tilde{x})}$ 倍，逐层累积引发精度暴跌（即\textbf{方差偏移}）。解决方案是：（1）将 Dropout 放在 BN 之后（或仅放在全连接分类头之前）；（2）现代大卷积与 Transformer 主干倾向于用权重衰减、数据增强与 DropPath 替代密集通道 Dropout。

\subsection*{习题 5.2.8\quad 低显存自适应优化器（Adafactor / 8-bit Adam / Muon）}
\textbf{显存账本}：设模型有 $P$ 个参数，在 BF16/FP16 混合精度训练中，参数权重占 $2P$ 字节，梯度占 $2P$ 字节，而标准 AdamW 需维护 FP32 主权重（$4P$）、一阶矩 $m_t$（$4P$）与二阶矩 $v_t$（$4P$）——仅 $m_t, v_t$ 两份优化器状态即占 $8P$ 字节，达到半精度模型权重显存（$2P$）的 $4\times$（亦为 FP32 权重的 $2\times$）。

\textbf{三类前沿降存与加速思路对比}：
\begin{itemize}
  \item \textbf{Adafactor}：对二维权重矩阵 $W \in \R^{m \times n}$ 的二阶矩矩阵 $V \in \R_+^{m \times n}$，仅存储行均值 $r \in \R^m$ 与列均值 $c \in \R^n$，以秩-1 外积 $\hat{V}_{ij} = r_i c_j / (\mathbf{1}^\top r)$ 近似 $V_{ij}$，将状态显存从 $O(mn)$ 压缩至 $O(m+n)$。
  \item \textbf{8-bit Adam}：采用分块动态指数非线性量化（Block-wise Dynamic Quantization）将 $m_t, v_t$ 压至 8-bit 整数，优化器显存直降 $75\%$ 且几乎无精度损失。
  \item \textbf{Muon（矩阵谱正交化优化器）}：只维护一阶动量矩阵 $M_t \in \R^{m\times n}$（省去 $v_t$），并借助 5 步低精度矩阵乘 Newton--Schulz 迭代 $X_{k+1} = a X_k + b (X_k X_k^\top) X_k + c (X_k X_k^\top)^2 X_k$ 近似计算 $M_t$ 的极分解正交因子 $U V^\top$（若 $M_t = U \Sigma V^\top$）。
\end{itemize}
\textbf{为何按奇异值谱而非逐元素缩放更具结构收益}：神经网络隐藏层权重 $W$ 本质上是线性算子，其输入输出放大倍数由奇异值谱 $\sigma_1(W) \ge \dots \ge \sigma_r(W)$ 而非单个矩阵元大小决定。Adam 的逐元素预条件 $\diag(v_t)^{-1/2} m_t$ 依赖于坐标基选择（旋转坐标系后更新方向随之改变），无法消除高度相关的特征方向间的谱条件数失衡；而 Muon 将更新矩阵的所有非零奇异值统一拉平为 $1$，相当于在谱范数球约束下执行最速下降，对左右正交基旋转严格不变，能让长尾低奇异值子空间以与主方向同等速率同步学习，在现代 LLM 隐藏层矩阵训练中展现出比 AdamW 更高的样本效率。

% =========================================================================
\section*{第 3 章习题解答（共 8 题）}
\addcontentsline{toc}{section}{第 3 章习题解答（共 8 题）}
% =========================================================================

\subsection*{习题 5.3.1\quad 多层卷积的特征图尺寸、参数量与感受野手算}
输入尺寸 $(C_0, H_0, W_0) = (3, 224, 224)$，初始感受野 $r_0 = 1$，累计步幅 $J_0 = 1$。逐层套用式~\eqref{eq:conv-out-size} 与式~\eqref{eq:receptive-field}：
\begin{enumerate}[label=(\arabic*), leftmargin=1.8em]
  \item \textbf{Conv1}（$C_{\text{out}}=64,\; k=7,\; S=2,\; P=3,\; d=1$）：
  \begin{itemize}
    \item 输出尺寸：$H_1 = \lfloor(224 + 2\times 3 - 7)/2\rfloor + 1 = \lfloor 223/2 \rfloor + 1 = 112$，故为 $(64, 112, 112)$；
    \item 参数量（\texttt{bias=False}）：$64 \times 3 \times 7 \times 7 = 9{,}408$；
    \item 感受野与累计步幅：$r_1 = 1 + (7 - 1)\times 1 = 7$，$J_1 = 1 \times 2 = 2$。
  \end{itemize}

  \item \textbf{MaxPool}（$k=3,\; S=2,\; P=1$）：
  \begin{itemize}
    \item 输出尺寸：$H_2 = \lfloor(112 + 2\times 1 - 3)/2\rfloor + 1 = \lfloor 111/2 \rfloor + 1 = 56$，故为 $(64, 56, 56)$；
    \item 参数量：$0$（无权重参数）；
    \item 感受野与累计步幅：$r_2 = 7 + (3 - 1)\times 2 = 11$，$J_2 = 2 \times 2 = 4$。
  \end{itemize}

  \item \textbf{Conv2}（$C_{\text{out}}=128,\; k=3,\; S=1,\; P=1,\; d=2$）：先求空洞卷积等效核尺寸 $k_{\text{eff}} = k + (k-1)(d-1) = 3 + 2\times 1 = 5$。
  \begin{itemize}
    \item 输出尺寸：$H_3 = \lfloor(56 + 2\times 1 - 5)/1\rfloor + 1 = 54$，故为 $(128, 54, 54)$（注：因 $P=1 < (k_{\text{eff}}-1)/2 = 2$，空间尺寸缩减 $2$ 像素）；
    \item 参数量（\texttt{bias=False}）：空洞不增加权重数，仍为 $128 \times 64 \times 3 \times 3 = 73{,}728$；
    \item 感受野与累计步幅：$r_3 = 11 + (5 - 1)\times 4 = 27$，$J_3 = 4 \times 1 = 4$。
  \end{itemize}
\end{enumerate}
综上，三层总参数量为 $9{,}408 + 0 + 73{,}728 = 83{,}136$，Conv2 输出每个像素在原图上的精确感受野为 $27 \times 27$。

\subsection*{习题 5.3.2\quad ResNet 瓶颈块 vs 基础块的参数效率对比}
取通道数 $C = 256$（不计 BN 与偏置，卷积单像素乘加量 MACs 在数值上恰好等于对应卷积层的权重参数量）：
\begin{enumerate}[label=(\alph*), leftmargin=1.8em]
  \item \textbf{基础块（BasicBlock：$256 \xrightarrow{3\times 3} 256 \xrightarrow{3\times 3} 256$）}：
  \[
    \text{Params}_{\text{Basic}} = 2 \times (3 \times 3 \times 256 \times 256) = 1{,}179{,}648 \approx 1.18\text{M}.
  \]
  单像素乘加量同为 $1{,}179{,}648\text{ MACs}$。

  \item \textbf{瓶颈块（Bottleneck：$256 \xrightarrow{1\times 1} 64 \xrightarrow{3\times 3} 64 \xrightarrow{1\times 1} 256$）}：
  \begin{align*}
    \text{Params}_{\text{Bottle}}
    &= (1 \times 1 \times 256 \times 64) + (3 \times 3 \times 64 \times 64) + (1 \times 1 \times 64 \times 256) \\
    &= 16{,}384 + 36{,}864 + 16{,}384 = 69{,}632 \approx 0.070\text{M}.
  \end{align*}
  单像素乘加量同为 $69{,}632\text{ MACs}$。
\end{enumerate}
\textbf{效率对比}：瓶颈块将参数量与计算量压缩为原来的 $\frac{69{,}632}{1{,}179{,}648} = \frac{17}{288} \approx 5.90\%$（即\textbf{节省 $94.10\%$，压缩约 $16.94\times$}），同时保持了相同的 $256$ 维主干道容量并额外增加了一层非线性激活，这正是 ResNet-50/101/152 得以堆叠百余层的关键工程设计。

\subsection*{习题 5.3.3\quad 一维卷积的反向传播与转置卷积}
\begin{proof}
设输入 $x \in \R^N$、卷积核 $w \in \R^k$、步幅 $S=1$ 无填充输出 $y \in \R^{N-k+1}$，前向互相关定义为 $y_i = \sum_{u=0}^{k-1} w_u x_{i+u}$（$i = 0, \dots, N-k$）。对标量损失 $\mathcal{L}(y)$ 取全微分：
\[
  d\mathcal{L} = \sum_{i=0}^{N-k} \bar{y}_i\, dy_i
  = \sum_{i=0}^{N-k} \bar{y}_i \sum_{u=0}^{k-1} (x_{i+u}\, dw_u + w_u\, dx_{i+u}).
\]
\begin{enumerate}[label=(\arabic*), leftmargin=1.8em]
  \item \textbf{权重梯度 $\bar{w}_u$}：交换第一项求和次序，读取 $dw_u$ 系数：
  \[
    \bar{w}_u = \sum_{i=0}^{N-k} \bar{y}_i\, x_{i+u}, \qquad u = 0, \dots, k-1,
  \]
  这恰好是输入 $x$ 与上游梯度 $\bar{y}$ 的有效互相关。

  \item \textbf{输入梯度 $\bar{x}_m$ 与翻转卷积核}：在第二项中作指标代换 $m = i + u$（即 $i = m - u$），合法范围要求 $0 \le m - u \le N - k$。若将 $\bar{y}$ 两端各补零 $k-1$ 格，定义扩展序列 $\tilde{y}$ 满足 $\tilde{y}_{t} = \bar{y}_{t-(k-1)}$（当 $k-1 \le t \le N-1$ 时），其余位置为 $0$，再定义空间翻转核 $\tilde{w}_v := w_{k-1-v}$（$v = 0, \dots, k-1$），则：
  \[
    \bar{x}_m = \sum_{u=0}^{k-1} \bar{y}_{m-u}\, w_u
    = \sum_{v=0}^{k-1} \bar{y}_{m - (k-1-v)}\, w_{k-1-v}
    = \sum_{v=0}^{k-1} \tilde{y}_{m+v}\, \tilde{w}_v, \qquad m = 0, \dots, N-1.
  \]
  右端恰好是将 $\bar{y}$ 两端补零 $k-1$ 格后，与翻转核 $\tilde{w} = \operatorname{flip}(w)$ 所做的标准一维互相关（即所谓“满卷积”或转置卷积）。
\end{enumerate}
\end{proof}

\subsection*{习题 5.3.4\quad 带缩放残差分支的方差增长与 $\frac{1}{\sqrt{L}}$ 缩放}
\begin{proof}
\begin{enumerate}[label=(\alph*), leftmargin=1.8em]
  \item \textbf{指数增长递推}：由 $x_{l+1} = x_l + \alpha\,\mathcal{F}_l(x_l)$ 及 $x_l$ 与 $\mathcal{F}_l(x_l)$ 不相关（$\Cov(x_l, \mathcal{F}_l(x_l)) = 0$），取方差得：
  \[
    \Var(x_{l+1}) = \Var(x_l) + \alpha^2 \Var(\mathcal{F}_l(x_l))
    = (1 + \alpha^2)\,\Var(x_l).
  \]
  从 $l = 0$ 递推至 $l = L$ 即证：
  \[
    \Var(x_L) = (1 + \alpha^2)^L \Var(x_0).
  \]

  \item \textbf{取 $\alpha = 1/\sqrt{L}$ 的收敛性}：当 $\alpha = 1$ 且 $L = 100$ 时，$(1 + 1)^{100} = 2^{100} \approx 1.27 \times 10^{30}$，前向信号必然溢出。若取 $\alpha = \frac{1}{\sqrt{L}}$，则 $\alpha^2 = \frac{1}{L}$，代入得：
  \[
    \Var(x_L) = \left(1 + \frac{1}{L}\right)^L \Var(x_0).
  \]
  由于数列 $(1 + 1/L)^L$ 严格单调递增且收敛于自然常数 $e \approx 2.7183$（例如 $L=100$ 时为 $2.7048$），因此对任意深度 $L \ge 1$ 均有 $\Var(x_L) < e\,\Var(x_0)$，彻底消除了方差随深度的指数爆炸。GPT-2 每个 Transformer 块含注意力与 MLP 两个残差子层，共 $2 N_{\text{layers}}$ 条残差支路，故对其输出投影矩阵按 $1/\sqrt{2 N_{\text{layers}}}$ 缩放，正是这一结论的直接工程落地。
\end{enumerate}
\end{proof}

\subsection*{习题 5.3.5\quad Post-Act 与 Pre-Act ResNet 的 50 层梯度传播}
扩展 \texttt{exp3\_resnet\_grad\_highway} 至 $L = 50$ 层（$D = 64$），需特别区分两种权重初始化尺度下的行为：
\begin{enumerate}[label=(\arabic*), leftmargin=1.8em]
  \item \textbf{压缩初始化（如原脚本 $\sigma_w = 0.65 / \sqrt{D}$ 配合 $\tanh$ / ReLU）}：
  \begin{itemize}
    \item \textbf{Plain 网络}：雅可比谱范数 $< 1$，经过 50 层连乘后梯度衰减超 10 个数量级（接近机器零）；
    \item \textbf{Post-Act ResNet}（外层套 $\ReLU$）：主干道每经过一层就被外层 ReLU 的导数 $\mathbb{I}(z > 0)$ 截断约一半负值通道，50 层后仍出现明显指数衰减；
    \item \textbf{Pre-Act ResNet}（洁净恒等加法 $h_l = h_{l-1} + \mathcal{F}(h_{l-1})$）：反向雅可比保留完整常数项 $I$，第 1 层与第 50 层梯度范数维持在同量级 $O(1)$。
  \end{itemize}

  \item \textbf{纯 Kaiming 初始化（$\sigma_w = \sqrt{2/D}$）配固定支路系数 $0.35$（未加 BN/LN）}：在无归一化与无 $1/\sqrt{L}$ 缩放时，每层残差支路方差约 $0.35^2 \approx 0.1225$，由习题 5.3.4 可知 50 层前向方差膨胀 $(1 + 0.1225)^{50} \approx 324$ 倍，从而在 \texttt{Pre-Act}（$\norm{\partial \mathcal{L}/\partial h_1} / \norm{\partial \mathcal{L}/\partial h_{50}} \approx 1.47\times 10^4$）与 \texttt{Post-Act}（$\approx 1.62\times 10^4$）上均呈现方差放大；一旦为残差支路加上 $1/\sqrt{50} \approx 0.141$ 缩放或 Pre-LN，Pre-Act 即展现出远优于 Post-Act 的洁净等幅梯度流。
\end{enumerate}

\subsection*{习题 5.3.6\quad 手工实现单步 LSTM 与长距离复制任务}
\textbf{核心实现（合并四门矩阵乘以提速）}：将 $[x_t, h_{t-1}]$ 拼接为维度 $d_{\text{in}} + d$ 的向量，经单个 \texttt{nn.Linear(d\_in + d, 4*d)} 投影后用 \texttt{.chunk(4, dim=-1)} 切分为四门：
\begin{lstlisting}
class MyLSTMCell(nn.Module):
    def __init__(self, d_in, d_hid, forget_bias=0.0):
        super().__init__()
        self.proj = nn.Linear(d_in + d_hid, 4 * d_hid)
        nn.init.constant_(self.proj.bias[d_hid:2*d_hid], forget_bias)

    def forward(self, x, state):
        h_prev, c_prev = state
        gates = self.proj(torch.cat([x, h_prev], dim=-1))
        i, f, g, o = gates.chunk(4, dim=-1)  # 顺序与 PyTorch 一致: i, f, g, o
        c = torch.sigmoid(f) * c_prev + torch.sigmoid(i) * torch.tanh(g)
        h = torch.sigmoid(o) * torch.tanh(c)
        return h, c
\end{lstlisting}
\textbf{与 \texttt{nn.LSTMCell} 精度互验}：PyTorch 内部将输入权与隐态权分别存为 \texttt{weight\_ih}、\texttt{weight\_hh}（按 \texttt{i, f, g, o} 排列），将二者沿列拼接赋给 \texttt{self.proj.weight}，偏置相加赋给 \texttt{self.proj.bias}。实测前向输出的 $h_t, c_t$ 最大逐元素绝对误差仅为 $5.96 \times 10^{-8} < 10^{-6}$。

\textbf{首词记忆任务（$T = 40$）预期对比}：
\begin{itemize}
  \item \textbf{Elman RNN}：梯度跨 40 步经 $\prod \diag(1-h_t^2) W_{hh}^\top$ 指数坍缩，训练损失长期停滞在随机猜想基线 $\ln K$ 附近，难以收敛；
  \item \textbf{MyLSTMCell（$b_f = 0$）}：初始遗忘门 $\sigma(0) = 0.5$，跨 40 步记忆衰减为 $0.5^{40} \approx 10^{-12}$，早期需要较长热身才能学会抬高遗忘门；
  \item \textbf{MyLSTMCell（$b_f = +2.0$）}：初始遗忘门 $\sigma(2) \approx 0.881$，细胞状态高速公路从第 1 步起即畅通传递远距梯度，收敛速度比 $b_f = 0$ 快数倍并迅速达到接近 $100\%$ 准确率（这正是 Jozefowicz et al., 2015 推荐将遗忘门偏置初始化为 $+1 \sim +2$ 的理论依据）。
\end{itemize}

\subsection*{习题 5.3.7\quad 旋转等变卷积（Group Equivariant CNN）}
设离散平面旋转平移群 $p4 = \mathbb{Z}^2 \rtimes C_4$，其每个群元 $g = (t, r)$ 由平移 $t \in \mathbb{Z}^2$ 与旋转 $r \in \{0^\circ, 90^\circ, 180^\circ, 270^\circ\}$ 构成。
\begin{enumerate}[label=(\arabic*), leftmargin=1.8em]
  \item \textbf{第一层（提升卷积，Lifting Convolution）}：输入图像 $f: \mathbb{Z}^2 \to \R^{C_{\text{in}}}$ 仅定义在平面网格上。给定单个 $3\times 3$ 卷积核 $\psi$，将其旋转 $4$ 个角度得到 $\{\psi_r\}_{r \in C_4}$（完全共享同一份 $9$ 个参数），分别与图像做标准 2D 互相关，输出一个定义在群 $p4$ 上的 4 通道组特征图 $F(t, r) = (f * \psi_r)(t)$。当输入图像旋转 $90^\circ$ 时，输出的 4 张特征图不仅自身旋转 $90^\circ$，还在 4 个角度通道间发生循环置换——严格满足群等变性。
  \item \textbf{深层（群卷积，Group Convolution）}：对定义在 $p4$ 上的特征图 $F(h)$ 与群卷积核 $\Psi(h)$（维度 $4 \times 3 \times 3$），定义群卷积：
  \[
    (F *_{p4} \Psi)(g) = \sum_{h \in p4} F(h)\,\Psi(g^{-1} h).
  \]
  验证左作用等变性：设群平移算子 $(L_u F)(h) := F(u^{-1} h)$，作变量替换 $h' = u^{-1} h$（则 $h = u h'$ 且 $g^{-1} h = (u^{-1} g)^{-1} h'$）：
  \[
    \big((L_u F) *_{p4} \Psi\big)(g)
    = \sum_{h \in p4} F(u^{-1} h)\,\Psi(g^{-1} h)
    = \sum_{h' \in p4} F(h')\,\Psi\big((u^{-1} g)^{-1} h'\big)
    = \big(L_u(F *_{p4} \Psi)\big)(g).
  \]
  \item \textbf{参数效率}：通过角度间权值共享，G-CNN 把原本需要靠数据增强和 $4\times$ 冗余卷积核才能死记硬背的四个方向模板，压缩为同一组核参数，在病理切片、遥感与分子点云等无绝对朝向的数据上大幅提升样本效率。
\end{enumerate}

\subsection*{习题 5.3.8\quad 线性递归的并行前缀扫描原理}
\begin{proof}
对任意三组数对 $(a_3, b_3),\; (a_2, b_2),\; (a_1, b_1)$，分别计算左右两种结合顺序：
\begin{align*}
  \big((a_3, b_3) \bullet (a_2, b_2)\big) \bullet (a_1, b_1)
  &= (a_3 a_2,\; a_3 b_2 + b_3) \bullet (a_1, b_1) \\
  &= \big((a_3 a_2) a_1,\; (a_3 a_2) b_1 + (a_3 b_2 + b_3)\big)
   = (a_3 a_2 a_1,\; a_3 a_2 b_1 + a_3 b_2 + b_3), \\
  (a_3, b_3) \bullet \big((a_2, b_2) \bullet (a_1, b_1)\big)
  &= (a_3, b_3) \bullet (a_2 a_1,\; a_2 b_1 + b_2) \\
  &= \big(a_3(a_2 a_1),\; a_3(a_2 b_1 + b_2) + b_3\big)
   = (a_3 a_2 a_1,\; a_3 a_2 b_1 + a_3 b_2 + b_3).
\end{align*}
两者恒等，故 $\bullet$ 严格满足\textbf{结合律}（更直观地，每个 $(a_t, b_t)$ 同构于上三角仿射矩阵 $\begin{psmallmatrix} a_t & b_t \\ 0 & 1 \end{psmallmatrix}$，$\bullet$ 即矩阵乘法，天然满足结合律）。
\end{proof}
\textbf{为何实现 $O(\log_2 T)$ 并行深度，而 $\tanh$ RNN 不能}：记初始状态对应 $(1, h_0)$，则任意时刻状态恰为前缀积 $(a_t, b_t) \bullet \cdots \bullet (a_1, b_1) \bullet (1, h_0)$ 的第二分量。因为结合律允许自由重排括号，Blelloch 并行前缀扫描（Parallel Scan）可在二叉树上先用 $\lceil \log_2 T \rceil$ 层上行规约、再用 $\lceil \log_2 T \rceil$ 层下行分发，以总工作量 $O(T)$、并行延迟仅 $O(\log_2 T)$ 在 GPU 上一次算出全部 $h_1, \dots, h_T$。相反，Elman RNN 中的非线性函数 $h_t = \tanh(W h_{t-1} + u_t)$ 在复合两步后变为 $\tanh(W \tanh(W h_{t-2} + u_{t-1}) + u_t)$，无法用有限维封闭参数表征任意步复合（即不构成封闭的结合半群），只能严格按 $t = 1, 2, \dots, T$ 串行计算。

% =========================================================================
\section*{第 4 章习题解答（共 8 题）}
\addcontentsline{toc}{section}{第 4 章习题解答（共 8 题）}
% =========================================================================

\subsection*{习题 5.4.1\quad 因果缩放点积注意力完整手算}
先计算未缩放点积 $Q K^\top$ 与缩放 logits $Z = \frac{1}{\sqrt{2}} Q K^\top$（记 $\alpha := \frac{1}{\sqrt{2}} \approx 0.7071$，则 $e^{\alpha} \approx 2.0281,\; e^{\sqrt{2}} \approx 4.1133$）：
\[
  Q K^\top = \begin{pmatrix} 1 & 1 & 0 \\ 0 & 1 & 1 \\ 1 & 2 & 1 \end{pmatrix}, \qquad
  Z = \frac{1}{\sqrt{2}} Q K^\top = \begin{pmatrix}
    \alpha & \alpha & 0 \\
    0 & \alpha & \alpha \\
    \alpha & 2\alpha & \alpha
  \end{pmatrix}
  \approx \begin{pmatrix}
    0.7071 & 0.7071 & 0 \\
    0 & 0.7071 & 0.7071 \\
    0.7071 & 1.4142 & 0.7071
  \end{pmatrix}.
\]
\begin{enumerate}[label=(\arabic*), leftmargin=1.8em]
  \item \textbf{无因果掩码（双向注意力）}：逐行计算 $\softmax$：
  \begin{itemize}
    \item 第 1 行分母 $2 e^{\alpha} + 1 \approx 5.0562$，权重 $(0.4011,\; 0.4011,\; 0.1978)$；
    \item 第 2 行对称得权重 $(0.1978,\; 0.4011,\; 0.4011)$；
    \item 第 3 行将整行同减 $\alpha$ 得 logits $(0, \alpha, 0)$，分母 $e^{\alpha} + 2 \approx 4.0281$，权重 $(0.2483,\; 0.5035,\; 0.2483)$。
  \end{itemize}
  乘以 $V = \begin{psmallmatrix} 2 & 0 \\ 0 & 2 \\ 1 & 1 \end{psmallmatrix}$（保留三位小数）：
  \[
    P_{\text{full}} \approx \begin{pmatrix}
      0.401 & 0.401 & 0.198 \\
      0.198 & 0.401 & 0.401 \\
      0.248 & 0.503 & 0.248
    \end{pmatrix}, \qquad
    O_{\text{full}} = P_{\text{full}} V \approx \begin{pmatrix}
      1.000 & 1.000 \\
      0.797 & 1.203 \\
      0.745 & 1.255
    \end{pmatrix}.
  \]
  （注意第 1 行输出精确为 $(2p_{11}+p_{13},\; 2p_{12}+p_{13}) = (1, 1)$，因 $p_{11}=p_{12}$ 且 $2p_{11}+p_{13}=1$！）

  \item \textbf{带因果掩码（式~\eqref{eq:causal-mask}，令 $j > i$ 处为 $-\infty$）}：
  \begin{itemize}
    \item 第 1 行仅允许 $j=1$：$P_{1,:} = (1,\; 0,\; 0)$，输出 $O_{1,:} = v_1 = (2.000,\; 0.000)$；
    \item 第 2 行 logits 为 $(0, \alpha, -\infty)$：分母 $1 + e^{\alpha} \approx 3.0281$，权重 $(0.3302,\; 0.6698,\; 0)$，输出 $(2\times 0.3302,\; 2\times 0.6698) \approx (0.660,\; 1.340)$；
    \item 第 3 行所有 $j \le 3$ 均可见，与无掩码完全一致：权重 $(0.248,\; 0.503,\; 0.248)$，输出 $(0.745,\; 1.255)$。
  \end{itemize}
  \[
    P_{\text{causal}} \approx \begin{pmatrix}
      1.000 & 0.000 & 0.000 \\
      0.330 & 0.670 & 0.000 \\
      0.248 & 0.503 & 0.248
    \end{pmatrix}, \qquad
    O_{\text{causal}} = P_{\text{causal}} V \approx \begin{pmatrix}
      2.000 & 0.000 \\
      0.660 & 1.340 \\
      0.745 & 1.255
    \end{pmatrix}.
  \]
\end{enumerate}
\textbf{第 1 行对比}：无掩码时第 1 行“偷看”了未来词元 $2, 3$，将 $v_1=(2,0)$ 混合成了 $(1.000, 1.000)$；加入因果掩码后第 1 行严格只能看见自身，权重坍缩为 $(1,0,0)$，输出恒等于 $v_1 = (2.000, 0.000)$，从根源阻断了未来信息泄露。

\subsection*{习题 5.4.2\quad 70B 大模型的参数结构与 KV Cache 精算}
已知 $L = 80,\; d = 8192,\; n_q = 64,\; n_{\text{kv}} = 8,\; d_h = 128,\; d_{\text{ff}} = 28672,\; |V| = 128{,}000$（不共享输入/输出嵌入）。
\begin{enumerate}[label=(\alph*), leftmargin=1.8em]
  \item \textbf{总可训练参数量精算}（略去极少量的 RMSNorm 缩放向量 $161 d \approx 1.32\text{M}$）：
  \begin{itemize}
    \item 单层 GQA 注意力投影（$n_{\text{kv}} d_h = 8 \times 128 = 1024$）：
    \[
      P_{\text{attn}} = \underbrace{d^2}_{W_Q} + \underbrace{d(n_{\text{kv}} d_h)}_{W_K} + \underbrace{d(n_{\text{kv}} d_h)}_{W_V} + \underbrace{d^2}_{W_O}
      = 8192 \times (8192 + 1024 + 1024 + 8192) = 150{,}994{,}944.
    \]
    \item 单层 SwiGLU 前馈网络（门控 $W_g$、上投影 $W_u$、下投影 $W_d$ 共 3 矩阵）：
    \[
      P_{\text{ffn}} = 3 \times d \times d_{\text{ff}} = 3 \times 8192 \times 28672 = 704{,}643{,}072.
    \]
    \item 单层合计 $P_{\text{layer}} = 855{,}638{,}016$，80 层合计 $80 \times P_{\text{layer}} = 68{,}451{,}041{,}280$（约 $68.45\text{B}$）。
    \item 独立输入嵌入与输出头：$P_{\text{emb}} = 2 \times |V| \times d = 2 \times 128{,}000 \times 8192 = 2{,}097{,}152{,}000$（约 $2.10\text{B}$）。
    \item \textbf{总参数量}：$P_{\text{total}} = 68{,}451{,}041{,}280 + 2{,}097{,}152{,}000 = 70{,}548{,}193{,}280 \approx \mathbf{70.55\text{B}}$（这正是 Llama-3 70B 的精确架构账本！）。
  \end{itemize}

  \item \textbf{BF16 下 KV Cache 显存占用（$T = 32{,}768,\; B = 16,\; b_{\text{elem}} = 2$）}：代入式~\eqref{eq:kv-cache-mem}（$1\text{ GiB} = 2^{30}\text{ B}$）：
  \begin{align*}
    \text{Mem}_{\mathrm{KV}}^{\mathrm{GQA}}
    &= 2 \times 80 \times 16 \times 32{,}768 \times 8 \times 128 \times 2\text{ B} \\
    &= 160 \times 2^{4} \times 2^{15} \times 2^{10} \times 2^{1}\text{ B}
     = 160 \times 2^{30}\text{ B} = \mathbf{160.00\text{ GiB}}.
  \end{align*}
  若改用标准 MHA（$n_{\text{kv}} = n_q = 64$，增为 $8$ 倍），则显存暴涨至：
  \[
    \text{Mem}_{\mathrm{KV}}^{\mathrm{MHA}} = 8 \times 160.00\text{ GiB} = \mathbf{1280.00\text{ GiB}} = 1.25\text{ TiB}.
  \]
  相比之下，模型自身 BF16 权重仅约 $131\text{ GiB}$。GQA 凭借 $8\times$ 缓存压缩，使得 32K 长上下文多并发推理能够在 4--8 卡 80GB GPU 节点上落地，而标准 MHA 在同并发下早已彻底显存爆炸。
\end{enumerate}

\subsection*{习题 5.4.3\quad Softmax 雅可比半正定性与缩放点积注意力 VJP}
\begin{proof}
\begin{enumerate}[label=(\alph*), leftmargin=1.8em]
  \item \textbf{对称半正定性与零空间}：显然 $J_{\softmax} = \diag(p) - p p^\top$ 对称。对任意 $u \in \R^T$，将二次型改写为以概率分布 $p$ 为测度的随机变量方差：
  \[
    u^\top J_{\softmax} u
    = \sum_{j=1}^T p_j u_j^2 - \left(\sum_{j=1}^T p_j u_j\right)^2
    = \E_{j\sim p}[u_j^2] - (\E_{j\sim p}[u_j])^2
    = \Var_{j\sim p}(u_j) \ge 0.
  \]
  （也可直接对向量 $(\sqrt{p_j} u_j)$ 与 $(\sqrt{p_j})$ 施用 Cauchy--Schwarz 不等式：$\big(\sum_j p_j u_j\big)^2 \le \big(\sum_j p_j u_j^2\big)\big(\sum_j p_j\big) = \sum_j p_j u_j^2$。）

  由于 Softmax 输出严格满足 $p_j > 0$（对所有 $j$），等号成立当且仅当 $\Var_{j\sim p}(u_j) = 0$，即所有分量相等：$u_1 = u_2 = \dots = u_T$，等价于 $u \in \operatorname{span}(\mathbf{1})$。又半正定矩阵满足 $u^\top J u = 0 \iff J u = 0$，因此 $\ker(J_{\softmax}) = \operatorname{span}(\mathbf{1})$。

  \item \textbf{前向平移不变性与反向行和为零}：给第 $i$ 行 logits 同加常数 $c$（$z \mapsto z + c\mathbf{1}$），分子分母同乘 $e^c$ 约去，故前向概率 $p$ 不变。反向传播中，第 $i$ 行上游伴随量 $\bar{p}_i$ 通过式~\eqref{eq:attn-softmax-vjp} 回传为 $\bar{Z}_{i,:}^\top = J_{\softmax}(p_i)^\top \bar{p}_i = J_{\softmax}(p_i)\bar{p}_i$。左乘全一行向量 $\mathbf{1}^\top$ 并利用 $\mathbf{1} \in \ker(J_{\softmax})$：
  \[
    \sum_{j=1}^T \bar{Z}_{ij}
    = \mathbf{1}^\top \bar{Z}_{i,:}^\top
    = \mathbf{1}^\top J_{\softmax}(p_i)\,\bar{p}_i
    = \big(J_{\softmax}(p_i)\mathbf{1}\big)^\top \bar{p}_i
    = 0^\top \bar{p}_i = 0.
  \]
  可见，前向对平移 $\mathbf{1}$ 的不变性在反向传播中精确对偶为行梯度和恒为零（亦保证了数值稳定技巧 $z_j \leftarrow z_j - \max_k z_k$ 对前后向均无任何偏差）。
\end{enumerate}
\end{proof}

\subsection*{习题 5.4.4\quad RoPE 编码的相对距离远程衰减上界}
\begin{proof}
记第 $r$ 个二维子平面上的复数乘积为 $h_r = \tilde{q}_r^* \tilde{k}_r$，相对位置 $\Delta = n - m$，部分和 $S_r := \sum_{j=1}^r e^{\mathrm{i}\Delta\theta_j}$（约定 $S_0 = 0$），则 $e^{\mathrm{i}\Delta\theta_r} = S_r - S_{r-1}$。
\begin{enumerate}[label=(\arabic*), leftmargin=1.8em]
  \item \textbf{Abel 分部求和与上界推导}：由定理~\ref{thm:rope}，RoPE 相对内积等于复和的实部：
  \begin{align*}
    \inner{R_\Theta(m)q}{R_\Theta(n)k}
    &= \operatorname{Re}\!\left(\sum_{r=1}^K h_r e^{\mathrm{i}\Delta\theta_r}\right)
     = \operatorname{Re}\!\left(\sum_{r=1}^K h_r (S_r - S_{r-1})\right) \\
    &= \operatorname{Re}\!\left(h_K S_K - \sum_{r=1}^{K-1} (h_{r+1} - h_r) S_r\right).
  \end{align*}
  取绝对值并依次利用 $|\operatorname{Re}(z)| \le |z|$ 与三角不等式：
  \begin{align*}
    \abs{\inner{R_\Theta(m)q}{R_\Theta(n)k}}
    &\le |h_K|\,|S_K| + \sum_{r=1}^{K-1} |h_{r+1} - h_r|\,|S_r| \\
    &\le |h_K|\,|S_K| + \left(\max_{1\le r \le K-1} |h_{r+1} - h_r|\right)\sum_{r=1}^{K-1} |S_r|.
  \end{align*}
  若约定 $h_{K+1} := 0$，可将边界项一并收入 $\big(\max_{1\le r \le K} |h_{r+1} - h_r|\big)\sum_{r=1}^K |S_r|$。

  \item \textbf{几何频率分布导致远程振荡衰减的机理}：当 $\Delta = 0$ 时所有复指数同相为 $1$，$|S_r(0)| = r$，上界系数达到峰值 $\sum_{r=1}^K r = \frac{K(K+1)}{2}$。当 $|\Delta|$ 增大时，几何级数频率 $\theta_r = 10000^{-2(r-1)/d_h}$ 在 $[10^{-4}, 1]$ 跨四个数量级展开，相邻相位 $\Delta\theta_r$ 迅速分散并绕单位圆多圈旋转，产生剧烈的\textbf{相消干涉}，使 $|S_r(\Delta)| \ll r$。以 $d_h = 64\; (K = 32)$ 数值实测为例，$\sum_{r=1}^{32}|S_r(\Delta)|$ 随距离呈显著下降趋势：
  \[
    \Delta=0:\; 528.0 \;\to\;
    \Delta=5:\; 333.7 \;\to\;
    \Delta=20:\; 256.3 \;\to\;
    \Delta=100:\; 238.9 \;\to\;
    \Delta=500:\; 97.4.
  \]
  这就是 RoPE 赋予 Transformer 温和“局部优先”归纳偏置的数学根源。
\end{enumerate}
\end{proof}

\subsection*{习题 5.4.5\quad 无 $\sqrt{d_k}$ 缩放消融实验}
在 \texttt{MiniGPT}（$d_{\text{model}} = 256,\; H = 2,\; d_h = 128$）上分别开启（除以 $\sqrt{128} \approx 11.31$）与关闭缩放训练序列反转任务，前 100 步实测对比如下：
\begin{center}
\footnotesize
\begin{tabularx}{\linewidth}{@{}lcccX@{}}
\toprule
实验配置 & 第 0 步 $H(P)$ & 前 100 步均值 & 第 0 步 $\norm{\nabla_{W_Q}\mathcal{L}}_F$ & 行为特征 \\
\midrule
开启 $1/\sqrt{d_h}$ 缩放 & $1.624$ & $0.508$ & $3.97 \times 10^{-2}$ & 分布平滑，多词元并行聚合 \\
关闭缩放（无 $\sqrt{128}$） & $0.571$ & $0.098$ & $3.77 \times 10^{-1}$ & 初始即逼近单热尖峰，熵骤降 $65\%$ \\
\bottomrule
\end{tabularx}
\end{center}
\textbf{结果分析}：未缩放时，点积方差比缩放组放大了 $d_h = 128$ 倍（标准差放大 $11.3\times$），第 0 步注意力熵便从接近均匀分布的 $1.624$ 坍缩至 $0.571$（近乎单热）；在训练早期未缩放的 $q_i^\top k_j$ 幅度大且无 $1/\sqrt{d_h}$ 阻尼，一方面少数未完全饱和的行产生尖刺梯度震荡，另一方面一旦某些位置 logits 差距超过 $15$，Softmax 雅可比 $p_j(\delta_{jk}-p_k)$ 便彻底进入饱和区使梯度消失。当序列更长（如 $T = 512$）或层数更深时，关闭缩放会迅速引发训练发散或停滞，与定理~\ref{thm:sqrt-dk} 严丝合缝。

\subsection*{习题 5.4.6\quad 为手写 \texttt{MiniGPT} 添加 RoPE 与增量 KV Cache}
\begin{enumerate}[label=(\alph*), leftmargin=1.8em]
  \item \textbf{$O(d_h)$ 逐元素 RoPE 实现与验证}：对张量 $q, k \in \R^{B \times H \times T \times d_h}$，预计算频率 $\theta_r = 10000^{-2r/d_h}$ 以及 $\cos(m\theta_r)$ 与 $\sin(m\theta_r)$，将最后一维拆为偶奇两半做二维旋转：
  \begin{lstlisting}
def apply_rope(x, cos, sin):
    # x: (B, H, T, dh); cos, sin: (1, 1, T, dh//2)
    x1, x2 = x[..., 0::2], x[..., 1::2]
    out1 = x1 * cos - x2 * sin
    out2 = x1 * sin + x2 * cos
    return torch.stack([out1, out2], dim=-1).flatten(-2)
  \end{lstlisting}
  验证相对位置不变性：任取位置 $m, n$ 及任意平移量 $\delta$，比较
  \[
    \inner{\mathrm{RoPE}(q, m)}{\mathrm{RoPE}(k, n)}
    \quad\text{与}\quad
    \inner{\mathrm{RoPE}(q, m+\delta)}{\mathrm{RoPE}(k, n+\delta)},
  \]
  两者在 FP32 下最大绝对误差小于 $10^{-6}$（FP64 下 $< 10^{-14}$）。

  \item \textbf{增量 KV Cache 推理与一致性检验}：为 \texttt{CausalSelfAttn.\allowbreak forward} 新增可选参数 \texttt{kv\_cache} 与位置偏移 \texttt{pos\_offset}：
  \begin{enumerate}[label=(\roman*), leftmargin=1.6em]
    \item 计算当前输入 $x$（长度 $T_{\text{new}}$）的 $q, k, v$，使用位置索引 \texttt{pos\_offset .. pos\_offset + T\_new - 1} 对 $q, k$ 施加 RoPE（切忌把增量步位置误设为 $0$！）；
    \item 若 \texttt{kv\_cache} 非空，沿序列维将新 $k, v$ 拼接到历史缓存后得到 $K_{\text{all}}, V_{\text{all}}$ 并返回更新后的缓存；
    \item 当 $T_{\text{new}} = 1$（单步解码）时，当前查询位于序列最末，可合法看见全部历史键，无需施加掩码；当 $T_{\text{new}} > 1$（预填充阶段）时照常施加下三角因果掩码。
  \end{enumerate}
  固定评估模式（\texttt{model.eval()}），对比“输入全前缀 $x_{0:t}$”与“逐步传入 $x_t$ + 缓存”所得最后一步 logits，最大绝对误差通常在 $10^{-6}$ 量级（远小于 $10^{-5}$ 验收线）。
\end{enumerate}

\subsection*{习题 5.4.7\quad 长上下文外推：位置内插（PI）与 YaRN}
从定理~\ref{thm:rope} 的“多频时钟表盘”视角看，RoPE 的第 $r$ 个二维子平面对应角频率 $\theta_r = 10000^{-2(r-1)/d_h}$、波长 $\lambda_r = 2\pi / \theta_r$：
\begin{itemize}
  \item \textbf{高频维度（小 $r$，秒针）}：$\theta_1 = 1,\; \lambda_1 \approx 6.28$ 个词元。在长度 $T_{\text{train}} = 4096$ 内已旋转 $650$ 余圈，模型早已见过所有相位角，负责精确分辨紧邻词元的前后顺序；
  \item \textbf{低频维度（大 $r$，时针）}：$\theta_K \approx 10^{-4},\; \lambda_K \approx 62{,}832$ 个词元。在 $T_{\text{train}} = 4096$ 内连一圈的 $1/15$ 都未走完。直接外推到 $T_{\text{test}} = 16384$ 时，低频相位首次闯入训练从未见过的分布外角度，导致注意力分布坍塌。
\end{itemize}
\textbf{均匀位置内插（PI）的缺陷}：将所有维度统一按 $m \to m / s$（此处 $s = 4$）压缩，虽把低频角度拉回训练区间，却同步把高频秒针的步进角缩小了 $4$ 倍——相邻两个词元 $m$ 与 $m+1$ 之间的相位差被压扁为原来的 $1/4$，破坏了局部高分辨率词序感知，需要大量微调才能部分恢复。

\textbf{NTK-Aware 与 YaRN 的分区缩放优势}：
\begin{itemize}
  \item \textbf{NTK-Aware} 通过增大底数（如 $10000 \to 10000 \cdot s^{d_h/(d_h-2)}$）自动减小低频角速度而几乎不改变高频 $\theta_1 \approx 1$；
  \item \textbf{YaRN} 进一步按每个维度在训练长度内转过的圈数 $r_d = T_{\text{train}} / \lambda_d$ 划定三个频段：高频段（$r_d > \beta$）\textbf{完全不插值}以保留精确局部词序，低频段（$r_d < \alpha$）\textbf{完全按 $1/s$ 内插}以防止分布外相位溢出，中频段平滑线性过渡，并辅以注意力温度修正因子 $\sqrt{1/t} = 0.1\ln s + 1$ 补偿长序列熵增。由此实现了极低微调成本下 $4\times \sim 16\times$ 的稳健上下文外推。
\end{itemize}

\subsection*{习题 5.4.8\quad 归纳头与两层残差流符号电路}
\begin{enumerate}[label=(\arabic*), leftmargin=1.8em]
  \item \textbf{为何单层注意力难以实现复杂多跳复制与序列反转}：在单层自注意力中，第 $t$ 个位置的查询 $q_t$ 只能直接与历史位置 $j \le t$ 的词元表征 $k_j$ 做一次内积匹配。它无法完成“\textbf{间接寻址}”（例如：找满足“其前驱词元等于当前词元 \texttt{A}”的位置 $j$，或者在可变长序列反转中动态配对对称位置）。要根据邻居内容或动态相对偏移定位目标，至少需要两次顺序信息路由。
  \item \textbf{两层残差流如何构成可编程符号电路}：把残差流 $x \in \R^{T \times d}$ 视作所有层共享读写的“高维通信总线”，各子空间正交分区存放不同寄存器字段：
  \begin{itemize}
    \item \textbf{第 1 层（写入上下文标签）}：第 1 层的前词头（Previous-Token Head）将位置 $j-1$ 的词元信息（或相对分隔符的距离编码）搬运并写入位置 $j$ 残差流的专属子空间，相当于为每个位置 $j$ 贴上了元数据标签“\texttt{prev\_tok = A}”；
    \item \textbf{第 2 层（基于新标签执行第二跳查询）}：第 2 层的匹配头（Induction Head）用当前词元 \texttt{A} 构造查询 $q_t$，专门读取第 1 层刚刚写入残差流的那段标签子空间作为键 $k_j$，从而精准锁定所有前驱为 \texttt{A} 的历史位置 $j$，将其值 $v_j = \texttt{B}$ 复制到当前位置完成补全 \texttt{[A][B] ... [A] -> [B]}。
  \end{itemize}
  \item \textbf{与第 4 章迷你 GPT 相变现象的呼应}：在第 4 章的 2 层迷你 GPT 反转实验中，前 20--40 步两层尚未建立“第 1 层编码相对分隔符偏移 $\to$ 第 2 层按镜像偏移精准抓取源词元”的跨层配合，整序精确率几乎停留在 $0\%$；一旦这两层电路通过梯度协同对齐，验证集整序精确率便在短短数十步内从接近 $0\%$ 陡峭跃迁至 $100\%$——这正是两层注意力电路形成“相变涌现”的缩影。
\end{enumerate}

% =========================================================================
\section*{第 5 章习题解答（共 8 题）}
\addcontentsline{toc}{section}{第 5 章习题解答（共 8 题）}
% =========================================================================

\subsection*{习题 5.5.1\quad BPE 分词手算与新词切分}
微型语料为：\texttt{hug}(10)、\texttt{pug}(5)、\texttt{pun}(12)、\texttt{bun}(4)、\texttt{hugs}(5)（词尾加 \texttt{</w>}）。
\begin{enumerate}[label=(\alph*), leftmargin=1.8em]
  \item \textbf{初始字符词表}（共 8 个基础符号）：
  \[
    V_0 = \{\texttt{b},\; \texttt{g},\; \texttt{h},\; \texttt{n},\; \texttt{p},\; \texttt{s},\; \texttt{u},\; \texttt{</w>}\}.
  \]
  初始分词状态：\texttt{h u g </w>}(10)、\texttt{p u g </w>}(5)、\texttt{p u n </w>}(12)、\texttt{b u n </w>}(4)、\texttt{h u g s </w>}(5)。
  \textbf{前 4 轮合并追踪}（遇到同频并列时按语料首次出现/字典序优先破平，并注明并列项）：
  \begin{enumerate}[label=\textbf{第 \arabic* 轮：}, leftmargin=4.2em]
    \item 统计高频对：\texttt{(u, g)} 出现 $10+5+5=20$ 次，\texttt{(p, u)} 出现 $17$ 次，\texttt{(u, n)} 出现 $16$ 次。合并 \texttt{(u, g) -> ug}（频次 \textbf{20}）。
    更新后：\texttt{h ug </w>}(10)、\texttt{p ug </w>}(5)、\texttt{p u n </w>}(12)、\texttt{b u n </w>}(4)、\texttt{h ug s </w>}(5)。
    \item 统计高频对：\texttt{(u, n)} 与 \texttt{(n, </w>)} 并列最高，均为 $12+4=16$ 次；其次 \texttt{(h, ug)} 与 \texttt{(ug, </w>)} 为 $15$ 次。取 \texttt{(u, n) -> un}（频次 \textbf{16}）。
    更新后：\texttt{h ug </w>}(10)、\texttt{p ug </w>}(5)、\texttt{p un </w>}(12)、\texttt{b un </w>}(4)、\texttt{h ug s </w>}(5)。
    \item 统计高频对：\texttt{(un, </w>)} 为 $12+4=16$ 次，高于 \texttt{(h, ug)}(15) 与 \texttt{(ug, </w>)}(15)。合并 \texttt{(un, </w>) -> un</w>}（频次 \textbf{16}）。
    更新后：\texttt{h ug </w>}(10)、\texttt{p ug </w>}(5)、\texttt{p un</w>}(12)、\texttt{b un</w>}(4)、\texttt{h ug s </w>}(5)。
    \item 统计高频对：\texttt{(h, ug)}（$10+5=15$）与 \texttt{(ug, </w>)}（$10+5=15$）并列最高，其次 \texttt{(p, un</w>)}(12)。取 \texttt{(h, ug) -> hug}（频次 \textbf{15}）。
  \end{enumerate}
  前 4 条有序合并规则为：
  \[
    \texttt{(u, g)->ug},\quad \texttt{(u, n)->un},\quad \texttt{(un, </w>)->un</w>},\quad \texttt{(h, ug)->hug}.
  \]

  \item \textbf{测试词切分}（按上述规则顺序逐条应用）：
  \begin{itemize}
    \item \texttt{bugs</w>}：初始 \texttt{[b, u, g, s, </w>]} $\xrightarrow{\text{规则 1}}$ \texttt{[b, ug, s, </w>]}，其余规则均不匹配，最终切分为 \texttt{['b', 'ug', 's', '</w>']}。
    \item \texttt{puns</w>}：初始 \texttt{[p, u, n, s, </w>]} $\xrightarrow{\text{规则 2}}$ \texttt{[p, un, s, </w>]}，因 \texttt{un} 后紧跟 \texttt{s} 而非 \texttt{</w>}，规则 3 不触发，最终切分为 \texttt{['p', 'un', 's', '</w>']}。
  \end{itemize}
  （注：即便第 2 轮破平改选 \texttt{(n, </w>)->n</w>}、第 3 轮选 \texttt{(u, n</w>)->un</w>}，则 \texttt{puns</w>} 会保持 \texttt{['p', 'u', 'n', 's', '</w>']}，由此可见词尾标记与并列破平策略对复数词形切分有微妙影响。）
\end{enumerate}

\subsection*{习题 5.5.2\quad 困惑度、交叉熵与 Chinchilla 算力预算换算}
\begin{enumerate}[label=(\alph*), leftmargin=1.8em]
  \item \textbf{困惑度与信息压缩比}：
  \begin{itemize}
    \item 困惑度：$\mathrm{PPL} = \exp(\mathcal{L}_{\mathrm{CE}}) = e^{2.3026} \approx \mathbf{10.0001} \approx 10$（因 $\ln 10 \approx 2.302585$），等价于模型每一步预测的不确定性仅相当于在 $10$ 个等概率词元中随机挑一个。
    \item 每词元比特数：$\mathrm{BPT} = \frac{\mathcal{L}_{\mathrm{CE}}}{\ln 2} = \frac{2.3026}{0.693147} \approx \mathbf{3.3219}\text{ bits/token}$。
    \item 词表 $|V| = 100{,}000$ 完全盲猜（均匀分布）的信息量为 $\log_2(10^5) = 5\log_2 10 \approx \mathbf{16.6096}\text{ bits/token}$。
    \item 压缩倍数：$\frac{16.6096}{3.3219} = \mathbf{5.00\text{ 倍}}$（也可直接由 $\frac{\ln 10^5}{\ln 10} = 5$ 看出）。
  \end{itemize}

  \item \textbf{Chinchilla 算力与最优 $(N^*, D^*)$ 换算}：
  \begin{itemize}
    \item 有效训练秒数：$14 \times 86{,}400 = 1{,}209{,}600\text{ s}$；
    \item 单卡有效算力：$1000\text{ TFLOPS} \times 40\% = 400\text{ TFLOPS} = 4 \times 10^{14}\text{ FLOP/s}$；
    \item 总有效算力：
    \[
      C = 64 \times (4 \times 10^{14}) \times 1{,}209{,}600 = \mathbf{3.096576 \times 10^{22}\text{ FLOPs}}.
    \]
    \item 代入 $C = 6 N^* D^* = 6 N^* (20 N^*) = 120 (N^*)^2$：
    \[
      N^* = \sqrt{\frac{3.096576 \times 10^{22}}{120}} = \sqrt{2.58048 \times 10^{20}} \approx \mathbf{1.606 \times 10^{10}} \approx \mathbf{16.06\text{B}},
    \]
    \[
      D^* = 20 N^* \approx \mathbf{3.213 \times 10^{11}} \approx \mathbf{321.3\text{B}\text{ tokens}}.
    \]
  \end{itemize}
\end{enumerate}

\subsection*{习题 5.5.3\quad InfoNCE 的硬负样本聚焦定理}
\begin{proof}
设负样本索引集为 $\mathcal{N}_i = \{1, \dots, N\} \setminus \{i\}$（基数 $N-1$），相对梯度权重为
\[
  w_j(\tau) = \frac{\exp(s_j / \tau)}{\sum_{m \in \mathcal{N}_i} \exp(s_m / \tau)}.
\]
\begin{enumerate}[label=(\alph*), leftmargin=1.8em]
  \item \textbf{$\tau \to 0^+$ 极限（退化为硬负样本三元组挖掘）}：设 $j^* = \argmax_{m \in \mathcal{N}_i} s_m$ 唯一，即对一切 $m \in \mathcal{N}_i \setminus \{j^*\}$ 均有严格间隙 $\Delta_m := s_{j^*} - s_m > 0$。分子分母同除以 $\exp(s_{j^*} / \tau)$：
  \[
    w_j(\tau) = \frac{\exp\!\big(-(s_{j^*} - s_j)/\tau\big)}{1 + \sum_{m \ne j^*} \exp(-\Delta_m / \tau)}.
  \]
  当 $\tau \to 0^+$ 时，对一切 $m \ne j^*$ 有 $-\Delta_m / \tau \to -\infty$，故 $\exp(-\Delta_m / \tau) \to 0$。从而：
  \[
    \lim_{\tau \to 0^+} w_j(\tau) = \begin{cases} 1, & j = j^*, \\ 0, & j \neq j^*. \end{cases}
  \]
  代入式~\eqref{eq:infonce-grad}，负样本排斥梯度 $\sum_{j \ne i} w_j(\tau)\hat{z}_j^B$ 全部集中于唯一最难负样本 $\hat{z}_{j^*}^B$ 上（同时 $\tau \ln \sum_m e^{s_m/\tau} \to \max_m s_m$），损失退化为只拉近正对、只推开最难负样本的\textbf{硬挖掘三元组损失}。

  \item \textbf{$\tau \to +\infty$ 极限（均匀排斥）}：记 $u = 1/\tau \to 0^+$，利用一阶泰勒展开 $e^{u s_j} = 1 + u s_j + O(u^2)$，记负样本平均相似度 $\bar{s} := \frac{1}{N-1}\sum_{m \in \mathcal{N}_i} s_m$：
  \begin{align*}
    w_j(\tau)
    &= \frac{1 + u s_j + O(u^2)}{(N-1)\big(1 + u\bar{s} + O(u^2)\big)} \\
    &= \frac{1}{N-1}\big(1 + u(s_j - \bar{s}) + O(u^2)\big)
     = \frac{1}{N-1} + \frac{s_j - \bar{s}}{(N-1)\tau} + O(\tau^{-2}).
  \end{align*}
  令 $\tau \to +\infty$ 立得 $\lim_{\tau \to +\infty} w_j(\tau) = \frac{1}{N-1}$，且一阶修正项按 $1/\tau$ 消失，所有负样本不论难易均施加几乎相等的排斥力，丧失硬负样本辨别精度。
\end{enumerate}
\end{proof}

\subsection*{习题 5.5.4\quad 含推理摊销成本的广义缩放最优解}
\begin{proof}
目标是在固定全生命周期总算力 $C_{\text{total}} = 6 N D + 2 N D_{\text{inf}}$ 下最小化式~\eqref{eq:chinchilla-loss} 的测试损失 $L(N, D) = E + A N^{-\alpha} + B D^{-\beta}$（$A, B, \alpha, \beta > 0$）。构造拉格朗日函数：
\[
  \mathcal{J}(N, D, \lambda) = E + A N^{-\alpha} + B D^{-\beta} + \lambda\big(6 N D + 2 N D_{\text{inf}} - C_{\text{total}}\big).
\]
令关于 $N$ 与 $D$ 的偏导数为零：
\begin{align*}
  \frac{\partial \mathcal{J}}{\partial N} &= -\alpha A N^{-\alpha - 1} + \lambda(6 D + 2 D_{\text{inf}}) = 0
  \;\Longrightarrow\; \alpha A N^{-\alpha} = \lambda N (6 D + 2 D_{\text{inf}}), \\
  \frac{\partial \mathcal{J}}{\partial D} &= -\beta B D^{-\beta - 1} + 6 \lambda N = 0
  \;\Longrightarrow\; \beta B D^{-\beta} = 6 \lambda N D.
\end{align*}
将两式相除消去拉格朗日乘子 $\lambda > 0$，得到广义最优均衡方程：
\[
  \frac{\alpha A N^{-\alpha}}{\beta B D^{-\beta}}
  = \frac{6 D + 2 D_{\text{inf}}}{6 D}
  = 1 + \frac{D_{\text{inf}}}{3 D}.
\]
\textbf{证明最优词元参数比严格增大}：对比纯训练最优（$D_{\text{inf}} = 0$）时的 Chinchilla 均衡条件 $\frac{\alpha A N_0^{-\alpha}}{\beta B D_0^{-\beta}} = 1$：只要上线存在推理需求 $D_{\text{inf}} > 0$，必有
\[
  \frac{\alpha A (N^*)^{-\alpha}}{\beta B (D^*)^{-\beta}} = 1 + \frac{D_{\text{inf}}}{3 D^*} > 1
  \;\Longrightarrow\;
  \frac{(D^*)^\beta}{(N^*)^\alpha} > \frac{\beta B}{\alpha A} = \frac{D_0^\beta}{N_0^\alpha}.
\]
换言之，在相同训练计算量或相同目标损失下，广义最优解必须\textbf{减小参数量 $N^*$ 并大幅增加训练词元数 $D^*$}[注意：每多训一个词元只付一次 $6N$ 的训练费，而每多加一个参数却要为未来全部 $D_{\text{inf}}$ 个推理词元持续缴纳 $2 D_{\text{inf}}$ 的“推理税”]。特别地，当 $\alpha \approx \beta$（如 Chinchilla 拟合值 $\alpha=0.34, \beta=0.28$）时，$D^*/N^*$ 严格高于纯训练最优的 $20$，且随预期服务量 $D_{\text{inf}}$ 增大而大幅飙升——这正是 Llama-3 8B（$D/N \approx 1875$）与 Qwen 系列小尺寸模型坚持“重度过训练”的经济学第一性原理。
\end{proof}

\subsection*{习题 5.5.5\quad 30 行纯 Python 从零实现微型 BPE 分词器}
\textbf{完整核心代码}（仅依赖标准库 \texttt{collections.Counter}）：
\begin{lstlisting}
from collections import Counter

def train_bpe(corpus_words, num_merges):
    # corpus_words: dict[str, int], e.g. {"hug": 10, "pug": 5, ...}
    vocab = {tuple(w) + ("</w>",): f for w, f in corpus_words.items()}
    merges = []
    for _ in range(num_merges):
        pairs = Counter()
        for syms, freq in vocab.items():
            for i in range(len(syms) - 1):
                pairs[(syms[i], syms[i + 1])] += freq
        if not pairs:
            break
        best = pairs.most_common(1)[0][0]
        merges.append(best)
        new_vocab = {}
        for syms, freq in vocab.items():
            out, i = [], 0
            while i < len(syms):
                if i < len(syms) - 1 and (syms[i], syms[i+1]) == best:
                    out.append(best[0] + best[1]); i += 2
                else:
                    out.append(syms[i]); i += 1
            new_vocab[tuple(out)] = freq
        vocab = new_vocab
    return merges

def encode(word, merges):
    syms = list(word) + ["</w>"]
    for a, b in merges:
        out, i = [], 0
        while i < len(syms):
            if i < len(syms) - 1 and syms[i] == a and syms[i+1] == b:
                out.append(a + b); i += 2
            else:
                out.append(syms[i]); i += 1
        syms = out
    return syms
\end{lstlisting}
\textbf{运行预期与校验基准}：在例~\ref{ex:bpe-trace} 语料上运行即可复现正文全步轨迹；在 500 词英文技术段落上运行 50 轮合并，前 15 条规则会优先捕获高频英语词缀与功能词（如 \texttt{th}, \texttt{in}, \texttt{er}, \texttt{on}, \texttt{the</w>}, \texttt{ing</w>}, \texttt{tion</w>}），全段总词元数通常可从约 $3000$ 字符压缩至 $1200 \sim 1400$ 子词（压缩率约 $55\% \sim 60\%$）。

\subsection*{习题 5.5.6\quad InfoNCE 温度 $\tau$ 与批大小 $N$ 消融实验}
严格复用 \texttt{vol5\_experiments.py} 中的 \texttt{exp5\_infonce\_contrastive}（$N = 128,\; D_{\text{in}} = D_{\text{emb}} = 16$，200 轮 AdamW \texttt{lr=5e-3}），两组消融实测基准如下：
\begin{center}
\small
\begin{tabular}{@{}lccc|lccc@{}}
\toprule
温度 $\tau$（$B=128$） & Top-1 & \texttt{pos\_cos} & \texttt{neg\_cos} & 批大小 $B$（$\tau=0.15$） & Top-1 & \texttt{pos\_cos} & \texttt{neg\_cos} \\
\midrule
$\tau = 0.02$ & $100.0\%$ & $0.651$ & $+0.025$ & $B = 8$ & $50.8\%$ & $0.721$ & $-0.005$ \\
$\tau = 0.15$ & $100.0\%$ & $0.973$ & $-0.008$ & $B = 32$ & $98.4\%$ & $0.890$ & $-0.006$ \\
$\tau = 1.00$ & $100.0\%$ & $0.997$ & $-0.008$ & $B = 128$ & $100.0\%$ & $0.973$ & $-0.008$ \\
$\tau = 5.00$ & $98.4\%$ & $0.998$ & $-0.008$ & 理论互信息上限 $\ln B$ & $2.08$ & $3.47$ & $4.85$ \\
\bottomrule
\end{tabular}
\end{center}
\textbf{定量结论}：
\begin{enumerate}[label=(\alph*), leftmargin=1.8em]
  \item \textbf{温度 $\tau$ 的双向权衡}：$\tau = 0.02$ 过于尖锐，一旦正样本余弦仅略高于负样本，Softmax 概率 $p_{ii}$ 便饱和至 $1$、梯度提前熄火，导致对齐度 \texttt{pos\_cos} 停留在 $0.651$ 且负样本均余弦偏高（$+0.025$）；相反 $\tau = 5.0$ 使梯度整体缩小 $1/\tau = 0.2$ 倍，200 轮内难以完全拉开难分样本，Top-1 降至 $98.4\%$；$\tau \in [0.07, 0.2]$ 是对齐与均匀性的黄金窗口。
  \item \textbf{批大小 $B$ 收紧互信息下界}：由定理~\ref{thm:infonce-mi}，$I(X; Y) \ge \ln B - \mathcal{L}_B$。当 $B=8$ 时每步仅见 $7$ 个负样本，互信息下界天花板被死锁在 $\ln 8 \approx 2.08$ nats，全库 128 路全局 Top-1 检索率暴跌至 $50.8\%$；增至 $B=32$（上限 $3.47$ nats）跃升至 $98.4\%$，$B=128$（上限 $4.85$ nats）直达 $100.0\%$。这精准解释了 CLIP/SimCLR 为何极度依赖超大对比批大小。
\end{enumerate}

\subsection*{习题 5.5.7\quad 从 Softmax InfoNCE 到 Sigmoid 逐对对比（SigLIP）}
\begin{enumerate}[label=(\arabic*), leftmargin=1.8em]
  \item \textbf{分布式通信与显存优势}：标准 CLIP 的 Softmax 分母 $\sum_{j=1}^N \exp(S_{ij}/\tau)$ 耦合了全局全部 $N$ 个样本，分布式训练时必须把所有卡上的特征全部 \texttt{all-gather} 并物化完整 $N\times N$ 矩阵，显存与同步开销高昂。SigLIP 彻底消除了跨列归一化，每个 $(i, j)$ 格子都是独立二元逻辑损失。借助分块环形通信（Chunked Ring Exchange），每张 GPU 每次仅持有一小块本地正负对计算梯度并立即释放，随后把键向量环传给下一张卡，将峰值激活显存从 $O(N^2)$ 压至 $O(N^2 / K^2)$（$K$ 为设备数），在小批量与百万级超大批下均表现优异。
  \item \textbf{为何必须将偏置初始化为 $b \approx -10$}：在 $N \times N$ 矩阵中，共有 $N$ 个正对（$y_{ij}=+1$）与 $N(N-1)$ 个负对（$y_{ij}=-1$），正负比例极端失衡（如 $N = 32{,}768$ 时高达 $1 : 32{,}767$）。若按常规初始化 $b = 0$，初始时所有格子预测概率均为 $\sigma(0) = 0.5$：负对产生的总梯度量级将是正对的 $32{,}767$ 倍，瞬间将所有特征粗暴推开、引发剧烈震荡。而将偏置初始化为 $b \approx -\ln N \approx -10.4$（实操取 $b = -10$）后，负对初始预测正概率降至 $\sigma(-10) \approx 4.5 \times 10^{-5}$，于是：
  \[
    \text{全部负对梯度总和量级} \approx (N - 1)\cdot \sigma(b) \approx N \cdot \frac{1}{N} = 1,
  \]
  恰好与单个正对的梯度量级 $1 - \sigma(b) \approx 1$ 完全平衡！这相当于把先验对数赔率 $\ln\frac{1}{N-1}$ 写入偏置，使训练从第 1 步起就平稳收敛。
\end{enumerate}

\subsection*{习题 5.5.8\quad 高质量数据饱和与“模型坍缩”}
\begin{enumerate}[label=(\arabic*), leftmargin=1.8em]
  \item \textbf{数据受限缩放与纯合成递归的坍缩机理}：Muennighoff 等人（2023）发现重复消费有限高质量语料在 4 个 epoch 内几乎无损，超过后有效数据量按 $D' = U_D + U_D R_D^*(1 - e^{-R_D / R_D^*})$ 饱和。若试图改用前代模型自由采样生成的文本递归训练下一代，会产生两重退化：（i）\textbf{有限样本统计截断}：设第 $k$ 代拟合高斯分布 $\mathcal{N}(\hat{\mu}_k, \hat{\sigma}_k^2)$，每次从 $n$ 个有限样本重估方差，其期望满足 $\E[\hat{\sigma}_{k+1}^2] = \frac{n-1}{n}\E[\hat{\sigma}_k^2] = (1 - 1/n)^k \sigma_0^2 \to 0$，低概率长尾事件一旦在某代未被抽到便永久灭绝；（ii）\textbf{函数逼近误差逐代累积}，最终导致多峰真实分布坍缩为单调乏味的窄峰（Model Collapse）。
  \item \textbf{可验证合成数据为何能打破诅咒}：单纯无监督自采样受限于信息论的\textbf{数据处理不等式}——闭环内部不会凭空产生超出初始分布的新信息。而引入外部确定性验证器（如 Lean 形式化证明检查器、代码编译器与单元测试集、数学符号验算器）后，系统本质上在执行\textbf{拒绝采样与搜索（Generate-then-Verify）}：生成器负责在广阔解空间做高熵探索，验证器则以极高信噪比把真正满足形式逻辑与事实约束的轨迹筛选出来，等价于向训练集持续注入来自形式公理体系与真实环境的\textbf{外部负熵（新信息）}。这正是当代推理大模型借助“可验证奖励强化学习（RLVR）与自举搜索”突破人类语料天花板的根本原因（也是下一卷第 6 卷的核心主题）。
\end{enumerate}

% =========================================================================
\section*{第 6 章习题解答（共 8 题）}
\addcontentsline{toc}{section}{第 6 章习题解答（共 8 题）}
% =========================================================================

\subsection*{习题 5.6.1\quad VAE 单样本损失与重参数化梯度完整手算}
已知 $x = 2.0,\; \mu = 1.0,\; s = \ln \sigma^2 = 0.0$（故 $\sigma^2 = e^s = 1.0,\; \sigma = e^{s/2} = 1.0$），外部噪声 $\varepsilon = -0.5$，解码器 $\hat{x} = z$。
\begin{enumerate}[label=(\alph*), leftmargin=1.8em]
  \item \textbf{前向损失计算}：
  \begin{align*}
    z &= \mu + \sigma \varepsilon = 1.0 + 1.0 \times (-0.5) = \mathbf{0.5}, \\
    \mathcal{L}_{\text{rec}} &= \frac{1}{2}(x - z)^2 = \frac{1}{2}(2.0 - 0.5)^2 = \frac{1}{2}\times 2.25 = \mathbf{1.125}, \\
    \DKL &= \frac{1}{2}(\mu^2 + e^s - s - 1) = \frac{1}{2}(1^2 + 1 - 0 - 1) = \mathbf{0.5}, \\
    \mathcal{L}_{\text{tot}} &= \mathcal{L}_{\text{rec}} + \DKL = 1.125 + 0.5 = \mathbf{1.625}.
  \end{align*}

  \item \textbf{重参数化精确偏导数}：先算重建项对 $z$ 的导数 $\frac{\partial \mathcal{L}_{\text{rec}}}{\partial z} = z - x = 0.5 - 2.0 = -1.5$。结合 $\frac{\partial z}{\partial \mu} = 1$ 与 $\frac{\partial z}{\partial s} = \frac{1}{2} e^{s/2}\varepsilon = \frac{1}{2}\sigma\varepsilon = -0.25$，由链式法则得：
  \begin{align*}
    \frac{\partial \mathcal{L}_{\text{tot}}}{\partial \mu}
    &= \frac{\partial \mathcal{L}_{\text{rec}}}{\partial z}\frac{\partial z}{\partial \mu} + \frac{\partial \DKL}{\partial \mu}
     = (z - x)\cdot 1 + \mu = -1.5 + 1.0 = \mathbf{-0.5}, \\
    \frac{\partial \mathcal{L}_{\text{tot}}}{\partial s}
    &= \frac{\partial \mathcal{L}_{\text{rec}}}{\partial z}\frac{\partial z}{\partial s} + \frac{\partial \DKL}{\partial s}
     = (z - x)\cdot\frac{1}{2}\sigma\varepsilon + \frac{1}{2}(e^s - 1) \\
    &= (-1.5)\times(-0.25) + \frac{1}{2}(1 - 1) = \mathbf{+0.375}.
  \end{align*}
  直觉检验：当前 $z = 0.5 < x = 2.0$，负梯度 $-\frac{\partial \mathcal{L}}{\partial \mu} = +0.5$ 把均值 $\mu$ 向上拉向 $x$；同时因本次噪声 $\varepsilon = -0.5 < 0$ 把 $z$ 拽离了 $x$，正梯度 $\frac{\partial \mathcal{L}}{\partial s} = +0.375$ 促使编码器缩小方差 $\sigma^2$。
\end{enumerate}

\subsection*{习题 5.6.2\quad 微型 $T = 3$ 步 DDPM 系数表与逆向去噪手算}
\begin{enumerate}[label=(\alph*), leftmargin=1.8em]
  \item \textbf{完整系数表}：给定 $\beta_1 = 0.1,\; \beta_2 = 0.2,\; \beta_3 = 0.3$，依次得 $\alpha_t = 1 - \beta_t$ 为 $0.9,\; 0.8,\; 0.7$；累积乘积 $\bar{\alpha}_1 = 0.9,\; \bar{\alpha}_2 = 0.9\times 0.8 = 0.72,\; \bar{\alpha}_3 = 0.72\times 0.7 = 0.504$。
  信噪比 $\mathrm{SNR}(t) = \frac{\bar{\alpha}_t}{1-\bar{\alpha}_t}$ 与后验方差 $\tilde{\beta}_t = \frac{1-\bar{\alpha}_{t-1}}{1-\bar{\alpha}_t}\beta_t$（式~\eqref{eq:ddpm-post-mean-x0}）为：
  \[
    \tilde{\beta}_2 = \frac{1 - 0.9}{1 - 0.72}\times 0.2 = \frac{0.02}{0.28} = \frac{1}{14} \approx 0.0714, \qquad
    \tilde{\beta}_3 = \frac{1 - 0.72}{1 - 0.504}\times 0.3 = \frac{0.084}{0.496} = \frac{21}{124} \approx 0.1694.
  \]
  汇总如下（保留四位小数）：
  \begin{center}
  \small
  \begin{tabular}{@{}cccccc@{}}
  \toprule
  步数 $t$ & $\beta_t$ & $\alpha_t$ & $\bar{\alpha}_t$ & $\mathrm{SNR}(t)$ & 后验方差 $\tilde{\beta}_t$ \\
  \midrule
  $t = 1$ & $0.1000$ & $0.9000$ & $0.9000$ & $9.0000$ & ---（$0$） \\
  $t = 2$ & $0.2000$ & $0.8000$ & $0.7200$ & $2.5714\;(\tfrac{18}{7})$ & $0.0714\;(\tfrac{1}{14}) < 0.2$ \\
  $t = 3$ & $0.3000$ & $0.7000$ & $0.5040$ & $1.0161\;(\tfrac{63}{62})$ & $0.1694\;(\tfrac{21}{124}) < 0.3$ \\
  \bottomrule
  \end{tabular}
  \end{center}
  可见每步严格满足 $\tilde{\beta}_t < \beta_t$（因为 $1 - \bar{\alpha}_{t-1} < 1 - \bar{\alpha}_t$）。

  \item \textbf{第 $t = 2$ 步去噪均值与估计分数}：已知 $x_2 = 1.50,\; \varepsilon_\theta(x_2, 2) = 0.40$，$\alpha_2 = 0.8,\; \beta_2 = 0.2,\; 1 - \bar{\alpha}_2 = 0.28$（$\sqrt{0.8} \approx 0.894427,\; \sqrt{0.28} \approx 0.529150$）。代入式~\eqref{eq:ddpm-mu-theta} 与式~\eqref{eq:score-eps-link}：
  \begin{align*}
    \mu_\theta(x_2, 2)
    &= \frac{1}{\sqrt{\alpha_2}}\!\left(x_2 - \frac{\beta_2}{\sqrt{1 - \bar{\alpha}_2}}\,\varepsilon_\theta(x_2, 2)\right) \\
    &= \frac{1}{\sqrt{0.8}}\!\left(1.50 - \frac{0.20}{\sqrt{0.28}}\times 0.40\right)
     \approx 1.118034 \times (1.50 - 0.151186) \approx \mathbf{1.5080}, \\
    s_\theta(x_2, 2)
    &= -\frac{\varepsilon_\theta(x_2, 2)}{\sqrt{1 - \bar{\alpha}_2}}
     = -\frac{0.40}{\sqrt{0.28}} \approx \mathbf{-0.7559}.
  \end{align*}
\end{enumerate}

\subsection*{习题 5.6.3\quad 一般多元高斯 KL 散度公式与标准正态特例}
\begin{proof}
写出 $q(z) = \mathcal{N}(\mu_q, \Sigma_q)$ 与 $p(z) = \mathcal{N}(\mu_p, \Sigma_p)$ 的对数密度差：
\begin{align*}
  \ln \frac{q(z)}{p(z)}
  &= \frac{1}{2}\ln\frac{\det \Sigma_p}{\det \Sigma_q}
   - \frac{1}{2}(z - \mu_q)^\top \Sigma_q^{-1}(z - \mu_q)
   + \frac{1}{2}(z - \mu_p)^\top \Sigma_p^{-1}(z - \mu_p).
\end{align*}
对两端关于 $z \sim q$ 取数学期望，并逐项利用标量的迹等于自身及迹的循环不变性 $\tr(ABC) = \tr(CAB)$：
\begin{enumerate}[label=(\arabic*), leftmargin=1.8em]
  \item 第二项：$\E_q[(z - \mu_q)^\top \Sigma_q^{-1}(z - \mu_q)] = \tr\!\big(\Sigma_q^{-1}\E_q[(z - \mu_q)(z - \mu_q)^\top]\big) = \tr(\Sigma_q^{-1}\Sigma_q) = \tr(I_d) = d$。
  \item 第三项：将 $z - \mu_p = (z - \mu_q) + (\mu_q - \mu_p)$ 展开，注意到交叉项因 $\E_q[z - \mu_q] = 0$ 而消失：
  \begin{align*}
    \E_q\big[(z - \mu_p)^\top \Sigma_p^{-1}(z - \mu_p)\big]
    &= \tr\!\big(\Sigma_p^{-1}\E_q[(z - \mu_q)(z - \mu_q)^\top]\big) + (\mu_q - \mu_p)^\top \Sigma_p^{-1}(\mu_q - \mu_p) \\
    &= \tr(\Sigma_p^{-1}\Sigma_q) + (\mu_p - \mu_q)^\top \Sigma_p^{-1}(\mu_p - \mu_q).
  \end{align*}
\end{enumerate}
合并三项即证一般公式：
\[
  \DKL(q \,\|\, p) = \frac{1}{2}\!\left[\tr(\Sigma_p^{-1}\Sigma_q) + (\mu_p - \mu_q)^\top \Sigma_p^{-1}(\mu_p - \mu_q) - d + \ln\frac{\det \Sigma_p}{\det \Sigma_q}\right].
\]
\textbf{退化验证}：令 $\mu_p = 0,\; \Sigma_p = I_d,\; \mu_q = \mu,\; \Sigma_q = \diag(\sigma_1^2, \dots, \sigma_d^2)$，则 $\tr(\Sigma_p^{-1}\Sigma_q) = \sum_{j=1}^d \sigma_j^2$，$(\mu_p - \mu_q)^\top \Sigma_p^{-1}(\mu_p - \mu_q) = \sum_{j=1}^d \mu_j^2$，且 $\ln\frac{\det I_d}{\det \Sigma_q} = -\sum_{j=1}^d \ln \sigma_j^2$，代入立得式~\eqref{eq:gaussian-kl-closed}：
\[
  \DKL(q \,\|\, p) = \frac{1}{2}\sum_{j=1}^d \left(\mu_j^2 + \sigma_j^2 - \ln \sigma_j^2 - 1\right).
\]
\end{proof}

\subsection*{习题 5.6.4\quad DDPM 三种等价参数化与流匹配时间坐标同构}
\begin{proof}
记 $a_t := \sqrt{\bar{\alpha}_t},\; b_t := \sqrt{1 - \bar{\alpha}_t}$，满足 $a_t^2 + b_t^2 = 1$。
\begin{enumerate}[label=(\alph*), leftmargin=1.8em]
  \item \textbf{$(x_0, \varepsilon, v_t)$ 双向正交换算}：由前向方程与 $v_t$ 定义构成 $2\times 2$ 正交旋转矩阵：
  \[
    \begin{pmatrix} x_t \\ v_t \end{pmatrix}
    = \begin{pmatrix} a_t & b_t \\ -b_t & a_t \end{pmatrix}
      \begin{pmatrix} x_0 \\ \varepsilon \end{pmatrix}.
  \]
  因系数矩阵正交（逆矩阵即转置），立即反解出以 $(x_t, v_t)$ 表示 $x_0$ 与 $\varepsilon$ 的精确公式：
  \[
    x_0 = a_t x_t - b_t v_t = \sqrt{\bar{\alpha}_t}\,x_t - \sqrt{1-\bar{\alpha}_t}\,v_t, \qquad
    \varepsilon = b_t x_t + a_t v_t = \sqrt{1-\bar{\alpha}_t}\,x_t + \sqrt{\bar{\alpha}_t}\,v_t.
  \]
  其余双向换算为：$x_0 = \frac{x_t - b_t \varepsilon}{a_t}$、$\varepsilon = \frac{x_t - a_t x_0}{b_t}$、$v_t = \frac{a_t x_t - x_0}{b_t} = \frac{\varepsilon - b_t x_t}{a_t}$。

  \textbf{高噪声步稳定性}：当 $\bar{\alpha}_t \to 0$（即 $a_t \to 0,\; b_t \to 1$）时，$x_t \approx \varepsilon$ 几乎不含任何信号。此时若预测 $x_0$，微小的信号估计误差在换算回 $\varepsilon$ 或分数时会被 $1/a_t \to \infty$ 放大；而按 $\varepsilon$ 预测在低噪声端 $b_t \to 0$ 时也会被 $1/b_t$ 放大。反观 $v_t = a_t \varepsilon - b_t x_0$，由于 $x_0$ 与 $\varepsilon$ 独立，当数据标准化至 $\Var(x_0) = 1$ 时，对\textbf{任意}时间步 $t$ 恒有
  \[
    \Var(v_t) = a_t^2 \Var(\varepsilon) + b_t^2 \Var(x_0) = \bar{\alpha}_t \cdot 1 + (1-\bar{\alpha}_t)\cdot 1 = 1,
  \]
  在 $t \to 0$ 时平滑过渡为预测 $\varepsilon$、在 $t \to T$ 时平滑过渡为预测 $-x_0$，全时间段方差守恒且梯度良态。

  \item \textbf{流匹配与 DDPM 的时间坐标同构}：对流匹配直线插值 $x_s = (1-s)x_0 + s x_1$（其中 $x_0 \sim \mathcal{N}(0, I)$ 为纯噪声，$x_1 \sim p_{\text{data}}$ 为数据，$s \in (0, 1)$），令归一化系数 $c(s) := \sqrt{(1-s)^2 + s^2} > 0$，则重归一化后的随机变量为：
  \[
    \tilde{x}_s = \frac{x_s}{c(s)}
    = \frac{s}{\sqrt{(1-s)^2 + s^2}}\, x_1 + \frac{1-s}{\sqrt{(1-s)^2 + s^2}}\, x_0.
  \]
  检验两系数平方和：
  \[
    \left(\frac{s}{c(s)}\right)^2 + \left(\frac{1-s}{c(s)}\right)^2 = \frac{s^2 + (1-s)^2}{(1-s)^2 + s^2} = 1.
  \]
  因此令严格单调的时间重参数化满足 $\sqrt{\bar{\alpha}_t} = \frac{s}{\sqrt{(1-s)^2 + s^2}}$，自动有 $\sqrt{1-\bar{\alpha}_t} = \frac{1-s}{\sqrt{(1-s)^2 + s^2}}$。由于 $x_1$ 对应 DDPM 中的干净数据、$x_0$ 对应标准高斯噪声 $\varepsilon$，上式与 DDPM 前向分布~\eqref{eq:ddpm-xt-reparam} 在边缘分布上\textbf{完全同分布}！区别仅在于流匹配在时间与空间尺度上选择了更直的确定性传输速度场。
\end{enumerate}
\end{proof}

\subsection*{习题 5.6.5\quad 流匹配 vs DDPM 少步采样正面对决}
严格复用代码~\ref{lst:ddpm-2d} 的 4 模态高斯环（半径 $2.0$，簇标准差 $0.18$，真实样本误差下限约 $0.226$）与完全相同的 3 层网络 \texttt{Linear(3,64)-SiLU-Linear(64,64)-SiLU-Linear(64,2)}，按式~\eqref{eq:cfm-loss} 以批大小 256 训练 600 步（Adam \texttt{lr=3e-3}），核心训练与欧拉采样代码如下：
\begin{lstlisting}
# 训练步：随机采样 t ~ U[0,1]，直线插值并回归条件速度 (x1 - x0)
x0 = torch.randn_like(x1)
t = torch.rand(x1.size(0), 1)
xt = (1.0 - t) * x0 + t * x1
loss = F.mse_loss(v_net(torch.cat([xt, t], dim=1)), x1 - x0)

# 确定性 K 步欧拉积分采样（t: 0 -> 1）
x = torch.randn(n, 2); dt = 1.0 / K
for k in range(K):
    t_col = torch.full((n, 1), k * dt)
    x = x + dt * v_net(torch.cat([x, t_col], dim=1))
\end{lstlisting}
\textbf{少步采样实测基准}（评估 1000 个生成点到最近模态中心的均值距离 \texttt{gen\_err}）：
\begin{center}
\small
\begin{tabular}{@{}ccccc@{}}
\toprule
采样步数 $K$ & $K = 5$ & $K = 10$ & $K = 20$ & $K = 40$ \\
\midrule
校正流（Rectified Flow）\texttt{gen\_err} & $0.377$ & $0.298$ & $0.275$ & $0.267$ \\
\bottomrule
\end{tabular}
\end{center}
仅用 $K = 10$ 步确定性欧拉积分即可取得 $0.298$ 的生成误差，与 40 步随机 DDPM 处于同一水准，而在 $K = 20 \sim 40$ 步时迅速逼近真实数据下限 $0.226$。原因在于校正流的传输轨迹远比方差保持扩散 SDE 笔直，低阶欧拉离散化截断误差大幅降低。

\subsection*{习题 5.6.6\quad 可控复现并修复 VAE 后验坍缩}
在 2D 四模态高斯环上构造隐维 $d_z = 4$ 的极简 VAE（编码器输出 $\mu \in \R^4$ 与 $\log\sigma^2 \in \R^4$，解码器还原 $\hat{x} \in \R^2$），使用 Adam（\texttt{lr=5e-3}）训练 300 轮，三组消融实测结果如下：
\begin{center}
\small
\begin{tabular}{@{}lccc@{}}
\toprule
实验配置 & 验证集重建损失 $\mathcal{L}_{\text{rec}}$ & 最大单维 $\max_j \DKL^{(j)}$ & 活跃隐维度（$\Var(\mu_j) > 0.01$） \\
\midrule
标准训练（$\beta = 1.0$） & $0.123$ & $0.406$ & $2 / 4$（精确匹配 2D 内蕴维度） \\
过重 KL 惩罚（$\beta = 10.0$） & $1.990$ & $0.0001 < 0.01$ & $0 / 4$（完全后验坍缩） \\
线性 KL 退火（前 $30\%$ 轮） & $0.154$ & $0.397$ & $2 / 4$（稳定避免坍缩） \\
\bottomrule
\end{tabular}
\end{center}
\textbf{现象与原理解析}：当 $\beta = 10.0$ 时，优化初期编码器尚未学出有效特征，减小 $\beta \DKL$ 成为最廉价的降损捷径，全部 4 个维度瞬间坍缩至先验 $\mathcal{N}(0, 1)$（$\max_j \DKL^{(j)} = 10^{-4}$），解码器被迫丢弃隐码、对任意输入均输出数据集均值 $(0, 0)$（重建误差恰好等于半径平方的一半 $\frac{1}{2}\times 2^2 = 2.0$，与实测 $1.990$ 吻合）。线性 KL 退火在训练早期把 $\beta$ 保持在接近 $0$ 的水平，让编码器先建立起高信息量隐空间表征，再逐步引入先验正则，从而在强解码器（如自回归或卷积解码器）场景下可靠治愈后验坍缩。

\subsection*{习题 5.6.7\quad Reflow（二次拉直）与一致性模型}
\begin{enumerate}[label=(\alph*), leftmargin=1.8em]
  \item \textbf{边缘分布不变性与传输代价单调不增证明}：
  \begin{proof}
  记初始独立耦合为 $(X_0, X_1)$，其线性插值 $X_t = (1-t)X_0 + t X_1$ 对应的精确条件期望速度场为 $u_t(x) = \E[X_1 - X_0 \mid X_t = x]$。由定理~\ref{thm:cfm}，沿 ODE $\frac{d Z_t}{dt} = u_t(Z_t)$（初值 $Z_0 = X_0$）演化所得随机过程 $\{Z_t\}_{t\in[0,1]}$ 在每个时刻 $t$ 均与 $X_t$ 满足相同的连续性方程，因而具有完全相同的边缘分布：$Z_t \stackrel{d}{=} X_t$，特别地终点 $\hat{X}_1 := Z_1$ 满足 $\hat{X}_1 \stackrel{d}{=} X_1 \sim p_1$，边缘分布严格保持不变。

  接着计算新配对 $(Z_0, Z_1)$ 的二次传输代价：由微积分基本定理 $Z_1 - Z_0 = \int_0^1 u_t(Z_t)\,dt$，对时间积分 $[0,1]$ 施用凸函数 $\norm{\cdot}_2^2$ 的 Jensen 不等式：
  \[
    \E\!\left[\norm{Z_1 - Z_0}_2^2\right]
    = \E\!\left[\norm{\int_0^1 u_t(Z_t)\,dt}_2^2\right]
    \le \int_0^1 \E\!\left[\norm{u_t(Z_t)}_2^2\right] dt.
  \]
  因 $Z_t \stackrel{d}{=} X_t$，有 $\E[\norm{u_t(Z_t)}_2^2] = \E[\norm{u_t(X_t)}_2^2] = \E\big[\norm{\E[X_1 - X_0 \mid X_t]}_2^2\big]$。再由条件期望的 $L^2$ 投影收缩性（或条件 Jensen 不等式 $\norm{\E[V \mid X_t]}_2^2 \le \E[\norm{V}_2^2 \mid X_t]$）：
  \[
    \E\!\left[\norm{u_t(X_t)}_2^2\right]
    \le \E\!\left[\E\!\left[\norm{X_1 - X_0}_2^2 \mid X_t\right]\right]
    = \E\!\left[\norm{X_1 - X_0}_2^2\right].
  \]
  代入积分即证单调不增性：
  \[
    \E\!\left[\norm{\hat{X}_1 - X_0}_2^2\right] \le \E\!\left[\norm{X_1 - X_0}_2^2\right].
  \]
  等号成立当且仅当路径几乎处处为直线（$u_t(Z_t)$ 沿整条轨迹为常数）且不同轨迹永不相交！因此迭代 Reflow 不断解开交叉路径，将随机配对逼近笔直的 2-Wasserstein 最优传输耦合。
  \end{proof}

  \item \textbf{Reflow vs 一致性模型（Consistency Models）的核心权衡}：
  \begin{itemize}
    \item \textbf{Reflow}：不改变流匹配标准训练范式，通过“生成合成配对 $(x_0, \hat{x}_1) \to$ 重新回归速度场”逐步拉直轨迹；\emph{优点}是训练极稳定、天然兼容任意 $K$ 步 ODE 求解器并可平滑权衡步数与画质；\emph{缺点}是每次迭代都要离线跑海量 ODE 生成合成数据，存在老师模型分布漂移误差累积。
    \item {\rightskip=0pt plus 1fil\relax \textbf{一致性模型}：直接学习自洽映射 $f_\theta(x_t, t) \approx x_1$，强制同一 ODE 轨迹上相邻两点满足 $f_\theta(x_{t_{n+1}}, t_{n+1}) = f_{\theta^-}(x_{t_n}, t_n)$ 及边界条件 $f_\theta(x_1, 1) = x_1$；\emph{优点}是原生支持 1--2 步直接映射且可脱离预训练教师从零训练（Consistency Training）；\emph{缺点}是自举目标存在类似强化学习 TD 的训练方差与不稳定性，需要精心设计时间步离散化调度、EMA 延迟目标与伪 Huber 损失。\par}
  \end{itemize}
\end{enumerate}

\subsection*{习题 5.6.8\quad 离散扩散语言模型（MDLM / LLaDA）与统一多模态之争}
\begin{enumerate}[label=(\alph*), leftmargin=1.8em]
  \item \textbf{掩码离散扩散 ELBO 为何化简为连续加权 BERT 损失}：在掩码离散扩散中，前向过程对每个词元独立以概率 $1 - \alpha_t$ 替换为吸收态 \texttt{[MASK]}（$\alpha_0 = 1$ 全干净，$\alpha_1 = 0$ 全掩码）。因为吸收态一旦变黑就不会再变回其他词元，后验分布 $q(x_s \mid x_t, x_0)$（$s < t$）具有极其干净的解析形式：
  \begin{itemize}
    \item 若位置 $i$ 在时刻 $t$ 尚未被掩码，则它在更早时刻 $s$ 必与 $x_{t,i}$ 相同（零 KL 贡献）；
    \item 若位置 $i$ 在时刻 $t$ 为 \texttt{[MASK]}，则它在时刻 $s$ 要么保持 \texttt{[MASK]}，要么以概率 $\frac{\alpha_s - \alpha_t}{1 - \alpha_t}$ 揭示出真实原词 $x_{0,i}$！
  \end{itemize}
  令时间步长 $t - s \to 0^+$，将全部时刻的离散 KL 项求和取连续极限，即得连续时间变分上界（负 ELBO）：
  \[
    \mathcal{L}_{\mathrm{MDLM}}
    = \int_0^1 \frac{-\alpha_t'}{1 - \alpha_t}\,
      \E_{q(x_t \mid x_0)}\!\left[\sum_{i:\, x_{t,i} = \texttt{[MASK]}} -\ln p_\theta(x_{0,i} \mid x_t)\right] dt.
  \]
  当采用线性掩码调度 $\alpha_t = 1 - t$ 时，权重简化为 $\frac{1}{t}$！对比第 5 章式~\eqref{eq:mlm-loss}，括号内正是标准的双向 BERT 掩码交叉熵损失——唯一本质差别在于 BERT 死守固定 $15\%$ 掩码率，无法从全 \texttt{[MASK]} 生成文本；而 MDLM / LLaDA 在全部掩码率 $t \in (0, 1]$ 上积分训练，从而具备了从 $100\%$ 全掩码画布出发，迭代揭晓最自信词元的生成式去噪能力。

  \item \textbf{相比因果自回归 LLM 的优势与工程硬伤}：
  \begin{itemize}
    \item \textbf{核心优势}：（1）\textbf{并行多词元解码与可控延迟}：单次前向可同时解出多个高置信度位置，有望打破自回归逐词生成的显存带宽墙；（2）\textbf{原生双向全局规划与任意位置填空}：可天然处理首尾约束、反转诅咒（Reversal Curse）与局部修改重绘；（3）\textbf{统一多模态范式}：文本与视觉/音频均可在同一“迭代去噪”框架下联合建模。
    \item \textbf{主要硬伤}：（1）\textbf{丧失精确 KV Cache 复用}：由于每轮去噪都是全序列双向注意力，任意新揭晓的词元都会改变全部其余位置的键和值，无法像因果 LLM 那样直接缓存并增量复用历史前缀（必须借助分块半自回归 Block Diffusion 或近似缓存折中）；（2）\textbf{并行解码的条件独立误差}：一步同时独立采样多个相关词元会忽略词元间的联合依存关系，容易产生局部语法冲突，需靠多步小批量揭示平衡吞吐与长程逻辑一致性；（3）\textbf{可变长生成不够自然}：需预先划定画布长度或依赖 \texttt{[EOS]} 填充机制。
  \end{itemize}
\end{enumerate}


\end{document}
